% =============================================================
% Bölüm 5 - Aritmetik İşlemler
% Bilgisayar Mimarisi Kitabı
% Oğuz Ergin
% =============================================================

\chapter{Aritmetik İşlemler}
\label{chap:aritmetik}
\bolumfoto[Divriği'nin taş oymalarındaki geometrik desenler, matematiğin sanata dönüşmüş halidir. İşlemcinin aritmetik birimi de benzer bir dönüşüm yapar: ikili sayılar üzerinde toplama, çarpma ve bölme işlemlerini donanıma dönüştürür.]{gorseller/bolum05-divrigi.jpg}{Divriği Ulu Camisi Taç Kapısı -- Sivas (Fotoğraf: Engin Mutlu, CC BY-SA 4.0)}

%% Bölüm açılış görseli
% \begin{figure}[htbp]
% \centering
% \includegraphics[width=0.6\textwidth]{sekiller/peri-bacalari.jpg}
% \caption*{Peri Bacaları -- Nevşehir}
% \end{figure}

\begin{tcolorbox}[
    colframe=sinyalyesil!60!black,
    title={\textbf{Öğrenme Hedefleri}},
    fonttitle=\bfseries,
    rounded corners,
    boxrule=0.8pt
]
Bu bölümü tamamladığınızda aşağıdakileri yapabileceksiniz:
\begin{itemize}[nosep, leftmargin=*]
    \item İkilik sayılarda toplama, çıkarma ve taşma denetimini gerçekleştirmek.
    \item Dalga eldeli ve önceden elde üreten toplayıcıların çalışma ilkesini ve hız farkını açıklamak.
    \item Mantık işlemlerini ve AMB'nin iç yapısını kavramak.
    \item Çarpma ve bölme algoritmalarını (ardışık, Booth, geri yüklemeli/geri yüklemesiz, SRT) adım adım uygulamak.
    \item IEEE~754 kayan nokta gösterimini, özel değerlerini ve yuvarlama yöntemlerini açıklamak.
    \item Kayan nokta toplama ve çarpmasının donanımda hangi adımlarla gerçekleştirildiğini betimlemek.
    \item BF16, TF32, FP8 gibi yapay zeka kayan nokta biçimlerinin neden geliştirildiğini ve karma duyarlık eğitiminin nasıl çalıştığını tartışmak.
\end{itemize}
\end{tcolorbox}

\section*{Giriş}

Bilgisayarlar özünde birer hesaplama makinesidir; ilk yapıldıkları günden bu yana ana amaçları, aritmetik işlemleri insanlardan çok daha hızlı ve hatasız biçimde gerçekleştirmektir. Günümüzün bilgisayarları dokunmatik ekranlar, sesli asistanlar ve yapay zeka uygulamalarıyla donatılmış olsa da temellerinde yatan şey hâlâ aynıdır: hesaplama. Bir bilgisayarın kalbi işlemcisiyse, işlemcinin kalbi \emph{aritmetik mantık birimi}\index{aritmetik işlemler} (AMB)'dir. Aslına bakılırsa, bir işlemcinin tüm derdi iki şeye indirgenebilir: bu işlem birimlerinin girişlerine veriyi zamanında getirmek ve birim zamanda olabildiğince çok işlem biriminin çalışmasını sağlamak. Kitabın ilerleyen bölümlerinde göreceğimiz boru hattı, önbellek ve sırasız yürütüm gibi tekniklerin hepsi, özünde bu iki hedefe hizmet eder.

Bu ilke, günümüzün yapay zeka çağında her zamankinden daha geçerlidir. Büyük dil modellerini eğiten ve çalıştıran grafik işlemcileri (GPU) ile yapay zeka hızlandırıcıları, özünde devasa toplayıcı ve çarpıcı dizilerinden oluşur. Örneğin NVIDIA'nın güncel GPU'ları on binlerce kayan nokta çarpma-toplama birimi barındırır; Google'ın TPU hızlandırıcısı ise 128$\times$128'lik çarpıcı dizileriyle tek saat darbesinde on binlerce çarpma-toplama işlemi gerçekleştirir. Bu bölümde öğreneceğimiz toplayıcı, çarpıcı ve kayan nokta birimi tasarımları, tüm bu devlerin temel yapı taşlarıdır.

Peki bir işlemci 0 ve 1'lerden ibaret bitlerle nasıl toplama, çarpma hatta ondalıklı sayılarla işlem yapabiliyor? Üstelik tüm bu farklı sayı türleri (işaretli tam sayılar\index{tam sayı}, kesirler, kayan nokta sayıları) sınırlı bir bit genişliğine sığdırılmak zorunda. Bu bölümde, bilgisayarın bu temel sorunu nasıl çözdüğünü adım adım göreceğiz: mantık kapılarından başlayarak toplama, çıkarma, çarpma ve bölme devrelerini tasarlayacak, sonra da kayan nokta gösterimine geçeceğiz.


% =============================================
\section{Toplama ve Çıkarma}
\label{sec:toplama-cikarma}

\subsection{Tek Bitlik Toplama}

Toplama\index{toplama}\index{addition|see{toplama}} en temel aritmetik işlemidir. Diğer tüm aritmetik işlemler toplama yapılarak gerçekleştirilebilir. Toplama işleminin ikilik tabandaki kuralları yalındır:

\begin{align*}
0 + 0 &= 0 \\
0 + 1 &= 1 \\
1 + 0 &= 1 \\
1 + 1 &= 0 \quad \text{(elde var 1)}
\end{align*}

İkilik tabanda 2 olmadığı için toplanan iki basamakta da 1 rakamı olduğunda sonuç 2 yerine 0 olarak yazılır ve bir sonraki basamağa 1 elde\index{elde}\index{carry|see{elde}} aktarılır. Her bir toplanan basamağa bir önceki basamağın toplamı sonucunda oluşan elde eklenir ve bu toplamdan oluşan elde bir sonraki basamağa geçer. RISC-V'te (RV32I) işlenen veriler 32 bit olduğu için bu toplama işlemi 32 basamağın her biri için ayrı ayrı yapılır.

\begin{ornek}[5.1]
\textbf{Soru:} $24 + 11 = 35$ işleminin ikilik tabanda nasıl yapıldığını gösterin.

\textbf{Çözüm:} 24 sayısının ikilik tabandaki karşılığı $(11000)_2$, 11 sayısının ikilik tabandaki karşılığı $(1011)_2$'dir. Bu iki sayının nasıl toplanacağı aşağıda gösterilmiştir.

\begin{center}
\begin{tabular}{c c c c c c c}
  & &   & $1$ & $1$ & & elde \\
\cline{1-6}
  & $1$ & $1$ & $0$ & $0$ & $0$ & (24) \\
$+$ & $0$ & $1$ & $0$ & $1$ & $1$ & (11) \\
\hline
$1$ & $0$ & $0$ & $0$ & $1$ & $1$ & (35) \\
\end{tabular}
\end{center}

Değeri 1 olan iki basamağın toplanmasından ortaya çıkan elde, soldaki basamağın üzerine eklenerek toplama işlemine devam edilir. En son ortaya çıkan $(100011)_2$ değerinin onluk tabandaki karşılığı 35'tir.
\end{ornek}

Aşağıda RISC-V'te (RV32I) 79 ve 73 değerlerinin nasıl toplandığı görülür:

\begin{center}
\small
\begin{tabular}{r c l}
$\underbrace{00000000000000000000000001001111}_{32 \text{ bit}}$ & $\rightarrow$ & $79$ \\
$+ \; \underbrace{00000000000000000000000001001001}_{32 \text{ bit}}$ & $\rightarrow$ & $+73$ \\
\hline
$\underbrace{00000000000000000000000010011000}_{32 \text{ bit}}$ & $\rightarrow$ & $152$ \\
\end{tabular}
\end{center}

\subsection{Çıkarma}

Çıkarma\index{çıkarma}\index{subtraction|see{çıkarma}} işlemi de aslında toplamanın aynısıdır; çıkarmanın toplamadan tek farkı toplama işleminden önce ikinci işlenenin (çıkarılanın) eksilisinin alınmasıdır. Eğer toplama işlemi $a + b = c$ şeklinde ifade edilirse çıkarma işlemi de toplama kullanılarak $a + (-b) = d$ şeklinde gösterilebilir. Böylece yalnızca toplama işlemi yapabilen bir sayısal devrede kolaylıkla çıkarma işlemi yapılabilir. İkinci işlenenin eksilisinin bulunması, kullanılan sayısal devrenin hangi tür sayı sistemi kullandığına bağlıdır. Güncel işlemciler büyük ölçüde 2'ye tümleyen\index{ikiye tümleyen}\index{two's complement|see{ikiye tümleyen}} kullandığı için bu işlemin ikiye tümleyen alınarak yapıldığı varsayılabilir.

2'ye tümleyen alma işlemi iki adımda gerçekleşir: önce sayının her biti ters çevrilir (0'lar 1, 1'ler 0 olur), ardından sonuca 1 eklenir. Bu yöntem donanım açısından son derece güçlüdür: çıkarma işlemi, mevcut toplayıcı donanımı tekrar kullanılarak yapılabilir. Donanımda $b$'nin her biti DEĞİL kapısıyla ters çevrilir ve toplayıcının giren elde girişi 1 olarak ayarlanır; böylece tek bir ek adımda hem bit terslemesi hem de +1 işlemi birleştirilmiş olur. Aynı toplayıcı devresi, giren elde sinyalinin değerine göre hem toplama hem çıkarma yapabilir. Bu tasarım ilkesinin donanımdaki somut gerçeklemesi Bölüm~\ref{sec:amb-ile-cikarma}'de ele alınmaktadır.

\begin{ornek}[5.2]
\textbf{Soru:} $24 - 11 = ?$ işlemini yapın.

\textbf{Çözüm:} 24 sayısının ikilik tabandaki karşılığı $(011000)_2$'dir. İlk olarak 2'ye tümleyen kullanılarak çıkarma, toplama yardımıyla yapılır:

\begin{center}
\begin{tabular}{r c c c c c c l}
  & $0$ & $1$ & $1$ & $0$ & $0$ & $0$ & (24) \\
$+$ & $1$ & $1$ & $0$ & $1$ & $0$ & $1$ & (-11) \\
\hline
  & $0$ & $0$ & $1$ & $1$ & $0$ & $1$ & (13) \\
\end{tabular}
\end{center}

Bir diğer örnek:
\begin{center}
\begin{tabular}{r c c c c l}
  & $0$ & $1$ & $0$ & $0$ & (9'un alt bitleri) \\
$+$ & $1$ & $1$ & $0$ & $1$ & (-5'in alt bitleri) \\
\hline
  & $0$ & $0$ & $1$ & $0$ & (4'ün alt bitleri) \\
\end{tabular}
\end{center}

Verilen çıkarma işlemlerinde birinci sayı aynen ve ikincisinin 2'ye tümleyeni alınarak çıkarma işlemleri toplama yardımıyla yapılmıştır.

İkilik düzende doğrudan çıkarma da yapılabilir. Çıkarılmak istenen sayı yeterli değilse bir sonraki basamaktan 1 alınır. 2'lik düzende basamaklar birbirinin 2 katı olarak ilerlediği için bir sonraki basamaktan alınan 1, bir önceki basamakta 2 değerindedir.
\end{ornek}

Şimdiye kadar yapılan işlemlerde hep büyük sayılardan küçük sayılar çıkarıldığı için sonuçlar artı çıkmıştır. Aşağıdaki işlemde küçük bir sayıdan daha büyük bir sayı çıkararak eksi sonuç elde edilmiştir:

\begin{center}
\begin{tabular}{r c c c c c c l}
  & $0$ & $1$ & $1$ & $0$ & $0$ & $0$ & (24)  \\
$+$ & $1$ & $0$ & $0$ & $0$ & $0$ & $0$ & (-32) \\
\hline
  & $1$ & $1$ & $1$ & $0$ & $0$ & $0$ & (-8) \\
\end{tabular}
\end{center}

Yukarıda verilen işlemde sonucun en soldaki bitinin 1 olmasından sayının eksi olduğu anlaşılır.

Bir başka örnekte:
\begin{center}
\begin{tabular}{r c c c c l}
  & $0$ & $1$ & $0$ & $1$ & (5)  \\
$+$ & $1$ & $1$ & $1$ & $1$ & (-1) \\
\hline
$1$ & $0$ & $1$ & $0$ & $0$ & (4) \\
\end{tabular}
\end{center}

sonuç toplanan iki sayıdan daha fazla bitten oluşur. İşlenenlerinden biri artı biri eksi olan bir çıkarma işleminde sonucun her iki sayıdan da büyük olması mümkün değildir. Bu nedenle taşan basamak dikkate alınmaz ve sonuç $(0100)_2 = 4$ olarak hesaplanır. İşaretsiz yapılan işlemlerde en anlamlı bitten dışarı çıkan elde taşmayı gösterir. 2'ye tümleyen gösteriminde ise bu elde biti taşma bilgisi taşımaz ve önemsenmez; 2'ye tümleyende taşma, farklı bir yöntemle (en anlamlı bite giren ve çıkan eldelerin XOR'u ile) saptanır (bkz.\ Bilgi Notu~\ref{bilginot:tasma-elde}, s.~\pageref{bilginot:tasma-elde}).

\subsection{32 Bit ile Çıkarma}

Toplama gibi çıkarma için de RISC-V (RV32I) üzerinde işlem yapılacak olan veriler 32 bittir:

\begin{center}
\small
\begin{tabular}{r c l}
$\underbrace{00000000000000000000000001001101}_{32 \text{ bit}}$ & $\rightarrow$ & $77$ \\
$+ \; \underbrace{11111111111111111111111111110001}_{32 \text{ bit}}$ & $\rightarrow$ & $-15$ \\
\hline
$1\underbrace{00000000000000000000000000111110}_{32 \text{ bit}}$ & $\rightarrow$ & $62$ \\
\end{tabular}
\end{center}

32 bitle yapılan örnekte yine artan bir elde biti ile karşılaşılır ancak bu bit önemsenmez ve işlemin sonucu 62 olarak bulunur. Bu fazladan ortaya çıkan bit taşmayı ifade etmez.

\subsection{Taşma}
\label{sec:tasma}

Taşmanın\index{taşma}\index{overflow|see{taşma}} belirli kuralları vardır:
\begin{itemize}
\item 2 tane artı sayının toplamı eksi ediyorsa,
\item 2 tane eksi sayının toplamı artı ediyorsa,
\end{itemize}
taşma oluşur. Aşağıda verilen toplama işlemi taşmaya bir örnektir:

\begin{center}
\begin{tabular}{r c c c c c c c c l}
  & $0$ & $1$ & $1$ & $0$ & $1$ & $0$ & $0$ & $0$ & (104) \\
$+$ & $0$ & $0$ & $1$ & $0$ & $1$ & $1$ & $0$ & $1$ & (45) \\
\hline
  & $1$ & $0$ & $0$ & $1$ & $0$ & $1$ & $0$ & $1$ & $(-107)$ \\
\end{tabular}
\end{center}

Eğer sayı gösterim metodu 2'ye tümleyen olmasaydı toplama işleminin sonucu doğru olacaktı (149). Ancak 2'ye tümleyen yönteminde en anlamlı biti 1 olan sayılar eksi olduğu için bu toplamanın sonucu eksi olmuştur. Belirtilen taşma kuralları gerçekleştiğinde sonuç, mevcut bit sayısıyla doğru biçimde gösterilemez.

Aynı şekilde 2 tane eksi sayının toplanmasından da artı bir sonuç elde edilmesi taşmaya sebep olur:

\begin{center}
\begin{tabular}{r c c c c c c c c l}
  & $1$ & $0$ & $0$ & $1$ & $1$ & $0$ & $0$ & $1$ & (-103) \\
$+$ & $1$ & $0$ & $1$ & $1$ & $1$ & $0$ & $1$ & $1$ & (-69) \\
\hline
$1$ & $0$ & $1$ & $0$ & $1$ & $0$ & $1$ & $0$ & $0$ & (84) \\
\end{tabular}
\end{center}

Elde biti önemsenmediği için ilk basamaktaki 1 iptal edilir ve sonuç 84 olarak bulunur.

Çıkarma işleminde de taşma kurallarının gerçekleşmesine sebep olabilecek durumlar vardır:
\begin{itemize}
\item 1. sayı artı, 2. sayı eksi iken çıkarmanın sonucu eksi ediyorsa,
\item 1. sayı eksi, 2. sayı artı iken çıkarmanın sonucu artı ediyorsa,
\end{itemize}
taşma olur. Dolayısıyla $(A-B)<0$ iken $A>0$, $B<0$ ve $(A-B)>0$ iken $A<0$, $B>0$ durumlarında taşma meydana gelir.

\subsection{RISC-V Aritmetik Buyrukları ve Taşma}

RISC-V buyruk kümesinde \texttt{add}, \texttt{sub} ve \texttt{addi} buyrukları taşma durumunda herhangi bir kural dışı durum (\emph{exception}) üretmez. MIPS'ten farklı olarak RISC-V'te işaretli ve işaretsiz toplama/çıkarma için ayrı buyruklar (örneğin \texttt{addu}, \texttt{addiu}) bulunmaz. Tüm aritmetik buyruklar taşma denetimi yapmadan sonucu hedef yazmacına yazar. Eğer taşma algılanması gerekiyorsa, yazılım tarafından açıkça denetlenmelidir. Bu buyrukların söz dizimi ve taşma davranışı Tablo~\ref{tab:ari-aritmetik-buyruklar}'de özetlenmiştir.

\begin{table}[htbp]
\centering
\caption{Aritmetik İşlem Buyrukları}
\label{tab:ari-aritmetik-buyruklar}
\small
\begin{tabularx}{\textwidth}{l l X c}
\toprule
\textbf{Buyruk} & \textbf{Söz Dizimi} & \textbf{Anlam} & \textbf{Taşma} \\
\midrule
Topla (\emph{Add}) & \texttt{add x1, x2, x3} & x1 = x2 + x3 & Yok \\
Çıkar (\emph{Subtract}) & \texttt{sub x1, x2, x3} & x1 = x2 $-$ x3 & Yok \\
Anlık değerle Topla (\emph{Add immediate}) & \texttt{addi x1, x2, imm} & x1 = x2 + Anlık değer & Yok \\
Çarp (\emph{Multiply}) & \texttt{mul x1, x2, x3} & x1 = (x2 $\times$ x3)[31:0] & Yok \\
Çarp Üst (\emph{Multiply High}) & \texttt{mulh x1, x2, x3} & x1 = (x2 $\times$ x3)[63:32] & Yok \\
Böl (\emph{Divide}) & \texttt{div x1, x2, x3} & x1 = x2 / x3 (bölüm) & Yok \\
Kalan (\emph{Remainder}) & \texttt{rem x1, x2, x3} & x1 = x2 \% x3 (kalan) & Yok \\
\bottomrule
\end{tabularx}
\end{table}

2'ye tümleyen yönteminde en anlamlı bitin (en soldaki bit) bir bakıma işaret biti\index{işaret biti}\index{sign bit|see{işaret biti}} olarak kullanılması buyruk kümesi tasarımını da yalınlaştırır ve daha az buyrukla daha çok iş yapılması sağlanır; bu durum RISC yapısını da destekler. RISC-V'te taşma algılaması donanım tarafından otomatik olarak yapılmaz; gerektiğinde yazılım tarafından işlenenlerin ve sonucun en anlamlı biti denetlenerek taşma olup olmadığı belirlenmelidir. Bölüm~\ref{chap:buyruk-kumesi}'te tanıtılan aritmetik buyrukların kodlama biçimi ve işaret genişletme kuralları bu davranışı doğrudan belirler.


% =============================================
\section{Mantık İşlemleri}
\label{sec:mantik-islemleri}

Temel mantık kapılarının (VE, VEYA, DEĞİL) yapısı ve çalışma ilkeleri Ek~\ref{ch:sayisal-tasarim}'de ayrıntılı olarak ele alınmıştır. Burada bu kapıların RISC-V buyruk kümesindeki karşılıklarına odaklanacağız.

Bir işlemci yalnızca toplama ve çıkarma yapmakla kalmaz; bitleri tek tek incelemek, maskelemek ya da karşılaştırmak gibi işlemlere de ihtiyaç duyar. Mantık işlemleri bu tür bit düzeyindeki işlemleri sağlar. Bellek adresleme, IP adresi maskeleme ve bayrak denetimi gibi pek çok temel işlem mantık buyrukları olmadan gerçekleştirilemez.

Temel mantık işlemleri VE\index{VE kapısı}\index{AND|see{VE kapısı}} (\emph{AND}), VEYA\index{VEYA kapısı}\index{OR|see{VEYA kapısı}} (\emph{OR}) ve DEĞİL\index{DEĞİL kapısı}\index{NOT|see{DEĞİL kapısı}} (\emph{NOT})'dir. Bu işlemler, elektronikte \emph{mantık kapıları}\index{mantık kapısı}\index{logic gate|see{mantık kapısı}} adı verilen temel devre elemanlarıyla gerçekleştirilir. Bu kapıların doğruluk tabloları ve devre sembolleri Ek~\ref{ch:sayisal-tasarim}'de Şekil~\ref{fig:B.1}'de verilmiştir. Kısaca hatırlatmak gerekirse: VE kapısı ancak tüm girişler 1 olduğunda 1 üretir; VEYA kapısı girişlerden en az biri 1 olduğunda 1 üretir; DEĞİL kapısı ise girişin tersini verir.

Mantık işlemlerinin girişleri 1 bit büyüklüğündeki verilerdir. Fakat RISC-V'teki (RV32I) tüm veriler 32 bit olduğu için her basamak bir giriş olarak alınır. Aşağıda 32 bit üzerinde yapılan VE, VEYA ve DEĞİL işlemleri gösterilir:

\begin{center}
\small
\begin{tabular}{r@{\;\;}l}
  \multicolumn{2}{l}{\textbf{VE İşlemi:}} \\[4pt]
  & \texttt{0000\,0000\,0000\,0000\,1010\,1101\,0100\,1101} \\
  VE & \texttt{0000\,0000\,0000\,0011\,0010\,1010\,0111\,0001} \\
  \midrule
  & \texttt{0000\,0000\,0000\,0000\,0010\,1000\,0100\,0001} \\[12pt]
  \multicolumn{2}{l}{\textbf{VEYA İşlemi:}} \\[4pt]
  & \texttt{0000\,0000\,0000\,0000\,1010\,1101\,0100\,1101} \\
  VEYA & \texttt{0000\,0000\,0000\,0011\,0010\,1010\,0111\,0001} \\
  \midrule
  & \texttt{0000\,0000\,0000\,0011\,1010\,1111\,0111\,1101} \\[12pt]
  \multicolumn{2}{l}{\textbf{DEĞİL İşlemi:}} \\[4pt]
  & \texttt{0000\,0000\,0000\,0011\,0010\,1010\,0111\,0001} \\
  \midrule
  DEĞİL & \texttt{1111\,1111\,1111\,1100\,1101\,0101\,1000\,1110} \\
\end{tabular}
\end{center}

Boole cebirinde VE işlemi ``$\cdot$'' (çarpma) ile, VEYA işlemi ``$+$'' (toplama) ile gösterilir. Bu gösterim rastlantısal değildir: VE işleminin çarpmayla, VEYA işleminin de toplamayla benzer davrandığını düşünürsek anlam kazanır. VE işleminde girişlerden herhangi biri 0 ise sonuç 0 olur; tıpkı çarpmada bir çarpanın sıfır olması gibi. VEYA işleminde ise girişlerden biri 1 ise sonuç 1 olur; toplama ile benzer bir sezgidir ($1 + 1 = 1$ olması dışında, çünkü mantık cebiri aritmetik değildir). DEĞİL işlemi ise giriş üzerine çizilen bir çizgi ($\overline{A}$) ile gösterilir. Programlama dillerinde farklı simgeler kullanılır: C dilinde VE için \texttt{\&\&}, VEYA için \texttt{||} ve DEĞİL için \texttt{!} yaygındır; bunlar mantık cebiri gösteriminden farklıdır.

Mantık işlemleri sayının ikilik düzeni (1 ve 0'lardan oluşan gösterimi) üzerinde yapılır.

\subsection{Kaydırma (Shift) İşlemi}

Kaydırma\index{kaydırma}\index{shift|see{kaydırma}} işlemi de mantık işlemleri içerisinde sayılır. RISC-V'te R türü buyruklarda 5 bitlik \texttt{shamt} alanı (funct7'nin alt bitleri) kaç basamak kaydırma yapılacağını gösterir. Sayıyı kaydırmak demek kaydırılmak istenen yönde sayının basamaklarının ilerletilmesidir. Boş kalan basamaklara mantıksal kaydırmada 0 değeri yazılır. Aritmetik kaydırmada\index{aritmetik kaydırma}\index{arithmetic shift|see{aritmetik kaydırma}} ise sayının işaretini gösteren ilk bit kopyalanır. RISC-V'teki mantıksal kaydırma buyrukları \texttt{slli} ve \texttt{srli}, aritmetik kaydırma buyruğu ise \texttt{srai}'dir; bu buyrukların söz dizimi Tablo~\ref{tab:ari-kaydirma-buyruklar}'de yer almaktadır.

\begin{center}\texttt{slli t1, s2, 3} \quad{\small\# s2'deki değeri sola 3 basamak kaydır (mantıksal)}\end{center}\vspace{-6pt}
\begin{center}
\small
\begin{tabular}{r@{\;\;}l}
  s2 & \texttt{0000\,0000\,0000\,0000\,0000\,0010\,1001\,1101} \\
  t1 & \texttt{0000\,0000\,0000\,0000\,0001\,0100\,1110\,1\underline{000}} \\
  & \hfill$\longleftarrow$ \textnormal{sola 3 kaydırma} \\
\end{tabular}
\end{center}

\begin{center}\texttt{srai t0, s0, 5} \quad{\small\# s0'daki değeri sağa 5 basamak kaydır (aritmetik)}\end{center}\vspace{-6pt}
\begin{center}
\small
\begin{tabular}{r@{\;\;}l}
  s0 & \texttt{1111\,1111\,1111\,1100\,1101\,0101\,1000\,1110} \\
  t0 & \texttt{\underline{1111\,1}111\,1111\,1111\,1110\,0110\,1010\,1100} \\
  & \hfill$\longrightarrow$ \textnormal{sağa 5 aritmetik kaydırma} \\
\end{tabular}
\end{center}

\begin{table}[htbp]
\centering
\caption{Bit Düzeyinde Kaydırma Buyrukları}
\label{tab:ari-kaydirma-buyruklar}
\begin{tabularx}{\textwidth}{l l X c}
\toprule
\textbf{Buyruk} & \textbf{Söz Dizimi} & \textbf{Anlam} & \textbf{Biçim} \\
\midrule
Sola Mantıksal Kaydır (\emph{slli}) & \texttt{slli x1, x2, DEĞER} & x1 = x2 $\ll$ Değer & I \\
Sağa Mantıksal Kaydır (\emph{srli}) & \texttt{srli x1, x2, DEĞER} & x1 = x2 $\gg$ Değer & I \\
Sağa Aritmetik Kaydır (\emph{srai}) & \texttt{srai x1, x2, DEĞER} & x1 = x2 $\gg$ (Değer \% 32) & I \\
\bottomrule
\end{tabularx}
\end{table}

\begin{table}[htbp]
\centering
\caption{Mantık Buyrukları}
\label{tab:ari-mantik-buyruklar}
\begin{tabularx}{\textwidth}{l l X c}
\toprule
\textbf{Buyruk} & \textbf{Söz Dizimi} & \textbf{Anlam} & \textbf{Biçim} \\
\midrule
VE (\emph{And}) & \texttt{and x1, x2, x3} & x1 = x2 \& x3 & R \\
Anlık değerle VE (\emph{Andi}) & \texttt{andi x1, x2, DEĞER} & x1 = x2 \& DEĞER & I \\
VEYA (\emph{Or}) & \texttt{or x1, x2, x3} & x1 = x2 $|$ x3 & R \\
Anlık değerle VEYA (\emph{Ori}) & \texttt{ori x1, x2, DEĞER} & x1 = x2 $|$ DEĞER & I \\
Dışlayan VEYA (\emph{Xor}) & \texttt{xor x1, x2, x3} & x1 = x2 $\oplus$ x3 & R \\
Anlık değerle Dışlayan VEYA (\emph{Xori}) & \texttt{xori x1, x2, DEĞER} & x1 = x2 $\oplus$ DEĞER & I \\
Küçükse Birle (\emph{Slt}) & \texttt{slt x1, x2, x3} & x1 = (x2 $<$ x3) & R \\
Anlık Değerle Küçükse Birle (\emph{Slti}) & \texttt{slti x1, x2, DEĞER} & x1 = (x2 $<$ DEĞER) & I \\
\bottomrule
\end{tabularx}
\end{table}

Kaydırma işleminin önemli bir aritmetik anlamı da vardır: bir sayıyı sola $k$ basamak kaydırmak, o sayıyı $2^k$ ile \emph{çarpmaya} eşdeğerdir; sağa $k$ basamak kaydırmak ise $2^k$'ye \emph{bölmeye} eşdeğerdir. Örneğin bir sayıyı sola 3 basamak kaydırmak (\texttt{slli}) o sayıyı $2^3 = 8$ ile çarpar; sağa 2 basamak kaydırmak (\texttt{srli} veya \texttt{srai}) ise 4'e bölmekle aynı sonucu verir. Bu ilişki, ikilik sayı sisteminin doğal bir sonucudur; tıpkı onluk tabanda bir sayıya sıfır eklemenin o sayıyı 10 ile çarpması gibi.

\begin{center}
\small
\begin{tabular}{l c l l}
\toprule
\textbf{İşlem} & \textbf{Buyruk} & \textbf{Aritmetik Karşılığı} & \textbf{Örnek} \\
\midrule
Sola $k$ kaydır & \texttt{slli rd, rs, k} & $\text{rs} \times 2^k$ & \texttt{slli t0, t1, 3} $\rightarrow$ $t1 \times 8$ \\
Sağa $k$ kaydır (mantıksal) & \texttt{srli rd, rs, k} & $\lfloor\text{rs} / 2^k\rfloor$ (işaretsiz) & \texttt{srli t0, t1, 2} $\rightarrow$ $t1 / 4$ \\
Sağa $k$ kaydır (aritmetik) & \texttt{srai rd, rs, k} & $\lfloor\text{rs} / 2^k\rfloor$ (işaretli) & \texttt{srai t0, t1, 1} $\rightarrow$ $t1 / 2$ \\
\bottomrule
\end{tabular}
\end{center}

Bu ilişki yalnızca kuramsal bir gözlem değildir; derleyiciler ve donanım tasarımcıları bunu sürekli kullanır. Çarpma ve bölme buyrukları birden fazla saat çevrimi gerektirirken, kaydırma tek bir çevrimde tamamlanır. Bu nedenle derleyiciler, 2'nin kuvvetiyle yapılan çarpma ve bölme işlemlerini otomatik olarak kaydırma buyruklarına dönüştürür. Örneğin C dilindeki \texttt{x * 16} ifadesi derleyici tarafından \texttt{slli} ile 4 basamak sola kaydırmaya, \texttt{x / 8} ifadesi ise \texttt{srai} (veya \texttt{srli}) ile 3 basamak sağa kaydırmaya çevrilir. Hatta $x \times 10$ gibi 2'nin kuvveti olmayan çarpımlar bile kaydırma ve toplama birleşimiyle gerçekleştirilebilir: $x \times 10 = x \times 8 + x \times 2 = (x \ll 3) + (x \ll 1)$.

\begin{bilginot}[RISC-V'te DEĞİL İşlemi]
Birçok buyruk kümesi mimarisinde bitleri terslemek için özel buyruklar bulunur. Örneğin MIPS'te \texttt{nor} (VEYA-DEĞİL) buyruğu vardır ve DEĞİL işlemi \texttt{nor \$t0, \$s0, \$zero} biçiminde (bir girişi sıfır yaparak) gerçekleştirilir. x86'da ise doğrudan \texttt{NOT} buyruğu mevcuttur. RISC-V ise yalınlık ilkesi gereği ayrı bir DEĞİL veya \texttt{nor} buyruğu tanımlamaz. Bunun yerine DEĞİL işlemi, mevcut Dışlayan VEYA buyruğu ile gerçekleştirilir: \texttt{xori rd, rs, -1} (montaj dilinde \texttt{not rd, rs} sözde buyruğu). $-1$ değeri ikilik tabanda tüm bitleri 1 olan sayıdır; herhangi bir bitin 1 ile Dışlayan VEYA'sı o bitin tersini verir. Böylece tek bir buyrukla tüm bitler terslenir.
\end{bilginot}


% =============================================
\section{AMB (Aritmetik Mantık Birimi)}
\label{sec:amb}

Bu bölümde mantık kapılarını (Ek~\ref{ch:sayisal-tasarim}) kullanarak aritmetik ve mantık işlemlerini gerçekleştirecek bir AMB tasarlayacağız.

Toplama, çıkarma, VE, VEYA: tüm bu işlemleri ayrı ayrı öğrendik. Şimdi sıra geldi hepsini tek bir birimde birleştirmeye. İşte \emph{Aritmetik Mantık Birimi}\index{AMB}\index{aritmetik mantık birimi|see{AMB}}\index{ALU|see{AMB}} (AMB), bir denetim sinyaline bakarak hangi işlemin yapılacağına karar veren ve sonucu üreten birimdir. Tasarım stratejimiz yalın ama güçlüdür: önce 1 bitlik bir AMB tasarlayacağız, sonra bu birimden 32 tanesini yan yana koyarak 32 bitlik bir AMB elde edeceğiz.

1 bitlik AMB'nin desteklemesi gereken temel aritmetik işlemler toplama ve çıkarma; mantık işlemleri ise VE ve VEYA'dır.

AMB'nin çalışma ilkesini anlamak için şu noktayı kavramak gerekir: AMB, hangi işlemin yapılacağına kendisi karar vermez; \emph{tüm işlemleri aynı anda koşut olarak gerçekleştirir}. Toplama birimi toplar, VE kapısı VE işlemini yapar, VEYA kapısı VEYA işlemini yapar; bunların hepsi eş zamanlı çalışır. AMB'nin çıkışındaki çoklayıcı, \emph{denetim sinyaline} bakarak doğru sonucu seçer. Bu denetim sinyali nereden gelir? Yürütülmekte olan buyruğun işlem kodu alanından (\emph{opcode} ve \emph{funct} bitleri) bir denetim birimi aracılığıyla üretilir. Böylece \texttt{add} buyruğu yürütülürken toplama sonucu, \texttt{and} buyruğu yürütülürken VE sonucu çoklayıcı tarafından seçilir; ama her iki durumda da diğer işlemler de yapılmış, sonuçları atılmıştır.

\subsection{AMB ile Toplama}

\subsubsection{Yarım Toplayıcı}

Bir bitlik toplama işleminin doğruluk tablosu\index{doğruluk tablosu}\index{truth table|see{doğruluk tablosu}} ve devre şeması aşağıda verilmiştir. Sonuç sütunu Dışlayan VEYA mantıksal kapısının doğruluk tablosu ve elde sütunu ise VE'nin doğruluk tablosuyla aynı olduğu için girişlerine $a$ ve $b$'nin bağlanacağı bir Dışlayan VEYA ve VE kapısıyla ``yarım toplayıcı\index{yarım toplayıcı}\index{half adder|see{yarım toplayıcı}}'' denilen 1 bitlik toplama işlemi yapan donanım birimi elde edilmiştir. Yarım toplayıcının giriş-çıkış ilişkisi Tablo~\ref{tab:yarim-toplayici}'de, devre şeması ise Şekil~\ref{fig:yarim-toplayici}'de gösterilmiştir.

\begin{table}[htbp]
\centering
\caption{Yarım toplayıcı doğruluk tablosu}
\label{tab:yarim-toplayici}
\begin{tabular}{c c | c c}
\toprule
\textbf{a} & \textbf{b} & \textbf{Toplam} & \textbf{Elde} \\
\midrule
0 & 0 & 0 & 0 \\
0 & 1 & 1 & 0 \\
1 & 0 & 1 & 0 \\
1 & 1 & 0 & 1 \\
\bottomrule
\end{tabular}
\end{table}

\begin{figure}[htbp]
\centering
\begin{tikzpicture}[kapiboy, scale=1.2, transform shape]
    % Kapılar — alt alta, yakın
    \node[xor gate US, draw, thick, fill=blokgri, inputs={nn}] (xor1) at (4, 0.8) {};
    \node[and gate US, draw, thick, fill=blokgri, inputs={nn}] (and1) at (4, -0.8) {};
    % a ve b giriş noktaları
    \coordinate (startA) at (0, 0.4);
    \coordinate (startB) at (0, -0.4);
    \node[left] at (startA) {$a$};
    \node[left] at (startB) {$b$};
    % a dallanma noktası
    \draw (startA) -- ++(1.8,0) coordinate (Asplit);
    \filldraw (Asplit) circle (2pt);
    \draw (Asplit) |- (xor1.input 1);
    \draw (Asplit) |- (and1.input 1);
    % b dallanma noktası
    \draw (startB) -- ++(1.3,0) coordinate (Bsplit);
    \filldraw (Bsplit) circle (2pt);
    \draw (Bsplit) |- (xor1.input 2);
    \draw (Bsplit) |- (and1.input 2);
    % Çıkışlar
    \draw[okstil] (xor1.output) -- ++(1.2,0) node[right, font=\footnotesize] {$T$ (Toplam)};
    \draw[okstil] (and1.output) -- ++(1.2,0) node[right, font=\footnotesize] {$e_{\text{çık}}$ (Elde)};
\end{tikzpicture}
\caption{Yar{\i}m toplay{\i}c{\i} devresi}
\label{fig:yarim-toplayici}
\end{figure}

Yarım toplayıcı tek başına kullanışsız görünebilir; elde girişi olmadan gerçek bir toplama yapılamaz. Ancak yarım toplayıcının asıl değeri, tam toplayıcının yapı taşı olmasıdır: Şekil~\ref{fig:tam-toplayici}'de görüleceği gibi, iki yarım toplayıcı ve bir VEYA kapısı birleştirilerek elde girişini de hesaba katan bir tam toplayıcı elde edilir.

\subsubsection{Tam Toplayıcı}

Toplama işlemi sadece bir bit ile gerçekleştirilmez. Bir önceki basamaktan gelen eldeyi de toplayabilecek bir ``tam toplayıcı\index{tam toplayıcı}\index{full adder|see{tam toplayıcı}}'' kullanılır. Tam toplayıcının olası giriş kombinasyonlarının tümü Tablo~\ref{tab:tam-toplayici}'de verilmiştir.

\begin{table}[htbp]
\centering
\caption{Tam toplayıcı doğruluk tablosu}
\label{tab:tam-toplayici}
\begin{tabular}{c c c | c c}
\toprule
\textbf{a} & \textbf{b} & \textbf{Giren Elde} & \textbf{Toplam} & \textbf{Çıkan Elde} \\
\midrule
0 & 0 & 0 & 0 & 0 \\
0 & 0 & 1 & 1 & 0 \\
0 & 1 & 0 & 1 & 0 \\
0 & 1 & 1 & 0 & 1 \\
1 & 0 & 0 & 1 & 0 \\
1 & 0 & 1 & 0 & 1 \\
1 & 1 & 0 & 0 & 1 \\
1 & 1 & 1 & 1 & 1 \\
\bottomrule
\end{tabular}
\end{table}

Tablodan yararlanarak sadece çıkış sütunları olan ``sonuç'' ve ``çıkan elde''ye bakılarak AMB tasarımı tek bir kapıyla gerçeklenemez. Bu nedenle fonksiyonların sadeleştirilmesi için Karno haritası\index{Karno haritası}\index{Karnaugh map|see{Karno haritası}} yöntemi kullanılır (bkz.\ Ek~B).

\begin{figure}[htbp]
\centering
\begin{tikzpicture}[kapiboy, scale=1.0, transform shape]
    % === Kapılar ===
    % Birinci yarım toplayıcı (kapılar yakın)
    \node[xor gate US, draw, thick, fill=blokgri, inputs={nn}] (xor1) at (3.2, 2.8) {};
    \node[and gate US, draw, thick, fill=blokgri, inputs={nn}] (and1) at (3.2, 1.2) {};
    % İkinci yarım toplayıcı
    \node[xor gate US, draw, thick, fill=blokgri, inputs={nn}] (xor2) at (7.2, 3.3) {};
    \node[and gate US, draw, thick, fill=blokgri, inputs={nn}] (and2) at (7.2, 1.7) {};
    % OR kapısı
    \node[or gate US,  draw, thick, fill=blokgri, inputs={nn}] (or1)  at (9.8, 0.2) {};
    % === Giriş etiketleri ===
    \coordinate (startA) at (0, 2.3);
    \coordinate (startB) at (0, 1.7);
    \coordinate (startC) at (0, -0.8);
    \node[left] at (startA) {$a$};
    \node[left] at (startB) {$b$};
    \node[left] at (startC) {$e_{\text{gir}}$};
    % === a bağlantıları ===
    \draw (startA) -- ++(1.5,0) coordinate (Asplit);
    \filldraw (Asplit) circle (2pt);
    \draw (Asplit) -- (Asplit |- xor1.input 1) -- (xor1.input 1);
    \draw (Asplit) -- (Asplit |- and1.input 1) -- (and1.input 1);
    % === b bağlantıları ===
    \draw (startB) -- ++(0.8,0) coordinate (Bsplit);
    \filldraw (Bsplit) circle (2pt);
    \draw (Bsplit) -- (Bsplit |- xor1.input 2) -- (xor1.input 2);
    \draw (Bsplit) -- (Bsplit |- and1.input 2) -- (and1.input 2);
    % === XOR1 çıkışı → XOR2 ve AND2 ===
    \draw (xor1.output) -- ++(1.4,0) coordinate (X1out);
    \filldraw (X1out) circle (2pt);
    \draw (X1out) -- (X1out |- xor2.input 2) -- (xor2.input 2);
    \draw (X1out) -- (X1out |- and2.input 1) -- (and2.input 1);
    % === c_in bağlantıları ===
    \draw (startC) -- ++(4.7,0) coordinate (Csplit);
    \draw (Csplit) -- (Csplit |- xor2.input 1) -- (xor2.input 1);
    \draw (Csplit) -- (Csplit |- and2.input 2) coordinate (Cdot) -- (and2.input 2);
    \filldraw (Cdot) circle (2pt);
    % === AND çıkışları → OR ===
    \draw (and1.output) -- ++(0.8,0) coordinate (A1out);
    \draw (A1out) -- (A1out |- or1.input 2) -- (or1.input 2);
    \draw (and2.output) -- ++(0.6,0) coordinate (A2out);
    \draw (A2out) -- (A2out |- or1.input 1) -- (or1.input 1);
    % === Çıkışlar ===
    \draw[okstil] (xor2.output) -- ++(1.5,0) node[right] {$T$ (Toplam)};
    \draw[okstil] (or1.output)  -- ++(1.2,0) node[right] {$e_{\text{çık}}$ (Çıkan Elde)};
    % === Yarım toplayıcı grupları ===
    \node[draw, dashed, rounded corners, inner sep=8pt,
          fit=(xor1)(and1),
          label={[font=\scriptsize]above:Yarım Toplayıcı 1}] {};
    \node[draw, dashed, rounded corners, inner sep=8pt,
          fit=(xor2)(and2),
          label={[font=\scriptsize]above:Yarım Toplayıcı 2}] {};
\end{tikzpicture}
\caption{Tam toplay{\i}c{\i} devre {\c{s}}emas{\i}}
\label{fig:tam-toplayici}
\end{figure}

Bir tane tam toplayıcı ile 1 bitlik toplama yapılabilir. Peki çok basamaklı sayıları nasıl toplarız? Elle toplama yaparken en sağdaki basamaktan (birler basamağı) başlayıp sola doğru ilerlediğimiz gibi, donanımda da aynı mantık geçerlidir.

$n$ bitlik bir sayının bitleri sağdan sola doğru 0'dan başlayarak numaralanır. En sağdaki bit (bit~0) \emph{en az anlamlı bit}\index{en az anlamlı bit}\index{LSB|see{en az anlamlı bit}} (\emph{Least Significant Bit}, LSB) olarak adlandırılır; çünkü temsil ettiği değer ($2^0 = 1$) en küçüktür ve sayının toplam değerine en az katkıyı sağlar. En soldaki bit (bit~$n{-}1$) ise \emph{en anlamlı bit}\index{en anlamlı bit}\index{MSB|see{en anlamlı bit}} (\emph{Most Significant Bit}, MSB) olarak adlandırılır; temsil ettiği değer ($2^{n-1}$) en büyüktür ve sayının değerine en çok katkıyı bu bit yapar.

Çok basamaklı toplama yapmak için $n$ tane tam toplayıcı zincirleme bağlanır: her tam toplayıcının çıkan eldesi ($C_{out}$), bir sonraki (solundaki) tam toplayıcının giren eldesine ($C_{in}$) bağlanır. Elde sinyali böylece en az anlamlı bitten en anlamlı bite doğru ``dalgalanarak'' ilerler; bu yapıya \emph{dalga eldeli toplayıcı}\index{dalga eldeli toplayıcı}\index{ripple carry adder|see{dalga eldeli toplayıcı}} (\emph{ripple carry adder}) denir. En sağdaki toplayıcının giren eldesi ($C_0$) toplama için 0, çıkarma için 1 yapılır. En soldaki toplayıcının çıkan eldesi ($C_n$) ise taşma bilgisi taşır.

Şekil~\ref{fig:1bit-toplayici}'de, Şekil~\ref{fig:tam-toplayici}'deki tam toplayıcı devresinin giriş ve çıkış sinyallerini özetleyen blok gösterimi verilmiştir. Bu blok gösterim, çok bitlik toplayıcı tasarımlarında her bir tam toplayıcıyı tek bir kutu olarak çizmeyi kolaylaştırır. Şekil~\ref{fig:4bit-toplayici}'de ise bu blokların zincirlenmesiyle oluşan 4 bitlik dalga eldeli toplayıcının yapısı gösterilmiştir.

\begin{figure}[htbp]
\centering
\begin{tikzpicture}[scale=0.65, transform shape]
    \node[draw, rectangle, minimum width=2.8cm, minimum height=2.2cm,
          fill=blokmavi, thick, font=\normalsize, align=center] (adder) {1 Bitlik\\Toplay{\i}c{\i}};
    % Inputs (top)
    \draw[okstil] (-0.7, 2.0) -- (-0.7, 1.1) node[pos=0, above] {$a$};
    \draw[okstil] (0.7, 2.0) -- (0.7, 1.1) node[pos=0, above] {$b$};
    % Carry in (left)
    \draw[okstil] (-2.2, 0) -- (-1.4, 0) node[pos=0, left] {Giren Elde ($e_{\text{gir}}$)};
    % Sum output (bottom)
    \draw[okstil] (0, -1.1) -- (0, -1.8) node[pos=1, below] {Sonu{\c{c}} ($T$)};
    % Carry out (right)
    \draw[okstil] (1.4, 0) -- (2.2, 0) node[pos=1, right] {{\c{C}}{\i}kan Elde ($e_{\text{çık}}$)};
\end{tikzpicture}
\caption[1 bitlik toplay{\i}c{\i} (blok g{\"o}sterimi)]{1 bitlik toplayıcının blok gösterimi. Şekil~\ref{fig:tam-toplayici}'deki tam toplayıcı devresinin giriş/çıkışlarını özetler.}
\label{fig:1bit-toplayici}
\end{figure}

\begin{figure}[htbp]
\centering
\begin{tikzpicture}[scale=0.85, transform shape]
    \tikzset{
        fa/.style={draw, rectangle, minimum width=2cm, minimum height=1.6cm,
                   fill=blokmavi, font=\small, thick, align=center}
    }
    % Full adder blocks — en anlamlı bit (3) solda, en az anlamlı bit (0) sağda
    \foreach \i in {0,...,3} {
        \pgfmathtruncatemacro{\x}{\i * 3}
        \pgfmathtruncatemacro{\bitidx}{3 - \i}
        \node[fa] (fa\i) at (\x, 0) {Tam\\Toplayıcı};
        % Input labels
        \node[above] at (\x - 0.5, 1.2) {$A_{\bitidx}$};
        \node[above] at (\x + 0.5, 1.2) {$B_{\bitidx}$};
        \draw[okstil] (\x - 0.5, 1.2) -- (\x - 0.5, 0.8);
        \draw[okstil] (\x + 0.5, 1.2) -- (\x + 0.5, 0.8);
        % Output labels
        \node[below] at (\x, -1.2) {$T_{\bitidx}$};
        \draw[okstil] (\x, -0.8) -- (\x, -1.2);
    }
    % Carry chain — elde sağdan sola akar (e_0 en sağda)
    \node[right] at (10.5, 0) {$e_0$};
    \draw[okstil] (10.5, 0) -- (fa3.east);
    \draw[okstil] (fa3.west) -- (fa2.east) node[midway, above, font=\scriptsize] {$e_1$};
    \draw[okstil] (fa2.west) -- (fa1.east) node[midway, above, font=\scriptsize] {$e_2$};
    \draw[okstil] (fa1.west) -- (fa0.east) node[midway, above, font=\scriptsize] {$e_3$};
    \draw[okstil] (fa0.west) -- ++(-1,0) node[left] {$e_4$};
\end{tikzpicture}
\caption[4 bitlik dalga eldeli toplay{\i}c{\i}]{4 bitlik dalga eldeli toplay{\i}c{\i}. Elde sinyali en az anlaml{\i} bitten (sa\u{g}) en anlaml{\i} bite (sol) do\u{g}ru dalgalanarak ilerler.}
\label{fig:4bit-toplayici}
\end{figure}

\FloatBarrier
\begin{bilginot}[Bit Numaralandırması: 0'dan mı 1'den mi?]
Şekil~\ref{fig:4bit-toplayici}'de en sağdaki bit (en az anlamlı bit) 0 numarasını taşır. Bu, RISC-V dahil çoğu çağdaş mimaride (MIPS, ARM, x86) kullanılan standarttır: bit~0 en az anlamlı bittir ve her bitin değeri $2^i$'dir ($i$ = bit numarası). Ancak IBM System/360 (1964) ve onun mirasçıları olan z/Architecture sistemlerinde numaralandırma \emph{tersine} çalışır: bit~0 en anlamlı bittir. Bu iki gelenek zaman zaman karışıklığa yol açar; bir belge okurken hangi numaralandırma kuralının kullanıldığına dikkat etmek gerekir.
\end{bilginot}

\subsection{AMB ile Çıkarma}
\label{sec:amb-ile-cikarma}

Çıkarma işlemi ikinci sayının eksisinin alınıp birinci sayıyla toplanmasıyla yapılır. Sayının 2'ye tümleyeninin alınması için giren elde biti 0 değil 1 alınır (hatırlatma: 2'ye tümleyen almak için tersi alınan sayıya 1 eklenir).

Şekil~\ref{fig:1bit-toplayici}'de gösterilen 1 bitlik toplayıcı bloğunun üzerine çıkarma desteği eklenmiş hali Şekil~\ref{fig:1bit-toplama-cikarma}'de verilmiştir.

Çok bitlik bir toplayıcı-çıkarıcı devresinde, tek bir \emph{Çıkar} denetim sinyali tüm işlemi yönetir. Bu sinyalin iki etkisi vardır: $b$'nin her bitini Dışlayan VEYA kapısından geçirir ve en az anlamlı bitin giren eldesini ayarlar.

\begin{itemize}
    \item \textbf{Çıkar = 0 (toplama):} Dışlayan VEYA kapıları $b$ bitlerini değiştirmeden geçirir; giren elde $e_{g}=0$ olur. Devre $a + b$ hesaplar.
    \item \textbf{Çıkar = 1 (çıkarma):} Dışlayan VEYA kapıları $b$ bitlerini tersler; giren elde $e_{g}=1$ olur. Devre $a + \overline{b} + 1$ hesaplar; bu ise 2'ye tümleyen tanımı gereği $a + (-b) = a - b$'dir.
\end{itemize}

\begin{figure}[htbp]
\centering
\begin{tikzpicture}[kapiboy, scale=1.0, transform shape]
    % === Tam Toplayıcı bloğu (sağda) ===
    \node[draw, rectangle, minimum width=2.4cm, minimum height=2.8cm,
          fill=blokmavi, thick, font=\small, align=center] (fa) at (8.5, 0.5) {Tam\\Toplayıcı};
    % === XOR kapısı (FA alt girişiyle aynı y → çıkış düz yatay) ===
    \node[xor gate US, draw, thick, fill=blokgri, scale=1.6] (xorB) at (4.2, -0.2) {};
    % === Giriş etiketleri ===
    \node[left] at (0, 1.2) {$a$};
    % === a → FA üst giriş (düz yatay, oklu) ===
    \draw[->] (0, 1.2) -- ([yshift=7mm]fa.west);
    \node[above left, font=\scriptsize] at ([yshift=7mm]fa.west) {$a$};
    % === b → XOR üst giriş (düz yatay, kapı girişiyle aynı y) ===
    \node[left] at (xorB.input 1 -| 0,0) {$b$};
    \draw (xorB.input 1 -| 0,0) -- (xorB.input 1);
    % === İşlem → dallanma noktası + XOR alt giriş + c_in ===
    \node[left, font=\small] at (xorB.input 2 -| 0,0) {İşlem};
    \draw (xorB.input 2 -| 0,0) -- ++(0.8,0) coordinate (Isplit);
    \filldraw (Isplit) circle (2pt);
    \draw (Isplit) -- (xorB.input 2);
    \draw[->] (Isplit) -- ++(0,-1.2) coordinate (Iturn) -- (Iturn -| fa.south) -- (fa.south);
    \node[below right, font=\scriptsize] at ([yshift=-1mm]fa.south) {$e_{\text{gir}}$};
    % === XOR çıkışı → FA alt giriş (düz yatay, oklu) ===
    \draw[->] (xorB.output) -- ([yshift=-7mm]fa.west);
    \node[above left, font=\scriptsize] at ([yshift=-7mm]fa.west) {$b/\overline{b}$};
    % === Çıkışlar (sağda) ===
    \draw[->] ([yshift=4mm]fa.east) -- ++(1.5,0) node[right, font=\small] {$e_{\text{çık}}$};
    \draw[->] ([yshift=-4mm]fa.east) -- ++(1.5,0) node[right, font=\small] {Toplam ($T$)};
\end{tikzpicture}
\caption{Toplama ve {\c{c}}{\i}karma yapabilen 1 bitlik AMB}
\label{fig:1bit-toplama-cikarma}
\end{figure}

Şekil~\ref{fig:1bit-toplama-cikarma}'de $b$ girişini terslemek için Dışlayan VEYA kapısı kullanılmıştır: İşlem sinyali 0 olduğunda bu kapı $b$'yi aynen geçirir, 1 olduğunda ise tersler. Gerçek AMB tasarımlarında ise $b$'yi terslemek için çoğunlukla DEĞİL kapısı tercih edilir; çünkü CMOS teknolojisinde DEĞİL kapısı yalnızca 2 transistörle gerçeklenirken Dışlayan VEYA kapısı 8--12 transistör gerektirir.

\begin{figure}[htbp]
\centering
\begin{tikzpicture}[kapiboy, scale=0.8, transform shape]
    \tikzset{
        fa/.style={draw, rectangle, minimum width=1.8cm, minimum height=1.6cm,
                   fill=blokmavi, font=\small, thick, align=center}
    }

    % ── Tam Toplayıcı blokları (soldan sağa: bit3 → bit0; S₀ en sağda) ─
    \node[fa] (fa3) at ( 1.5, 0) {Tam\\Toplayıcı};
    \node[fa] (fa2) at ( 5.5, 0) {Tam\\Toplayıcı};
    \node[fa] (fa1) at ( 9.5, 0) {Tam\\Toplayıcı};
    \node[fa] (fa0) at (13.5, 0) {Tam\\Toplayıcı};

    % ── XOR kapıları (çıkış aşağı bakıyor, b girişi solda) ────────────
    \node[xor gate US, draw, thick, fill=blokgri, rotate=-90] (xor3) at ( 1.1, 2.5) {};
    \node[xor gate US, draw, thick, fill=blokgri, rotate=-90] (xor2) at ( 5.1, 2.5) {};
    \node[xor gate US, draw, thick, fill=blokgri, rotate=-90] (xor1) at ( 9.1, 2.5) {};
    \node[xor gate US, draw, thick, fill=blokgri, rotate=-90] (xor0) at (13.1, 2.5) {};

    % ── Adlandırılmış koordinatlar (anchor hesapları için) ─────────────
    \coordinate (xor3i1) at (xor3.input 1);  \coordinate (xor3i2) at (xor3.input 2);
    \coordinate (xor2i1) at (xor2.input 1);  \coordinate (xor2i2) at (xor2.input 2);
    \coordinate (xor1i1) at (xor1.input 1);  \coordinate (xor1i2) at (xor1.input 2);
    \coordinate (xor0i1) at (xor0.input 1);  \coordinate (xor0i2) at (xor0.input 2);

    % ── b girişleri → XOR input 2 (rotate=-90 sonrası sol/üst giriş) ────
    \draw[->] (xor3i2 |- 0,5.5) node[above]{$b_3$} -- (xor3i2);
    \draw[->] (xor2i2 |- 0,5.5) node[above]{$b_2$} -- (xor2i2);
    \draw[->] (xor1i2 |- 0,5.5) node[above]{$b_1$} -- (xor1i2);
    \draw[->] (xor0i2 |- 0,5.5) node[above]{$b_0$} -- (xor0i2);

    % ── a girişleri (FA'ya doğrudan, XOR'un sağından geçer) ───────────
    \draw[->] ( 1.9, 5.5) node[above]{$a_3$} -- ( 1.9, 0.8);
    \draw[->] ( 5.9, 5.5) node[above]{$a_2$} -- ( 5.9, 0.8);
    \draw[->] ( 9.9, 5.5) node[above]{$a_1$} -- ( 9.9, 0.8);
    \draw[->] (13.9, 5.5) node[above]{$a_0$} -- (13.9, 0.8);

    % ── Çıkar denetim sinyali (kapıların üstünde, etiket sağda) ──────
    \draw (17.5, 4.5) node[right]{Çıkar} -- (-0.5, 4.5);
    % Çıkar → XOR input 1 (rotate=-90 sonrası sağ/üst giriş)
    \filldraw (xor3i1 |- 0,4.5) circle (2pt);
    \draw[->] (xor3i1 |- 0,4.5) -- (xor3i1);
    \filldraw (xor2i1 |- 0,4.5) circle (2pt);
    \draw[->] (xor2i1 |- 0,4.5) -- (xor2i1);
    \filldraw (xor1i1 |- 0,4.5) circle (2pt);
    \draw[->] (xor1i1 |- 0,4.5) -- (xor1i1);
    \filldraw (xor0i1 |- 0,4.5) circle (2pt);
    \draw[->] (xor0i1 |- 0,4.5) -- (xor0i1);

    % ── XOR çıkışları → FA üst girişi ─────────────────────────────────
    \draw[->] (xor3.output) -- (xor3.output |- 0,0.8);
    \draw[->] (xor2.output) -- (xor2.output |- 0,0.8);
    \draw[->] (xor1.output) -- (xor1.output |- 0,0.8);
    \draw[->] (xor0.output) -- (xor0.output |- 0,0.8);
    % b' etiketleri
    \node[left, font=\footnotesize] at ( 1.1, 1.4) {$b_3'$};
    \node[left, font=\footnotesize] at ( 5.1, 1.4) {$b_2'$};
    \node[left, font=\footnotesize] at ( 9.1, 1.4) {$b_1'$};
    \node[left, font=\footnotesize] at (13.1, 1.4) {$b_0'$};

    % ── c_in: Çıkar → en sağdaki FA (bit 0) elde girişi ──────────────
    \draw[->] (15.5, 4.5) -- (15.5, 0) -- (fa0.east);
    \node[right, font=\footnotesize] at (15.5, 2.25) {$e_{\text{gir}}$};

    % ── Elde zinciri (sağdan sola: bit0 → bit3) ───────────────────────
    \draw[->] (fa0.west) -- (fa1.east)
              node[midway, above, font=\footnotesize]{$e_1$};
    \draw[->] (fa1.west) -- (fa2.east)
              node[midway, above, font=\footnotesize]{$e_2$};
    \draw[->] (fa2.west) -- (fa3.east)
              node[midway, above, font=\footnotesize]{$e_3$};
    \draw[->] (fa3.west) -- ++(-1.2,0) node[left, font=\footnotesize]{$e_{\text{çık}}$};

    % ── Sonuç çıkışları (T₀ en sağda) ─────────────────────────────────
    \draw[->] (13.5, -0.8) -- (13.5, -1.5) node[below]{$T_0$};
    \draw[->] ( 9.5, -0.8) -- ( 9.5, -1.5) node[below]{$T_1$};
    \draw[->] ( 5.5, -0.8) -- ( 5.5, -1.5) node[below]{$T_2$};
    \draw[->] ( 1.5, -0.8) -- ( 1.5, -1.5) node[below]{$T_3$};

\end{tikzpicture}
\caption[4 bitlik toplay{\i}c{\i}-{\c{c}}{\i}kar{\i}c{\i}]{4 bitlik toplayıcı-çıkarıcı devresi. Çıkar sinyali 0 olduğunda $b$ değerleri aynen geçer ve toplama yapılır; 1 olduğunda $b$ terslenir ve giren elde 1 yapılarak çıkarma (2'ye tümleyen) gerçekleştirilir.}
\label{fig:toplayici-cikarici}
\end{figure}

Bu tasarımda ek bir toplama devresi gerekmez: aynı donanım üzerinde yalnızca tek bir denetim sinyaliyle hem toplama hem çıkarma işlemi gerçekleştirilir. 4 bitlik toplayıcı-çıkarıcı devresinin şeması Şekil~\ref{fig:toplayici-cikarici}'de verilmiştir.

\subsection{AMB ile Mantık İşlemleri}

Mantık işlemleri için temel kapıların kullanılması yeterlidir. Bu durumda sadece mantık işlemleri yapacak bir birim VE ve VEYA kapıları ve bir çoklayıcı ile yapılabilir; bu birimin devre şeması Şekil~\ref{fig:ve-veya-amb}'de gösterilmiştir. RISC-V buyruk kümesindeki mantık buyrukları Tablo~\ref{tab:ari-mantik-buyruklar}'de listelenmiştir.

\begin{figure}[htbp]
\centering
\begin{tikzpicture}[kapiboy, scale=1.3, transform shape]
    % === MUX bloğu ===
    \node[draw, rectangle, minimum width=1.4cm, minimum height=3.2cm,
          fill=blokgri, thick, font=\scriptsize, align=center] (mux) at (6.5, 0) {2$\times$1\\Çoklayıcı};
    % === Kapılar (çıkışları MUX girişleriyle aynı y) ===
    \pgfmathsetmacro{\andY}{0.8}   % MUX üst giriş y
    \pgfmathsetmacro{\orY}{-0.8}   % MUX alt giriş y
    \node[and gate US, draw, thick, fill=blokgri] (andg) at (4.3, \andY) {};
    \node[or gate US, draw, thick, fill=blokgri]  (org)  at (4.3, \orY) {};
    % === Giriş etiketleri (kapı girişleriyle aynı y → düz yatay) ===
    \node[left, font=\small] at (andg.input 2 -| 0.8,0) {$a$};
    \node[left, font=\small] at (org.input 2 -| 0.8,0) {$b$};
    % === a bağlantıları ===
    \draw (andg.input 2 -| 0.8,0) -- ++(1.2,0) coordinate (Asplit);
    \filldraw (Asplit) circle (2pt);
    \draw (Asplit) -- (Asplit |- andg.input 1) -- (andg.input 1);
    \draw (Asplit) -- (Asplit |- org.input 1) -- (org.input 1);
    % === b bağlantıları (b → OR alt girişle aynı y, nokta sağda) ===
    \draw (org.input 2 -| 0.8,0) -- ++(1.8,0) coordinate (Bsplit);
    \filldraw (Bsplit) circle (2pt);
    \draw (Bsplit) -- (Bsplit |- andg.input 2) -- (andg.input 2);
    \draw (Bsplit) -- (org.input 2);
    % === Kapı çıkışları → MUX (düz yatay) ===
    \draw (andg.output) -- ([yshift=8mm]mux.west);
    \node[right, font=\scriptsize] at ([yshift=8mm, xshift=1mm]mux.west) {0};
    \draw (org.output) -- ([yshift=-8mm]mux.west);
    \node[right, font=\scriptsize] at ([yshift=-8mm, xshift=1mm]mux.west) {1};
    % === İşlem denetim sinyali ===
    \node[above, font=\scriptsize] at (6.5, 2.3) {İşlem};
    \draw[->] (6.5, 2.3) -- (6.5, 1.6);
    % === Çıkış ===
    \draw[->] (mux.east) -- ++(1.2,0) node[right, font=\scriptsize] {Sonuç};
\end{tikzpicture}
\caption{VE ve VEYA i{\c{s}}lemleri yapabilen AMB (1 bitlik)}
\label{fig:ve-veya-amb}
\end{figure}

\subsection{Bir Bitlik AMB}
\label{sec:1bit-amb}

Mantık işlemleri yapan AMB ile aritmetik işlemler yapan AMB birleştirildiğinde ve sonuç seçimi için $4 \times 1$'lik çoklayıcı eklendiğinde 1 bitlik AMB elde edilir; bu birimin tam devre şeması Şekil~\ref{fig:1bit-amb}'de verilmiştir.

\begin{figure}[htbp]
\centering
\begin{tikzpicture}[kapiboy, scale=1.15, transform shape]
    % === SÜTUN 3: 4×1 Çoklayıcı ===
    \node[draw, rectangle, minimum width=2cm, minimum height=4.5cm,
          fill=blokgri, thick, font=\footnotesize, align=center] (mux4) at (9.5, 3.2) {4$\times$1\\Çoklayıcı};
    % === SÜTUN 2: AND, OR, FA (büyütülmüş) ===
    % OR çıkışı MUX 01 girişiyle aynı y (=3.6) olacak şekilde
    \node[and gate US, draw, thick, fill=blokgri, scale=1.3] (andg) at (6, 4.8) {};
    \node[or gate US, draw, thick, fill=blokgri, scale=1.3]  (org)  at (6, 3.6) {};
    \node[draw, rectangle, minimum width=2.2cm, minimum height=1.8cm,
          fill=blokmavi, thick, font=\footnotesize, align=center] (fa) at (6, 1.8) {Tam\\Toplayıcı};
    % === SÜTUN 1: XOR (büyütülmüş, FA alt girişiyle aynı y) ===
    \node[xor gate US, draw, thick, fill=blokgri, scale=1.3] (xorFark) at (2, 1.5) {};
    % === SOL: Giriş etiketleri (kapılara yakın) ===
    \node[left] at (andg.input 1 -| 1.5,0) {$a$};
    \node[left] at (xorFark.input 1 -| 0.5,0) {$b$};
    \node[left, font=\footnotesize] at (xorFark.input 2 -| 0.5,0) {Çıkar};
    % === b → XOR üst giriş (düz yatay) ===
    \draw (xorFark.input 1 -| 0.5,0) -- (xorFark.input 1);
    % === Çıkar → XOR alt giriş (düz yatay) ===
    \draw (xorFark.input 2 -| 0.5,0) -- (xorFark.input 2);
    % === a dallanması (tek nokta, kapılara yakın) ===
    \draw (andg.input 1 -| 1.5,0) -- ++(1.5,0) coordinate (Asplit);
    \filldraw (Asplit) circle (2pt);
    % a → AND üst giriş (düz yatay)
    \draw (Asplit) -- (andg.input 1);
    % a → OR üst giriş (dikey aşağı, yatay sağa)
    \draw (Asplit) |- (org.input 1);
    % a → FA üst giriş (dikey aşağı, yatay sağa)
    \draw (Asplit) |- ([yshift=3mm]fa.west);
    % === XOR çıkışı (b') dallanması ===
    % b' hattı: XOR'dan FA girişine kadar, dallanma noktası FA hizasında
    \draw (xorFark.output) -- ([yshift=-3mm]fa.west) coordinate[pos=0.7] (Bout);
    \filldraw (Bout) circle (2pt);
    % b' → AND alt giriş (dikey yukarı, yatay sağa)
    \draw (Bout) -- (Bout |- andg.input 2) -- (andg.input 2);
    % b' → OR alt giriş (dikey yukarı, yatay sağa — noktayla bağlantı)
    \draw (Bout) -- (Bout |- org.input 2) coordinate (Bor);
    \filldraw (Bor) circle (2pt);
    \draw (Bor) -- (org.input 2);
    % === e_g (alttan FA'ya) ===
    \node[below] at (6, 0.2) {$e_{\text{gir}}$};
    \draw (6, 0.2) -- (fa.south);
    % === e_ç (FA sağından) ===
    \draw[->] (fa.east) -- ++(0.5,0) node[right, font=\footnotesize] {$e_{\text{çık}}$};
    % === Kapı/FA çıkışları → Çoklayıcı (yatay hizalı) ===
    \draw (andg.output) -- ([yshift=16mm]mux4.west);
    \node[right, font=\scriptsize] at ([yshift=16mm, xshift=1mm]mux4.west) {00};
    \draw (org.output) -- ([yshift=4mm]mux4.west);
    \node[right, font=\scriptsize] at ([yshift=4mm, xshift=1mm]mux4.west) {01};
    \draw ([yshift=3mm]fa.east) -- ++(0.3,0) |- ([yshift=-11mm]mux4.west);
    \node[right, font=\scriptsize] at ([yshift=-11mm, xshift=1mm]mux4.west) {10};
    % === AMBİşlem denetim sinyali ===
    \node[above, font=\footnotesize] at (9.5, 6.5) {AMBİşlem};
    \draw[->] (9.5, 6.5) -- (9.5, 5.45);
    % === Sonuç çıkışı ===
    \draw[->] (mux4.east) -- ++(1.2,0) node[right, font=\footnotesize] {Sonuç};
\end{tikzpicture}
\caption{1 bitlik AMB}
\label{fig:1bit-amb}
\end{figure}

Örneğin çıkarma işlemi yapmak için: İşlem denetimi = 10, Çıkar = 1, Giren elde = 1 olmalıdır. Toplama yapmak için ise: İşlem denetimi = 10, Çıkar = 0, Giren elde = 0 alınır. Çıkarma işlemi yapılırken ``çıkar'' ve ``giren elde'' girişlerinin ikisi de 1 alındığı için bu iki girişe birlikte ``çıkar'' girişi adı verilip tek bir değer gönderilebilir.

\begin{bilginot}[Gerçek İşlemcilerde VEYA-DEĞİL ve VE-DEĞİL Kapıları]
Bu bölümde AMB'yi VE, VEYA ve Dışlayan VEYA kapılarıyla tasarladık. Ancak gerçek CMOS devrelerinde bu kapılar doğrudan kullanılmaz. CMOS teknolojisinde \emph{VEYA-DEĞİL} (NOR) ve \emph{VE-DEĞİL} (NAND) kapıları, VE ve VEYA kapılarından daha az transistörle gerçeklenir ve daha hızlı çalışır. Bir CMOS VE kapısı aslında bir VE-DEĞİL kapısı ile ardından gelen bir DEĞİL kapısından (evirici) oluşur; iki aşama gerektirir. VE-DEĞİL kapısı ise tek aşamada sonuç üretir. Bu nedenle yüksek başarımlı AMB tasarımlarında tüm mantık VE-DEĞİL ve VEYA-DEĞİL kapılarıyla kurulur; çıkıştaki evirici gereken yerde eklenir.
\end{bilginot}

\subsection{32 Bitlik AMB}
\label{sec:32bit-amb}

32 bitlik AMB yapmak için 32 tane 1 bitlik AMB birbirine bağlanır; bu yapı Şekil~\ref{fig:32bit-amb}'de gösterilmiştir. 1 bitlik AMB bazı temel işlemleri gerçekleştirir, fakat RISC-V'te bu AMB'nin yapabildiği temel işlerden daha fazla aritmetik işlem buyruğu vardır. Bunlardan birisi \texttt{slt} (\emph{set on less than}) ``küçükse birle'' buyruğudur.

\begin{lstlisting}[style=riscv]
slt t1, s0, s2   # s0 < s2 ise t1 = 1
\end{lstlisting}

$\text{s0} = a$ ve $\text{s2} = b$ olarak alınırsa $a$, $b$'den küçükse sonuç (t1) yazmacının değerinin 1 yapılacağı anlamına gelir. Peki AMB bu karşılaştırmayı nasıl yapar?

Anahtar gözlem şudur: $a < b$ olup olmadığını anlamak için $a - b$ çıkarması yapılır. Eğer sonuç eksiyse $a < b$ demektir. 2'ye tümleyen düzeninde eksi bir sayının en anlamlı biti (en soldaki bit) her zaman 1'dir. Demek ki çıkarma sonucunun en anlamlı bitine bakmak, karşılaştırmanın sonucunu verir:
\begin{itemize}
  \item $(a - b)$'nin en anlamlı biti $= 1$ $\Rightarrow$ sonuç eksi $\Rightarrow$ $a < b$ $\Rightarrow$ çıkış 1 olmalı
  \item $(a - b)$'nin en anlamlı biti $= 0$ $\Rightarrow$ sonuç artı veya sıfır $\Rightarrow$ $a \geq b$ $\Rightarrow$ çıkış 0 olmalı
\end{itemize}

Ancak burada bir incelik vardır: \texttt{slt} buyruğunun sonucu 32 bitlik bir yazmaca yazılır ve bu yazmacın değeri ya 0 ya da 1 olmalıdır. ``1'' değeri ikilik tabanda \texttt{00...001}'dir; en az anlamlı bit 1, diğer 31 bit 0. Oysa çıkarma sonucunun işaret bilgisi \emph{en anlamlı} bittedir. Bu bilginin en az anlamlı bite taşınması gerekir.

Çözüm olarak 1 bitlik AMB'ye ``küçük'' adında dördüncü bir giriş eklenir (Şekil~\ref{fig:1bit-amb-kucuk}). 32 bitlik AMB'de bu giriş şöyle kullanılır:
\begin{enumerate}
  \item Önce $a - b$ çıkarması yapılır (AMBİşlem $= 10$, Çıkar $= 1$).
  \item En anlamlı biti hesaplayan son AMB birimi (bit~31), çıkarma sonucunun işaret bitini üretir.
  \item Bu işaret biti, \emph{en az anlamlı bitin} (bit~0) AMB biriminin ``küçük'' girişine bağlanır.
  \item Bit~0'ın çoklayıcısı AMBİşlem $= 11$ olduğunda bu ``küçük'' girişini seçer.
  \item Diğer tüm bitlerin (bit~1 -- bit~31) ``küçük'' girişi 0'a bağlıdır; böylece bu bitlerin çıkışı 0 olur.
\end{enumerate}
Sonuçta yazmaç ya \texttt{00...001} ($a < b$) ya da \texttt{00...000} ($a \geq b$) değerini alır.

\begin{figure}[htbp]
\centering
\begin{tikzpicture}[kapiboy, scale=1.15, transform shape]
    % === SÜTUN 3: 4×1 Çoklayıcı (Küçük girişi dahil) ===
    \node[draw, rectangle, minimum width=2cm, minimum height=5.5cm,
          fill=blokgri, thick, font=\footnotesize] (mux4) at (9.5, 3.0) {};
    \node[font=\footnotesize, align=center] at (9.5, 4.3) {4$\times$1\\Çoklayıcı};
    % === SÜTUN 2: AND, OR, FA (büyütülmüş) ===
    % VE yukarı → çıkışı MUX 00 ile hizalı; VEYA aşağı → çıkışı MUX 01 ile hizalı
    \node[and gate US, draw, thick, fill=blokgri, scale=1.3] (andg) at (6, 5.0) {};
    \node[or gate US, draw, thick, fill=blokgri, scale=1.3]  (org)  at (6, 3.4) {};
    \node[draw, rectangle, minimum width=2.2cm, minimum height=1.8cm,
          fill=blokmavi, thick, font=\footnotesize, align=center] (fa) at (6, 1.8) {Tam\\Toplayıcı};
    % === SÜTUN 1: XOR (büyütülmüş, FA alt girişiyle aynı y) ===
    \node[xor gate US, draw, thick, fill=blokgri, scale=1.3] (xorFark) at (2, 1.5) {};
    % === SOL: Giriş etiketleri (kapılara yakın) ===
    \node[left] at (andg.input 1 -| 1.5,0) {$a$};
    \node[left] at (xorFark.input 1 -| 0.5,0) {$b$};
    \node[left, font=\footnotesize] at (xorFark.input 2 -| 0.5,0) {Çıkar};
    % === b → XOR üst giriş (düz yatay) ===
    \draw (xorFark.input 1 -| 0.5,0) -- (xorFark.input 1);
    % === Çıkar → XOR alt giriş (düz yatay) ===
    \draw (xorFark.input 2 -| 0.5,0) -- (xorFark.input 2);
    % === a dallanması (tek nokta, kapılara yakın) ===
    \draw (andg.input 1 -| 1.5,0) -- ++(1.5,0) coordinate (Asplit);
    \filldraw (Asplit) circle (2pt);
    % a → AND üst giriş (düz yatay)
    \draw (Asplit) -- (andg.input 1);
    % a → OR üst giriş (dikey aşağı, yatay sağa)
    \draw (Asplit) |- (org.input 1);
    % a → FA üst giriş (dikey aşağı, yatay sağa)
    \draw (Asplit) |- ([yshift=3mm]fa.west);
    % === XOR çıkışı (b') dallanması ===
    \draw (xorFark.output) -- ([yshift=-3mm]fa.west) coordinate[pos=0.7] (Bout);
    \filldraw (Bout) circle (2pt);
    % b' → AND alt giriş (dikey yukarı, yatay sağa)
    \draw (Bout) -- (Bout |- andg.input 2) -- (andg.input 2);
    % b' → OR alt giriş (dikey yukarı, yatay sağa — noktayla bağlantı)
    \draw (Bout) -- (Bout |- org.input 2) coordinate (Bor);
    \filldraw (Bor) circle (2pt);
    \draw (Bor) -- (org.input 2);
    % === e_g (alttan FA'ya) ===
    \node[below] at (6, 0.2) {$e_{\text{gir}}$};
    \draw (6, 0.2) -- (fa.south);
    % === e_ç (FA sağından) ===
    \draw[->] (fa.east) -- ++(0.5,0) node[right, font=\footnotesize] {$e_{\text{çık}}$};
    % === Kapı/FA çıkışları → Çoklayıcı (yatay hizalı) ===
    \draw (andg.output) -- ([yshift=20mm]mux4.west);
    \node[right, font=\scriptsize] at ([yshift=20mm, xshift=1mm]mux4.west) {00};
    \draw (org.output) -- ([yshift=4mm]mux4.west);
    \node[right, font=\scriptsize] at ([yshift=4mm, xshift=1mm]mux4.west) {01};
    \draw ([yshift=3mm]fa.east) -- ++(0.3,0) |- ([yshift=-9mm]mux4.west);
    \node[right, font=\scriptsize] at ([yshift=-9mm, xshift=1mm]mux4.west) {10};
    % === Küçük girişi (slt için, alttan) ===
    \node[below, font=\footnotesize] at (8, 0.2) {Küçük};
    \draw (8, 0.2) |- ([yshift=-20mm]mux4.west);
    \node[right, font=\scriptsize] at ([yshift=-20mm, xshift=1mm]mux4.west) {11};
    % === AMBİşlem denetim sinyali ===
    \node[above, font=\footnotesize] at (9.5, 6.5) {AMBİşlem};
    \draw[->] (9.5, 6.5) -- (9.5, 5.75);
    % === Sonuç çıkışı ===
    \draw[->] (mux4.east) -- ++(1.2,0) node[right, font=\footnotesize] {Sonuç};
\end{tikzpicture}
\caption[``K{\"u}{\c{c}}{\"u}k'' giri{\c{s}}i eklenmi{\c{s}} 1 bitlik AMB]{``K{\"u}{\c{c}}{\"u}k'' giri{\c{s}}i eklenmi{\c{s}} 1 bitlik AMB.}
\label{fig:1bit-amb-kucuk}
\end{figure}

Özetle: AMB zaten çıkarma yapabildiği için küçüklük karşılaştırması için yeni bir devre tasarlamaya gerek yoktur; yalnızca çıkarma sonucunun işaret bitini doğru yere yönlendirmek yeterlidir. Küçükse birle buyruğunun işleme konulması için AMBİşlem denetim sinyalinin ``$(11)_2 = 3$'' olması gerekir.

\subsection{Dallanma Buyruklarının Desteklenmesi}

RISC-V buyruk kümesi mimarisi için AMB'nin yerine getirmesi gereken buyruklardan birisi de dallanma buyruklarıdır. Örneğin, \texttt{bne} (\emph{branch not equal}) eşit değilse dallan buyruğunun çalıştırılabilmesi için bir eşitlik testine ihtiyaç vardır. Bu eşitlik testi,
\[
a = b \implies a - b = 0
\]
eşitliği ile sağlanabilir. Eşitliği test etmek için $a$'dan $b$ çıkarılır ve sonucun 0 olup olmadığı araştırılır. Yazmacın değerinin 0 olması demek tüm bitlerinin 0 olması anlamına gelir. Tüm 1 bitlik AMB'lerden çıkan sonuçlar VEYA kapısına sokulduğunda ve tümü 0 olduğunda kapının çıkışı 0 olacaktır. Sıfır çıkışı girdilerin eşit olup olmadığını test etmeyi sağladığı için değerler eşitse (VEYA'nın çıkışı 0 olduğunda) değer 1, çıkış 1 ise değer 0 olmalıdır. Bu nedenle VEYA kapısının çıkışına bir tersleyici bağlanır.

\begin{figure}[htbp]
\centering
\begin{tikzpicture}[scale=0.75, transform shape]
    \tikzset{
        ambblok/.style={draw, rectangle, minimum width=2.4cm, minimum height=1.6cm,
                       fill=blokmavi, thick, font=\small, align=center}
    }

    % ── AMB blokları (soldan sağa: bit31 → bit0; S₀ en sağda) ─────
    \node[ambblok] (alu31) at (0, 0) {AMB$_{31}$};
    \node at (2.8, 0) {\large$\cdots$};
    \node[ambblok] (alu2) at (5.5, 0) {AMB$_2$};
    \node[ambblok] (alu1) at (9, 0) {AMB$_1$};
    \node[ambblok] (alu0) at (12.5, 0) {AMB$_0$};

    % ── a ve b girişleri (üstten) ─────────────────────────────────
    \foreach \x/\lbl in {0/31, 5.5/2, 9/1, 12.5/0} {
        \draw[->] (\x-0.4, 3.8) node[above]{$a_{\lbl}$} -- (\x-0.4, 0.8);
        \draw[->] (\x+0.4, 3.8) node[above]{$b_{\lbl}$} -- (\x+0.4, 0.8);
    }

    % ── AMBİşlem denetim hattı (sağdan giriş, yatay) ──────────────
    \draw[denetimstil] (15, 2.8) node[right, text=black]{AMB{\.I}{\c{s}}lem} -- (-2.5, 2.8);
    \foreach \x in {0, 5.5, 9, 12.5} {
        \filldraw[sinyalkirmizi] (\x, 2.8) circle (2pt);
        \draw[denetimstil] (\x, 2.8) -- (\x, 0.8);
    }

    % ── Elde zinciri (sağdan sola) ────────────────────────────────
    \draw[okstil] (14.5, 0) -- (alu0.east);
    \node[right, font=\footnotesize, align=left] at (14.5, 0) {$e_0$\\[-2pt]{\scriptsize(Giren Elde)}};
    \draw[okstil] (alu0.west) -- (alu1.east)
              node[midway, above, font=\footnotesize]{$e_1$};
    \draw[okstil] (alu1.west) -- (alu2.east)
              node[midway, above, font=\footnotesize]{$e_2$};
    \draw[okstil] (alu2.west) -- (3.3, 0)
              node[midway, above, font=\footnotesize]{$e_3$};
    \draw[okstil] (2.3, 0) -- (alu31.east)
              node[midway, above, font=\footnotesize]{$e_{31}$};
    \draw[okstil] (alu31.west) -- ++(-1.5, 0) node[left, font=\footnotesize]{$e_{\text{çık}}$};

    % ── Sonuç çıkışları (alttan) ──────────────────────────────────
    \draw[->] (0, -0.8) -- (0, -1.8) node[below]{$T_{31}$};
    \draw[->] (5.5, -0.8) -- (5.5, -1.8) node[below]{$T_2$};
    \draw[->] (9, -0.8) -- (9, -1.8) node[below]{$T_1$};
    \draw[->] (12.5, -0.8) -- (12.5, -1.8) node[below]{$T_0$};

    % ── Taşma saptama (c₃₁ ⊕ c_out, XOR kapısı ile) ──────────────
    \node[draw, circle, inner sep=2pt, fill=blokyesil, font=\small] (xor) at (-2, 1.8) {$\oplus$};
    % c_out bağlantı noktası
    \filldraw (-2, 0) circle (2pt);
    \draw (-2, 0) -- (xor.south);
    % c₃₁ bağlantı noktası (c₃₁ yazısıyla çakışmasın diye sağda)
    \filldraw (2.2, 0) circle (2pt);
    \draw (2.2, 0) -- (2.2, 1.8) -- (xor.east);
    % Taşma çıkışı
    \draw[okstil] (xor.west) -- (-3.5, 1.8) node[left, font=\footnotesize]{Ta{\c{s}}ma};

    % ── Sıfır saptama (VEYA-DEĞİL kapısı ile) ──────────────────
    \node[draw, rounded corners=2pt, fill=blokyesil, font=\scriptsize,
          minimum width=1.8cm, minimum height=0.5cm] (nor) at (14.8, -2.8) {VEYA-DE{\u{G}}{\.I}L};
    \draw[thick] (-0.5, -2.8) -- (nor.west);
    \foreach \x in {0, 5.5, 9, 12.5} {
        \draw (\x, -2.3) -- (\x, -2.8);
        \filldraw (\x, -2.8) circle (2pt);
    }
    \node at (2.8, -2.5) {$\cdots$};
    \draw[okstil] (nor.east) -- ++(1.2, 0)
        node[right, font=\footnotesize]{S{\i}f{\i}r};

    % ── Küçük geri besleme (AMB₃₁ → AMB₀, işaret biti) ──────────
    % AMB₃₁'den aşağı, yatay hat, AMB₁-AMB₀ arasından yukarı, AMB₀ sol girişi
    \draw[okstil, sinyalkirmizi] (0.9, -0.8) -- (0.9, -3.8)
        -- (10.7, -3.8) -- (10.7, -0.3) -- ([yshift=-0.3cm]alu0.west);
    \node[font=\footnotesize, text=sinyalkirmizi, below] at (5.8, -3.8)
        {K{\"u}{\c{c}}{\"u}k: {\.I}{\c{s}}aret biti ($T_{31}$)};

\end{tikzpicture}
\caption{32 Bitlik AMB}
\label{fig:32bit-amb}
\end{figure}

\begin{table}[htbp]
\centering
\caption{32 Bitlik AMB İşlem Tablosu}
\label{tab:amb-islem}
\begin{tabular}{c c l}
\toprule
\textbf{Çıkar} & \textbf{İşlem Denetimi} & \textbf{İşlem} \\
\midrule
0 & 00 & VE \\
0 & 01 & VEYA \\
0 & 10 & Toplama \\
1 & 10 & Çıkarma \\
1 & 11 & Küçükse Birle \\
\bottomrule
\end{tabular}
\end{table}

Bu AMB'ye ayrıca taşma denetimi birimi de eklenmiştir. Taşma denetiminin yapılacağı birime hangi işlemin yapılacağı, işlenecek değerler ve çıkan elde değeri bağlanır. Bölüm~\ref{sec:tasma}'te anlatılan taşma kurallarının gerçekleşip gerçekleşmediği sınanır. AMB'nin desteklediği işlemler ve bu işlemlere karşılık gelen denetim sinyalleri Tablo~\ref{tab:amb-islem}'de özetlenmiştir. AMB'nin tek çevrimli ve çok çevrimli veri yolundaki rolü Bölüm~\ref{chap:islemci-tasarimi}'da ele alınacaktır.

\begin{bilginot}[Taşma Denetiminde Elde Kuralı]\label{bilginot:tasma-elde}
Toplamada en anlamlı bitlerin bulunduğu kolona giren elde ($e_{31}$) ve çıkan elde ($e_{\text{çık}}$) birbirinden farklıysa taşma olur. Bu kural, Şekil~\ref{fig:32bit-amb}'deki dışlayan VEYA kapısının ($\oplus$) temelini oluşturur:

\begin{itemize}
\item $e_{31} = 0$, $e_{\text{çık}} = 1$: İki artı sayı toplanmış, elde en anlamlı bitten dışarı taşmıştır; sonuç yanlışlıkla eksi çıkar.
\item $e_{31} = 1$, $e_{\text{çık}} = 0$: İki eksi sayı toplanmış, en anlamlı bite elde gelmesine karşın dışarı elde çıkmamıştır; sonuç yanlışlıkla artı çıkar.
\item $e_{31} = e_{\text{çık}}$ olduğunda ise taşma yoktur.
\end{itemize}

Çıkarma ($A - B$), donanımda $A + \overline{B} + 1$ biçiminde gerçekleştirildiği için aynı elde kuralı çıkarma taşmasını da doğru biçimde saptar.
\end{bilginot}

\subsection{Eldenin Önceden Üretilmesi (Carry Lookahead)}\index{elden önceden üretme}\index{carry lookahead|see{elden önceden üretme}}
\label{sec:carry-lookahead}

1 bitlik AMB, 32 bitlik AMB'nin temel birimidir ve 32 bitlik AMB yapımında 32 tane 1 bitlik AMB kullanılır. AMB'nin bu şekilde tasarlanması pek çok sorunu da beraberinde getirmiştir. Bunlardan en önemlisi elde zincirinin yarattığı gecikmedir: her 1 bitlik toplayıcı, çıkan eldesini hesaplayabilmek için bir önceki basamaktan gelen eldeyi beklemek zorundadır.

\subsubsection*{Dalga Eldeli Toplayıcının Gecikme Çözümlemesi}
Her 1 bitlik toplayıcının elde çıkışını üretmesi 1 kapı gecikmesi (1G) sürsün. 4 bitlik bir dalga eldeli toplayıcıda gecikme şöyle birikir:

\begin{itemize}
\item AMB$_0$: $e_0$ hazırdır, 1G sonra $e_1$ üretilir.
\item AMB$_1$: $e_1$'i bekler, 1G daha harcar $\rightarrow$ $e_2$ toplam 2G'de hazır olur.
\item AMB$_2$: $e_2$'yi bekler $\rightarrow$ $e_3$ toplam 3G'de hazır olur.
\item AMB$_3$: $e_3$'ü bekler $\rightarrow$ $e_4$ toplam 4G'de hazır olur.
\end{itemize}

Genel olarak $n$ bitlik bir dalga eldeli toplayıcının elde gecikmesi $n$G'dir. 32 bitlik bir toplayıcıda son elde ($e_{32}$) ancak 32G sonra hazır olur; bu süre boyunca sonuç bitleri de güvenilir değildir. Bu yöntemle yapılmış olan toplayıcılara \emph{dalga eldeli toplayıcı}\index{dalga eldeli toplayıcı}\index{ripple carry adder|see{dalga eldeli toplayıcı}} denir ve elde sinyalinin basamaktan basamağa dalga gibi ilerlemesinden adını alır.

\subsubsection*{Geçirme ve Üretme Kavramları}
Dalga eldeli toplayıcının temel sorunu, her basamağın bir öncekinin eldesine bağlı olmasıdır. Peki bir basamağın elde üretip üretmeyeceği önceden bilinebilir mi? Bu sorunun yanıtı, her basamaktaki $a_i$ ve $b_i$ değerlerine bakılarak verilebilir.

Elde değeri şöyle ifade edilir:
\begin{equation}
e_{i+1} = a_i b_i + a_i e_i + b_i e_i = a_i b_i + e_i(a_i + b_i)
\label{eq:elde}
\end{equation}

Bu denklemde iki ayrı durum gizlidir. $a_i$ ve $b_i$'nin ikisi de 1 ise giren eldeden bağımsız olarak elde \emph{üretilir}. $a_i$ ve $b_i$'den yalnızca biri 1 ise giren elde varsa çıkan elde de olur; başka bir deyişle elde \emph{geçirilir}. Bu iki davranışı ayrı ayrı tanımlayan $p_i$ ve $g_i$ eşitliklerinden yararlanılır:
\begin{itemize}
\item $p_i$ (\emph{propagate} -- geçirme)\index{geçirme}\index{propagate|see{geçirme}}: $p_i = a_i \oplus b_i$
\item $g_i$ (\emph{generate} -- üretme)\index{üretme}\index{generate|see{üretme}}: $g_i = a_i \cdot b_i$
\end{itemize}

Bitlerin ikisi de 1 ise elde üretilir ($g_i$), ikisi de farklı ise elde geçirilir ($p_i$). Bir sonraki basamakta elde oluşması için bir önceki basamaktan gelmiş olan eldenin geçirilmesi ya da elde üretilmesi gerekir:
\begin{equation}
e_{i+1} = g_i + e_i \cdot p_i
\end{equation}

Tüm basamaklardaki elde durumları:
\begin{align}
e_1 &= g_0 + e_0 p_0 \label{eq:c1} \\
e_2 &= g_1 + e_1 p_1 \label{eq:c2} \\
e_3 &= g_2 + e_2 p_2 \label{eq:c3} \\
e_4 &= g_3 + e_3 p_3 \label{eq:c4}
\end{align}

Eşitlikler birbirinin yerine konularak açılırsa:
\begin{align}
e_1 &= g_0 + p_0 e_0 \\
e_2 &= g_1 + g_0 p_1 + p_0 p_1 e_0 \\
e_3 &= g_2 + g_1 p_2 + g_0 p_2 p_1 + p_0 p_1 p_2 e_0 \\
e_4 &= g_3 + g_2 p_3 + g_1 p_3 p_2 + g_0 p_3 p_2 p_1 + p_0 p_1 p_2 p_3 e_0
\end{align}
Bu denklemlerin mantık kapılarıyla doğrudan gerçekleştirilmesi Şekil~\ref{fig:onceden-elde}'de gösterilmiştir.

\begin{figure}[htbp]
\centering
\begin{tikzpicture}[kapiboy, scale=0.78, transform shape]

%% ---------------------------------------------------------------
%% SOL SÜTUN — üst: c₁ = g₀ + p₀·c₀
%% ---------------------------------------------------------------
\begin{scope}[shift={(0,0)}]
  \node[font=\small] at (4.5, 1.2) {$e_1 = g_0 + p_0 \cdot e_0$};

  \node[and gate US, draw, thick, fill=blokgri, inputs={nn}] (a1) at (3, -0.3) {};
  \node[or gate US, draw, thick, fill=blokgri, inputs={nn}]  (o1) at (6, 0.3) {};

  % VE girişleri
  \draw (a1.input 1) -- ++(-2.0,0) node[left, font=\footnotesize] {$p_0$};
  \draw (a1.input 2) -- ++(-2.0,0) node[left, font=\footnotesize] {$e_0$};

  % g₀ → VEYA üst giriş (VE kapısının üstünden, kesişme yok)
  \draw (o1.input 1) -- ++(-5.0,0) node[left, font=\footnotesize] {$g_0$};

  % VE → VEYA alt giriş
  \draw[okstil] (a1.output) -- ++(0.3,0) |- (o1.input 2);

  % Çıkış
  \draw[okstil] (o1.output) -- ++(0.7,0) node[right, font=\footnotesize] {$e_1$};
\end{scope}

%% ---------------------------------------------------------------
%% SOL SÜTUN — alt: c₂ = g₁ + g₀·p₁ + p₀·p₁·c₀
%% ---------------------------------------------------------------
\begin{scope}[shift={(0,-4.5)}]
  \node[font=\small, align=center] at (4.5, 1.5)
      {$e_2 = g_1 + g_0 p_1 + p_0 p_1 e_0$};

  \node[and gate US, draw, thick, fill=blokgri, inputs={nn}]  (a20) at (3, -0.3) {};
  \node[and gate US, draw, thick, fill=blokgri, inputs={nnn}] (a21) at (3, -2.2) {};
  \node[or gate US, draw, thick, fill=blokgri, inputs={nnn}]  (o2) at (6.3, 0.3) {};

  % a20 girişleri (g₀, p₁)
  \draw (a20.input 1) -- ++(-2.0,0) node[left, font=\footnotesize] {$g_0$};
  \draw (a20.input 2) -- ++(-2.0,0) node[left, font=\footnotesize] {$p_1$};

  % a21 girişleri (p₀, p₁, c₀)
  \draw (a21.input 1) -- ++(-2.0,0) node[left, font=\footnotesize] {$p_0$};
  \draw (a21.input 2) -- ++(-2.0,0) node[left, font=\footnotesize] {$p_1$};
  \draw (a21.input 3) -- ++(-2.0,0) node[left, font=\footnotesize] {$e_0$};

  % g₁ → VEYA üst giriş (VE kapılarının üstünden, kesişme yok)
  \draw (o2.input 1) -- ++(-5.3,0) node[left, font=\footnotesize] {$g_1$};

  % VE → VEYA (kademeli dirsekler)
  \draw[okstil] (a20.output) -- ++(0.3,0) |- (o2.input 2);
  \draw[okstil] (a21.output) -- ++(0.5,0) |- (o2.input 3);

  % Çıkış
  \draw[okstil] (o2.output) -- ++(0.7,0) node[right, font=\footnotesize] {$e_2$};
\end{scope}

%% ---------------------------------------------------------------
%% SAĞ SÜTUN — c₄ = g₃ + g₂·p₃ + g₁·p₂·p₃ + g₀·p₁·p₂·p₃ + p₀·p₁·p₂·p₃·c₀
%% ---------------------------------------------------------------
\begin{scope}[shift={(9,0)}]
  \node[font=\small, align=center] at (5.0, 1.2)
      {$e_4 = g_3 + g_2 p_3 + g_1 p_2 p_3 + g_0 p_1 p_2 p_3 + p_0 p_1 p_2 p_3 e_0$};

  % 5-girişli VEYA kapısı
  \node[or gate US, draw, thick, fill=blokgri, inputs={nnnnn}] (o4) at (7, 0.0) {};

  % VE kapıları (tümü VEYA üst girişinin altında → g₃ ile kesişme yok)
  \node[and gate US, draw, thick, fill=blokgri, inputs={nn}]    (a40) at (3, -0.3) {};
  \node[and gate US, draw, thick, fill=blokgri, inputs={nnn}]   (a41) at (3, -2.2) {};
  \node[and gate US, draw, thick, fill=blokgri, inputs={nnnn}]  (a42) at (3, -4.5) {};
  \node[and gate US, draw, thick, fill=blokgri, inputs={nnnnn}] (a43) at (3, -7.2) {};

  % g₃ → VEYA üst giriş (düz yatay hat, tüm VE kapılarının üstünde)
  \draw (o4.input 1) -- ++(-6.0,0) node[left, font=\footnotesize] {$g_3$};

  % --- a40 girişleri (g₂, p₃) ---
  \draw (a40.input 1) -- ++(-2.0,0) node[left, font=\footnotesize] {$g_2$};
  \draw (a40.input 2) -- ++(-2.0,0) node[left, font=\footnotesize] {$p_3$};

  % --- a41 girişleri (g₁, p₂, p₃) ---
  \draw (a41.input 1) -- ++(-2.0,0) node[left, font=\footnotesize] {$g_1$};
  \draw (a41.input 2) -- ++(-2.0,0) node[left, font=\footnotesize] {$p_2$};
  \draw (a41.input 3) -- ++(-2.0,0) node[left, font=\footnotesize] {$p_3$};

  % --- a42 girişleri (g₀, p₁, p₂, p₃) ---
  \draw (a42.input 1) -- ++(-2.0,0) node[left, font=\footnotesize] {$g_0$};
  \draw (a42.input 2) -- ++(-2.0,0) node[left, font=\footnotesize] {$p_1$};
  \draw (a42.input 3) -- ++(-2.0,0) node[left, font=\footnotesize] {$p_2$};
  \draw (a42.input 4) -- ++(-2.0,0) node[left, font=\footnotesize] {$p_3$};

  % --- a43 girişleri (p₀, p₁, p₂, p₃, c₀) ---
  \draw (a43.input 1) -- ++(-2.0,0) node[left, font=\footnotesize] {$p_0$};
  \draw (a43.input 2) -- ++(-2.0,0) node[left, font=\footnotesize] {$p_1$};
  \draw (a43.input 3) -- ++(-2.0,0) node[left, font=\footnotesize] {$p_2$};
  \draw (a43.input 4) -- ++(-2.0,0) node[left, font=\footnotesize] {$p_3$};
  \draw (a43.input 5) -- ++(-2.0,0) node[left, font=\footnotesize] {$e_0$};

  % --- VE → VEYA (artan x-ofseti: alt kapılar dışarıda → kesişme yok) ---
  \draw[okstil] (a40.output) -- ++(0.15,0) |- (o4.input 2);
  \draw[okstil] (a41.output) -- ++(0.35,0) |- (o4.input 3);
  \draw[okstil] (a42.output) -- ++(0.55,0) |- (o4.input 4);
  \draw[okstil] (a43.output) -- ++(0.75,0) |- (o4.input 5);

  % Çıkış
  \draw[okstil] (o4.output) -- ++(0.7,0) node[right, font=\footnotesize] {$e_4$};
\end{scope}

\end{tikzpicture}
\caption[Önceden elde üreten devre]{Önceden elde üreten devrenin mantık kapılarıyla gerçekleştirilmesi ($e_3$ devresi $e_2$ ile benzer yapıdadır).}
\label{fig:onceden-elde}
\end{figure}

4 bitteki eldeler hesaplanırken hiçbir elde değeri bir öncekine bağlantılı değildir. Yalnızca hepsi ilk giren eldeye ($e_0$) bağlıdır ve ilk giren elde işlemin başında hazır bulunacağı için gecikmeye sebep olmaz.

Bu yaklaşımın hız kazancı sağlaması için önemli bir ön koşul vardır: çok girişli kapıların (örneğin 5 girişli VE veya VEYA kapısının) gecikmesi, aynı işlevi zincirleme bağlanmış çok sayıda 2 girişli kapıyla gerçekleştirmenin gecikmesinden \emph{az} olmalıdır. CMOS teknolojisinde çok girişli bir kapının gecikmesi giriş sayısıyla birlikte artar; ancak bu artış, aynı sonucu elde etmek için gereken zincirleme 2 girişli kapıların toplam gecikmesinden çok daha düşüktür. Örneğin 5 girişli bir VE kapısı tek bir kapı gecikmesinde sonuç üretirken, aynı işlevi dört adet zincirleme bağlanmış 2 girişli VE kapısıyla yapmak dört kapı gecikmesi gerektirir. Bu fark, önceden elde üretiminin temel üstünlüğünü oluşturur.

Ancak bir kapının giriş sayısı sınırsız değildir. Bir mantık kapısının sahip olabileceği en fazla giriş sayısı \emph{kapı giriş sayısı kısıtı}\index{kapı giriş sayısı kısıtı}\index{fan-in|see{kapı giriş sayısı kısıtı}} (\emph{fan-in}) olarak adlandırılır. CMOS teknolojisinde giriş sayısı arttıkça kapının gecikmesi ve alan tüketimi belirgin biçimde artar; belirli bir eşiğin ötesinde kapı güvenilir çalışmaz. Bu nedenle önceden elde üretimi doğrudan 32 veya 64 bit için uygulanamaz; $e_{32}$ denklemi 33 girişli kapılar gerektirirdi ki bu uygulamada gerçekleştirilebilir değildir. Çözüm olarak önceden elde üretimi 4'er bitlik gruplar halinde yapılır ve gruplar arasında hiyerarşik bir yapı kurulur.

Bu yöntem pek çok mantıksal kapının kullanılmasını gerektirir ve donanımı karmaşıklaştırır.

Bu nedenle önceden elde üretici devreyi 4'er bitlik gruplar için yaparak gruplardan çıkan değerleri diğer 4'lü gruplara yönlendirecek bloklara ihtiyaç vardır; bu yapı Şekil~\ref{fig:4bit-onceden-elde}'de gösterilmiştir. Blok düzeysinde fonksiyonlar tanımlanır:
\begin{align}
P_0 &= p_0 p_1 p_2 p_3 \\
G_0 &= g_3 + g_2 p_3 + g_1 p_3 p_2 + g_0 p_3 p_2 p_1
\end{align}

16 bitte işlem yapan bir toplayıcı için:
\begin{align}
P_1 &= p_4 p_5 p_6 p_7, \quad G_1 = g_7 + g_6 p_7 + g_5 p_7 p_6 + g_4 p_7 p_6 p_5 \\
P_2 &= p_8 p_9 p_{10} p_{11}, \quad G_2 = g_{11} + g_{10} p_{11} + g_9 p_{11} p_{10} + g_8 p_{11} p_{10} p_9 \\
P_3 &= p_{12} p_{13} p_{14} p_{15}, \quad G_3 = g_{15} + g_{14} p_{15} + g_{13} p_{15} p_{14} + g_{12} p_{15} p_{14} p_{13}
\end{align}

\begin{align}
e_4 &= G_0 + P_0 e_0 \\
e_8 &= G_1 + P_1 e_4 \\
e_{12} &= G_2 + P_2 e_8 \\
e_{16} &= G_3 + P_3 e_{12}
\end{align}

\begin{figure}[htbp]
\centering
\begin{tikzpicture}[scale=0.78, transform shape]
    \tikzset{
        alubox/.style={draw, rectangle, minimum width=2cm, minimum height=1.5cm,
                       fill=blokmavi, thick, font=\small, align=center},
        clabox/.style={draw, rectangle, fill=blokyesil, thick, font=\small, align=center}
    }

    % ── CLA bloğu (üstte) ────────────────────────────────────────
    \node[clabox, minimum width=15cm, minimum height=1.4cm] (cla) at (6, 4)
        {{\"O}nceden Elde {\"U}retici (CLA)};

    % ── AMB blokları (altta, soldan sağa: bit3 → bit0) ──────────
    \node[alubox] (amb3) at (-1,   0) {AMB$_3$};
    \node[alubox] (amb2) at (3.5,  0) {AMB$_2$};
    \node[alubox] (amb1) at (8,    0) {AMB$_1$};
    \node[alubox] (amb0) at (12.5, 0) {AMB$_0$};

    % ── a, b girişleri (soldan AMB'ye giriyor, ok sağa bakıyor) ──
    \foreach \amb/\lbl in {amb3/3, amb2/2, amb1/1, amb0/0} {
        \draw[okstil] ([yshift=3mm]\amb.west) ++(-1.2,0)
            node[left, font=\scriptsize] {$a_{\lbl}$} -- ([yshift=3mm]\amb.west);
        \draw[okstil] ([yshift=-3mm]\amb.west) ++(-1.2,0)
            node[left, font=\scriptsize] {$b_{\lbl}$} -- ([yshift=-3mm]\amb.west);
    }

    % ── T çıkışları (alttan) ─────────────────────────────────────
    \foreach \amb/\lbl in {amb3/3, amb2/2, amb1/1, amb0/0} {
        \draw[okstil] (\amb.south) -- ++(0,-0.8)
            node[below, font=\scriptsize] {$T_{\lbl}$};
    }

    % ── p, g sinyalleri (AMB üstünden CLA altına) ───────────────
    \foreach \x/\lbl in {-1/3, 3.5/2, 8/1, 12.5/0} {
        \draw[okstil, sinyalmavi] (\x-0.3, 0.75) -- (\x-0.3, 3.3)
            node[pos=0.5, left, font=\scriptsize] {$p_{\lbl}$};
        \draw[okstil, sinyalkirmizi] (\x+0.3, 0.75) -- (\x+0.3, 3.3)
            node[pos=0.5, right, font=\scriptsize] {$g_{\lbl}$};
    }

    % ── e₀ girişi (sağdan, CLA'ya ve AMB₀'a dallanır) ──────────
    \coordinate (c0in) at (15, 4);
    \node[right, font=\scriptsize] at (c0in) {$e_0$};
    \draw[okstil] (c0in) -- (cla.east);
    % e₀ aynı zamanda AMB₀'ın elde girişi
    \draw[okstil, sinyalyesil] (15, 3.3) -- (15, 0) -- (amb0.east)
        node[pos=0.15, right, font=\scriptsize, xshift=2pt] {$e_0$};
    \filldraw (15, 4) circle (2pt);
    \draw (15, 4) -- (15, 3.3);

    % ── e₄ çıkışı (CLA solundan) ─────────────────────────────────
    \draw[okstil] (cla.west) -- ++(-1.5, 0)
        node[left, font=\scriptsize] {$e_4$};

    % ── Elde sinyalleri (CLA'dan aşağı, AMB sağ tarafına) ───────
    % e₃: CLA'dan aşağı, AMB₃ sağ kenarına
    \draw[okstil, sinyalyesil] (0.5, 3.3) -- (0.5, 0) -- (amb3.east)
        node[pos=0.15, right, font=\scriptsize] {$e_3$};
    % e₂: CLA'dan aşağı, AMB₂ sağ kenarına
    \draw[okstil, sinyalyesil] (5, 3.3) -- (5, 0) -- (amb2.east)
        node[pos=0.15, right, font=\scriptsize] {$e_2$};
    % e₁: CLA'dan aşağı, AMB₁ sağ kenarına
    \draw[okstil, sinyalyesil] (9.5, 3.3) -- (9.5, 0) -- (amb1.east)
        node[pos=0.15, right, font=\scriptsize] {$e_1$};

    % P₀, G₀ bu şekilde gösterilmiyor; hiyerarşik CLA'da ayrıca ele alınacak.

\end{tikzpicture}
\caption{4 Bitlik {\"o}nceden elde {\"u}retici}
\label{fig:4bit-onceden-elde}
\end{figure}

Yapılacak olan AMB'nin eldeyi önceden üreten devrelerle kurulması, AMB'ye hız kazandırır. Çünkü dalga eldeli toplayıcıdaki gibi her basamakta elde değerinin hesaplanmasının bekleme süresi yoktur. Bu gecikmelerin boru hattı tasarımına etkisi Bölüm~\ref{chap:boruhatti}'de incelenecektir.

\subsection{Elde ile Seçmeli Toplayıcı (Carry Select Adder)}
\label{sec:carry-select}
\index{elde ile seçmeli toplayıcı}\index{carry select adder|see{elde ile seçmeli toplayıcı}}

Önceden elde üreten toplayıcı, elde sinyalini mantıksal denklemlerle \emph{önceden üreterek} hızlanma sağlar. Elde ile seçmeli toplayıcı ise farklı bir strateji izler: sonucu \emph{önceden hesaplayarak} hızlanır. Temel fikir şudur: bir sonraki blok, önceki bloğun eldesini beklemek yerine her iki olasılığı da ($e_{g} = 0$ ve $e_{g} = 1$) koşut olarak hesaplar. Önceki bloğun eldesi hazır olduğunda bir çoklayıcıyla doğru sonuç seçilir.

Şekil~\ref{fig:carry-select-8bit}, 8 bitlik bir elde ile seçmeli toplayıcıyı göstermektedir. İlk 4 bitlik blok normal çalışır. İkinci blok ise iki kopya halinde işlem yapar: biri $e_4 = 0$, diğeri $e_4 = 1$ varsayarak. İlk bloğun $e_4$ çıkışı hazır olduğunda çoklayıcı doğru toplamı ve eldeyi seçer.

\begin{figure}[htbp]
\centering
\begin{tikzpicture}[scale=0.78, transform shape,
    blok/.style={draw, rectangle, fill=blokmavi, thick, font=\footnotesize,
                 align=center, minimum width=2.8cm, minimum height=2cm},
    muxstil/.style={draw, trapezium, trapezium angle=70, shape border rotate=270,
                    fill=bloksari, thick, font=\scriptsize,
                    minimum width=0.8cm, minimum height=0.6cm}]

    % ══════════════════════════════════════════════════════════════
    % Yerleşim: MSB (bit [7:4]) solda, LSB (bit [3:0]) sağda.
    % Blok 1a/1b üstte yan yana, çoklayıcılar altlarında.
    % ══════════════════════════════════════════════════════════════

    % ── Toplayıcı blokları (üst sıra) ─────────────────────────────
    \node[blok, fill=blokmavi!50] (blk1a) at (0, 0) {4 bit\\Toplayıcı};
    \node[blok, fill=blokmavi!50] (blk1b) at (4.5, 0) {4 bit\\Toplayıcı};
    \node[blok] (blk0) at (11, 0) {4 bit\\Toplayıcı};

    \node[font=\scriptsize, above left=-1pt and 3pt of blk1a.north] {Bit [7:4]};
    \node[font=\scriptsize, above left=-1pt and 3pt of blk1b.north] {Bit [7:4]};
    \node[font=\scriptsize, above left=-1pt and 3pt of blk0.north] {Bit [3:0]};

    % ── a,b girişleri (üstten) ────────────────────────────────────
    % a[7:4], b[7:4] → her iki bloğa
    \draw[->] (0, 2.5) node[above, font=\footnotesize] {$a_{[7:4]},\; b_{[7:4]}$} -- (blk1a.north);
    \filldraw (0, 2) circle (2pt);
    \draw[->] (0, 2) -- (4.5, 2) -- (blk1b.north);

    % a[3:0], b[3:0] → Blok 0
    \draw[->] (11, 2.5) node[above, font=\footnotesize] {$a_{[3:0]},\; b_{[3:0]}$} -- (blk0.north);

    % ── c₀ girişi (sağdan) ────────────────────────────────────────
    \draw[->] (13.5, 0) node[right, font=\footnotesize] {$e_0$} -- (blk0.east);

    % ── c_in girişleri (sağdan, kutunun kenarında c_in etiketi) ───
    \draw[->] (2.2, 0) -- (blk1a.east);
    \node[font=\tiny, left] at ([xshift=-1mm]blk1a.east) {$e_{\text{gir}}$};
    \node[font=\scriptsize, right] at (2.2, 0) {$0$};
    \draw[->] (6.7, 0) -- (blk1b.east);
    \node[font=\tiny, left] at ([xshift=-1mm]blk1b.east) {$e_{\text{gir}}$};
    \node[font=\scriptsize, right] at (6.7, 0) {$1$};

    % ── Blok 0: S[3:0] çıkışı (aşağı) ────────────────────────────
    \draw[->] (blk0.south) -- ++(0, -1.5) node[below, font=\footnotesize] {$T_{[3:0]}$};

    % ── Blok 0: c₄ çıkışı (soldan, elde zinciri) ─────────────────
    \draw (blk0.west) -- ++(-0.5, 0) coordinate (c4);
    \node[above, font=\scriptsize] at (c4) {$e_4$};

    % ── Çoklayıcılar (her toplayıcının tam altında) ────────────────
    % MUX-S: blk1a'nın altında (sonuç seçimi)
    \node[draw, rectangle, fill=bloksari, thick, font=\scriptsize,
          minimum width=2.4cm, minimum height=1.3cm, align=center]
          (muxS) at (0, -3.2) {Çkl ($e$)};
    % MUX-C: blk1b'nin altında (toplam seçimi)
    \node[draw, rectangle, fill=bloksari, thick, font=\scriptsize,
          minimum width=2.4cm, minimum height=1.3cm, align=center]
          (muxC) at (4.5, -3.2) {Çkl ($T$)};

    % ── Blk1a → MUX-S (düz aşağı, giriş 0) ─────────────────────
    \draw[->] ([xshift=-3mm]blk1a.south) -- ([xshift=-3mm]muxS.north);
    \node[font=\scriptsize] at ([xshift=-3mm, yshift=3mm]blk1a.south) {$e$};
    \node[font=\tiny] at ([xshift=-3mm, yshift=-3mm]muxS.north) {0};

    % ── Blk1b → MUX-S (aşağı→sola→aşağı, y=-1.7'de yatay) ──────
    \draw[->] (4.2, -1) -- (4.2, -1.7) -- (0.3, -1.7) -- (0.3, -2.55);
    \node[font=\scriptsize] at ([xshift=-3mm, yshift=3mm]blk1b.south) {$e$};
    \node[font=\tiny] at ([xshift=3mm, yshift=-3mm]muxS.north) {1};

    % ── Blk1a → MUX-C (aşağı→sağa→aşağı→sağa→aşağı) ──────────
    %  Erken dirsek y=-1.2, e telini (y=-1.7) geçtikten sonra y=-2.0'da tekrar sağa
    \draw[->] (0.3, -1) -- (0.3, -1.2) -- (2.25, -1.2)
        -- (2.25, -2.0) -- (4.2, -2.0) -- (4.2, -2.55);
    \node[font=\scriptsize] at ([xshift=3mm, yshift=3mm]blk1a.south) {$T$};
    \node[font=\tiny] at ([xshift=-3mm, yshift=-3mm]muxC.north) {0};

    % ── Blk1b → MUX-C (düz aşağı, giriş 1) ─────────────────────
    \draw[->] ([xshift=3mm]blk1b.south) -- ([xshift=3mm]muxC.north);
    \node[font=\scriptsize] at ([xshift=3mm, yshift=3mm]blk1b.south) {$T$};
    \node[font=\tiny] at ([xshift=3mm, yshift=-3mm]muxC.north) {1};

    % ── c₄ → çoklayıcı seçim girişleri (alttan dolanıp sağ kenara) ─
    % c₄'ten aşağı, çoklayıcıların altına iniyor
    \draw (c4) -- ++(0, -4.7) coordinate (selbot);
    % sola yatay hat (çoklayıcıların sağ kenarı + 5mm hizası)
    \coordinate (muxCr) at ([xshift=5mm]muxC.east);
    \coordinate (muxSr) at ([xshift=5mm]muxS.east);
    \draw (selbot) -- (selbot -| muxSr);
    % muxC: noktadan yukarı, sola dirsek yaparak sağ kenara
    \filldraw (selbot -| muxCr) circle (2pt);
    \draw[->] (selbot -| muxCr) -- (muxCr) -- (muxC.east);
    \node[font=\tiny, above=1pt] at (muxCr) {Seç};
    % muxS: doğrudan yukarı, sola dirsek yaparak sağ kenara
    \draw[->] (selbot -| muxSr) -- (muxSr) -- (muxS.east);
    \node[font=\tiny, above=1pt] at (muxSr) {Seç};

    % ── Çoklayıcı çıkışları (alttan aşağı) ──────────────────────
    \draw[->] (muxS.south) -- ++(0, -1.3) node[below, font=\footnotesize] {$e_8$};
    \draw[->] (muxC.south) -- ++(0, -1.3) node[below, font=\footnotesize] {$T_{[7:4]}$};

\end{tikzpicture}
\caption[8 bitlik elde ile seçmeli toplayıcı]{8 bitlik elde ile seçmeli toplayıcı. İkinci blok, elde girişinin her iki olasılığını koşut hesaplar; $e_4$ hazır olduğunda çoklayıcı doğru sonucu seçer.}
\label{fig:carry-select-8bit}
\end{figure}

\begin{bilginot}[Güncel İşlemcilerde ve Yapay Zeka Hızlandırıcılarında Hızlı Toplayıcılar]
Bu bölümde ele aldığımız önceden elde üreten ve elde ile seçmeli toplayıcıların ötesinde, daha da hızlı toplayıcı yapıları vardır. \emph{Koşut önek toplayıcıları}\index{koşut önek toplayıcı}\index{parallel prefix adder|see{koşut önek toplayıcı}} (\emph{parallel prefix adder}), elde hesaplamasını bir ağaç yapısında koşut olarak gerçekleştirir ve $O(\log_2 n)$ gecikme sağlar. Bunların en bilinenleri:
\begin{itemize}
  \item \emph{Kogge-Stone}\index{Kogge-Stone toplayıcı}: En düşük gecikmeyi sağlar; 32 bit için yalnızca 5 düzey yeterlidir. Ancak $O(n \log n)$ kablo gerektirir; bu da yonga üzerinde yoğun kablo trafiğine yol açar.
  \item \emph{Brent-Kung}\index{Brent-Kung toplayıcı}: Kogge-Stone'a göre daha az kablo kullanır ($O(n)$) ama gecikme biraz artar ($2\log_2 n - 1$ düzey). Alan ve hız arasında daha dengeli bir ödünleşme sunar.
  \item \emph{Han-Carlson}\index{Han-Carlson toplayıcı}: Kogge-Stone ve Brent-Kung arasında bir uzlaşı sunar; hem kablo yoğunluğunu düşük tutar hem de gecikmeyi kabul edilebilir düzeyde bırakır. Uygulamada sıkça tercih edilen tasarımdır.
\end{itemize}
Bu toplayıcılar günümüzün yüksek başarımlı işlemcilerinde, grafik işlemcilerinde (GPU) ve yapay zeka hızlandırıcılarında kullanılan güncel tasarımlardır. Bir GPU'daki on binlerce çarpma-toplama (FMA) biriminin her birinde, son toplama aşaması koşut önek toplayıcıyla gerçekleştirilir. Çarpıcılarda ise değiştirilmiş Booth kodlaması (radix-4) ile kısmi çarpım sayısı yarıya indirilir ve Wallace ya da Dadda ağaçlarıyla sıkıştırılır. Düşük hassasiyetli yapay zeka iş yüklerinde (BF16, FP8, INT8) toplayıcılar daha kısa bit genişliğinde çalıştığından aynı yonga alanına çok daha fazla birim sığdırılabilir; NVIDIA H100 GPU'nun 16.896 FMA birimi barındırabilmesinin temel nedeni budur.
\end{bilginot}

\subsubsection*{Gecikme Karşılaştırması}

Gecikme çözümlemesi için bir kapı gecikmesini $t_g$ ile gösterelim. $n$ bitlik bir dalga eldeli (ripple carry) toplayıcıda en kötü durum gecikmesi elde zincirinden gelir:
\[
T_{\text{dalga}} = 2n \cdot t_g
\]
Çünkü her tam toplayıcı eldeyi 2 kapı gecikmesinde ($2t_g$) üretir ve bu elde bir sonraki basamağa aktarılır.

Elde ile seçmeli toplayıcıda ise $k$ bitlik bloklar kullanıldığında ilk blok $2k \cdot t_g$'de tamamlanır, ardından yalnızca bir çoklayıcı gecikmesi ($t_{mux}$) eklenir:
\[
T_{\text{seçmeli}} = 2k \cdot t_g + \left(\frac{n}{k} - 1\right) \cdot t_{mux}
\]
$t_{mux} \approx t_g$ kabul edilirse 8 bitlik toplayıcı için:
\begin{itemize}
  \item Dalga eldeli: $T = 2 \times 8 \times t_g = 16\,t_g$
  \item Elde ile seçmeli (4'erli blok): $T = 2 \times 4 \times t_g + 1 \times t_g = 9\,t_g$
\end{itemize}
Gecikme neredeyse yarıya indi; ancak bunun bir bedeli vardır: ikinci blok iki kez kurulduğu için donanım alanı yaklaşık 1{,}5 katına çıkar. Bu, tasarım ilkelerinden \emph{ödünleşme}nin güzel bir örneğidir: alandan vererek hız kazanılır. Bu toplayıcı türlerinin gecikme, alan ve karmaşıklık açısından karşılaştırması Tablo~\ref{tab:toplayici-karsilastirma}'de verilmiştir.

\begin{table}[htbp]
\centering
\caption[Toplayıcı türlerinin karşılaştırması]{Toplayıcı türlerinin gecikme, alan ve karmaşıklık açısından karşılaştırması ($n$ bit genişlik, $k$ bit blok boyutu).}
\label{tab:toplayici-karsilastirma}
\small
\begin{tabularx}{\textwidth}{l l X c c}
\toprule
\textbf{Türkçe Ad} & \textbf{İngilizce} & \textbf{Strateji} & \textbf{Gecikme} & \textbf{Alan} \\
\midrule
Dalga eldeli & Ripple Carry & Elde bir sonraki basamağa dalgalanır & $O(n)$ & $O(n)$ \\
Önceden elde üreten & Carry Lookahead & Elde mantıksal denklemlerle önceden üretilir & $O(\log n)$ & $O(n)$ \\
Elde ile seçmeli & Carry Select & Her iki elde olasılığı koşut hesaplanır, sonuç seçilir & $O(\sqrt{n})$ & $O(1{,}5n)$ \\
Kogge-Stone & Kogge-Stone & Koşut önek ağacı, en düşük gecikme & $O(\log n)$ & $O(n \log n)$ \\
Brent-Kung & Brent-Kung & Koşut önek ağacı, az kablo & $O(\log n)$ & $O(n)$ \\
\bottomrule
\end{tabularx}
\end{table}

\begin{ornek}[5.3]
\textbf{Soru:} 8 bitlik elde ile seçmeli toplayıcının gecikme ve alan karşılaştırmasını yapın. Dalga eldeli toplayıcıya göre ne kadar hızlanma sağlanır? Alan artışı ne kadardır?

\textbf{Çözüm:}

\emph{Gecikme:}
\begin{itemize}
  \item Dalga eldeli (8 bit): $T = 2 \times 8 \times t_g = 16\,t_g$
  \item Elde ile seçmeli (4+4): İlk blok $= 8\,t_g$, çoklayıcı $= 1\,t_g$ $\Rightarrow$ $T = 9\,t_g$
  \item Hızlanma: $16/9 \approx 1{,}78\times$
\end{itemize}

\emph{Alan (tam toplayıcı sayısı):}
\begin{itemize}
  \item Dalga eldeli: 8 tam toplayıcı
  \item Elde ile seçmeli: $4 + 4 + 4 = 12$ tam toplayıcı $+$ 2 çoklayıcı
  \item Alan artışı: $\approx$1{,}5$\times$ (çoklayıcılar dahil $\approx$1{,}6$\times$)
\end{itemize}
\end{ornek}

% =============================================
\section{Çarpma}
\label{sec:carpma}

AMB'miz toplamayı ve çıkarmayı öğrendi; şimdi sıra çarpmaya\index{çarpma}\index{multiplication|see{çarpma}} geldi. Bir çarpma işleminde üç temel bileşen vardır:
%
\begin{itemize}
\item \textbf{Çarpılan} (\emph{multiplicand}): çarpma işleminde tekrarlanan sayı,
\item \textbf{Çarpan} (\emph{multiplier}): çarpılanın kaç kez ve hangi konumda toplanacağını belirleyen sayı,
\item \textbf{Çarpım} (\emph{product}): işlemin sonucu.
\end{itemize}

Çarpmayı anlamanın en kolay yolu, ilkokulda öğrendiğimiz elde çarpma yöntemini ikili tabana uygulamaktır. Onluk tabanda çarpanın her bir basamağını çarpılanla çarpar ve sonucu uygun konumda alt alta yazarız; sonra tüm kısmi çarpımları toplarız. İkili tabanda bu işlem çok daha basittir: çarpanın her bir biti yalnızca 0 veya 1 olduğundan, kısmi çarpım ya sıfırdır (bit = 0) ya da çarpılanın kendisidir (bit = 1). Aşağıdaki örnekte $7 \times 5$ çarpımı ikili tabanda gösterilmiştir:

\begin{center}
\begin{tikzpicture}[scale=1.0]
  \node[font=\ttfamily, anchor=east, vurgumavi] at (0, 0) {0111};
  \node[anchor=east, vurgumavi] at (-2.0, 0) {Çarpılan:};
  \node[font=\ttfamily, anchor=east, vurgukirmizi] at (0, -0.6) {$\times$\; 0101};
  \node[anchor=east, vurgukirmizi] at (-2.0, -0.6) {Çarpan:};
  \draw[thick] (-1.0, -0.9) -- (0.2, -0.9);
  \node[font=\ttfamily, anchor=east, vurgumavi] at (0, -1.3) {0111};
  \node[font=\small, anchor=west] at (0.3, -1.3) {$\leftarrow$ \textcolor{vurgukirmizi}{bit 0 = 1}, çarpılanı yaz};
  \node[font=\ttfamily, anchor=east, gray] at (0, -1.8) {0000\phantom{0}};
  \node[font=\small, anchor=west] at (0.3, -1.8) {$\leftarrow$ \textcolor{vurgukirmizi}{bit 1 = 0}, sıfır yaz};
  \node[font=\ttfamily, anchor=east, vurgumavi] at (0, -2.3) {0111\phantom{00}};
  \node[font=\small, anchor=west] at (0.3, -2.3) {$\leftarrow$ \textcolor{vurgukirmizi}{bit 2 = 1}, çarpılanı yaz};
  \node[font=\ttfamily, anchor=east, gray] at (0, -2.8) {0000\phantom{000}};
  \node[font=\small, anchor=west] at (0.3, -2.8) {$\leftarrow$ \textcolor{vurgukirmizi}{bit 3 = 0}, sıfır yaz};
  \draw[thick] (-1.8, -3.1) -- (0.2, -3.1);
  \node[font=\ttfamily\bfseries, anchor=east, vurguyesil] at (0, -3.5) {00100011};
  \node[anchor=east, vurguyesil] at (-2.0, -3.5) {Çarpım:};
  \node[font=\small, anchor=west, vurguyesil] at (0.3, -3.5) {$= 35$};
\end{tikzpicture}
\end{center}

Görüldüğü gibi ikili çarpma, özünde \emph{kaydır ve topla} işleminden ibarettir: çarpanın her biti için çarpılanı uygun miktarda sola kaydırarak kısmi çarpımları oluşturur ve bunları toplarız. Bu yalın ilke, donanımda çarpma biriminin temelini oluşturur.

\subsubsection*{Çarpma Donanımına Neden İhtiyaç Var?}

Bir işlemci toplama yapabiliyorsa, çarpmayı yazılımla da gerçekleştirebilir: çarpma, özünde arka arkaya toplama işlemine indirgenebilir. Nitekim RISC-V'in temel buyruk kümesi RV32I'da çarpma buyruğu yoktur; çarpma desteği isteğe bağlı \textbf{M uzantısı}\index{M uzantısı} ile eklenir. Peki bu durumda özelleşmiş çarpma donanımına neden ihtiyaç duyulur? Aşağıdaki iki RISC-V programını karşılaştıralım:

\begin{table}[htbp]
\centering
\caption[Yazılım ve donanım çarpma karşılaştırması]{Aynı çarpma işlemi ($7 \times 5$): yalnızca toplama ile (yazılım) ve \texttt{mul} buyruğu ile (donanım).}
\label{tab:carpma-yazilim-donanim}
\small
\begin{tabular}{p{6.8cm}p{6.8cm}}
\toprule
\textbf{Yalnızca toplama ile (M uzantısı yok)} & \textbf{\texttt{mul} buyruğu ile (M uzantısı var)} \\
\midrule
\begin{minipage}[t]{6.5cm}
\begin{lstlisting}[style=riscv, xleftmargin=0pt]
# 7 x 5 = 7+7+7+7+7
    li   t0, 0       # sonuc = 0
    li   t1, 7       # carpilan
    li   t2, 5       # sayac (carpan)
dongu:
    beqz t2, son     # sayac=0? bitir
    add  t0, t0, t1  # sonuc += 7
    addi t2, t2, -1  # sayac--
    j    dongu       # tekrarla
son:
\end{lstlisting}
\end{minipage}
&
\begin{minipage}[t]{6.5cm}
\begin{lstlisting}[style=riscv, xleftmargin=0pt]
    li   t1, 7       # carpilan
    li   t2, 5       # carpan
    mul  t0, t1, t2  # sonuc = 7x5
\end{lstlisting}
\end{minipage}
\\
\midrule
Durağan buyruk sayısı: 7 & Durağan buyruk sayısı: 3 \\
Devingen buyruk sayısı: 24 & Devingen buyruk sayısı: 3 \\
Çevrim sayısı (BBÇ${}\approx 1$): $\sim$24 & Çevrim sayısı: 3--5 \\
\bottomrule
\end{tabular}
\end{table}

Soldaki programda döngü 5 kez yinelenir (her yinelemede 4 buyruk) ve başlangıç buyrukları ile birlikte toplam 24 buyruk yürütülür. Çarpan büyüdükçe döngü sayısı doğrusal olarak artar; örneğin $1000 \times 1000$ için 4004, 32~bitlik en kötü durumda ise milyarlarca buyruk gerekir. Sağdaki programda ise aynı işlem \emph{tek bir buyrukla} tamamlanır; \texttt{mul} buyruğu birkaç saat çevriminde sonucu üretir. Bölüm~\ref{chap:basarim}'deki yürütme zamanı formülünü hatırlayalım:
%
\[
T_{\text{yürütme}} = N_{\text{buyruk}} \times \text{BBÇ} \times T_{\text{çevrim}}
\]
%
Buyruk sayısının $24$'ten $3$'e düşmesi bile \textbf{8 kat} iyileşme demektir; büyük sayılarda bu fark astronomik boyutlara ulaşır. Bu dramatik fark, çarpma donanımının neden hemen her çağdaş işlemcide bulunduğunu açıklar: \emph{buyruk kümesine eklenen tek bir buyruk, sorumluluğu yazılımdan donanıma aktararak hem programcıyı rahatlatır hem de sistem başarımını büyük ölçüde artırır.}

Aynı durum bölme için de geçerlidir ve Bölüm~\ref{sec:bolme}'de ele alınacaktır.

Çarpmada temel bir kural vardır: \emph{$m$ bitlik bir sayı ile $n$ bitlik bir sayının çarpımı en fazla $m+n$ bit uzunluğundadır}. RISC-V'te (RV32I) veriler 32 bit olduğuna göre çarpım sonucu 64 bite kadar çıkabilir. Bu bölümde önce artı sayıların çarpımını göreceğiz; eksi sayılarla çarpma ise ikiye tümleyen dönüşümü veya Booth algoritması ile yapılır.

Çarpma yöntemlerini verimsizden verimliye doğru adım adım geliştireceğiz.

\subsection{Kaydır-Topla Çarpma (İşaretsiz Sayılar)}
\label{sec:kaydir-topla-carpma}

İlk yöntemimiz ``kaydır ve topla'' (\emph{shift-and-add}) ilkesine dayanır. Çarpanın her bir bitini sırayla sınayarak çarpılanı uygun konumda çarpıma ekleriz; tıpkı elde kâğıt üzerinde yapılan çarpmada olduğu gibi. Bu temel fikir üzerinde donanımı adım adım iyileştiren üç yol geliştireceğiz.

\subsubsection{1. Yol}

İlk yol, onluk tabanda yapılan normal çarpma işleminin doğrudan gerçekleştirilmesidir. Bu yolda 64 bitlik çarpılan ve çarpım yazmaçları ile 32 bitlik çarpan yazmacı kullanılır. Her adımda çarpanın en sağ bitine bakılır; bit 1 ise çarpılan çarpıma eklenir, 0 ise toplama atlanır. Ardından çarpılan 1 bit sola, çarpan 1 bit sağa kaydırılır. Çarpılanın sola kaydırılması, her adımda çarpılanın bir üst basamağa hizalanmasını sağlar; tıpkı elde çarpma yönteminde kısmi çarpımları bir basamak sola kaydırarak yazdığımız gibi.

\begin{figure}[htbp]
\centering
\begin{tikzpicture}[node distance=0.8cm, scale=0.72, transform shape]
    % Start
    \node[akisbaslangic] (start) {Ba{\c{s}}la};
    % Initialize
    \node[akisislem, below=of start, minimum width=3.2cm] (init)
        {{\c{C}}arp{\i}m $\leftarrow$ 0, Say $\leftarrow$ 0};
    % Decision: LSB of multiplier
    \node[akiskarar, below=of init, minimum width=2.5cm] (check)
        {{\c{C}}arpan{\i}n\\en sa{\u{g}} biti = 1?};
    % Add
    \node[akisislem, right=2cm of check, minimum width=3.2cm] (add)
        {{\c{C}}arp{\i}m $\leftarrow$ {\c{C}}arp{\i}m + {\c{C}}arp{\i}lan};
    % Shift multiplicand left
    \node[akisislem, below=1.3cm of check, minimum width=3.2cm] (shiftl)
        {{\c{C}}arp{\i}lan{\i} 1 bit sola kayd{\i}r};
    % Shift multiplier right
    \node[akisislem, below=of shiftl, minimum width=3.2cm] (shiftr)
        {{\c{C}}arpan{\i} 1 bit sa{\u{g}}a kayd{\i}r};
    % Increment counter
    \node[akisislem, below=of shiftr, minimum width=3.2cm] (inc)
        {Say $\leftarrow$ Say + 1};
    % Check counter
    \node[akiskarar, below=of inc, minimum width=2cm] (countchk)
        {Say = 32?};
    % End
    \node[akisbaslangic, below=1cm of countchk] (stop) {Dur};
    % Arrows
    \draw[akisok] (start) -- (init);
    \draw[akisok] (init) -- (check);
    \draw[akisok] (check) -- node[above, font=\scriptsize] {Evet} (add);
    \draw[akisok] (check) -- node[left, font=\scriptsize] {Hay{\i}r} (shiftl);
    \draw[akisok] (add) |- (shiftl);
    \draw[akisok] (shiftl) -- (shiftr);
    \draw[akisok] (shiftr) -- (inc);
    \draw[akisok] (inc) -- (countchk);
    \draw[akisok] (countchk) -- node[left, font=\scriptsize] {Evet} (stop);
    \coordinate (dongu-sag) at ($(add.east)+(0.8,0)$);
    \draw[akisdonus] (countchk.east)
        -- node[above, font=\scriptsize, pos=0.2] {Hay{\i}r}
        (dongu-sag |- countchk.east)
        |- ($(check.north)+(0,0.5)$)
        -- (check.north);
\end{tikzpicture}
\caption[Kayd{\i}r-Topla {\c{C}}arpma (1.\ Yol) ak{\i}{\c{s}} {\c{s}}emas{\i}]{Kayd{\i}r-Topla {\c{C}}arpma (1.\ Yol) ak{\i}{\c{s}} {\c{s}}emas{\i}.}
\label{fig:carpma-1yol-algo}
\end{figure}

1. yolda kaydırma ve toplama işlemleri kullanılmıştır; algoritmanın akış şeması Şekil~\ref{fig:carpma-1yol-algo}'da, buna karşılık gelen donanım tasarımı ise Şekil~\ref{fig:carpma-1surum-donanim}'de gösterilmiştir. Donanım kaydırıcı ve toplayıcıların dışında, 2 tane 64 bitlik ve bir tane de 32 bitlik yazmaçtan oluşur. Toplayıcının kullanılacağı AMB de 64 bit ile işlem yapabilmek için 64 bitlik bir AMB olmalıdır. 64 bitlik AMB, 32 bitlik olandan çok daha karmaşıktır ve yavaştır.

Her yinelemede yapılan işlemler şunlardır: (1) çarpanın en sağ bitini sına, (2) bit 1 ise çarpılanı çarpıma ekle, (3) çarpılanı sola, çarpanı sağa kaydır. Burada iki kaydırma birbirinden bağımsızdır (biri çarpılan yazmacını, diğeri çarpan yazmacını etkiler); bu yüzden donanımda \emph{koşut} (aynı anda) yapılabilirler. Dolayısıyla her yineleme aslında üç aşamadan oluşur: test, toplama ve kaydırma. Denetim birimi içindeki bir sayaç yineleme sayısını takip eder: 32 bitlik bir RISC-V çarpmasında bu döngü 32 kez tekrarlanır ve sayaç 32'ye ulaştığında işlem sonlandırılır. Dolayısıyla çarpma işlemi toplam \textbf{32 saat vuruşu} sürer (toplamayı, testi ve kaydırmayı tek bir vuruşta yapabildiğimiz varsayımıyla).

\begin{figure}[H]
\centering
\begin{tikzpicture}[scale=0.72, transform shape]

    % ========== DÜĞÜMLER ==========

    % ── Çarpılan Yazmacı (64 bit) ──────────────────────────────
    \node[draw, rectangle, fill=blokmavi, thick, font=\footnotesize, align=center,
          minimum width=4.8cm, minimum height=1.3cm] (carpilan) at (0, 5)
          {{\c{C}}arp{\i}lan Yazmac{\i} {\scriptsize(64 bit)}};
    % Sola Kaydır yazısı kutunun üstünde (siyah — yazmacın özelliği)
    \node[font=\scriptsize, above=2pt of carpilan]
          {$\longleftarrow$ Sola Kayd{\i}r};

    % ── Çarpan Yazmacı (32 bit, en sağ bit ayrı bölmede) ────────
    \node[draw, rectangle, fill=blokmavi, thick, font=\footnotesize,
          minimum width=5.4cm, minimum height=1.3cm] (carpan) at (8, 5) {};
    \node[font=\footnotesize, anchor=west] at ([xshift=2mm]carpan.west)
          {{\c{C}}arpan Yazmac{\i} {\scriptsize(32 bit)}};
    % Sağa Kaydır yazısı kutunun üstünde (siyah — yazmacın özelliği)
    \node[font=\scriptsize, above=2pt of carpan]
          {Sa{\u{g}}a Kayd{\i}r $\longrightarrow$};
    % En sağ bit bölmesi (kutunun sağ ucunda dikey çizgiyle ayrılmış)
    \draw ([xshift=-14mm]carpan.north east) -- ([xshift=-14mm]carpan.south east);
    \node[font=\tiny, align=center] at ([xshift=-7mm]carpan.east)
          {en sa{\u{g}}\\bit};

    % ── 64 bit AMB (trapezoid) ─────────────────────────────────
    \node[draw, trapezium, trapezium angle=70, fill=bloksari, thick,
          shape border rotate=180, font=\footnotesize, align=center,
          minimum width=3cm, minimum height=1.6cm] (amb) at (0, 2.2)
          {64 bit AMB};

    % ── Çarpım Yazmacı (64 bit) ───────────────────────────────
    \node[draw, rectangle, fill=blokmavi, thick, font=\footnotesize, align=center,
          minimum width=4.8cm, minimum height=1.3cm] (carpim) at (0, -0.5)
          {{\c{C}}arp{\i}m Yazmac{\i} {\scriptsize(64 bit)}};

    % ── Denetim birimi ─────────────────────────────────────────
    \node[draw, rectangle, fill=blokkirmizi-acik, thick, font=\footnotesize, align=center,
          minimum width=2.8cm, minimum height=0.8cm] (denetim) at (7.5, 2.2)
          {Denetim};

    % ========== VERİ YOLLARI (kalın düz çizgi) ==========

    % Çarpılan → AMB
    \draw[okstil] (carpilan.south) -- (amb.north);

    % AMB → Çarpım (yaz)
    \draw[okstil] (amb.south) -- (carpim.north)
          node[midway, right=2pt, font=\scriptsize] {Yaz};

    % Çarpım → AMB geri besleme (sol taraftan)
    \draw[okstil] (carpim.west) -- ++(-1.8, 0) coordinate (fb1)
          -- (fb1 |- amb.west) -- (amb.west);

    % Çarpan en sağ bit bölmesinden → Denetim (aşağı, dirsekli)
    \draw[okstil] ([xshift=-7mm]carpan.south east) -- ++(0, -0.5)
          -| (denetim.north);

    % ========== DENETİM SİNYALLERİ (kırmızı kesikli) ==========

    % Denetim → Çarpılan (sola kaydır): denetim solundan, yukarı, çarpılana
    \draw[denetimstil] ([yshift=1mm]denetim.west) -- (4, 2.3)
          -- (4, 5) -- (carpilan.east);
    \node[font=\scriptsize, sinyalkirmizi, fill=white, inner sep=1pt]
          at (4, 3.6) {Kayd{\i}r};

    % Denetim → Çarpan (sağa kaydır): denetim sağından, yukarı, çarpana
    \draw[denetimstil] (denetim.east) -- (11.2, 2.2)
          -- (11.2, 5) -- (carpan.east);
    \node[font=\scriptsize, sinyalkirmizi, fill=white, inner sep=1pt]
          at (11.2, 3.6) {Kayd{\i}r};

    % Denetim → AMB (topla): tam yatay
    \draw[denetimstil] (denetim.west) -- (amb.east);
    \node[font=\scriptsize, sinyalkirmizi, fill=white, inner sep=1pt]
          at (3.8, 1.9) {Topla};

    % Denetim → Çarpım (yaz): aşağı, sola, çarpıma
    \draw[denetimstil] (denetim.south) -- ++(0, -1.5)
          -| ([xshift=8mm]carpim.east) -- (carpim.east);
    \node[font=\scriptsize, sinyalkirmizi, fill=white, inner sep=1pt]
          at (5.5, -0.2) {Yaz};

\end{tikzpicture}
\caption{Kayd{\i}r-Topla {\c{C}}arpma (1.\ Yol) donan{\i}m{\i}}
\label{fig:carpma-1surum-donanim}
\end{figure}

1.\ yolun asıl sorunu toplama adımının 64 bitlik AMB ile yapılmak zorunda olmasıdır: 64 bitlik bir toplama, 32 bitlik bir toplamaya göre belirgin şekilde daha yavaştır, daha fazla alan kaplar ve daha çok güç tüketir. Yöntemin iyileştirilmesi gerekir.

\begin{ornek}[5.4]
\textbf{Soru:} $7 \times 5 = ?$ \quad $(0111 \times 0101 = ?)$

\textbf{Çözüm:} Adımları tek tek izleyelim (Tablo~\ref{tab:carpma1-ornek}):

\begin{itemize}[nosep]
\item \textbf{Başlangıç:} Çarpılan $= \textcolor{vurgumavi}{0111}$, çarpan $= \textcolor{vurgukirmizi}{0101}$, çarpım $= \textcolor{gray}{0}$.
\item \textbf{Adım 1:} Çarpanın en sağ biti $= \textcolor{vurgukirmizi}{1}$ $\Rightarrow$ çarpılanı çarpıma ekle: $\textcolor{gray}{0} + \textcolor{vurgumavi}{0111} = \textcolor{vurguyesil}{0111}$. Çarpılanı sola, çarpanı sağa kaydır.
\item \textbf{Adım 2:} Çarpanın en sağ biti $= \textcolor{vurgukirmizi}{0}$ $\Rightarrow$ toplama yok. Çarpılanı sola, çarpanı sağa kaydır.
\item \textbf{Adım 3:} Çarpanın en sağ biti $= \textcolor{vurgukirmizi}{1}$ $\Rightarrow$ çarpılanı çarpıma ekle: $\textcolor{gray}{0000}\textcolor{vurguyesil}{0111} + \textcolor{gray}{000}\textcolor{vurgumavi}{0111}\textcolor{gray}{00} = \textcolor{gray}{00}\textcolor{vurguyesil}{100011}$. Kaydır.
\item \textbf{Adım 4:} Çarpanın en sağ biti $= \textcolor{vurgukirmizi}{0}$ $\Rightarrow$ toplama yok. Kaydır. Çarpım değişmez.
\end{itemize}

\begin{table}[H]
\centering
\caption{Örnek 5.4 -- 1.\ yol çarpma algoritması}
\label{tab:carpma1-ornek}
\small
\begin{tabular}{c l l l}
\toprule
\textbf{Adım} & \textbf{\textcolor{vurgumavi}{Çarpılan}} & \textbf{\textcolor{vurgukirmizi}{Çarpan}} & \textbf{\textcolor{vurguyesil}{Çarpım}} \\
\midrule
Başlangıç & \textcolor{gray}{0000}\textcolor{vurgumavi}{0111} & \textcolor{vurgukirmizi}{0101} & \textcolor{gray}{00000000} \\
1 (Topla + Kaydır) & \textcolor{gray}{000}\textcolor{vurgumavi}{0111}\textcolor{gray}{0} & \textcolor{gray}{0}\textcolor{vurgukirmizi}{010} & \textcolor{gray}{0000}\textcolor{vurguyesil}{0111} \\
2 (Kaydır) & \textcolor{gray}{00}\textcolor{vurgumavi}{0111}\textcolor{gray}{00} & \textcolor{gray}{00}\textcolor{vurgukirmizi}{01} & \textcolor{gray}{0000}\textcolor{vurguyesil}{0111} \\
3 (Topla + Kaydır) & \textcolor{gray}{0}\textcolor{vurgumavi}{0111}\textcolor{gray}{000} & \textcolor{gray}{000}\textcolor{vurgukirmizi}{0} & \textcolor{gray}{00}\textcolor{vurguyesil}{100011} \\
4 (Kaydır) & \textcolor{vurgumavi}{0111}\textcolor{gray}{0000} & \textcolor{gray}{0000} & \textcolor{gray}{00}\textcolor{vurguyesil}{100011} \\
\bottomrule
\end{tabular}
\end{table}

Çarpım $= (00100011)_2 = 35$
\end{ornek}

\subsubsection{2. Yol}

1.\ yoldaki donanıma dikkatli bakıldığında bir verimsizlik göze çarpar: çarpılan yazmacı 64 bit genişliğinde olmasına karşın, herhangi bir anda yalnızca 32 biti anlamlı veri taşır; kalan bitler sıfırdır. Çarpılanın sola kaydırılması, AMB'nin de 64 bit genişliğinde olmasını zorunlu kılmaktadır; oysa her toplama işleminde yalnızca 32 bitlik bir değer toplanır. Buradaki kilit gözlem şudur: çarpılanı sola kaydırmak yerine, çarpımı sağa kaydırmak aynı hizalama etkisini yaratır. Çarpım sağa kaydırıldığında, çarpılanın her zaman çarpım yazmacının üst yarısıyla (sol 32 bit) toplanması yeterli olur.

\begin{figure}[htbp]
\centering
\begin{tikzpicture}[scale=0.72, transform shape]

    % ========== D{\"U}{\u{G}}{\"U}MLER ==========

    % ---- {\c{C}}arp{\i}lan Yazma{\c{c}}{\i} (32 bit, sol {\"u}st, sabit, AMB ile hizal{\i}) ----
    % (Sabit) — kayd{\i}rma yok, bu nedenle {\"u}ste kayd{\i}rma bilgisi yaz{\i}lmaz
    \node[draw, rectangle, fill=blokmavi, thick, font=\footnotesize, align=center,
          minimum width=3.4cm, minimum height=1.3cm] (carpilan) at (-1.1, 5)
          {{\c{C}}arp{\i}lan Yazmac{\i} {\scriptsize(32 bit)}\\{\scriptsize(Sabit)}};

    % ---- {\c{C}}arpan Yazma{\c{c}}{\i} (32 bit, sa{\u{g}} {\"u}st, sa{\u{g}}a kayd{\i}r{\i}l{\i}r) ----
    % Kayd{\i}rma bilgisi kutunun {\"u}st{\"u}nde, siyah
    % En sa{\u{g}} bit b{\"o}lmesi (5.15'teki gibi dikey {\c{c}}izgiyle ayr{\i}lm{\i}{\c{s}})
    \node[draw, rectangle, fill=blokmavi, thick, font=\footnotesize,
          minimum width=5.4cm, minimum height=1.3cm] (carpan) at (7, 5) {};
    \node[font=\footnotesize, anchor=west] at ([xshift=2mm]carpan.west)
          {{\c{C}}arpan Yazmac{\i} {\scriptsize(32 bit)}};
    \node[font=\scriptsize, above=2pt of carpan]
          {Sa{\u{g}}a Kayd{\i}r $\longrightarrow$};
    % En sa{\u{g}} bit b{\"o}lmesi (kutunun sa{\u{g}} ucunda dikey {\c{c}}izgiyle ayr{\i}lm{\i}{\c{s}})
    \draw ([xshift=-14mm]carpan.north east) -- ([xshift=-14mm]carpan.south east);
    \node[font=\tiny, align=center] at ([xshift=-7mm]carpan.east)
          {en sa{\u{g}}\\bit};

    % ---- AMB (32 bit Toplay{\i}c{\i}, sol yar{\i}yla hizal{\i}, 5.15'in yar{\i}s{\i} boyutunda) ----
    \node[draw, trapezium, trapezium angle=70, fill=bloksari, thick,
          shape border rotate=180, font=\footnotesize, align=center,
          minimum width=1.5cm, minimum height=1.0cm] (amb) at (-1.1, 2.2)
          {32 bit AMB};

    % ---- {\c{C}}arp{\i}m Yazma{\c{c}}{\i} (64 bit, altta ortada, sa{\u{g}}a kayd{\i}r{\i}l{\i}r) ----
    % Kayd{\i}rma bilgisi kutunun {\"u}st{\"u}nde, siyah
    \node[draw, rectangle, fill=blokmavi, thick, font=\footnotesize, align=center,
          minimum width=4.4cm, minimum height=1.3cm] (carpim) at (0, -0.5)
          {{\c{C}}arp{\i}m Yazmac{\i} {\scriptsize(64 bit)}};
    % Ortaya kesikli {\c{c}}izgi (sol/sa{\u{g}} yar{\i} ayr{\i}m{\i})
    \draw[dashed, thick, gray] (carpim.north) -- (carpim.south);
    \node[font=\scriptsize, above=2pt of carpim, xshift=8mm]
          {Sa{\u{g}}a Kayd{\i}r $\longrightarrow$};

    % ---- Denetim birimi (sa{\u{g}} tarafta) ----
    \node[draw, rectangle, fill=blokkirmizi-acik, thick, font=\footnotesize, align=center,
          minimum width=2.8cm, minimum height=0.8cm] (denetim) at (7, 2.2)
          {Denetim};

    % ========== VER{\.I} YOLLARI (kal{\i}n d{\"u}z {\c{c}}izgi) ==========

    % {\c{C}}arp{\i}lan $\rightarrow$ AMB ({\"u}stten giri{\c{s}})
    \draw[okstil] (carpilan.south) -- (amb.north);

    % AMB $\rightarrow$ {\c{C}}arp{\i}m sol yar{\i}s{\i}n{\i}n ortas{\i}na ({\c{c}}{\i}k{\i}{\c{s}} a{\c{s}}a{\u{g}}{\i})
    \draw[okstil] (amb.south) -- ([xshift=-1.1cm]carpim.north)
          node[midway, right=2pt, font=\scriptsize] {Yaz [63:32]};

    % {\c{C}}arp{\i}m {\"u}st yar{\i} $\rightarrow$ AMB geri besleme (sol taraftan dola{\c{s}}arak)
    \draw[okstil] (carpim.west) -- ++(-1.5, 0) coordinate (fb2)
          -- (fb2 |- amb.west) -- (amb.west)
          node[pos=0.15, left=3pt, font=\scriptsize] {[63:32]};

    % {\c{C}}arpan en sa{\u{g}} bit b{\"o}lmesinden $\rightarrow$ Denetim (a{\c{s}}a{\u{g}}{\i}, dirsekli)
    \draw[okstil] ([xshift=-7mm]carpan.south east) -- ++(0, -0.5)
          -| (denetim.north);

    % ========== DENET{\.I}M S{\.I}NYALLER{\.I} (k{\i}rm{\i}z{\i} kesikli) ==========

    % Denetim $\rightarrow$ {\c{C}}arpan (sa{\u{g}}a kayd{\i}r sinyali)
    \draw[denetimstil] (denetim.east) -- (10.2, 2.2)
          -- (10.2, 5) -- (carpan.east);
    \node[font=\scriptsize, sinyalkirmizi, fill=white, inner sep=1pt]
          at (10.2, 3.6) {Kayd{\i}r};

    % Denetim $\rightarrow$ {\c{C}}arp{\i}m (kayd{\i}r / yaz sinyali)
    \draw[denetimstil] (denetim.south) -- ++(0, -0.8)
          -| ([xshift=10mm]carpim.east)
          -- (carpim.east);
    \node[font=\scriptsize, sinyalkirmizi, fill=white, inner sep=1pt, right=1pt]
          at (5.5, 0.7) {Kayd{\i}r/Yaz};

    % Denetim $\rightarrow$ AMB (topla sinyali)
    \draw[denetimstil] (denetim.west) -- (amb.east);
    \node[font=\scriptsize, sinyalkirmizi, fill=white, inner sep=1pt]
          at ($(denetim.west)!0.5!(amb.east)$) {Topla};

\end{tikzpicture}
\caption{Kayd{\i}r-Topla {\c{C}}arpma (2.\ Yol) donan{\i}m{\i}}
\label{fig:carpma-2surum-donanim}
\end{figure}

\begin{figure}[htbp]
\centering
\begin{tikzpicture}[node distance=0.8cm, scale=0.72, transform shape]
    % Start
    \node[akisbaslangic] (start) {Ba{\c{s}}la};
    % Initialize
    \node[akisislem, below=of start, minimum width=3.6cm] (init)
        {{\c{C}}arp{\i}m $\leftarrow$ 0, Say $\leftarrow$ 0};
    % Decision: LSB of multiplier
    \node[akiskarar, below=of init, minimum width=2.5cm] (check)
        {{\c{C}}arpan{\i}n\\en sa{\u{g}} biti = 1?};
    % Add
    \node[akisislem, right=2cm of check, minimum width=3.6cm] (add)
        {{\c{C}}arp{\i}m[63:32] $\leftarrow$\\{\c{C}}arp{\i}m[63:32] + {\c{C}}arp{\i}lan};
    % Shift product right
    \node[akisislem, below=1.3cm of check, minimum width=3.6cm] (shiftp)
        {{\c{C}}arp{\i}m{\i} 1 bit sa{\u{g}}a kayd{\i}r};
    % Shift multiplier right
    \node[akisislem, below=of shiftp, minimum width=3.6cm] (shiftr)
        {{\c{C}}arpan{\i} 1 bit sa{\u{g}}a kayd{\i}r};
    % Increment counter
    \node[akisislem, below=of shiftr, minimum width=3.6cm] (inc)
        {Say $\leftarrow$ Say + 1};
    % Check counter
    \node[akiskarar, below=of inc, minimum width=2cm] (countchk)
        {Say = 32?};
    % End
    \node[akisbaslangic, below=1cm of countchk] (stop) {Dur};
    % Arrows
    \draw[akisok] (start) -- (init);
    \draw[akisok] (init) -- (check);
    \draw[akisok] (check) -- node[above, font=\scriptsize] {Evet} (add);
    \draw[akisok] (check) -- node[left, font=\scriptsize] {Hay{\i}r} (shiftp);
    \draw[akisok] (add) |- (shiftp);
    \draw[akisok] (shiftp) -- (shiftr);
    \draw[akisok] (shiftr) -- (inc);
    \draw[akisok] (inc) -- (countchk);
    \draw[akisok] (countchk) -- node[left, font=\scriptsize] {Evet} (stop);
    \coordinate (dongu-sag) at ($(add.east)+(0.8,0)$);
    \draw[akisdonus] (countchk.east)
        -- node[above, font=\scriptsize, pos=0.2] {Hay{\i}r}
        (dongu-sag |- countchk.east)
        |- ($(check.north)+(0,0.5)$)
        -- (check.north);
\end{tikzpicture}
\caption[Kayd{\i}r-Topla {\c{C}}arpma (2.\ Yol) ak{\i}{\c{s}} {\c{s}}emas{\i}]{Kayd{\i}r-Topla {\c{C}}arpma (2.\ Yol) ak{\i}{\c{s}} {\c{s}}emas{\i}.}
\label{fig:carpma-2yol-algo}
\end{figure}

2.\ yol, bu gözlemden yararlanarak donanımı sadeleştirir (Şekil~\ref{fig:carpma-2surum-donanim}, akış şeması Şekil~\ref{fig:carpma-2yol-algo}): çarpılan yazmacı artık 32 bitte kalır, AMB 32 bite iner ve toplama sonucu çarpım yazmacının üst yarısına yazılır. Her adımda çarpım yazmacının tamamı sağa kaydırılarak kısmi çarpım doğru konuma yerleşir. Sonuç olarak, aynı işlevsellik yarı boyuttaki bir AMB ile sağlanmış olur; bu da hem alan hem de güç tüketimi açısından önemli bir kazanımdır. Her yinelemede yapılan işlem sayısı (1 test, 1 toplama, kaydırmalar) 1.\ yol ile benzerdir ve döngü yine 32 kez tekrarlanır; ancak toplama artık 32 bitlik AMB ile yapıldığından her adım daha hızlıdır ve donanım daha az alan kaplar.

\begin{ornek}[5.5]
\textbf{Soru:} $7 \times 7 = ?$ \quad $(0111 \times 0111 = ?)$

\textbf{Çözüm:} 2.\ yolda çarpılan (32 bit) sabittir; her adımda çarpanın en sağ bitine bakılır, bit 1 ise çarpılan çarpımın üst yarısına eklenir. Ardından çarpım yazmacının tamamı 1 bit sağa kaydırılır. Adımları izleyelim (Tablo~\ref{tab:carpma2-ornek}):

\begin{itemize}[nosep]
\item \textbf{Başlangıç:} Çarpılan $= \textcolor{vurgumavi}{0111}$, çarpan $= \textcolor{vurgukirmizi}{0111}$, çarpım $= \textcolor{gray}{0}$.
\item \textbf{Adım 1:} Çarpanın en sağ biti $= \textcolor{vurgukirmizi}{1}$ $\Rightarrow$ çarpılanı üst yarıya ekle: $\textcolor{gray}{0000} + \textcolor{vurgumavi}{0111} = \textcolor{vurguyesil}{0111}$. Sağa kaydır: $\textcolor{vurguyesil}{00111}\textcolor{gray}{000}$.
\item \textbf{Adım 2:} En sağ bit $= \textcolor{vurgukirmizi}{1}$ $\Rightarrow$ üst yarıya ekle: $\textcolor{vurguyesil}{0011} + \textcolor{vurgumavi}{0111} = \textcolor{vurguyesil}{1010}$. Sağa kaydır: $\textcolor{vurguyesil}{010101}\textcolor{gray}{00}$.
\item \textbf{Adım 3:} En sağ bit $= \textcolor{vurgukirmizi}{1}$ $\Rightarrow$ üst yarıya ekle: $\textcolor{vurguyesil}{0101} + \textcolor{vurgumavi}{0111} = \textcolor{vurguyesil}{1100}$. Sağa kaydır: $\textcolor{vurguyesil}{0110001}\textcolor{gray}{0}$.
\item \textbf{Adım 4:} En sağ bit $= \textcolor{vurgukirmizi}{0}$ $\Rightarrow$ toplama yok. Sağa kaydır: $\textcolor{vurguyesil}{00110001}$.
\end{itemize}

\begin{table}[H]
\centering
\caption{Örnek 5.5 -- 2.\ yol çarpma algoritması}
\label{tab:carpma2-ornek}
\small
\begin{tabular}{c l l l}
\toprule
\textbf{Adım} & \textbf{\textcolor{vurgumavi}{Çarpılan}} & \textbf{\textcolor{vurgukirmizi}{Çarpan}} & \textbf{\textcolor{vurguyesil}{Çarpım}} \\
\midrule
Başlangıç & \textcolor{vurgumavi}{0111} & \textcolor{vurgukirmizi}{0111} & \textcolor{gray}{00000000} \\
1 (Topla + Kaydır) & \textcolor{vurgumavi}{0111} & \textcolor{gray}{0}\textcolor{vurgukirmizi}{011} & \textcolor{vurguyesil}{00111}\textcolor{gray}{000} \\
2 (Topla + Kaydır) & \textcolor{vurgumavi}{0111} & \textcolor{gray}{00}\textcolor{vurgukirmizi}{01} & \textcolor{vurguyesil}{010101}\textcolor{gray}{00} \\
3 (Topla + Kaydır) & \textcolor{vurgumavi}{0111} & \textcolor{gray}{000}\textcolor{vurgukirmizi}{0} & \textcolor{vurguyesil}{0110001}\textcolor{gray}{0} \\
4 (Kaydır) & \textcolor{vurgumavi}{0111} & \textcolor{gray}{0000} & \textcolor{vurguyesil}{00110001} \\
\bottomrule
\end{tabular}
\end{table}

Çarpım $= (00110001)_2 = 49$
\end{ornek}

\subsubsection{3. Yol}

2.\ yoldaki çarpma sürecini adım adım izlediğimizde ilginç bir örüntü ortaya çıkar: çarpım yazmacının 64 bitlik alanı başlangıçta tamamen sıfırdır; her adımda toplama sonucu üst yarıya (sol 32 bit) yazılıp yazmaç sağa kaydırıldıkça, alt yarı (sağ 32 bit) yavaş yavaş kısmi çarpım bitleriyle dolmaya başlar. Öte yandan çarpan yazmacı da her adımda sağa kaydırılarak tüketilir: en sağdaki bit denetim birimine gönderilir ve kaydırmayla atılır; dolayısıyla çarpandaki anlamlı bit sayısı her adımda bir azalır.

\begin{figure}[htbp]
\centering
\begin{tikzpicture}[scale=0.72, transform shape]

    % ========== D{\"U}{\u{G}}{\"U}MLER ==========

    % ---- {\c{C}}arp{\i}lan Yazma{\c{c}}{\i} (32 bit, sol {\"u}st, sabit, sol yar{\i}yla hizal{\i}) ----
    % (Sabit) — kayd{\i}rma yok, bu nedenle {\"u}ste kayd{\i}rma bilgisi yaz{\i}lmaz
    \node[draw, rectangle, fill=blokmavi, thick, font=\footnotesize, align=center,
          minimum width=3.0cm, minimum height=1.3cm] (carpilan) at (-1.5, 5)
          {{\c{C}}arp{\i}lan Yazmac{\i} {\scriptsize(32 bit)}\\{\scriptsize(Sabit)}};

    % ---- AMB (32 bit Toplay{\i}c{\i}, sol yar{\i}yla hizal{\i}, 5.15'in yar{\i}s{\i} boyutunda) ----
    \node[draw, trapezium, trapezium angle=70, fill=bloksari, thick,
          shape border rotate=180, font=\footnotesize, align=center,
          minimum width=1.5cm, minimum height=1.0cm] (amb) at (-1.5, 3.0)
          {32 bit AMB};

    % ---- {\c{C}}arp{\i}m|{\c{C}}arpan birle{\c{s}}ik yazma{\c{c}} (64 bit, ikiye b{\"o}l{\"u}nm{\"u}{\c{s}}, altta) ----
    % Kayd{\i}rma bilgisi kutunun {\"u}st{\"u}nde, siyah
    % Sol yar{\i}: {\c{C}}arp{\i}m [63:32]
    \node[draw, rectangle, fill=blokmavi, thick, font=\footnotesize, align=center,
          minimum width=3.0cm, minimum height=1.3cm] (carpimhi) at (-1.5, 0.5)
          {{\c{C}}arp{\i}m};
    % Sa{\u{g}} yar{\i}: {\c{C}}arpan [31:0] — e{\c{s}}it boy, yaz{\i} sola kayd{\i}r{\i}lm{\i}{\c{s}}
    \node[draw, rectangle, fill=blokmavi!60, thick, font=\footnotesize,
          minimum width=3.0cm, minimum height=1.3cm] (carpimlo) at (1.5, 0.5) {};
    \node[font=\footnotesize] at ([xshift=-8mm]carpimlo.center) {{\c{C}}arpan};
    % D{\i}{\c{s}} {\c{c}}er{\c{c}}eve
    \draw[thick] (carpimhi.north west) rectangle (carpimlo.south east);
    % En sa{\u{g}} bit b{\"o}lmesi (dikey {\c{c}}izgiyle ayr{\i}lm{\i}{\c{s}})
    \draw ([xshift=-14mm]carpimlo.north east) -- ([xshift=-14mm]carpimlo.south east);
    \node[font=\tiny, align=center] at ([xshift=-7mm]carpimlo.east)
          {en sa{\u{g}}\\bit};
    % Bit etiketleri
    \node[font=\scriptsize, above=1pt of carpimhi.north west, anchor=south west] {[63:32]};
    \node[font=\scriptsize, above=1pt of carpimlo.north east, anchor=south east] {[31:0]};
    \node[font=\scriptsize, below=3pt] at (0, -0.15) {64 bit};
    % Kayd{\i}rma y{\"o}n g{\"o}stergesi: kutunun {\"u}st{\"u}nde, siyah
    \node[font=\scriptsize, above=14pt of carpimlo.north east, anchor=south east]
          {Sa{\u{g}}a Kayd{\i}r $\longrightarrow$};

    % ---- Denetim birimi (sa{\u{g}} tarafta) ----
    \node[draw, rectangle, fill=blokkirmizi-acik, thick, font=\footnotesize, align=center,
          minimum width=2.8cm, minimum height=0.8cm] (denetim) at (7, 3.0)
          {Denetim};

    % ========== VER{\.I} YOLLARI (kal{\i}n d{\"u}z {\c{c}}izgi) ==========

    % {\c{C}}arp{\i}lan $\rightarrow$ AMB ({\"u}stten giri{\c{s}})
    \draw[okstil] (carpilan.south) -- (amb.north);

    % AMB $\rightarrow$ {\c{C}}arp{\i}m {\"u}st yar{\i} [63:32] ({\c{c}}{\i}k{\i}{\c{s}} a{\c{s}}a{\u{g}}{\i})
    \draw[okstil] (amb.south) -- (carpimhi.north)
          node[midway, left=2pt, font=\scriptsize] {Yaz [63:32]};

    % {\c{C}}arp{\i}m {\"u}st yar{\i} $\rightarrow$ AMB geri besleme (sol taraftan dola{\c{s}}arak)
    \draw[okstil] (carpimhi.west) -- ++(-1.0, 0) coordinate (fb3)
          -- (fb3 |- amb.west) -- (amb.west)
          node[pos=0.0, left=2pt, font=\scriptsize] {[63:32]};

    % {\c{C}}arpan en sa{\u{g}} bit (sa{\u{g}} yar{\i} LSB) $\rightarrow$ Denetim (sağa kaydırılmış)
    \draw[okstil] ([xshift=-7mm]carpimlo.south east) -- ++(0, -0.6)
          -| ([xshift=4mm]denetim.south);

    % ========== DENET{\.I}M S{\.I}NYALLER{\.I} (k{\i}rm{\i}z{\i} kesikli) ==========

    % Denetim $\rightarrow$ {\c{C}}arp{\i}m|{\c{C}}arpan (kayd{\i}r sinyali, sola kaydırılmış)
    \draw[denetimstil] ([xshift=-4mm]denetim.south) -- ++(0, -1.2)
          -| ([xshift=8mm]carpimlo.east)
          -- (carpimlo.east);
    \node[font=\scriptsize, sinyalkirmizi, fill=white, inner sep=1pt, right=1pt]
          at ([xshift=4mm]carpimlo.east) {Kayd{\i}r/Yaz};

    % Denetim $\rightarrow$ AMB (topla sinyali)
    \draw[denetimstil] (denetim.west) -- (amb.east);
    \node[font=\scriptsize, sinyalkirmizi, fill=white, inner sep=1pt]
          at ($(denetim.west)!0.5!(amb.east)$) {Topla};

\end{tikzpicture}
\caption{Kayd{\i}r-Topla {\c{C}}arpma (3.\ Yol) donan{\i}m{\i}}
\label{fig:carpma-3surum-donanim}
\end{figure}

Buradaki kritik gözlem şudur: çarpım yazmacının alt yarısı dolarken, çarpan yazmacı aynı hızda boşalmaktadır. Her ikisi de sağa kaydırıldığından, çarpanın kalan bitleri çarpım yazmacının sağ yarısında tutulabilir. Böylece ayrı bir çarpan yazmacına gerek kalmaz; çarpan, çarpım yazmacının başlangıçta boş olan alt yarısına yerleştirilir ve her adımda birlikte sağa kaydırılır. Bu tasarım, 32 bitlik bir yazmacı ortadan kaldırarak donanımı daha da yalınlaştırır; 3.\ yolun donanım yapısı Şekil~\ref{fig:carpma-3surum-donanim}'de, akış şeması ise Şekil~\ref{fig:carpma-3yol-algo}'da gösterilmiştir.

\begin{figure}[htbp]
\centering
\begin{tikzpicture}[node distance=0.8cm, scale=0.72, transform shape]
    % Start
    \node[akisbaslangic] (start) {Ba{\c{s}}la};
    % Initialize
    \node[akisislem, below=of start, minimum width=4cm] (init)
        {{\c{C}}arp{\i}m[63:32] $\leftarrow$ 0\\{\c{C}}arp{\i}m[31:0] $\leftarrow$ {\c{C}}arpan\\Say $\leftarrow$ 0};
    % Decision: LSB of product
    \node[akiskarar, below=1cm of init, minimum width=2.5cm] (check)
        {{\c{C}}arp{\i}m{\i}n\\en sa{\u{g}} biti = 1?};
    % Add
    \node[akisislem, right=2cm of check, minimum width=4cm] (add)
        {{\c{C}}arp{\i}m[63:32] $\leftarrow$\\{\c{C}}arp{\i}m[63:32] + {\c{C}}arp{\i}lan};
    % Shift product right
    \node[akisislem, below=1.3cm of check, minimum width=4cm] (shiftp)
        {{\c{C}}arp{\i}m{\i} 1 bit sa{\u{g}}a kayd{\i}r};
    % Increment counter
    \node[akisislem, below=of shiftp, minimum width=4cm] (inc)
        {Say $\leftarrow$ Say + 1};
    % Check counter
    \node[akiskarar, below=of inc, minimum width=2cm] (countchk)
        {Say = 32?};
    % End
    \node[akisbaslangic, below=1cm of countchk] (stop) {Dur};
    % Arrows
    \draw[akisok] (start) -- (init);
    \draw[akisok] (init) -- (check);
    \draw[akisok] (check) -- node[above, font=\scriptsize] {Evet} (add);
    \draw[akisok] (check) -- node[left, font=\scriptsize] {Hay{\i}r} (shiftp);
    \draw[akisok] (add) |- (shiftp);
    \draw[akisok] (shiftp) -- (inc);
    \draw[akisok] (inc) -- (countchk);
    \draw[akisok] (countchk) -- node[left, font=\scriptsize] {Evet} (stop);
    \coordinate (dongu-sag) at ($(add.east)+(0.8,0)$);
    \draw[akisdonus] (countchk.east)
        -- node[above, font=\scriptsize, pos=0.2] {Hay{\i}r}
        (dongu-sag |- countchk.east)
        |- ($(check.north)+(0,0.5)$)
        -- (check.north);
\end{tikzpicture}
\caption[Kayd{\i}r-Topla {\c{C}}arpma (3.\ Yol) ak{\i}{\c{s}} {\c{s}}emas{\i}]{Kayd{\i}r-Topla {\c{C}}arpma (3.\ Yol) ak{\i}{\c{s}} {\c{s}}emas{\i}.}
\label{fig:carpma-3yol-algo}
\end{figure}

\FloatBarrier
\begin{ornek}[5.6]
\textbf{Soru:} $6 \times 3 = ?$ \quad $(0110 \times 0011 = ?)$

\textbf{Çözüm:} 3.\ yolda çarpan, çarpım yazmacının alt yarısında [31:0] tutulur. Her adımda birleşik yazmacın en sağ bitine (çarpanın en sağ biti) bakılır; bit 1 ise çarpılan üst yarıya [63:32] eklenir. Ardından yazmacın tamamı 1 bit sağa kaydırılır. Adımları izleyelim (Tablo~\ref{tab:carpma3-ornek}):

\begin{itemize}[nosep]
\item \textbf{Başlangıç:} Çarpılan $= \textcolor{vurgumavi}{0110}$. Birleşik yazmaç: üst yarı $= \textcolor{gray}{0000}$, alt yarı $= \textcolor{vurgukirmizi}{0011}$ (çarpan).
\item \textbf{Adım 1:} En sağ bit $= \textcolor{vurgukirmizi}{1}$ $\Rightarrow$ üst yarıya ekle: $\textcolor{gray}{0000} + \textcolor{vurgumavi}{0110} = \textcolor{vurguyesil}{0110}$. Tümünü sağa kaydır: $\textcolor{vurguyesil}{0011}\,|\,\textcolor{vurguyesil}{0}\textcolor{vurgukirmizi}{001}$.
\item \textbf{Adım 2:} En sağ bit $= \textcolor{vurgukirmizi}{1}$ $\Rightarrow$ üst yarıya ekle: $\textcolor{vurguyesil}{0011} + \textcolor{vurgumavi}{0110} = \textcolor{vurguyesil}{1001}$. Sağa kaydır: $\textcolor{vurguyesil}{0100}\,|\,\textcolor{vurguyesil}{10}\textcolor{vurgukirmizi}{00}$.
\item \textbf{Adım 3:} En sağ bit $= \textcolor{vurgukirmizi}{0}$ $\Rightarrow$ toplama yok. Sağa kaydır: $\textcolor{vurguyesil}{0010}\,|\,\textcolor{vurguyesil}{010}\textcolor{vurgukirmizi}{0}$.
\item \textbf{Adım 4:} En sağ bit $= \textcolor{vurgukirmizi}{0}$ $\Rightarrow$ toplama yok. Sağa kaydır: $\textcolor{vurguyesil}{0001}\,|\,\textcolor{vurguyesil}{0010}$. Çarpan tükendi.
\end{itemize}

\begin{table}[H]
\centering
\caption{Örnek 5.6 -- 3.\ yol çarpma algoritması}
\label{tab:carpma3-ornek}
\small
\begin{tabular}{c l l}
\toprule
\textbf{Adım} & \textbf{\textcolor{vurgumavi}{Çarpılan}} & \textbf{\textcolor{vurguyesil}{Çarpım [63:32]}\,|\,\textcolor{vurgukirmizi}{[31:0]=Çarpan}} \\
\midrule
Başlangıç & \textcolor{vurgumavi}{0110} & \textcolor{vurguyesil}{0000}\,|\,\textcolor{vurgukirmizi}{0011} \\
1 (Topla + Kaydır) & \textcolor{vurgumavi}{0110} & \textcolor{vurguyesil}{0011}\,|\,\textcolor{vurguyesil}{0}\textcolor{vurgukirmizi}{001} \\
2 (Topla + Kaydır) & \textcolor{vurgumavi}{0110} & \textcolor{vurguyesil}{0100}\,|\,\textcolor{vurguyesil}{10}\textcolor{vurgukirmizi}{00} \\
3 (Kaydır) & \textcolor{vurgumavi}{0110} & \textcolor{vurguyesil}{0010}\,|\,\textcolor{vurguyesil}{010}\textcolor{vurgukirmizi}{0} \\
4 (Kaydır) & \textcolor{vurgumavi}{0110} & \textcolor{vurguyesil}{0001}\,|\,\textcolor{vurguyesil}{0010} \\
\bottomrule
\end{tabular}
\end{table}

Çarpım $= (00010010)_2 = 18$
\end{ornek}

3.\ yol, kaydır-topla çarpma için en yalın donanım tasarımını sunar: yalnızca bir adet 32 bitlik çarpılan yazmacı, bir adet 32 bitlik AMB ve bir adet 64 bitlik birleşik çarpım/çarpan yazmacı yeterlidir. Yineleme başına yapılan işlemler 2.\ yol ile aynıdır; kazanım çevrim sayısında değil, bir yazmaç daha az kullanarak donanımın yalınlaşmasındadır.

\subsubsection{Karşılaştırma ve RISC-V Çarpma Buyrukları}

Üç yolun karşılaştırması Tablo~\ref{tab:carpma-yol-karsilastirma}'da özetlenmiştir. Her üç yolda da döngü 32 kez tekrarlanır ve her yineleme bir saat vuruşunda tamamlanabilir; dolayısıyla 32 bitlik bir çarpma \textbf{32 saat vuruşu} sürer. Yollar arasındaki temel fark, her vuruşun ne kadar hızlı olabileceği (AMB genişliği) ve donanımın ne kadar yer kapladığıdır (yazmaç sayısı).

\begin{table}[htbp]
\centering
\caption{Kaydır-topla çarpma yollarının karşılaştırması (32 bitlik RISC-V)}
\label{tab:carpma-yol-karsilastirma}
\small
\begin{tabular}{l c c c}
\toprule
& \textbf{AMB genişliği} & \textbf{Yazmaç sayısı} & \textbf{Yineleme başına işlem} \\
\midrule
1.\ Yol & 64 bit & 3 (64+32+64 bit) & test + toplama + 2 kaydırma$^*$ \\
2.\ Yol & 32 bit & 3 (32+32+64 bit) & test + toplama + 2 kaydırma$^*$ \\
3.\ Yol & 32 bit & 2 (32+64 bit) & test + toplama + kaydırma \\
\bottomrule
\multicolumn{4}{l}{\footnotesize $^*$İki kaydırma birbirinden bağımsızdır ve koşut yapılabilir.}
\end{tabular}
\end{table}

Bölüm~\ref{sec:carpma-bolme-sorunu}'de gördüğümüz gibi, RV32I'de yazmaçlar 32 bit genişliğindedir; oysa iki 32 bitlik sayının çarpımı 64 bite kadar çıkabilir. 64 bitlik sonucu tek bir 32 bitlik yazmaca yazmak mümkün olmadığından, RISC-V M uzantısı çarpma sonucunu \emph{iki ayrı buyrukla} elde eder: \texttt{mul} buyruğu sonucun alt 32 bitini (çarpım$[31{:}0]$), \texttt{mulh} buyruğu ise üst 32 bitini (çarpım$[63{:}32]$) hedef yazmacına yazar. Her iki buyruk da donanımda bağımsız birer çarpma işlemi gerçekleştirir; dolayısıyla tam 64 bitlik sonuca ihtiyaç duyan bir program bu iki buyruğu art arda çalıştırmalıdır ve toplam maliyet iki çarpma süresi olur. İşaretsiz çarpmanın üst yarısı için \texttt{mulhu}, bir işaretli bir işaretsiz çarpmanın üst yarısı için ise \texttt{mulhsu} buyruğu kullanılır (bu buyrukların kullanım örneği için bkz.\ Örnek~\ref{orn:carpma-64bit}).

\subsection{Booth Algoritması (İşaretli Sayılar)}
\label{sec:booth}

Kaydır-topla çarpma yalnızca işaretsiz sayılarla çalışır. İşaretli sayıları çarpmak için bir yol, işaretleri ayrı tutup işaretsiz çarpma yapmak ve sonuca doğru işareti eklemektir; ancak bu ek adımlar gerektirir. Booth algoritması\index{Booth algoritması}\index{Booth's algorithm|see{Booth algoritması}}\index{Booth, Andrew} çok daha zarif bir çözüm sunar: ikiye tümleyen gösterimindeki sayıları (ister artı ister eksi) \emph{hiçbir ayrım yapmadan} doğrudan çarpar. Ayrı bir işaret düzeltmesine veya farklı bir donanıma gerek kalmaz.

Booth algoritmasının çıkış noktası, ikili sayılardaki \emph{ardışık bit örüntüleridir}. Kaydır-topla yönteminde çarpanın her 1 biti bir toplama gerektirir. Ancak gerçek sayılarda ardışık 1'ler sıkça görülür; örneğin:
%
\begin{itemize}[nosep]
\item $30 = \underbrace{0\textbf{1111}0}_{\text{dört ardışık 1}}$; kaydır-topla ile 4 toplama gerekir,
\item $126 = 0\underbrace{\textbf{1111111}}_{\text{yedi ardışık 1}}0$; kaydır-topla ile 7 toplama gerekir.
\end{itemize}
%
Booth'un gözlemi şudur: ardışık $k$ tane 1 içeren bir sayı, yalnızca \emph{iki} işlemle (bir çıkarma ve bir toplama ile) ifade edilebilir. Örneğin $0111110_2 = 1000000_2 - 0000010_2$ olduğundan, 5 toplama yerine 1 çıkarma + 1 toplama yeterlidir. Algoritma, çarpanın bitlerini ikişerli çiftler halinde tarayarak ardışık 1'lerin \emph{başladığı} yerde çıkarma, \emph{bittiği} yerde toplama yapar; ardışık 0'lar için ise hiçbir işlem gerekmez. Bu şekilde toplama/çıkarma sayısı önemli ölçüde azalır.

Booth algoritmasının işlem sayısı, çarpanın bit örüntüsüne doğrudan bağlıdır:
%
\begin{itemize}[nosep]
\item \textbf{En iyi durum:} Çarpanda uzun ardışık 0 veya uzun ardışık 1 blokları varsa, toplama/çıkarma sayısı çok düşer. Örneğin $01111111110_2$ gibi bir çarpan için yalnızca 2 işlem (1 çıkarma + 1 toplama) yeterlidir; kaydır-topla ile 9 toplama gerekirdi.
\item \textbf{En kötü durum:} Çarpanın bitleri sürekli değişiyorsa ($\ldots 010101 \ldots$), her bit geçişi bir işlem tetikler ve toplama/çıkarma sayısı kaydır-topla yönteminden bile fazla olabilir. Örneğin $0101010101_2$ çarpanı için her adımda bir işlem gerekir.
\end{itemize}
%
Bu nedenle Booth algoritmasının başarımı \emph{veriye bağlıdır} (\emph{data-dependent}): aynı donanım, çarpanın değerine göre farklı sayıda aritmetik işlem yapar. Bu durum, sabit gecikmeli tasarım gerektiren boru hatlı işlemcilerde (bkz.\ Bölüm~\ref{chap:boruhatti}) ek zorluklar yaratabilir. Algoritmanın kuralları Tablo~\ref{tab:booth-kurallar}'de listelenmiştir, akış şeması ise Şekil~\ref{fig:booth-algo}'da gösterilmiştir.

\begin{figure}[htbp]
\centering
\begin{tikzpicture}[node distance=0.8cm, scale=0.72, transform shape]
    % Start
    \node[akisbaslangic] (start) {Ba{\c{s}}la};
    % Initialize
    \node[akisislem, below=of start, minimum width=4cm] (init)
        {{\c{C}}arp{\i}m $\leftarrow$ 0, Say $\leftarrow$ 0\\Son bit sa{\u{g}}{\i}nda 0 varsay};
    % Decision: check last 2 bits
    \node[akiskarar, below=1cm of init, minimum width=2.5cm, aspect=2.5] (check)
        {Son 2 bit?};
    % 10: Subtract
    \node[akisislem, left=2cm of check, minimum width=3.2cm] (sub)
        {{\c{C}}arp{\i}m $\leftarrow$\\{\c{C}}arp{\i}m $-$ {\c{C}}arp{\i}lan};
    % 01: Add
    \node[akisislem, right=2cm of check, minimum width=3.2cm] (add)
        {{\c{C}}arp{\i}m $\leftarrow$\\{\c{C}}arp{\i}m + {\c{C}}arp{\i}lan};
    % Shift right (arithmetic)
    \node[akisislem, below=2.2cm of check, minimum width=4cm] (shift)
        {{\c{C}}arp{\i}m{\i} 1 bit sa{\u{g}}a\\aritmetik kayd{\i}r};
    % Increment
    \node[akisislem, below=of shift, minimum width=4cm] (inc)
        {Say $\leftarrow$ Say + 1};
    % Check counter
    \node[akiskarar, below=of inc, minimum width=2cm] (countchk)
        {Say = 32?};
    % End
    \node[akisbaslangic, below=1cm of countchk] (stop) {Dur};
    % Arrows
    \draw[akisok] (start) -- (init);
    \draw[akisok] (init) -- (check);
    \draw[akisok] (check) -- node[above, font=\scriptsize] {10} (sub);
    \draw[akisok] (check) -- node[above, font=\scriptsize] {01} (add);
    \draw[akisok] (check) -- node[left, font=\scriptsize, align=center] {00\\veya\\11} (shift);
    \draw[akisok] (sub) |- (shift);
    \draw[akisok] (add) |- (shift);
    \draw[akisok] (shift) -- (inc);
    \draw[akisok] (inc) -- (countchk);
    \draw[akisok] (countchk) -- node[left, font=\scriptsize] {Evet} (stop);
    \coordinate (dongu-sag) at ($(add.east)+(0.8,0)$);
    \draw[akisdonus] (countchk.east)
        -- node[above, font=\scriptsize, pos=0.2] {Hay{\i}r}
        (dongu-sag |- countchk.east)
        |- ($(check.north)+(0,0.5)$)
        -- (check.north);
\end{tikzpicture}
\caption[Booth {\c{C}}arp{\i}m Algoritmas{\i} ak{\i}{\c{s}} {\c{s}}emas{\i}]{Booth {\c{C}}arp{\i}m Algoritmas{\i} ak{\i}{\c{s}} {\c{s}}emas{\i}.}
\label{fig:booth-algo}
\end{figure}

\begin{bilginot}[Booth Algoritmasının Tarihçesi]
Andrew D. Booth\index{Booth, Andrew} bu algoritmayı 1950 yılında, Birkbeck Koleji'nde masaüstü hesap makineleri üzerinde çalışırken geliştirmiştir. Orijinal yöntem ardışık 1 dizilerindeki gereksiz toplama işlemlerini ortadan kaldırarak çarpma süresini kısaltır.
\end{bilginot}

\begin{table}[htbp]
\centering
\caption{Booth Algoritması Kuralları}
\label{tab:booth-kurallar}
\begin{tabular}{c l}
\toprule
\textbf{Bit Çifti} & \textbf{İşlem} \\
\midrule
10 & 1 dizisi başlar $\rightarrow$ çarpılan çıkarılır, sağa kaydırma \\
11 & 1 dizisinin ortası $\rightarrow$ sağa kaydırma \\
01 & 1 dizisi biter $\rightarrow$ çarpılan toplanır, sağa kaydırma \\
00 & 0 dizisinin ortası $\rightarrow$ sağa kaydırma \\
\bottomrule
\end{tabular}
\end{table}

Algoritmanın en önemli özelliği işaretli sayılarla işlem yapabilmesidir. Bunun için kaydır-topla yöntemindeki \emph{mantıksal} sağa kaydırma yerine \emph{aritmetik} sağa kaydırma\index{aritmetik kaydırma}\index{arithmetic shift|see{aritmetik kaydırma}} kullanılır. Mantıksal kaydırmada soldan her zaman 0 girer; aritmetik kaydırmada ise en soldaki bit (işaret biti) korunarak tekrarlanır. Böylece eksi bir kısmi çarpım sağa kaydırıldığında işareti bozulmaz: örneğin $\textcolor{vurguyesil}{1101}$ aritmetik sağa kaydırılırsa $\textcolor{vurguyesil}{1110}$ olur; en soldaki 1 korunmuştur. Bu sayede ikiye tümleyen gösterimindeki eksi ara sonuçlar kaydırma sırasında doğru kalır. Ek olarak, en sağ basamağın da sağında 0 olduğu farz edilir.

Booth algoritması donanım açısından çarpma 3.\ yolunun (Şekil~\ref{fig:carpma-3surum-donanim}) üzerine kuruludur: aynı birleşik çarpım/çarpan yazmacı, aynı 32 bitlik çarpılan yazmacı ve aynı AMB kullanılır. İki temel fark vardır: (1)~AMB yalnızca toplama değil, çıkarma da yapabilmelidir ve (2)~denetim birimi çarpanın en sağ bitini tek başına değil, bir önceki bitle birlikte iki bitlik çift olarak değerlendirir (Tablo~\ref{tab:booth-kurallar}). Döngü sayısı değişmez: $n$ bitlik çarpan için yine $n$ adım gerekir.

\begin{ornek}[5.7]
\textbf{Soru:} $3 \times (-4) = ?$ \quad Booth algoritmasıyla 4 bitlik işaretli sayılarla çözün.

\textbf{Çözüm:} İkiye tümleyen gösterimde çarpılan $= \textcolor{vurgumavi}{0011}$ (3), çarpan $= \textcolor{vurgukirmizi}{1100}$ ($-4$). Booth algoritmasında çarpanın sağına başlangıçta $0$ eklenir; her adımda en sağdaki iki bit (mevcut bit ve önceki bit) incelenir. Adımları izleyelim:

\begin{itemize}[nosep]
\item \textbf{Başlangıç:} Çarpım üst yarı $= \textcolor{gray}{0000}$, çarpan $= \textcolor{vurgukirmizi}{1100}\mathbf{0}$ (eklenen bit).
\item \textbf{Adım 1:} En sağ iki bit $= \textcolor{vurgukirmizi}{00}$ $\Rightarrow$ işlem yok (Tablo~\ref{tab:booth-kurallar}). Aritmetik sağa kaydır.
\item \textbf{Adım 2:} En sağ iki bit $= \textcolor{vurgukirmizi}{00}$ $\Rightarrow$ işlem yok. Aritmetik sağa kaydır.
\item \textbf{Adım 3:} En sağ iki bit $= \textcolor{vurgukirmizi}{10}$ $\Rightarrow$ ardışık 1'lerin başlangıcı, çarpılanı \emph{çıkar}: $\textcolor{gray}{0000} - \textcolor{vurgumavi}{0011} = \textcolor{vurguyesil}{1101}$. Aritmetik sağa kaydır.
\item \textbf{Adım 4:} En sağ iki bit $= \textcolor{vurgukirmizi}{11}$ $\Rightarrow$ ardışık 1'lerin ortası, işlem yok. Aritmetik sağa kaydır.
\end{itemize}

\begin{center}
\small
\begin{tabular}{c c c c c l}
\toprule
\textbf{Adım} & \textbf{\textcolor{vurgumavi}{Çarpılan}} & \textbf{\textcolor{vurguyesil}{Çarpım (üst)}} & \textbf{Çarpan} & \textbf{Bit çifti} & \textbf{İşlem} \\
\midrule
Başla & \textcolor{vurgumavi}{0011} & $\textcolor{gray}{0000}$ & $\textcolor{vurgukirmizi}{1100}\mathbf{0}$ & --- & --- \\
1 & \textcolor{vurgumavi}{0011} & $\textcolor{gray}{0000}$ & $\textcolor{vurgukirmizi}{1100}\mathbf{0}$ & $00$ & İşlem yok, aritmetik sağa kaydır \\
  & & $\textcolor{gray}{0000}$ & $\textcolor{vurguyesil}{0}\textcolor{vurgukirmizi}{110}\mathbf{0}$ & & \\
2 & \textcolor{vurgumavi}{0011} & $\textcolor{gray}{0000}$ & $\textcolor{vurguyesil}{0}\textcolor{vurgukirmizi}{110}\mathbf{0}$ & $00$ & İşlem yok, aritmetik sağa kaydır \\
  & & $\textcolor{gray}{0000}$ & $\textcolor{vurguyesil}{00}\textcolor{vurgukirmizi}{11}\mathbf{0}$ & & \\
3 & \textcolor{vurgumavi}{0011} & $\textcolor{gray}{0000}$ & $\textcolor{vurguyesil}{00}\textcolor{vurgukirmizi}{11}\mathbf{0}$ & $10$ & Çıkar: $0000 - \textcolor{vurgumavi}{0011} = \textcolor{vurguyesil}{1101}$, kaydır \\
  & & $\textcolor{vurguyesil}{1110}$ & $\textcolor{vurguyesil}{100}\textcolor{vurgukirmizi}{1}\mathbf{1}$ & & \\
4 & \textcolor{vurgumavi}{0011} & $\textcolor{vurguyesil}{1110}$ & $\textcolor{vurguyesil}{100}\textcolor{vurgukirmizi}{1}\mathbf{1}$ & $11$ & İşlem yok, aritmetik sağa kaydır \\
  & & $\textcolor{vurguyesil}{1111}$ & $\textcolor{vurguyesil}{0100}\mathbf{1}$ & & \\
\bottomrule
\end{tabular}
\end{center}

Sonuç: $\textcolor{vurguyesil}{11110100}_2 = -12_{10}$. Doğrulama: $3 \times (-4) = -12$\;\checkmark. Booth algoritması \emph{aritmetik} sağa kaydırma (işaret bitini koruyarak) kullandığından, ikiye tümleyen sayılarla doğrudan çarpım yapabilir; ayrı bir işaret düzeltmesine gerek kalmaz.
\end{ornek}

\subsection{Değiştirilmiş Booth Kodlaması (Radix-4)}
\label{sec:radix4-booth}
\index{değiştirilmiş Booth kodlaması}\index{modified Booth encoding|see{değiştirilmiş Booth kodlaması}}\index{radix-4 Booth|see{değiştirilmiş Booth kodlaması}}

Bir önceki bölümde gördüğümüz gibi, orijinal Booth algoritmasının başarımı veriye bağlıdır: çarpandaki bitler sürekli 0-1 arasında değişiyorsa ($\ldots 010101 \ldots$), her adımda bir toplama veya çıkarma tetiklenir ve kaydır-topla yönteminden bile yavaş kalınabilir. Üstelik Booth algoritması kısmi çarpım sayısını azaltsa da, $n$ bitlik bir çarpan yine $n$ adım gerektirir; her adımda yalnızca 1 bit işlenir.

\begin{table}[htbp]
\centering
\caption{Radix-4 Booth kodlama tablosu ($M$ = çarpılan)}
\label{tab:radix4-booth}
\begin{tabular}{c c l}
\toprule
\textbf{Bit üçlüsü} & \textbf{Kısmi çarpım} & \textbf{Açıklama} \\
\midrule
000 & $0$ & Toplama yok \\
001 & $+M$ & Çarpılanı ekle \\
010 & $+M$ & Çarpılanı ekle \\
011 & $+2M$ & Çarpılanı 1 bit sola kaydırıp ekle \\
100 & $-2M$ & Çarpılanı 1 bit sola kaydırıp çıkar \\
101 & $-M$ & Çarpılanı çıkar \\
110 & $-M$ & Çarpılanı çıkar \\
111 & $0$ & Toplama yok \\
\bottomrule
\end{tabular}
\end{table}

\emph{Değiştirilmiş Booth kodlaması}\index{değiştirilmiş Booth kodlaması}\index{modified Booth encoding|see{değiştirilmiş Booth kodlaması}}\index{radix-4 Booth|see{değiştirilmiş Booth kodlaması}} (\emph{modified Booth encoding}, radix-4) bu iki sorunu birden çözer. Her adımda çarpanın 1 biti yerine \textbf{2 bitini} birden işleyerek adım sayısını yarıya indirir: 32 bitlik bir çarpan artık 32 değil \textbf{16 adımda} tamamlanır. Bunu başarmak için çarpanın bitleri 3'lük örtüşen gruplar halinde incelenir. Her üçlü, Tablo~\ref{tab:radix4-booth}'de gösterildiği gibi beş olası kısmi çarpımdan birini üretir.

Bu kodlamanın en büyük kazanımı, $n$ bitlik bir çarpan için yalnızca $n/2$ kısmi çarpım üretmesidir; orijinal Booth'a göre adım sayısı yarıya iner. $+2M$ ve $-2M$ değerleri kolay bir sola kaydırma ile elde edilebildiğinden ek donanım maliyeti düşüktür.

Günümüzde neredeyse tüm yüksek başarımlı çarpma birimleri radix-4 Booth kodlamasını kısmi çarpım üretimi için kullanır. Daha yüksek tabanlı kodlamalar (radix-8, radix-16) da mevcuttur ancak $\pm 3M$ gibi değerlerin üretilmesi ek toplayıcı gerektirdiğinden radix-4 en yaygın tercih olmaya devam eder.

\begin{ornek}[5.8]
\textbf{Soru:} Radix-4 Booth kodlamasını kullanarak $13 \times 11$ çarpımını 6 bitlik işaretli sayılarla gerçekleştirin. ($13 = 001101_2$, $11 = 001011_2$)

\textbf{Çözüm:} Çarpan $= 001011_2$ (11), çarpılan $M = 001101_2$ (13).

\textbf{Adım 1 — Çarpanı 3'erli örtüşen gruplara ayır.} 6 bitlik çarpan için $n/2 = 3$ grup oluşturulur. Her grup, önceki grubun en düşük anlamlı bitiyle örtüşür; başlangıçta sağa sıfır (0) eklenir:

Çarpan bitleri: $b_5 b_4 b_3 b_2 b_1 b_0 = 0\,0\,1\,0\,1\,1$, $b_{-1} = 0$ (eklenen sıfır).

\begin{center}
\small
\begin{tabular}{c c c c}
\toprule
\textbf{Grup} & \textbf{Bit üçlüsü} ($b_{2i+1}\,b_{2i}\,b_{2i-1}$) & \textbf{Kısmi çarpım} & \textbf{Ağırlık} \\
\midrule
Grup 0 ($i=0$) & $b_1\,b_0\,b_{-1} = 1\,1\,0$ & $-M = -13$ & $2^0 = 1$ \\
Grup 1 ($i=1$) & $b_3\,b_2\,b_1 = 1\,0\,1$ & $-M = -13$ & $2^2 = 4$ \\
Grup 2 ($i=2$) & $b_5\,b_4\,b_3 = 0\,0\,1$ & $+M = +13$ & $2^4 = 16$ \\
\bottomrule
\end{tabular}
\end{center}

Kodlama tablosuna göre: üçlü $110 \to -M$, üçlü $101 \to -M$, üçlü $001 \to +M$.

\textbf{Adım 2 — Kısmi çarpımları hizalayarak topla.} Her kısmi çarpım, grubun ağırlığıyla çarpılarak hizalanır:

\begin{center}
\small
\begin{tabular}{r r l}
& $-13$ & (Grup 0: $-M \times 2^0 = -13$) \\
$+$ & $-52$ & (Grup 1: $-M \times 2^2 = -13 \times 4 = -52$) \\
$+$ & $+208$ & (Grup 2: $+M \times 2^4 = +13 \times 16 = +208$) \\
\hline
\multicolumn{2}{r}{$\mathbf{143}$} & \\
\end{tabular}
\end{center}

\textbf{Sonuç:} $-13 + (-52) + 208 = 143 = 13 \times 11$. \checkmark

Radix-4 Booth, 6 bitlik çarpan için yalnızca 3 kısmi çarpım üretmiştir; orijinal Booth'un 6 adımı yerine 3 adım yeterli olmuştur.
\end{ornek}


\subsection{Koşut Çarpıcılar}
\label{sec:paralel-carpicilar}
\index{koşut çarpıcı}\index{parallel multiplier|see{koşut çarpıcı}}

Şimdiye kadar incelediğimiz tüm çarpma yöntemleri \emph{yinelemeli} yöntemlerdir: yazmaçlar, saat sinyali ve denetim birimi içeren \emph{ardışıl devrelerdir} (\emph{sequential circuits}); her saat çevriminde bir kısmi çarpım işlenir, $n$ bitlik çarpma için $n$ (veya radix-4 ile $n/2$) çevrim gerekir.

Bu bölümde ele alacağımız koşut çarpıcılar ise temelden farklıdır: bunlar saat sinyali gerektirmeyen \emph{birleşik devrelerdir} (\emph{combinational circuits}). Girdiler (çarpılan ve çarpan) devreye verildiği anda, çıkış (çarpım) belirli bir \emph{yayılım gecikmesi} sonrasında hazır olur; herhangi bir saat vuruşu beklenmez. Gecikme yalnızca kapıların ve bağlantıların fiziksel yayılım süresine bağlıdır.

Birleşik çarpıcının toplam gecikmesi tek bir saat çevrimine sığıyorsa çarpma \emph{tek saat vuruşunda} tamamlanır. Sığmıyorsa (ki yüksek saat hızlarında bu sık görülür), devrenin arasına yazmaçlar eklenerek birden fazla aşamaya bölünür ve çarpma birkaç saat vuruşunda tamamlanır. Bu teknik, Bölüm~\ref{chap:boruhatti}'de ayrıntılı olarak ele alınacak olan \emph{boru hattı} yapısının bir uygulamasıdır.

\begin{terimnotu}
\textbf{birleşik devre} (\emph{combinational circuit})\,: Çıkışı yalnızca o anki girdilerin \emph{birleşimine} (kombinasyonuna) bağlı olan devredir. Bellek öğesi (yazmaç, flip-flop) içermez; saat sinyali gerektirmez. Girdiler değiştiğinde çıkış, kapıların yayılım gecikmesi kadar bir süre sonra yeni değerine ulaşır. Toplayıcılar, çoğullayıcılar ve bu bölümde incelenen koşut çarpıcılar birleşik devrelere örnektir.

\textbf{ardışıl devre} (\emph{sequential circuit})\,: Çıkışı hem o anki girdilere hem de devrenin \emph{önceki durumuna} (belleğine) bağlı olan devredir. Flip-flop veya yazmaç gibi bellek öğeleri içerir ve bir saat sinyaliyle yönetilir. Her saat vuruşunda devre bir sonraki durumuna geçer. Kaydır-topla çarpıcılar, sayaçlar ve işlemcinin denetim birimi ardışıl devrelere örnektir.
\end{terimnotu}

\subsubsection{Dizi Çarpıcı}
\label{sec:dizi-carpici}
\index{dizi çarpıcı}\index{array multiplier|see{dizi çarpıcı}}

Dizi çarpıcı (\emph{array multiplier}), kaydır-topla algoritmasının doğrudan donanıma çevrilmiş halidir. Her kısmi çarpım biti, çarpılanın ilgili biti ile çarpanın ilgili bitinin VE işlemiyle üretilir; ardından bu kısmi çarpımlar bir ızgara düzeninde tam toplayıcılarla birbirine eklenir. $n \times n$ bitlik bir çarpma için $n^2$ adet VE kapısı ve yaklaşık $n \times (n-1)$ adet tam toplayıcı kullanılır.

\begin{figure}[htbp]
\centering
\begin{tikzpicture}[scale=0.75, transform shape,
    cell/.style={draw, rectangle, minimum width=1.6cm, minimum height=1.6cm,
                 fill=bloksari!30, thick, rounded corners=2pt, font=\footnotesize, align=center},
]
% =============================================================
% 4×4 dizi çarpıcı — birleşik hücreler (her hücre = VE + TT)
% Hücreler ızgara düzeninde; bağlantılar yalnızca komşu hücreler arasında.
% =============================================================
\def\cs{2.4}   % sütun aralığı
\def\rs{2.4}   % satır aralığı

% ===== Sütun ve satır başlıkları =====
\foreach \c/\ai in {0/3, 1/2, 2/1, 3/0} {
    \node[font=\normalsize, anchor=south] at (\c*\cs, 1.2) {$a_{\ai}$};
}
\foreach \r/\bi in {0/0, 1/1, 2/2, 3/3} {
    \node[font=\normalsize, anchor=east] at (-1.4, -\r*\rs) {$b_{\bi}$};
}

% ===== Hücreleri yerleştir =====
\foreach \r in {0,1,2,3} {
    \foreach \c/\ai in {0/3, 1/2, 2/1, 3/0} {
        \node[cell] (h\r\c) at (\c*\cs, -\r*\rs)
            {$a_{\ai}{\cdot}b_{\r}$};
    }
}

% ===== Toplam bağlantıları (mavi, aşağı-sağa çapraz) =====
% Her hücrenin toplam çıkışı, bir alt satırda bir sağdaki hücreye girer.
% h(r,c).south → h(r+1,c+1).north — ama c+1 <= 3
\foreach \r/\rr in {0/1, 1/2, 2/3} {
    \foreach \c/\cc in {0/1, 1/2, 2/3} {
        \draw[->, thick, sinyalmavi]
            ([xshift=4mm]h\r\c.south) -- ([xshift=-4mm]h\rr\cc.north);
    }
}

% ===== Elde bağlantıları (kırmızı, sola) =====
% Her hücrenin elde çıkışı, aynı satırda soldaki hücreye girer.
% h(r,c).west → h(r,c-1).east — ama c-1 >= 0
\foreach \r in {1,2,3} {
    \foreach \c/\cc in {1/0, 2/1, 3/2} {
        \draw[->, thick, sinyalkirmizi] (h\r\c.west) -- (h\r\cc.east);
    }
}

% ===== Çıkışlar =====
% p0: h(0,3) sağ alt — en sağ üst hücre
\draw[->, thick] (h03.east) -- ++(0.8, 0) node[right, font=\small] {$p_0$};

% p1–p3: her satırın en sağ hücresinin toplam çıkışı
\draw[->, thick] (h13.east) -- ++(0.8, 0) node[right, font=\small] {$p_1$};
\draw[->, thick] (h23.east) -- ++(0.8, 0) node[right, font=\small] {$p_2$};
\draw[->, thick] (h33.east) -- ++(0.8, 0) node[right, font=\small] {$p_3$};

% p4–p7: son satırın altından
\draw[->, thick] (h33.south) -- ++(0, -0.7) node[below, font=\small] {$p_4$};
\draw[->, thick] (h32.south) -- ++(0, -0.7) node[below, font=\small] {$p_5$};
\draw[->, thick] (h31.south) -- ++(0, -0.7) node[below, font=\small] {$p_6$};
\draw[->, thick] (h30.south) -- ++(0, -0.7) node[below, font=\small] {$p_7$};

% ===== Satır 0 özel: ilk satırda TT yok, yalnızca VE =====
% İlk satır hücrelerini farklı renkle vurgula
\foreach \c in {0,1,2,3} {
    \node[cell, fill=blokyesil!20] (h0\c) at (\c*\cs, 0)
        {$a_{\ifcase\c 3\or 2\or 1\or 0\fi}{\cdot}b_{0}$\\[-1pt]\scriptsize(yalnızca VE)};
}

% Diğer hücreleri yeniden çiz (VE+TT olarak etiketle)
\foreach \r in {1,2,3} {
    \foreach \c/\ai in {0/3, 1/2, 2/1, 3/0} {
        \node[cell] (h\r\c) at (\c*\cs, -\r*\rs)
            {$a_{\ai}{\cdot}b_{\r}$\\[-1pt]\scriptsize VE+TT};
    }
}

% ===== İnset: tek hücrenin iç yapısı (sağda, ortada) =====
\begin{scope}[shift={(3*\cs+6.0, -3.2)}]
    % Çerçeve (VE ve TT'yi tam saran kutu — önce çiz, elemanlar üste binsin)
    \draw[thick, rounded corners=4pt, fill=white, draw=gray!70]
        (-1.6, -1.5) rectangle (2.0, 1.95);
    %
    % Başlık (kutunun dışında, belirgin boşlukla üstte)
    \node[font=\footnotesize\bfseries, anchor=south] at (0.2, 3.15) {Hücre iç yapısı};
    %
    % --- VE kapısı (üstte, ortada) ---
    \node[draw, circle, minimum size=0.8cm, fill=blokyesil!60,
          font=\small, thick] (ive) at (0.2, 1.2) {$\wedge$};
    \node[font=\scriptsize, anchor=west] at (0.9, 1.2) {VE};
    %
    % --- TT (tam toplayıcı, altta) ---
    \node[draw, rectangle, minimum width=2.4cm, minimum height=1.0cm,
          fill=blokturuncu!50, font=\small, thick, rounded corners=2pt]
        (itt) at (0.2, -0.6) {TT};
    %
    % ====== GİRİŞLER ======
    %
    % Giriş: a_j ve b_i → VE (üstten, kutunun dışından)
    \draw[->, thick] (-0.2, 2.8) -- (-0.2, 1.6);
    \node[font=\scriptsize, anchor=south] at (-0.2, 2.85) {$a_j$};
    \draw[->, thick] (0.6, 2.8) -- (0.6, 1.6);
    \node[font=\scriptsize, anchor=south] at (0.6, 2.85) {$b_i$};
    %
    % VE çıkışı → TT üst giriş (kısmi çarpım, gri — kutu içi)
    \draw[->, very thick, gray] (ive.south) -- (itt.north);
    \node[font=\scriptsize, gray, anchor=west] at (0.7, 0.3) {$a_j{\cdot}b_i$};
    %
    % Giren toplam: yukarıdan çapraz gelip TT'ye üstten dik girer (mavi)
    \draw[->, very thick, sinyalmavi, rounded corners=4pt]
        (-2.8, 1.6) -- (-2.8, 0.4) -- ([xshift=-5mm]itt.north |- {0,0.4}) -- ([xshift=-5mm]itt.north);
    \node[font=\scriptsize, sinyalmavi, anchor=south] at (-2.8, 1.65)
        {giren toplam};
    %
    % Giren elde: sağdan → TT sağ giriş (kırmızı)
    \draw[<-, very thick, sinyalkirmizi]
        (itt.east) -- (3.4, -0.6);
    \node[font=\scriptsize, sinyalkirmizi, anchor=west] at (3.5, -0.6)
        {giren elde};
    %
    % ====== ÇIKIŞLAR ======
    %
    % Çıkan elde: TT'den yatay sola (kırmızı)
    \draw[->, very thick, sinyalkirmizi]
        (itt.west) -- (-2.8, -0.6);
    \node[font=\scriptsize, sinyalkirmizi, anchor=east] at (-2.9, -0.6)
        {çıkan elde};
    %
    % Çıkan toplam: TT'den aşağı (mavi)
    \draw[->, very thick, sinyalmavi]
        (itt.south) -- ++(0, -1.3);
    \node[font=\scriptsize, sinyalmavi, anchor=west] at (0.4, -2.1)
        {çıkan toplam};
\end{scope}

\end{tikzpicture}
\caption[Dizi \c{c}arp{\i}c{\i}: 4 bit]{4 bitlik dizi çarpıcının ızgara yapısı. Her hücrede bir $\wedge$ (VE kapısı) kısmi çarpım bitini üretir; altındaki TT (tam toplayıcı) bu biti, üst satırdan gelen toplamla ve sağ komşudan gelen eldeyle toplar. Her hücrenin iç yapısında \textcolor{gray}{gri ok} VE kapısından tam toplayıcıya giden kısmi çarpım bitini, \textcolor{sinyalmavi}{mavi oklar} üst satırdan gelen toplamı, \textcolor{sinyalkirmizi}{kırmızı oklar} ise elde yayılımını (sağdan sola) gösterir.}
\label{fig:dizi-carpici-4bit}
\end{figure}

Şekil~\ref{fig:dizi-carpici-4bit}, 4 bitlik bir dizi çarpıcının yapısını göstermektedir. En üst satırda kısmi çarpım bitleri ($a_i \cdot b_0$) doğrudan çıkışa verilir. Her sonraki satırda bir üst satırdan gelen toplam ve elde bitleri, o satırın kısmi çarpım bitleriyle ($a_i \cdot b_j$) tam toplayıcılarda toplanır. Eldeler bir sonraki sütuna, toplamlar aşağıya aktarılır.

Dizi çarpıcının üstünlüğü yapısının düzenli ve tasarımının yalın olmasıdır; VLSI yerleşiminde düzgün bir ızgara oluşturur. Ancak \emph{gecikme}, eldeler en üst satırdan en alt satıra doğru yayıldığından $O(n)$ mertebesindedir ve donanım maliyeti $O(n^2)$'dir. 32 bitlik bir dizi çarpıcının gecikme süresi kabaca 32 tam toplayıcının seri gecikmesine eşittir; bu da çağdaş işlemciler için kabul edilemez derecede yavaştır.

\subsubsection{Wallace Ağacı Çarpıcı}
\label{sec:wallace-agaci}
\index{Wallace ağacı}\index{Wallace tree|see{Wallace ağacı}}

Wallace ağacı\index{Wallace, Chris} (\emph{Wallace tree}), 1964 yılında Chris Wallace tarafından önerilen ve dizi çarpıcının gecikme sorununu çözen bir yapıdır. Dizi çarpıcıda kısmi çarpımlar bir ızgarada satır satır toplanır ve eldeler en üst satırdan en alt satıra seri olarak yayılır. Wallace ağacının temel fikri, bu sıralı toplamayı ortadan kaldırıp kısmi çarpımları \emph{ağaç yapısında koşut olarak} indirgemektir.

Yöntemin anahtar bileşeni \emph{3:2 sıkıştırıcı} (carry-save adder, CSA) adı verilen devredir. Bir 3:2 sıkıştırıcı, aslında her bit sütununa bağımsız olarak uygulanan bir tam toplayıcıdır: üç sayıyı girdi olarak alır ve iki sayı (toplam $T$ ve elde $C$) üretir. Kritik nokta şudur: elde bitleri bir sonraki sütuna taşınır ancak sütunlar arası seri elde yayılımı \emph{yoktur}. Her sütundaki tam toplayıcı bağımsız çalıştığından bir CSA'nın gecikmesi yalnızca tek bir tam toplayıcı gecikmesine eşittir; girdi sayısından bağımsızdır.

Wallace ağacı bu sıkıştırıcıları şu şekilde kullanır:
\begin{enumerate}[nosep]
    \item $n$ adet kısmi çarpım 3'erli gruplara ayrılır.
    \item Her grupta bir 3:2 CSA ile 3 sayı $\to$ 2 sayıya (toplam + elde) indirgenir.
    \item Geriye kalan sayılar yeniden 3'erli gruplara ayrılarak işlem yinelenir.
    \item Geriye yalnızca 2 sayı kalınca bunlar hızlı bir önceden elde üreten toplayıcıyla (CLA) toplanır ve sonuç elde edilir.
\end{enumerate}

Her aşamada sayı adedi kabaca $\frac{2}{3}$'e düştüğünden toplam aşama sayısı $O(\log_{1{,}5} n) \approx O(\log n)$'dir; dizi çarpıcının $O(n)$ gecikmesiyle kıyaslanınca büyük bir kazanımdır. Donanım maliyeti dizi çarpıcıyla benzerdir ($O(n^2)$ tam toplayıcı); ancak bağlantı düzeni ızgara yerine ağaç biçiminde olduğundan daha karmaşıktır. 32 bitlik bir çarpma örneğinde: 32 kısmi çarpım $\to$ 22 $\to$ 15 $\to$ 10 $\to$ 7 $\to$ 5 $\to$ 4 $\to$ 3 $\to$ 2 sayı; yalnızca 8 sıkıştırma aşaması yeterlidir.

\paragraph{3:2 Sıkıştırıcı Nasıl Çalışır?}\mbox{}\\
3:2 sıkıştırıcının (CSA) çalışma mantığını küçük bir örnekle görelim. Üç adet 4 bitlik sayıyı toplamak istediğimizi düşünelim: $A = 1101$ (13), $B = 0110$ (6) ve $C = 0011$ (3). Gerçek toplam $13+6+3 = 22$ olmalıdır.

Normal bir toplayıcıda elde bitleri sütun sütun sağdan sola yayılır ve her sütun, soldaki sütunun eldesini bekler. CSA ise her bit sütununa \emph{bağımsız} birer tam toplayıcı uygular; sütunlar birbirini beklemez:

\begin{center}
\small
\begin{tabular}{l c c c c l}
              & Bit 3 & Bit 2 & Bit 1 & Bit 0 & \\
\midrule
$A$           & 1 & 1 & 0 & 1 & (13) \\
$B$           & 0 & 1 & 1 & 0 & (6) \\
$C$           & 0 & 0 & 1 & 1 & (3) \\
\midrule
Toplam ($T$)  & 1 & 0 & 0 & 0 & $= 8$ \\
Elde ($C$)    & 0\;1 & 1 & 1 & 0 & $= 14$ \quad {\scriptsize (1 konum sola kaydırılmış)} \\
\end{tabular}
\end{center}

Her sütundaki tam toplayıcı, o sütundaki üç biti toplar ve bir toplam biti ile bir elde biti üretir. Elde biti bir sonraki (soldaki) sütuna kaydırılır; ancak \emph{sütunlar arası seri elde yayılımı yoktur}; dört sütundaki dört tam toplayıcı eş zamanlı çalışır.

Sonuçta $T = 1000_2 = 8$ ve $C = 01110_2 = 14$ elde edilir. Dikkat edin: $T + C = 8 + 14 = 22$; doğru sonucu verir! Ancak $T$ ve $C$ henüz \emph{toplanmamıştır}; yalnızca üç sayı iki sayıya \emph{indirgenmiştir}. Gerçek toplama hâlâ yapılmamıştır çünkü gerçek toplama elde yayılımı gerektirir ve bu işlem zaman alır. CSA'nın tüm hız kazancı buradan gelir: elde yayılımından kaçınarak yalnızca sayı adedini azaltır.

Wallace ağacında bu indirgeme işlemi yinelenerek sayı adedi her aşamada $\frac{2}{3}$'e düşürülür. En sonda kalan iki sayı için kaçınılmaz olarak \emph{bir kez} gerçek toplama (CLA ile) yapılır. Böylece $n$ kez seri elde yayılımı yerine, $O(\log n)$ adet eldesiz sıkıştırma $+$ tek bir hızlı toplama gerçekleştirilir.

\begin{figure}[htbp]
\centering
\begin{tikzpicture}[
    pp/.style={draw, rectangle, minimum width=1.6cm, minimum height=0.7cm,
               fill=blokmavi, font=\small, thick, align=center, rounded corners=2pt},
    csa/.style={draw, rectangle, minimum width=3.0cm, minimum height=0.8cm,
                fill=blokturuncu!60, font=\small, thick, align=center, rounded corners=2pt},
    son/.style={draw, rectangle, minimum width=3.0cm, minimum height=0.8cm,
                fill=blokturuncu!30, font=\small, thick, align=center, rounded corners=2pt},
    etiket/.style={font=\small, fill=white, inner sep=2pt},
    okstil/.style={->, >=stealth, thick}
]
% ===== Aşama 0: 4 kısmi çarpım =====
\node[pp] (pp0) at (0, 0) {$PP_0$};
\node[pp] (pp1) at (2.5, 0) {$PP_1$};
\node[pp] (pp2) at (5.0, 0) {$PP_2$};
\node[pp] (pp3) at (7.5, 0) {$PP_3$};

% Aşama 1 etiketi ve ayraç
\draw[dashed, gray!50] (-2.0, -0.7) -- (9.5, -0.7);
\node[font=\footnotesize, gray!80, anchor=east] at (-0.3, -1.0) {Aşama 1};

% ===== Aşama 1: PP0 + PP1 + PP2 → T1, C1  (PP3 bekler) =====
\node[csa] (csa1) at (2.5, -1.8) {3:2 CSA};

% PP0, PP1, PP2 → CSA1
\draw[okstil] (pp0.south) -- ++(0, -0.4) -| ([xshift=-8mm]csa1.north);
\draw[okstil] (pp1.south) -- (csa1.north);
\draw[okstil] (pp2.south) -- ++(0, -0.4) -| ([xshift=8mm]csa1.north);

% T1 ve C1 çıkışları
\draw[okstil] ([xshift=-8mm]csa1.south) -- ++(0, -0.6)
    node[etiket, anchor=east] {$T_1$} coordinate (t1end);
\draw[okstil] ([xshift=8mm]csa1.south) -- ++(0, -0.6)
    node[etiket, anchor=west] {$C_1$} coordinate (c1end);

% Aşama 2 etiketi ve ayraç
\draw[dashed, gray!50] (-2.0, -3.1) -- (9.5, -3.1);
\node[font=\footnotesize, gray!80, anchor=east] at (-0.3, -3.4) {Aşama 2};

% ===== Aşama 2: T1 + C1 + PP3 → T2, C2 =====
\node[csa] (csa2) at (3.75, -4.2) {3:2 CSA};

% PP3 bypass oku — aşağı inip Aşama 2 CSA'ya bağlanır
\draw[okstil, rounded corners=4pt] (pp3.south) -- ++(0, -0.4)
    -- (7.5, -0.4) -- (7.5, -3.6) -| ([xshift=8mm]csa2.north);

% T1, C1 → CSA2
\draw[okstil, rounded corners=3pt] (t1end) -- ++(0, -0.2) -| ([xshift=-8mm]csa2.north);
\draw[okstil, rounded corners=3pt] (c1end) -- ++(0, -0.2) -| (csa2.north);

% T2 ve C2 çıkışları
\draw[okstil] ([xshift=-8mm]csa2.south) -- ++(0, -0.6)
    node[etiket, anchor=east] {$T_2$} coordinate (t2end);
\draw[okstil] ([xshift=8mm]csa2.south) -- ++(0, -0.6)
    node[etiket, anchor=west] {$C_2$} coordinate (c2end);

% Son toplama etiketi ve ayraç
\draw[dashed, gray!50] (-2.0, -5.5) -- (9.5, -5.5);
\node[font=\footnotesize, gray!80, anchor=east] at (-0.3, -5.8) {Son Toplama};

% ===== Son aşama: CLA =====
\node[son] (cla) at (3.75, -6.5) {CLA};

% T2, C2 → CLA
\draw[okstil, rounded corners=3pt] (t2end) -- ++(0, -0.2) -| ([xshift=-8mm]cla.north);
\draw[okstil, rounded corners=3pt] (c2end) -- ++(0, -0.2) -| ([xshift=8mm]cla.north);

% Sonuç
\draw[okstil] (cla.south) -- ++(0, -0.7) node[below, font=\small] {Sonuç ($P = A \times B$)};

% Sağ tarafta açıklama kutuları
\node[font=\scriptsize, gray, text width=2.8cm, anchor=west] at (8.5, -1.8)
    {3 sayı $\to$ 2 sayıya\\indirgenir};
\node[font=\scriptsize, gray, text width=2.8cm, anchor=west] at (8.5, -4.2)
    {3 sayı $\to$ 2 sayıya\\indirgenir};
\node[font=\scriptsize, gray, text width=2.8cm, anchor=west] at (8.5, -6.5)
    {2 sayı $\to$ son\\toplam};

\end{tikzpicture}
\caption[Wallace a\u{g}ac{\i} \c{c}arp{\i}c{\i}: 4 k{\i}smi \c{c}arp{\i}m]{Dört k{\i}smi \c{c}arp{\i}m{\i}n Wallace a\u{g}ac{\i}yla indirgenmesi: iki 3:2 CSA a{\c{s}}amas{\i} dört girdi say{\i}s{\i}n{\i} ikiye d{\"u}{\c{s}}{\"{u}}r{\"u}r; ardından önceden elde üreten bir toplayıcı son toplamı üretir.}
\label{fig:wallace-agaci-4pp}
\end{figure}

\begin{bilginot}[Dadda Ağacı]
Luigi Dadda\index{Dadda, Luigi} 1965'te Wallace ağacına benzer ancak daha az donanım kullanan bir yöntem önermiştir. Dadda ağacı her aşamada yalnızca gerekli minimum sıkıştırmayı yapar ve gereksiz tam toplayıcılardan kaçınır. Uygulamada Dadda ağacı Wallace ağacıyla aynı gecikmeye sahiptir ancak daha az donanım kullanır. Günümüzde her iki yapı da ``Wallace-tarzı ağaç çarpıcı'' olarak anılır.
\end{bilginot}

Şekil~\ref{fig:wallace-agaci-4pp}, dört kısmi çarpımın Wallace ağacıyla iki sayıya nasıl indirildiğini göstermektedir. Her 3:2 sıkıştırıcı aşaması koşut çalışır; gecikme, sıkıştırıcı aşama sayısıyla orantılıdır.

\begin{ornek}[5.9]
\textbf{Soru:} $1101_2 \times 1111_2$ ($13 \times 15$) çarpımını Wallace ağacı yöntemiyle gerçekleştirin.

\textbf{Çözüm:} $A = 1101_2 = 13_{10}$ ve $B = 1111_2 = 15_{10}$ sayıları için $b_0{=}1$, $b_1{=}1$, $b_2{=}1$, $b_3{=}1$ olduğundan dört kısmi çarpımın tamamı sıfırdan farklıdır:

\begin{center}
\begin{tabular}{l l r}
\toprule
\textbf{Kısmi Çarpım} & \textbf{İkili (8 bit)} & \textbf{Onlu} \\
\midrule
$PP_0 = A \cdot b_0$           & $0000\;1101$ & 13  \\
$PP_1 = A \cdot b_1$ (1 sola)  & $0001\;1010$ & 26  \\
$PP_2 = A \cdot b_2$ (2 sola)  & $0011\;0100$ & 52  \\
$PP_3 = A \cdot b_3$ (3 sola)  & $0110\;1000$ & 104 \\
\bottomrule
\end{tabular}
\end{center}

Wallace ağacının iki sıkıştırma aşaması şöyle işler (CSA her bit sütununda bağımsız üç girdi toplar, elde bitleri bir konum sola kaydırılır):

\begin{enumerate}[nosep]
    \item \textit{Aşama~1:} $PP_0$, $PP_1$, $PP_2$ bir 3:2 CSA'ya girer:
        \[ T_1 = 0010\;0011_2 = 35, \quad C_1 = 0011\;1000_2 = 56. \]
        Doğrulama: $T_1 + C_1 = 35 + 56 = 91 = 13 + 26 + 52$. \checkmark
    \item \textit{Aşama~2:} $T_1$, $C_1$, $PP_3$ ikinci 3:2 CSA'ya girer:
        \[ T_2 = 0111\;0011_2 = 115, \quad C_2 = 0101\;0000_2 = 80. \]
        Doğrulama: $T_2 + C_2 = 115 + 80 = 195 = 91 + 104$. \checkmark
    \item \textit{Son toplama (önceden elde üreten toplayıcı):} $T_2 + C_2 = 115 + 80 = 195 = 1100\;0011_2$. \\
        Doğrulama: $13 \times 15 = 195$. \checkmark
\end{enumerate}

Gerçek devrede işlemler bit sütunları üzerinden yürütülür: CSA her sütunu bağımsız toplar, eldeler kaydırılarak bir sonraki sütuna taşınır. Tüm sütunlar \emph{koşut} olarak işlendiğinden herhangi bir seri elde yayılımı oluşmaz.
\end{ornek}

İki koşut çarpıcı yapısının karşılaştırması Tablo~\ref{tab:dizi-vs-wallace}'da özetlenmiştir.

\begin{table}[htbp]
\centering
\caption{Dizi çarpıcı ve Wallace ağacı karşılaştırması ($n$ bitlik çarpma)}
\label{tab:dizi-vs-wallace}
\small
\begin{tabular}{l c c c}
\toprule
& \textbf{Gecikme} & \textbf{Donanım} & \textbf{Düzenlilik} \\
\midrule
Dizi çarpıcı & $O(n)$ & $O(n^2)$ & Yüksek (ızgara) \\
Wallace ağacı & $O(\log n)$ & $O(n^2)$ & Düşük (karmaşık bağlantı) \\
\bottomrule
\end{tabular}
\end{table}

\FloatBarrier
\begin{bilginot}[Çağdaş \.{I}{\c{s}}lemcilerde \c{C}arpma Birimi]
Çağdaş y{\"u}ksek ba{\c{s}}ar{\i}ml{\i} i{\c{s}}lemcilerde \c{c}arpma birimi tipik olarak {\c{s}}u bile{\c{s}}enlerden olu{\c{s}}ur:
\begin{enumerate}[nosep]
    \item \textbf{Radix-4 Booth kodlayıcı:} Çarpandan $n/2$ adet kısmi çarpım üretir.
    \item \textbf{Wallace/Dadda ağacı:} Kısmi çarpımları $O(\log n)$ aşamada 2 sayıya indirger.
    \item \textbf{Önceden elde üreten toplayıcı:} Son iki sayıyı toplar.
\end{enumerate}
Bu mimari sayesinde 64 bitlik bir çarpma işlemi tipik olarak 3--5 saat çevriminde tamamlanır. Çarpma birimi genellikle boru hatlı (\emph{pipelined}) olarak tasarlanır: her çevrimde yeni bir çarpma başlatılabilir, her ne kadar tek bir çarpmanın sonucu birkaç çevrim sonra hazır olsa da.
\end{bilginot}

\subsection{İşlenenlerden Birinin Sıfır Olması}

Çarpma işleminde işlenenlerden birinin sıfır olduğu durumda sonucun da sıfır olacağı önceden bellidir. Kaydır-topla algoritmasında çarpan sıfırsa hiçbir kısmi çarpım üretilmez; çarpılan sıfırsa üretilen her kısmi çarpım zaten sıfırdır. Basit bir donanım bu durumu fark etmeden 32 yinelemenin tamamını çalıştırır; ancak bu gereksiz bir zaman kaybıdır.

Çağdaş işlemcilerde çarpma birimleri genellikle bir \emph{erken sonlandırma}\index{erken sonlandırma}\index{early termination|see{erken sonlandırma}} mekanizması içerir: çarpan yazmacının içeriği her sağa kaydırma sonrasında denetlenir; kalan bitler tümüyle sıfırsa başka kısmi çarpım oluşmayacağı garanti olduğundan döngü erkenden bitirilir. Bu denetim yalnızca bir VEYA kapısı ağacıyla (tüm bitlerin 0 olup olmadığını sınayan) gerçekleştirilebilir.

Erken sonlandırmanın yararı sıfır işlenenle sınırlı değildir. Çarpanın üst bitlerinin büyük bölümü sıfır olan küçük sayılarda da (örneğin $3 = 00000000\ldots011_2$) döngü birkaç adımda tamamlanır. Araştırmalar, gerçek programlarda işlenenlerin önemli bir bölümünün az sayıda anlamlı bite sahip \emph{dar değerler} (\emph{narrow values})\index{dar değer}\index{narrow value|see{dar değer}} olduğunu göstermiştir~\cite{ergin2004}; bu nedenle erken sonlandırma ortalama çarpma süresini belirgin biçimde kısaltır.


% =============================================
\section{Bölme}
\label{sec:bolme}

Çarpma toplama üzerine kuruluysa, bölme\index{bölme}\index{division|see{bölme}} de çıkarma üzerine kuruludur. Bölme, bir sayının içinde diğerinden kaç tane olduğunu bulmaktır. Tek bir ayrıklık vardır: bölen sıfır olamaz; tam sayı aritmetiğinde sıfıra bölmenin geçerli bir sonucu yoktur. Bölme işleminde dört temel bileşen vardır:
%
\begin{itemize}
\item \textbf{Bölünen} (\emph{dividend}): bölünecek sayı,
\item \textbf{Bölen} (\emph{divisor}): bölen sayı,
\item \textbf{Bölüm} (\emph{quotient}): bölme sonucu (tam kısım),
\item \textbf{Kalan} (\emph{remainder}): artık değer.
\end{itemize}

Aralarındaki ilişki şöyle ifade edilir:
\[
\text{Bölünen} = \text{Bölüm} \times \text{Bölen} + \text{Kalan}
\]

Çarpmada $n$ bitlik iki sayının çarpımının $2n$ bite kadar çıkabildiğini görmüştük. Bölme, çarpmanın tersi olduğuna göre, işlenen boyutları da bu durumla uyumlu olmalıdır. $n$ bitlik iki sayıyı birbirine böldüğümüzde bölüm en fazla $n$ bit, kalan ise en fazla $n$ bit uzunluğundadır. Ancak çarpma sonucunda elde edilen $2n$ bitlik bir değeri orijinal çarpanlardan birine bölmek istediğimizde, bölünenin çift genişlikte olması doğal bir gereksinimdir. Bu durumda bölüm de $2n$ bite kadar çıkabilir: örneğin $2n$ bitlik bir bölünen, 1 değerindeki bir bölene bölündüğünde bölüm bölünenin kendisine eşittir. Kalan ise her zaman bölenden küçük olacağından en fazla $n$ bit kalır. RISC-V'te \texttt{div} ve \texttt{rem} buyrukları 32 bitlik iki yazmaç üzerinde çalışır ve 32 bitlik sonuç üretir; bu nedenle $2n$ bitlik bölünen doğrudan desteklenmez. Bununla birlikte, ileride göreceğimiz bölme donanımında kalan yazmacının neden $2n$ bit genişliğinde tutulduğu, bu boyut ilişkisiyle doğrudan bağlantılıdır.

Çarpmada olduğu gibi, ikili tabanda bölme işlemi de onluk tabandaki elde bölmeye benzer; ancak çok daha basittir, çünkü her adımda bölümün ilgili basamağı yalnızca 0 veya 1 olabilir: bölen çıkarılabiliyorsa 1, çıkarılamıyorsa 0 yazılır. Aşağıdaki örnekte $35 \div 5$ işlemi ikili tabanda gösterilmiştir:

\begin{center}
\begin{tikzpicture}[scale=1.0]
  % Sol taraf: Bölünen
  \node[font=\ttfamily, vurgumavi, anchor=east] at (0, 0) {100011};
  \node[font=\small, anchor=west] at (-2.8, 0) {$(35)$};
  \node[font=\scriptsize, gray, anchor=south] at (-1.0, 0.15) {Bölünen};
  % Dikey ayırıcı çizgi
  \draw[thick] (0.15, 0.3) -- (0.15, -1.2);
  % Sağ taraf: Bölen (üstte)
  \node[font=\ttfamily, vurgukirmizi, anchor=west] at (0.35, 0) {101};
  \node[font=\small, anchor=west] at (1.5, 0) {$(5)$};
  \node[font=\scriptsize, gray, anchor=south] at (1.0, 0.15) {Bölen};
  % Bölenin altında yatay çizgi
  \draw[thick] (0.35, -0.25) -- (2.0, -0.25);
  % Bölüm (bölenin altında)
  \node[font=\ttfamily\bfseries, vurguyesil, anchor=west] at (0.35, -0.6) {000111};
  \node[font=\small, anchor=west] at (2.1, -0.6) {$(7)$};
  \node[font=\scriptsize, gray, anchor=north] at (1.2, -0.85) {Bölüm};
  % Sol tarafta: Kalan (bölünenin altında)
  \draw[thick] (-2.0, -0.7) -- (0, -0.7);
  \node[font=\ttfamily, gray, anchor=east] at (0, -1.0) {000};
  \node[font=\scriptsize, gray, anchor=north] at (-1.0, -1.15) {Kalan};
\end{tikzpicture}
\end{center}

Doğrulama: $\textcolor{vurguyesil}{7} \times \textcolor{vurgukirmizi}{5} + \textcolor{gray}{0} = 35 = \textcolor{vurgumavi}{35}$\;\checkmark. İkili bölme de tıpkı çarpma gibi \emph{kaydır ve çıkar} mantığına dayanır: her adımda bölenin çıkarılıp çıkarılamayacağı sınanır, bölüm biti belirlenir ve kalan bir basamak kaydırılır. Bu ilke, bölme donanımının temelini oluşturur.

\subsubsection*{Bölme Donanımına Neden İhtiyaç Var?}

Çarpma bölümünde gördüğümüz gibi, çıkarma yapabilen her işlemci bölmeyi de yazılımla gerçekleştirebilir: bölünen sayıdan böleni arka arkaya çıkararak bölümü bulmak yeterlidir. Ancak bu yaklaşım büyük sayılarda son derece yavaştır:

\begin{table}[htbp]
\centering
\caption[Yazılım ve donanım bölme karşılaştırması]{Aynı bölme işlemi ($35 \div 5$): yalnızca çıkarma ile (yazılım) ve \texttt{div} buyruğu ile (donanım).}
\label{tab:bolme-yazilim-donanim}
\small
\begin{tabular}{p{6.8cm}p{6.8cm}}
\toprule
\textbf{Yalnızca çıkarma ile (M uzantısı yok)} & \textbf{\texttt{div} buyruğu ile (M uzantısı var)} \\
\midrule
\begin{minipage}[t]{6.5cm}
\begin{lstlisting}[style=riscv, xleftmargin=0pt]
    li   t0, 0        # bolum = 0
    li   t1, 35       # bolunen
    li   t2, 5        # bolen
dongu:
    blt  t1, t2, son  # bolunen<bolen? bitir
    sub  t1, t1, t2   # bolunen -= bolen
    addi t0, t0, 1    # bolum++
    j    dongu        # tekrarla
son:  # t0=bolum, t1=kalan
\end{lstlisting}
\end{minipage}
&
\begin{minipage}[t]{6.5cm}
\begin{lstlisting}[style=riscv, xleftmargin=0pt]
    li   t1, 35       # bolunen
    li   t2, 5        # bolen
    div  t0, t1, t2   # bolum = 35/5
    rem  t3, t1, t2   # kalan = 35%5
\end{lstlisting}
\end{minipage}
\\
\midrule
Durağan buyruk sayısı: 7 & Durağan buyruk sayısı: 4 \\
Devingen buyruk sayısı: 32 & Devingen buyruk sayısı: 4 \\
Çevrim sayısı (BBÇ${}\approx 1$): $\sim$32 & Çevrim sayısı: 20--40 \\
\bottomrule
\end{tabular}
\end{table}

Soldaki programda döngü, bölüm değeri kadar yinelenir ($35 \div 5 = 7$ yineleme, toplam 32 buyruk). Bölünen büyüdükçe döngü sayısı doğrusal olarak artar; örneğin $10^9 \div 1$ için yaklaşık $4 \times 10^9$ buyruk gerekir; bu da saniyeler süren bir işlem demektir. Sağdaki programda ise \texttt{div} ve \texttt{rem} buyrukları aynı sonucu yalnızca birkaç düzine saat çevriminde üretir. Çarpmada olduğu gibi, bölme donanımı da \emph{sorumluluğu yazılımdan donanıma aktararak} buyruk sayısını düşürür ve başarımı büyük ölçüde artırır.

Çarpmada olduğu gibi, bölme algoritmaları da verimsizden verimliye doğru geliştirilmiştir.

\subsection{Geri Yüklemeli Bölme}
\label{sec:geri-yuklemeli-bolme}
\index{geri yüklemeli bölme}\index{restoring division|see{geri yüklemeli bölme}}

Temel bölme algoritması \emph{geri yüklemeli bölme} (\emph{restoring division}) olarak adlandırılır: bölen, kalanın üst yarısından çıkarılır; sonuç eksiyse çıkarma geri alınarak (kalan ``geri yüklenerek'') önceki değerine döndürülür. Çarpmada olduğu gibi, donanımı adım adım iyileştiren üç yol geliştireceğiz.

Donanım gerçekleştirmesinde dikkat çekici bir durum gözlemlenir: bölünen, işlem boyunca yazmaçlarda yavaş yavaş \emph{tüketilir} ve yerini iki yeni değere (bölüm ve kalana) bırakır. Her adımda bölen, kalan ara değerinden çıkarılmaya çalışılır; başarılı her çıkarma bölümün bir bitini belirlerken bölünenin o kısmını da dönüştürür. İşlem sonunda başlangıçtaki bölünen ortadan kalkmış, bilgi bölüm ve kalan arasında paylaşılmış olur. Bu dönüşüm yazmaçlarda fiziksel olarak da gözlemlenecektir: bölünen bitleri adım adım erirken, bölüm bitleri sağdan dolmaya başlar.

\subsubsection{1. Yol}

Bölmenin birinci yolunda çarpmada olduğu gibi bölme de en temel haliyle, arka arkaya çıkarma ve kaydırmalarla yapılır. İlk olarak bölünecek sayının tümü kalan yazmacına atılır. Verilerin uzunluğunun $n$ bit olduğu kabul edilirse kalan yazmacının büyüklüğü $2n$'dir. Bölen yazmacı da $2n$ büyüklüğünde alınır.

Bölen yazmacındaki değer kalandan çıkarılır. Sonuç eksi ise sayının içinde bölen yoktur ve işlem geriye alınır (kalan yazmacı eski haline getirilir). Bölüm yazmacı bir bit sola kaydırılır ve en sağ bitine 0 yazılır. Eğer kalandan bölen çıkarıldığında sonuç artı çıkarsa, fark kalan yazmacına yazılır ve bölüm 1 bit sola kaydırılarak en sağ bitine 1 yazılır. Döngü $n+1$ kez devam eder.

\begin{figure}[htbp]
\centering
\begin{tikzpicture}[node distance=0.8cm, scale=0.72, transform shape]
    % Start
    \node[akisbaslangic] (start) {Ba{\c{s}}la};
    % Initialize
    \node[akisislem, below=of start, minimum width=4cm] (init)
        {Kalan $\leftarrow$ B{\"o}l{\"u}nen\\B{\"o}l{\"u}m $\leftarrow$ 0\\Say $\leftarrow$ 0};
    % Subtract
    \node[akisislem, below=0.6cm of init, minimum width=4cm] (sub)
        {Kalan $\leftarrow$ Kalan $-$ B{\"o}len};
    % Decision: remainder >= 0?
    \node[akiskarar, below=0.6cm of sub, minimum width=2.5cm] (check)
        {Kalan $\geq$ 0?};
    % If yes: set quotient bit 1
    \node[akisislem, left=1.6cm of check, minimum width=3.6cm] (setbit)
        {B{\"o}l{\"u}m{\"u} sola kayd{\i}r\\en sa{\u{g}} bit $\leftarrow$ 1};
    % If no: restore remainder
    \node[akisislem, right=1.6cm of check, minimum width=3.6cm] (restore)
        {Kalan{\i} geri y{\"u}kle\\B{\"o}l{\"u}m{\"u} sola kayd{\i}r\\en sa{\u{g}} bit $\leftarrow$ 0};
    % Shift divisor right
    \node[akisislem, below=1.4cm of check, minimum width=4cm] (shiftd)
        {B{\"o}leni 1 bit sa{\u{g}}a kayd{\i}r};
    % Increment
    \node[akisislem, below=0.5cm of shiftd, minimum width=4cm] (inc)
        {Say $\leftarrow$ Say + 1};
    % Check counter
    \node[akiskarar, below=0.5cm of inc, minimum width=2cm] (countchk)
        {Say = 33?};
    % End
    \node[akisbaslangic, below=0.7cm of countchk] (stop) {Dur};
    % Arrows
    \draw[akisok] (start) -- (init);
    \draw[akisok] (init) -- (sub);
    \draw[akisok] (sub) -- (check);
    \draw[akisok] (check) -- node[above, font=\scriptsize] {Evet} (setbit);
    \draw[akisok] (check) -- node[above, font=\scriptsize] {Hay{\i}r} (restore);
    \draw[akisok] (setbit) |- (shiftd);
    \draw[akisok] (restore) |- (shiftd);
    \draw[akisok] (shiftd) -- (inc);
    \draw[akisok] (inc) -- (countchk);
    \draw[akisok] (countchk) -- node[left, font=\scriptsize] {Evet} (stop);
    \coordinate (dongu-sol) at ($(setbit.west)+(-0.8,0)$);
    \draw[akisdonus] (countchk.west)
        -- node[above, font=\scriptsize, pos=0.2] {Hay{\i}r}
        (dongu-sol |- countchk.west)
        |- ($(sub.north)+(0,0.5)$)
        -- (sub.north);
\end{tikzpicture}
\caption[Geri Y{\"u}klemeli B{\"o}lme (1.\ Yol) ak{\i}{\c{s}} {\c{s}}emas{\i}]{Geri Y{\"u}klemeli B{\"o}lme (1.\ Yol) ak{\i}{\c{s}} {\c{s}}emas{\i}.}
\label{fig:bolme-1yol-algo}
\end{figure}

Algoritmanın akış şeması Şekil~\ref{fig:bolme-1yol-algo}'da gösterilmiştir. Her yinelemede yapılan işlemler şunlardır: (1)~kalan yazmacından bölen yazmacını çıkar, (2)~sonuç eksiyse kalanı geri yükle ve bölümün en sağ bitine 0 yaz; artıysa farkı kalana yaz ve bölümün en sağ bitine 1 yaz, (3)~böleni 1 bit sağa, bölümü 1 bit sola kaydır. Bu döngü $n{+}1$ kez tekrarlanır; 4 bitlik veriler için 5, RISC-V (RV32) gibi 32 bitlik bir mimaride ise 33 adım gerekir.

\begin{figure}[htbp]
\centering
\begin{tikzpicture}[scale=0.72, transform shape]

    % ========== D{\"U}{\u{G}}{\"U}MLER ==========

    % ---- B{\"o}len Yazma{\c{c}}{\i} (64 bit, sol {\"u}st, sa{\u{g}}a kayd{\i}r{\i}l{\i}r) ----
    % Kayd{\i}rma bilgisi kutunun {\"u}st{\"u}nde, siyah
    \node[draw, rectangle, fill=blokmavi, thick, font=\footnotesize, align=center,
          minimum width=4.4cm, minimum height=1.3cm] (bolen) at (0, 5)
          {B{\"o}len Yazmac{\i} {\scriptsize(64 bit)}};
    \node[font=\scriptsize, above=2pt of bolen]
          {Sa{\u{g}}a Kayd{\i}r $\longrightarrow$};

    % ---- B{\"o}l{\"u}m Yazma{\c{c}}{\i} (32 bit, sa{\u{g}} {\"u}st, sola kayd{\i}r{\i}l{\i}r) ----
    % Kayd{\i}rma bilgisi kutunun {\"u}st{\"u}nde, siyah
    \node[draw, rectangle, fill=blokmavi, thick, font=\footnotesize, align=center,
          minimum width=3cm, minimum height=1.3cm] (bolum) at (7, 5)
          {B{\"o}l{\"u}m Yazmac{\i} {\scriptsize(32 bit)}};
    \node[font=\scriptsize, above=2pt of bolum]
          {$\longleftarrow$ Sola Kayd{\i}r};

    % ---- AMB (64 bit {\c{C}}{\i}kar{\i}c{\i}, ortada) ----
    \node[draw, trapezium, trapezium angle=70, fill=bloksari, thick,
          shape border rotate=180, font=\footnotesize, align=center,
          minimum width=3cm, minimum height=1.6cm] (amb) at (0, 2.2)
          {64 bit AMB\\({\c{C}}{\i}kar{\i}c{\i})};

    % ---- Kalan Yazma{\c{c}}{\i} (64 bit, altta ortada) ----
    \node[draw, rectangle, fill=blokmavi, thick, font=\footnotesize, align=center,
          minimum width=4.4cm, minimum height=1.3cm] (kalan) at (0, -0.5)
          {Kalan Yazmac{\i} {\scriptsize(64 bit)}};

    % ---- Denetim birimi (sa{\u{g}} tarafta) ----
    \node[draw, rectangle, fill=blokkirmizi-acik, thick, font=\footnotesize, align=center,
          minimum width=2.8cm, minimum height=0.8cm] (denetim) at (7, 2.2)
          {Denetim};

    % ========== VER{\.I} YOLLARI (kal{\i}n d{\"u}z {\c{c}}izgi) ==========

    % B{\"o}len $\rightarrow$ AMB ({\"u}stten giri{\c{s}})
    \draw[okstil] (bolen.south) -- (amb.north);

    % AMB $\rightarrow$ Kalan ({\c{c}}{\i}k{\i}{\c{s}} a{\c{s}}a{\u{g}}{\i})
    \draw[okstil] (amb.south) -- (kalan.north)
          node[midway, right=2pt, font=\scriptsize] {Yaz};

    % Kalan $\rightarrow$ AMB geri besleme (sol taraftan dola{\c{s}}arak)
    \draw[okstil] (kalan.west) -- ++(-1.8, 0) coordinate (fb1)
          -- (fb1 |- amb.west) -- (amb.west);

    % AMB sonu{\c{c}} i{\c{s}}areti $\rightarrow$ Denetim
    \draw[okstil] (amb.east) -- (denetim.west)
          node[midway, above, font=\scriptsize] {Sonu{\c{c}} i{\c{s}}areti};

    % ========== DENET{\.I}M S{\.I}NYALLER{\.I} (k{\i}rm{\i}z{\i} kesikli) ==========

    % Denetim $\rightarrow$ B{\"o}len (sa{\u{g}}a kayd{\i}r sinyali)
    \draw[denetimstil] (denetim.north) -- ++(0, 0.8) coordinate (den-bol-kose)
          -| ([xshift=10mm]bolen.east) coordinate (bol-sag-kose)
          -- (bolen.east);
    \node[font=\scriptsize, sinyalkirmizi, fill=white, inner sep=1pt]
          at ($(den-bol-kose)!0.5!(bol-sag-kose |- den-bol-kose)$) {Kayd{\i}r};

    % Denetim $\rightarrow$ B{\"o}l{\"u}m (sola kayd{\i}r / yaz sinyali)
    \draw[denetimstil] ([xshift=5mm]denetim.north) -- ++(0, 1.3) coordinate (den-blm-kose)
          -| ([xshift=5mm]bolum.east) coordinate (blm-sag-kose)
          -- (bolum.east);
    \node[font=\scriptsize, sinyalkirmizi, fill=white, inner sep=1pt]
          at ($(den-blm-kose)!0.5!(blm-sag-kose |- den-blm-kose)$) {Kayd{\i}r/Yaz};

    % Denetim $\rightarrow$ Kalan (yaz sinyali)
    \draw[denetimstil] (denetim.south) -- ++(0, -0.8) coordinate (den-kal-kose)
          -| ([xshift=10mm]kalan.east) coordinate (kal-sag-kose)
          -- (kalan.east);
    \node[font=\scriptsize, sinyalkirmizi, fill=white, inner sep=1pt]
          at ($(den-kal-kose)!0.5!(kal-sag-kose |- den-kal-kose)$) {Yaz};

    % Denetim $\rightarrow$ AMB ({\c{c}}{\i}kar sinyali)
    \draw[denetimstil] ([yshift=-3mm]denetim.west) -- ([yshift=-3mm]amb.east);
    \node[font=\scriptsize, sinyalkirmizi, fill=white, inner sep=1pt, below=1pt]
          at ($(denetim.west)!0.5!(amb.east)+(0,-3mm)$) {{\c{C}}{\i}kar};

\end{tikzpicture}
\caption{Geri Y{\"u}klemeli B{\"o}lme (1.\ Yol) donan{\i}m{\i}}
\label{fig:bolme-1surum-donanim}
\end{figure}

Buna karşılık gelen donanım tasarımı Şekil~\ref{fig:bolme-1surum-donanim}'de gösterilmiştir. Donanım, 64 bitlik bir bölen yazmacı, 64 bitlik bir kalan yazmacı, 32 bitlik bir bölüm yazmacı ve 64 bitlik bir AMB'den (çıkarıcı) oluşur. Çarpmadaki 1.\ yola benzer biçimde, bölen yazmacının sol yarısına yerleştirilen değer her adımda sağa kaydırılarak kalanla hizalanır. Denetim birimi içindeki bir sayaç yineleme sayısını takip eder ve $n{+}1$ adım tamamlandığında işlemi sonlandırır. 64 bitlik AMB gerekliliği bu yolun temel dezavantajıdır: çarpmada olduğu gibi, 64 bitlik çıkarma 32 bitlik çıkarmaya göre daha yavaş, daha fazla alan kaplayan ve daha çok güç tüketen bir işlemdir.

\begin{ornek}[5.10]
\textbf{Soru:} $5 \div 3 = ?$ \quad $(0101 \div 0011 = ?)$

\textbf{Çözüm:} 1.\ yol algoritmasıyla (4 bit işlenenler, 8 bit bölen ve kalan yazmaçları). Başlangıçta bölünen $\textcolor{vurgumavi}{0101}$ kalan yazmacının alt yarısına, bölen $\textcolor{vurgukirmizi}{0011}$ ise bölen yazmacının üst yarısına yerleştirilir. Tabloda \textcolor{vurgukirmizi}{kırmızı} bitler bölenin her adımda sağa kayan anlamlı konumunu, \textcolor{vurgumavi}{mavi} bitler kalanı, \textcolor{vurguyesil}{yeşil} bitler ise sağdan dolarak oluşan bölümü gösterir; \textcolor{gray}{gri} bitler henüz anlam taşımayan sıfırlardır. Adım adım çözüm Tablo~\ref{tab:bolme1-ornek}'de gösterilmiştir.

\begin{table}[H]
\centering
\caption[Örnek 5.10 -- 1.\ yol geri yüklemeli bölme]{Örnek 5.10 -- 1.\ yol geri yüklemeli bölme algoritması ($0101 \div 0011$).}
\label{tab:bolme1-ornek}
\small
\begin{tabular}{c l c c c}
\toprule
\textbf{Adım} & \textbf{İşlem} & \textbf{\textcolor{vurgukirmizi}{Bölen} (8 bit)} & \textbf{\textcolor{vurgumavi}{Kalan} (8 bit)} & \textbf{\textcolor{vurguyesil}{Bölüm}} \\
\midrule
Başlangıç & --- & \textcolor{vurgukirmizi}{0011}\ \textcolor{gray}{0000} & \textcolor{gray}{0000}\ \textcolor{vurgumavi}{0101} & \textcolor{gray}{0000} \\
\midrule
1 & Kalan $-$ Bölen $\to$ eksi & \textcolor{vurgukirmizi}{0011}\ \textcolor{gray}{0000} & \textcolor{gray}{0000}\ \textcolor{vurgumavi}{0101} & --- \\
  & Geri yükle; Bölüm sola kaydır, bit $\leftarrow$ 0 & \textcolor{gray}{0}\textcolor{vurgukirmizi}{0011}\textcolor{gray}{000} & \textcolor{gray}{0000}\ \textcolor{vurgumavi}{0101} & \textcolor{gray}{000}\textcolor{vurguyesil}{0} \\
\midrule
2 & Kalan $-$ Bölen $\to$ eksi & \textcolor{gray}{0}\textcolor{vurgukirmizi}{0011}\textcolor{gray}{000} & \textcolor{gray}{0000}\ \textcolor{vurgumavi}{0101} & --- \\
  & Geri yükle; Bölüm sola kaydır, bit $\leftarrow$ 0 & \textcolor{gray}{00}\textcolor{vurgukirmizi}{0011}\textcolor{gray}{00} & \textcolor{gray}{0000}\ \textcolor{vurgumavi}{0101} & \textcolor{gray}{00}\textcolor{vurguyesil}{00} \\
\midrule
3 & Kalan $-$ Bölen $\to$ eksi & \textcolor{gray}{00}\textcolor{vurgukirmizi}{0011}\textcolor{gray}{00} & \textcolor{gray}{0000}\ \textcolor{vurgumavi}{0101} & --- \\
  & Geri yükle; Bölüm sola kaydır, bit $\leftarrow$ 0 & \textcolor{gray}{000}\textcolor{vurgukirmizi}{0011}\textcolor{gray}{0} & \textcolor{gray}{0000}\ \textcolor{vurgumavi}{0101} & \textcolor{gray}{0}\textcolor{vurguyesil}{000} \\
\midrule
4 & Kalan $-$ Bölen $\to$ eksi & \textcolor{gray}{000}\textcolor{vurgukirmizi}{0011}\textcolor{gray}{0} & \textcolor{gray}{0000}\ \textcolor{vurgumavi}{0101} & --- \\
  & Geri yükle; Bölüm sola kaydır, bit $\leftarrow$ 0 & \textcolor{gray}{0000}\ \textcolor{vurgukirmizi}{0011} & \textcolor{gray}{0000}\ \textcolor{vurgumavi}{0101} & \textcolor{vurguyesil}{0000} \\
\midrule
5 & Kalan $-$ Bölen $= \textcolor{vurgumavi}{0000\ 0010} \geq 0$ & \textcolor{gray}{0000}\ \textcolor{vurgukirmizi}{0011} & \textcolor{gray}{0000}\ \textcolor{vurgumavi}{0010} & --- \\
  & Bölüm sola kaydır, bit $\leftarrow$ 1 & \textcolor{gray}{0000}\ \textcolor{vurgukirmizi}{0011} & \textcolor{gray}{0000}\ \textcolor{vurgumavi}{0010} & \textcolor{vurguyesil}{0001} \\
\bottomrule
\end{tabular}
\end{table}

Bölüm $= \textcolor{vurguyesil}{(0001)_2 = 1}$, \quad Kalan $= \textcolor{vurgumavi}{(00000010)_2 = 2}$
\end{ornek}

\subsubsection{2. Yol}

1. yolda $2n$ bitlik yazmaçlar sorun yaratır. 2. yolda bölen yazmacı $n$ bit alınmış ve bölenin sağa kaydırılması yerine kalanın sola kaydırılarak işlemin daha verimli yapılması sağlanmıştır; bu yolun akış şeması Şekil~\ref{fig:bolme-2yol-algo}'da, donanım yapısı Şekil~\ref{fig:bolme-2surum-donanim}'de gösterilmiştir.

\begin{figure}[htbp]
\centering
\begin{tikzpicture}[node distance=0.8cm, scale=0.72, transform shape]
    % Start
    \node[akisbaslangic] (start) {Ba{\c{s}}la};
    % Initialize
    \node[akisislem, below=of start, minimum width=4cm] (init)
        {Kalan $\leftarrow$ B{\"o}l{\"u}nen\\B{\"o}l{\"u}m $\leftarrow$ 0\\Say $\leftarrow$ 0};
    % Shift remainder left
    \node[akisislem, below=1cm of init, minimum width=4cm] (shiftrem)
        {Kalan{\i} 1 bit sola kayd{\i}r};
    % Subtract
    \node[akisislem, below=of shiftrem, minimum width=4.4cm] (sub)
        {Kalan[63:32] $\leftarrow$\\Kalan[63:32] $-$ B{\"o}len};
    % Decision: remainder >= 0?
    \node[akiskarar, below=1cm of sub, minimum width=2.5cm] (check)
        {Kalan[63:32]\\$\geq$ 0?};
    % If yes: set quotient bit 1
    \node[akisislem, left=1.6cm of check, minimum width=3.6cm] (setbit)
        {B{\"o}l{\"u}m{\"u} sola kayd{\i}r\\en sa{\u{g}} bit $\leftarrow$ 1};
    % If no: restore remainder
    \node[akisislem, right=1.6cm of check, minimum width=3.6cm] (restore)
        {Kalan[63:32] geri y{\"u}kle\\B{\"o}l{\"u}m{\"u} sola kayd{\i}r\\en sa{\u{g}} bit $\leftarrow$ 0};
    % Increment
    \node[akisislem, below=2.5cm of check, minimum width=4cm] (inc)
        {Say $\leftarrow$ Say + 1};
    % Check counter
    \node[akiskarar, below=of inc, minimum width=2cm] (countchk)
        {Say = 33?};
    % End
    \node[akisbaslangic, below=1cm of countchk] (stop) {Dur};
    % Arrows
    \draw[akisok] (start) -- (init);
    \draw[akisok] (init) -- (shiftrem);
    \draw[akisok] (shiftrem) -- (sub);
    \draw[akisok] (sub) -- (check);
    \draw[akisok] (check) -- node[above, font=\scriptsize] {Evet} (setbit);
    \draw[akisok] (check) -- node[above, font=\scriptsize] {Hay{\i}r} (restore);
    \draw[akisok] (setbit) |- (inc);
    \draw[akisok] (restore) |- (inc);
    \draw[akisok] (inc) -- (countchk);
    \draw[akisok] (countchk) -- node[left, font=\scriptsize] {Evet} (stop);
    \coordinate (dongu-sol) at ($(setbit.west)+(-0.8,0)$);
    \draw[akisdonus] (countchk.west)
        -- node[above, font=\scriptsize, pos=0.2] {Hay{\i}r}
        (dongu-sol |- countchk.west)
        |- ($(shiftrem.north)+(0,0.5)$)
        -- (shiftrem.north);
\end{tikzpicture}
\caption[Geri Y{\"u}klemeli B{\"o}lme (2.\ Yol) ak{\i}{\c{s}} {\c{s}}emas{\i}]{Geri Y{\"u}klemeli B{\"o}lme (2.\ Yol) ak{\i}{\c{s}} {\c{s}}emas{\i}.}
\label{fig:bolme-2yol-algo}
\end{figure}

Algoritmanın akış şeması Şekil~\ref{fig:bolme-2yol-algo}'da gösterilmiştir. Her yinelemede yapılan işlemler şunlardır: (1)~kalanı 1 bit sola kaydır, (2)~kalanın üst yarısından böleni çıkar, (3)~sonuç eksiyse kalanı geri yükle ve bölümün en sağ bitine 0 yaz; artıysa farkı kalanın üst yarısına yaz ve bölümün en sağ bitine 1 yaz. 1.\ yoldan farklı olarak bölen sabit kalır ve kalan sola kaydırılarak hizalama sağlanır. Denetim birimi içindeki sayaç yineleme sayısını takip eder: $n$ bitlik veriler için döngü $n$ kez tekrarlanır; RISC-V (RV32) gibi 32 bitlik bir mimaride 32 adım gerekir.

\begin{figure}[htbp]
\centering
\begin{tikzpicture}[scale=0.72, transform shape]

    % ========== D{\"U}{\u{G}}{\"U}MLER ==========

    % ---- B{\"o}len Yazma{\c{c}}{\i} (32 bit, sol {\"u}st, sabit) ----
    % (Sabit) — kayd{\i}rma yok, sol yar{\i}yla hizal{\i}
    \node[draw, rectangle, fill=blokmavi, thick, font=\footnotesize, align=center,
          minimum width=3.0cm, minimum height=1.3cm] (bolen) at (-1.1, 5)
          {B{\"o}len Yazmac{\i} {\scriptsize(32 bit)}\\{\scriptsize(Sabit)}};

    % ---- B{\"o}l{\"u}m Yazma{\c{c}}{\i} (32 bit, sa{\u{g}} {\"u}st, sola kayd{\i}r{\i}l{\i}r) ----
    % Kayd{\i}rma bilgisi kutunun {\"u}st{\"u}nde, siyah
    \node[draw, rectangle, fill=blokmavi, thick, font=\footnotesize, align=center,
          minimum width=3cm, minimum height=1.3cm] (bolum) at (7, 5)
          {B{\"o}l{\"u}m Yazmac{\i} {\scriptsize(32 bit)}};
    \node[font=\scriptsize, above=2pt of bolum]
          {$\longleftarrow$ Sola Kayd{\i}r};

    % ---- AMB (32 bit {\c{C}}{\i}kar{\i}c{\i}, sol yar{\i}yla hizal{\i}, k{\"u}{\c{c}}{\"u}k) ----
    \node[draw, trapezium, trapezium angle=70, fill=bloksari, thick,
          shape border rotate=180, font=\footnotesize, align=center,
          minimum width=1.5cm, minimum height=1.0cm] (amb) at (-1.1, 2.5)
          {32 bit AMB};

    % ---- Kalan Yazma{\c{c}}{\i} (64 bit, altta ortada, sola kayd{\i}r{\i}l{\i}r) ----
    \node[draw, rectangle, fill=blokmavi, thick, font=\footnotesize, align=center,
          minimum width=4.4cm, minimum height=1.3cm] (kalan) at (0, -0.5)
          {Kalan Yazmac{\i} {\scriptsize(64 bit)}};
    % Ortaya kesikli {\c{c}}izgi (sol/sa{\u{g}} yar{\i} ayr{\i}m{\i})
    \draw[dashed, thick, gray] (kalan.north) -- (kalan.south);
    % Bit etiketleri
    \node[font=\scriptsize, above=1pt of kalan.north west, anchor=south west] {[63:32]};
    \node[font=\scriptsize, above=1pt of kalan.north east, anchor=south east] {[31:0]};
    % Kayd{\i}rma y{\"o}n g{\"o}stergesi
    \node[font=\scriptsize, above=14pt of kalan.north east, anchor=south east]
          {$\longleftarrow$ Sola Kayd{\i}r};

    % ---- Denetim birimi (sa{\u{g}} tarafta) ----
    \node[draw, rectangle, fill=blokkirmizi-acik, thick, font=\footnotesize, align=center,
          minimum width=2.8cm, minimum height=0.8cm] (denetim) at (7, 2.5)
          {Denetim};

    % ========== VER{\.I} YOLLARI (kal{\i}n d{\"u}z {\c{c}}izgi) ==========

    % B{\"o}len $\rightarrow$ AMB ({\"u}stten giri{\c{s}})
    \draw[okstil] (bolen.south) -- (amb.north);

    % AMB $\rightarrow$ Kalan sol yar{\i} [63:32] ({\c{c}}{\i}k{\i}{\c{s}} a{\c{s}}a{\u{g}}{\i}, dik ok)
    \draw[okstil] (amb.south) -- ([xshift=-1.1cm]kalan.north)
          node[midway, left=2pt, font=\scriptsize] {Yaz};

    % Kalan sol yar{\i} [63:32] $\rightarrow$ AMB geri besleme (sol taraftan dola{\c{s}}arak)
    \draw[okstil] (kalan.west) -- ++(-1.5, 0) coordinate (fb2)
          -- (fb2 |- amb.west) -- (amb.west)
          node[pos=0.0, above, font=\scriptsize] {[63:32]};

    % AMB sonu{\c{c}} i{\c{s}}areti $\rightarrow$ Denetim
    \draw[okstil] (amb.east) -- (denetim.west)
          node[midway, above, font=\scriptsize] {Sonu{\c{c}} i{\c{s}}areti};

    % ========== DENET{\.I}M S{\.I}NYALLER{\.I} (k{\i}rm{\i}z{\i} kesikli) ==========

    % Denetim $\rightarrow$ B{\"o}l{\"u}m (sola kayd{\i}r / yaz sinyali)
    \draw[denetimstil] ([xshift=5mm]denetim.north) -- ++(0, 1.3)
          -| ([xshift=5mm]bolum.east)
          -- (bolum.east);
    \node[font=\scriptsize, sinyalkirmizi, fill=white, inner sep=1pt]
          at ([xshift=5mm, yshift=5mm]bolum.east |- denetim.north) {Kayd{\i}r/Yaz};

    % Denetim $\rightarrow$ Kalan (kayd{\i}r / yaz sinyali)
    \draw[denetimstil] (denetim.south) -- ++(0, -0.8)
          -| ([xshift=10mm]kalan.east)
          -- (kalan.east);
    \node[font=\scriptsize, sinyalkirmizi, fill=white, inner sep=1pt, above=1pt]
          at ([xshift=14mm]kalan.east) {Kayd{\i}r/Yaz};

    % Denetim $\rightarrow$ AMB ({\c{c}}{\i}kar sinyali)
    \draw[denetimstil] ([yshift=-3mm]denetim.west) -- ([yshift=-3mm]amb.east);
    \node[font=\scriptsize, sinyalkirmizi, fill=white, inner sep=1pt]
          at ($(denetim.west)!0.5!(amb.east)+(0,-3mm)$) {{\c{C}}{\i}kar};

\end{tikzpicture}
\caption{Geri Y{\"u}klemeli B{\"o}lme (2.\ Yol) donan{\i}m{\i}}
\label{fig:bolme-2surum-donanim}
\end{figure}

Buna karşılık gelen donanım tasarımı Şekil~\ref{fig:bolme-2surum-donanim}'de gösterilmiştir. Donanım, 32 bitlik sabit bir bölen yazmacı, 64 bitlik sola kaydırılan bir kalan yazmacı, 32 bitlik bir bölüm yazmacı ve 32 bitlik bir AMB'den (çıkarıcı) oluşur. AMB artık yalnızca kalanın üst yarısı (32 bit) ile bölen (32 bit) arasında çıkarma yaptığından, 1.\ yolun 64 bitlik AMB'sine göre yarı boyuttadır; bu da her adımı daha hızlı ve donanımı daha küçük kılar. Denetim birimi her iki yazmaca da ``Kaydır/Yaz'' sinyali gönderir: kalan yazmacı her adımın başında 1 bit sola kaydırılır; bölüm yazmacı ise çıkarma sonucuna göre 1 bit sola kaydırılıp en sağ bitine 0 veya 1 yazılır.

\begin{ornek}[5.11]
\textbf{Soru:} $7 \div 2 = ?$ \quad $(0111 \div 0010 = ?)$

\textbf{Çözüm:} 2.\ yol algoritmasıyla (4 bit işlenenler, 8 bit kalan yazmacı, bölen 4 bit sabittir). Başlangıçta bölünen $\textcolor{vurguturuncu}{0111}$ kalan yazmacının alt yarısına yerleştirilir. Her sola kaydırmada bölünenin bir biti üst yarıya taşınarak kalan hesabına katılır (\textcolor{vurgumavi}{mavi}); alt yarıdaki \textcolor{vurguturuncu}{turuncu} bitler adım adım tükenir, sağdan giren \textcolor{gray}{gri} sıfırlar ise boşalan konumları doldurur. Adım adım çözüm Tablo~\ref{tab:bolme2-ornek}'de izlenebilir.

\begin{center}
\captionof{table}[Örnek 5.11 -- 2.\ yol geri yüklemeli bölme]{Örnek 5.11 -- 2.\ yol geri yüklemeli bölme algoritması ($0111 \div 0010$).}
\label{tab:bolme2-ornek}
\footnotesize
\resizebox{\linewidth}{!}{%
\begin{tabular}{c l c c c}
\toprule
\textbf{Adım} & \textbf{İşlem} & \textbf{\textcolor{vurgukirmizi}{Bölen} (4 bit)} & \textbf{\textcolor{vurgumavi}{Kalan} (8 bit)} & \textbf{\textcolor{vurguyesil}{Bölüm} (4 bit)} \\
\midrule
Başlangıç & --- & \textcolor{vurgukirmizi}{0010} & \textcolor{gray}{0000}\ \textcolor{vurguturuncu}{0111} & \textcolor{gray}{0000} \\
\midrule
1 & Kalan'ı 1 bit sola kaydır & \textcolor{vurgukirmizi}{0010} & \textcolor{gray}{000}\textcolor{vurgumavi}{0}\ \textcolor{vurguturuncu}{111}\textcolor{gray}{0} & \textcolor{gray}{0000} \\
  & Kalan[üst] $-$ Bölen $= 0000 - \textcolor{vurgukirmizi}{0010} \to$ eksi & \textcolor{vurgukirmizi}{0010} & \textcolor{gray}{000}\textcolor{vurgumavi}{0}\ \textcolor{vurguturuncu}{111}\textcolor{gray}{0} & --- \\
  & Geri yükle; Bölüm sola kaydır, bit $\leftarrow$ 0 & \textcolor{vurgukirmizi}{0010} & \textcolor{gray}{000}\textcolor{vurgumavi}{0}\ \textcolor{vurguturuncu}{111}\textcolor{gray}{0} & \textcolor{gray}{000}\textcolor{vurguyesil}{0} \\
\midrule
2 & Kalan'ı 1 bit sola kaydır & \textcolor{vurgukirmizi}{0010} & \textcolor{gray}{00}\textcolor{vurgumavi}{01}\ \textcolor{vurguturuncu}{11}\textcolor{gray}{00} & \textcolor{gray}{000}\textcolor{vurguyesil}{0} \\
  & Kalan[üst] $-$ Bölen $= 0001 - \textcolor{vurgukirmizi}{0010} \to$ eksi & \textcolor{vurgukirmizi}{0010} & \textcolor{gray}{00}\textcolor{vurgumavi}{01}\ \textcolor{vurguturuncu}{11}\textcolor{gray}{00} & --- \\
  & Geri yükle; Bölüm sola kaydır, bit $\leftarrow$ 0 & \textcolor{vurgukirmizi}{0010} & \textcolor{gray}{00}\textcolor{vurgumavi}{01}\ \textcolor{vurguturuncu}{11}\textcolor{gray}{00} & \textcolor{gray}{00}\textcolor{vurguyesil}{00} \\
\midrule
3 & Kalan'ı 1 bit sola kaydır & \textcolor{vurgukirmizi}{0010} & \textcolor{gray}{0}\textcolor{vurgumavi}{011}\ \textcolor{vurguturuncu}{1}\textcolor{gray}{000} & \textcolor{gray}{00}\textcolor{vurguyesil}{00} \\
  & Kalan[üst] $-$ Bölen $= \textcolor{vurgumavi}{0011} - \textcolor{vurgukirmizi}{0010} = \textcolor{vurgumavi}{0001} \geq 0$ & \textcolor{vurgukirmizi}{0010} & \textcolor{gray}{0}\textcolor{vurgumavi}{001}\ \textcolor{vurguturuncu}{1}\textcolor{gray}{000} & --- \\
  & Bölüm sola kaydır, bit $\leftarrow$ 1 & \textcolor{vurgukirmizi}{0010} & \textcolor{gray}{0}\textcolor{vurgumavi}{001}\ \textcolor{vurguturuncu}{1}\textcolor{gray}{000} & \textcolor{gray}{0}\textcolor{vurguyesil}{001} \\
\midrule
4 & Kalan'ı 1 bit sola kaydır & \textcolor{vurgukirmizi}{0010} & \textcolor{vurgumavi}{0011}\ \textcolor{gray}{0000} & \textcolor{gray}{0}\textcolor{vurguyesil}{001} \\
  & Kalan[üst] $-$ Bölen $= \textcolor{vurgumavi}{0011} - \textcolor{vurgukirmizi}{0010} = \textcolor{vurgumavi}{0001} \geq 0$ & \textcolor{vurgukirmizi}{0010} & \textcolor{vurgumavi}{0001}\ \textcolor{gray}{0000} & --- \\
  & Bölüm sola kaydır, bit $\leftarrow$ 1 & \textcolor{vurgukirmizi}{0010} & \textcolor{vurgumavi}{0001}\ \textcolor{gray}{0000} & \textcolor{vurguyesil}{0011} \\
\bottomrule
\end{tabular}}%
\end{center}

Bölüm $= \textcolor{vurguyesil}{(0011)_2 = 3}$, \quad Kalan $= \textcolor{vurgumavi}{(0001)_2 = 1}$

Doğrulama: $7 = \textcolor{vurguyesil}{3} \times \textcolor{vurgukirmizi}{2} + \textcolor{vurgumavi}{1}$. \checkmark
\end{ornek}

\subsubsection{3. Yol}

\begin{figure}[htbp]
\centering
\begin{tikzpicture}[scale=0.72, transform shape]

    % ========== D{\"U}{\u{G}}{\"U}MLER ==========

    % ---- B{\"o}len Yazma{\c{c}}{\i} (32 bit, sol {\"u}st, sabit, sol yar{\i}yla hizal{\i}) ----
    \node[draw, rectangle, fill=blokmavi, thick, font=\footnotesize, align=center,
          minimum width=3.0cm, minimum height=1.3cm] (bolen) at (-1.3, 5)
          {B{\"o}len Yazmac{\i} {\scriptsize(32 bit)}\\{\scriptsize(Sabit)}};

    % ---- AMB (32 bit {\c{C}}{\i}kar{\i}c{\i}, sol yar{\i}yla hizal{\i}, k{\"u}{\c{c}}{\"u}k) ----
    \node[draw, trapezium, trapezium angle=70, fill=bloksari, thick,
          shape border rotate=180, font=\footnotesize, align=center,
          minimum width=1.5cm, minimum height=1.0cm] (amb) at (-1.3, 2.5)
          {32 bit AMB};

    % ---- Kalan|B{\"o}l{\"u}m birle{\c{s}}ik yazma{\c{c}} (64 bit, ikiye b{\"o}l{\"u}nm{\"u}{\c{s}}, altta) ----
    % Sol yar{\i}: Kalan [63:32]
    \node[draw, rectangle, fill=blokmavi, thick, font=\footnotesize, align=center,
          minimum width=2.6cm, minimum height=1.3cm] (kalanhi) at (-1.3, -0.5)
          {Kalan};
    % Sa{\u{g}} yar{\i}: B{\"o}l{\"u}m [31:0]
    \node[draw, rectangle, fill=blokmavi!60, thick, font=\footnotesize, align=center,
          minimum width=2.6cm, minimum height=1.3cm] (kalanlo) at (1.3, -0.5)
          {B{\"o}l{\"u}m};
    % D{\i}{\c{s}} {\c{c}}er{\c{c}}eve
    \draw[thick] (kalanhi.north west) rectangle (kalanlo.south east);
    % Bit etiketleri
    \node[font=\scriptsize, above=1pt of kalanhi.north west, anchor=south west] {[63:32]};
    \node[font=\scriptsize, above=1pt of kalanlo.north east, anchor=south east] {[31:0]};
    \node[font=\scriptsize, below=3pt of kalanhi.south east, anchor=north] {64 bit};
    % Kayd{\i}rma y{\"o}n g{\"o}stergesi: kutunun sa{\u{g}} {\"u}st{\"u}nde
    \node[font=\scriptsize, above=14pt of kalanlo.north east, anchor=south east]
          {$\longleftarrow$ Sola Kayd{\i}r};

    % ---- Denetim birimi (sa{\u{g}} tarafta) ----
    \node[draw, rectangle, fill=blokkirmizi-acik, thick, font=\footnotesize, align=center,
          minimum width=2.8cm, minimum height=0.8cm] (denetim) at (7, 2.5)
          {Denetim};

    % ========== VER{\.I} YOLLARI (kal{\i}n d{\"u}z {\c{c}}izgi) ==========

    % B{\"o}len $\rightarrow$ AMB ({\"u}stten giri{\c{s}}, dik ok)
    \draw[okstil] (bolen.south) -- (amb.north);

    % AMB $\rightarrow$ Kalan sol yar{\i} [63:32] ({\c{c}}{\i}k{\i}{\c{s}} a{\c{s}}a{\u{g}}{\i}, dik ok)
    \draw[okstil] (amb.south) -- (kalanhi.north)
          node[midway, left=2pt, font=\scriptsize] {Yaz};

    % Kalan sol yar{\i} [63:32] $\rightarrow$ AMB geri besleme (sol taraftan dola{\c{s}}arak)
    \draw[okstil] (kalanhi.west) -- ++(-1.5, 0) coordinate (fb3)
          -- (fb3 |- amb.west) -- (amb.west)
          node[pos=0.0, above, font=\scriptsize] {[63:32]};

    % AMB sonu{\c{c}} i{\c{s}}areti $\rightarrow$ Denetim
    \draw[okstil] (amb.east) -- (denetim.west)
          node[midway, above, font=\scriptsize] {Sonu{\c{c}} i{\c{s}}areti};

    % ========== DENET{\.I}M S{\.I}NYALLER{\.I} (k{\i}rm{\i}z{\i} kesikli) ==========

    % Denetim $\rightarrow$ Kalan|B{\"o}l{\"u}m (kayd{\i}r / yaz sinyali)
    \draw[denetimstil] ([xshift=5mm]denetim.south) -- ++(0, -1.2)
          -| ([xshift=10mm]kalanlo.east)
          -- (kalanlo.east);
    \node[font=\scriptsize, sinyalkirmizi, fill=white, inner sep=1pt, above=1pt]
          at ([xshift=14mm]kalanlo.east) {Kayd{\i}r/Yaz};

    % Denetim $\rightarrow$ AMB ({\c{c}}{\i}kar sinyali)
    \draw[denetimstil] ([yshift=-3mm]denetim.west) -- ([yshift=-3mm]amb.east);
    \node[font=\scriptsize, sinyalkirmizi, fill=white, inner sep=1pt]
          at ($(denetim.west)!0.5!(amb.east)+(0,-3mm)$) {{\c{C}}{\i}kar};

\end{tikzpicture}
\caption{Geri Y{\"u}klemeli B{\"o}lme (3.\ Yol) donan{\i}m{\i}}
\label{fig:bolme-3surum-donanim}
\end{figure}

\begin{figure}[htbp]
\centering
\begin{tikzpicture}[node distance=0.8cm, scale=0.72, transform shape]
    % Start
    \node[akisbaslangic] (start) {Ba{\c{s}}la};
    % Initialize
    \node[akisislem, below=of start, minimum width=4.4cm] (init)
        {Kalan[63:32] $\leftarrow$ 0\\Kalan[31:0] $\leftarrow$ B{\"o}l{\"u}nen\\Say $\leftarrow$ 0};
    % Shift remainder left
    \node[akisislem, below=1cm of init, minimum width=4.4cm] (shiftrem)
        {Kalan{\i} 1 bit sola kayd{\i}r};
    % Subtract
    \node[akisislem, below=of shiftrem, minimum width=4.4cm] (sub)
        {Kalan[63:32] $\leftarrow$\\Kalan[63:32] $-$ B{\"o}len};
    % Decision: remainder >= 0?
    \node[akiskarar, below=1cm of sub, minimum width=2.5cm] (check)
        {Kalan[63:32]\\$\geq$ 0?};
    % If yes
    \node[akisislem, left=1.6cm of check, minimum width=3.6cm] (setbit)
        {Kalan{\i}n en sa{\u{g}}\\bitine 1 yaz};
    % If no: restore
    \node[akisislem, right=1.6cm of check, minimum width=3.6cm] (restore)
        {Kalan[63:32] geri\\y{\"u}kle, en sa{\u{g}} bit $\leftarrow$ 0};
    % Increment
    \node[akisislem, below=2.5cm of check, minimum width=4.4cm] (inc)
        {Say $\leftarrow$ Say + 1};
    % Check counter
    \node[akiskarar, below=of inc, minimum width=2cm] (countchk)
        {Say = 33?};
    % Final shift
    \node[akisislem, below=1cm of countchk, minimum width=4.4cm] (finalshift)
        {Kalan[63:32] sa{\u{g}}a kayd{\i}r};
    % End
    \node[akisbaslangic, below=of finalshift] (stop) {Dur};
    % Arrows
    \draw[akisok] (start) -- (init);
    \draw[akisok] (init) -- (shiftrem);
    \draw[akisok] (shiftrem) -- (sub);
    \draw[akisok] (sub) -- (check);
    \draw[akisok] (check) -- node[above, font=\scriptsize] {Evet} (setbit);
    \draw[akisok] (check) -- node[above, font=\scriptsize] {Hay{\i}r} (restore);
    \draw[akisok] (setbit) |- (inc);
    \draw[akisok] (restore) |- (inc);
    \draw[akisok] (inc) -- (countchk);
    \draw[akisok] (countchk) -- node[left, font=\scriptsize] {Evet} (finalshift);
    \coordinate (dongu-sol) at ($(setbit.west)+(-0.8,0)$);
    \draw[akisdonus] (countchk.west)
        -- node[above, font=\scriptsize, pos=0.2] {Hay{\i}r}
        (dongu-sol |- countchk.west)
        |- ($(shiftrem.north)+(0,0.5)$)
        -- (shiftrem.north);
    \draw[akisok] (finalshift) -- (stop);
\end{tikzpicture}
\caption[Geri Y{\"u}klemeli B{\"o}lme (3.\ Yol) ak{\i}{\c{s}} {\c{s}}emas{\i}]{Geri Y{\"u}klemeli B{\"o}lme (3.\ Yol) ak{\i}{\c{s}} {\c{s}}emas{\i}.}
\label{fig:bolme-3yol-algo}
\end{figure}

2.\ yolda kalan ve bölüm yazmaçları aynı yönde (sola) kaydırılır. İkisi ayrı yazmaçlarda tutularak 2 ayrı kaydırma yapılması yerine bölüm, kalan yazmacının alt yarısına gömülür ve tek bir kaydırma yeterli olur; tıpkı çarpmadaki 3.\ yolda çarpanın çarpım yazmacına gömülmesi gibi. Böylece ayrı bir bölüm yazmacına gerek kalmaz; donanım yalnızca 32 bitlik sabit bir bölen yazmacı, 64 bitlik birleşik bir kalan/bölüm yazmacı ve 32 bitlik bir AMB'den oluşur (Şekil~\ref{fig:bolme-3surum-donanim}). Algoritmanın akış şeması Şekil~\ref{fig:bolme-3yol-algo}'da gösterilmiştir.

Her yinelemede birleşik yazmaç 1 bit sola kaydırılır, kalanın üst yarısından bölen çıkarılır ve sonuca göre en sağ bite 0 veya 1 yazılır. Bu en sağ bit artık bölüm yazmacının değil, birleşik yazmacın en sağ bitidir; böylece bölüm bitleri her adımda otomatik olarak alt yarıda birikir. Denetim birimi içindeki sayaç $n{+}1$ adımı takip eder; RISC-V (RV32) gibi 32 bitlik bir mimaride 33 adım gerekir. Döngü bittikten sonra, kalanın üst yarısındaki değer fazladan bir kaydırma ile düzeltilir.

\begin{ornek}[5.12]
\textbf{Soru:} $6 \div 2 = ?$ \quad $(0110 \div 0010 = ?)$

\textbf{Çözüm:} 3.\ yol algoritmasıyla (4 bit işlenenler, 8 bitlik birleşik kalan/bölüm yazmacı, bölen 4 bit sabittir). Başlangıçta bölünen birleşik yazmaca yerleştirilir: bölünen $n$ bit ise alt yarıya [31:0] yazılıp üst yarı [63:32] sıfırlanır; $2n$ bit ise tüm yazmaca yazılır. Bu örnekte bölünen $\textcolor{vurguturuncu}{0110}$ yalnızca 4 bit olduğundan alt yarıya yerleştirilmiştir. Her sola kaydırmada bölünenin bir biti üst yarıya taşınarak kalan hesabına katılır (\textcolor{vurgumavi}{mavi}); alt yarıda \textcolor{vurguturuncu}{turuncu} bölünen bitleri tükenirken, sağdan giren \textcolor{vurguyesil}{yeşil} bölüm bitleri dolar. Son adımda üst yarı fazladan bir kaydırmayla düzeltilir. Adım adım çözüm Tablo~\ref{tab:bolme3-ornek}'de gösterilmiştir.

\begin{center}
\captionof{table}[Örnek 5.12 -- 3.\ yol geri yüklemeli bölme]{Örnek 5.12 -- 3.\ yol geri yüklemeli bölme algoritması ($0110 \div 0010$).}
\label{tab:bolme3-ornek}
\footnotesize
\resizebox{\linewidth}{!}{%
\begin{tabular}{c l c c c}
\toprule
\textbf{Adım} & \textbf{İşlem} & \textbf{\textcolor{vurgukirmizi}{Bölen} (4 bit)} & \textbf{\textcolor{vurgumavi}{Kalan[üst]} (4 bit)} & \textbf{Kalan[alt]/\textcolor{vurguyesil}{Bölüm} (4 bit)} \\
\midrule
Başlangıç & --- & \textcolor{vurgukirmizi}{0010} & \textcolor{gray}{0000} & \textcolor{vurguturuncu}{0110} \\
\midrule
1 & Sola kaydır & \textcolor{vurgukirmizi}{0010} & \textcolor{gray}{0000} & \textcolor{vurguturuncu}{1100} \\
  & Kalan[üst] $-$ Bölen $\to$ eksi; geri yükle, bit $\leftarrow$ 0 & \textcolor{vurgukirmizi}{0010} & \textcolor{gray}{0000} & \textcolor{vurguturuncu}{110}\textcolor{vurguyesil}{0} \\
\midrule
2 & Sola kaydır & \textcolor{vurgukirmizi}{0010} & \textcolor{gray}{000}\textcolor{vurgumavi}{1} & \textcolor{vurguturuncu}{10}\textcolor{vurguyesil}{00} \\
  & Kalan[üst] $-$ Bölen $\to$ eksi; geri yükle, bit $\leftarrow$ 0 & \textcolor{vurgukirmizi}{0010} & \textcolor{gray}{000}\textcolor{vurgumavi}{1} & \textcolor{vurguturuncu}{1}\textcolor{vurguyesil}{000} \\
\midrule
3 & Sola kaydır & \textcolor{vurgukirmizi}{0010} & \textcolor{gray}{00}\textcolor{vurgumavi}{11} & \textcolor{vurguyesil}{0000} \\
  & $\textcolor{vurgumavi}{0011} - \textcolor{vurgukirmizi}{0010} = \textcolor{vurgumavi}{0001} \geq 0$; bit $\leftarrow$ 1 & \textcolor{vurgukirmizi}{0010} & \textcolor{gray}{00}\textcolor{vurgumavi}{01} & \textcolor{vurguyesil}{0001} \\
\midrule
4 & Sola kaydır & \textcolor{vurgukirmizi}{0010} & \textcolor{gray}{0}\textcolor{vurgumavi}{010} & \textcolor{vurguyesil}{0010} \\
  & $\textcolor{vurgumavi}{0010} - \textcolor{vurgukirmizi}{0010} = \textcolor{vurgumavi}{0000} \geq 0$; bit $\leftarrow$ 1 & \textcolor{vurgukirmizi}{0010} & \textcolor{gray}{0}\textcolor{vurgumavi}{000} & \textcolor{vurguyesil}{0011} \\
\midrule
5 & Sola kaydır & \textcolor{vurgukirmizi}{0010} & \textcolor{vurgumavi}{0000} & \textcolor{vurguyesil}{0110} \\
  & Kalan[üst] $-$ Bölen $\to$ eksi; geri yükle, bit $\leftarrow$ 0 & \textcolor{vurgukirmizi}{0010} & \textcolor{vurgumavi}{0000} & \textcolor{vurguyesil}{0110} \\
\midrule
Son & Kalan[üst] sağa kaydır (düzeltme) & \textcolor{vurgukirmizi}{0010} & \textcolor{vurgumavi}{0000} & \textcolor{vurguyesil}{0011} \\
\bottomrule
\end{tabular}}%
\end{center}

Bölüm $= \textcolor{vurguyesil}{(0011)_2 = 3}$, \quad Kalan $= \textcolor{vurgumavi}{(0000)_2 = 0}$. Doğrulama: $\textcolor{vurguyesil}{3} \times \textcolor{vurgukirmizi}{2} + \textcolor{vurgumavi}{0} = 6$. \checkmark
\end{ornek}

3.\ yol, geri yüklemeli bölme için en yalın donanım tasarımını sunar: yalnızca 1 adet 32 bitlik bölen yazmacı, 1 adet 64 bitlik birleşik kalan/bölüm yazmacı ve 32 bitlik bir AMB yeterlidir. Bu yapı RISC-V'e (RV32I) doğrudan uygulanabilir.

\begin{bilginot}[İşaretli Bölme İşlemi]
Yukarıdaki algoritmalar işaretsiz sayılar için tasarlanmıştır. İşaretli sayılarla bölme yapmak için en yaygın yöntem şudur: bölünen ve bölenin işaretleri saklanır, her iki sayı da mutlak değerine dönüştürülür ve işaretsiz bölme yapılır. İşlem sonunda bölümün ve kalanın işaretleri ayrıca belirlenir:
\begin{itemize}[nosep]
  \item \emph{Bölümün işareti:} Bölünen ile bölenin işaretleri farklıysa bölüm eksidir (tıpkı çarpmada olduğu gibi).
  \item \emph{Kalanın işareti:} Kalan her zaman bölünenle aynı işareti taşır. Örneğin $(-7) \div 2 = -3$ kalan $-1$'dir; çünkü $-7 = (-3) \times 2 + (-1)$.
\end{itemize}
RISC-V bu kuralı doğrudan destekler: \texttt{div}/\texttt{rem} buyrukları işaretli, \texttt{divu}/\texttt{remu} buyrukları ise işaretsiz bölme yapar. İşaretli sürümler yukarıdaki kalan işareti kuralına uyar.
\end{bilginot}

\subsection{Geri Yüklemesiz Bölme}
\label{sec:geri-yuklemesiz-bolme}
\index{geri yüklemesiz bölme}\index{non-restoring division|see{geri yüklemesiz bölme}}

Geri yüklemeli bölmede, çıkarma sonucu eksiyse kalanı eski değerine döndürmek (geri yüklemek) için ek bir toplama yapılır. Bu geri yükleme adımı zaman kaybına neden olur. \emph{Geri yüklemesiz bölme} (\emph{non-restoring division}) bu sorunu ortadan kaldırır.

Temel gözlem şudur: Geri yüklemeli bölmede, eksi bir kalanı geri yükledikten sonra bir sonraki adımda yine bölen çıkarılır. Dolayısıyla sırasıyla yapılan iki işlem ``$+$Bölen'' (geri yükleme) ve ``$-$Bölen'' (yeni çıkarma) birleştirilirse net etki ``$+$Bölen $-$ Bölen $= 0$'' yerine doğrudan bir sonraki adıma geçilmesidir. Daha doğrusu, eksi kalan elde edildiğinde geri yüklemek yerine bir sonraki adımda \emph{bölen eklenir} (çıkarmak yerine).

Algoritmanın kuralları:
\begin{itemize}[nosep]
    \item Kalanı sola kaydır.
    \item Kalan $\geq 0$ ise: bölenden çıkar, bölüm bitine 1 yaz.
    \item Kalan $< 0$ ise: bölen ekle, bölüm bitine 0 yaz.
\end{itemize}

Geri yüklemesiz bölme, her adımda tam olarak bir aritmetik işlem (toplama veya çıkarma) ve bir kaydırma yapar; koşullu geri yükleme yoktur. Bu nedenle geri yüklemeli bölmeden daha hızlıdır ve donanım denetimi daha basittir. İşlem sonunda kalanın eksi olması mümkündür; bu durumda son adımda bölen bir kez eklenerek kalan artı yapılır. Sonuç, geri yüklemeli bölmeyle matematiksel olarak aynıdır.

Donanım açısından geri yüklemesiz bölme, geri yüklemeli bölmenin 3.\ yol donanımını (Şekil~\ref{fig:bolme-3surum-donanim}) kullanır. Tek fark, AMB'nin yalnızca çıkarma değil toplama da yapabilmesidir; tıpkı Booth çarpma algoritmasında AMB'nin hem toplama hem çıkarma yapması gibi. Denetim birimi kalanın işaret bitine bakarak AMB'ye toplama veya çıkarma sinyali gönderir; geri yükleme adımı ortadan kalktığından her yineleme sabit sürede tamamlanır.

\begin{ornek}[5.13]
\textbf{Soru:} $7 \div 3 = ?$ \quad $(0111 \div 0011 = ?)$ işlemini geri yüklemesiz bölme ile çözün (4 bit işlenenler, 8 bit kalan yazmacı).

\textbf{Çözüm:} 3.\ yol donanımıyla (8 bitlik birleşik kalan/bölüm yazmacı, bölen 4 bit sabittir). Başlangıçta bölünen $\textcolor{vurguturuncu}{0111}$ birleşik yazmacın alt yarısına yerleştirilir, üst yarı sıfırlanır. Geri yüklemeli bölmeden farklı olarak, burada kalanın işaret bitine bakılır: kalan artıysa bölen \emph{çıkarılır}, eksiyse bölen \emph{eklenir}. Tabloda \textcolor{vurgukirmizi}{kırmızı} bölen, \textcolor{vurgumavi}{mavi} kalan (üst yarı), \textcolor{vurguturuncu}{turuncu} tükenen bölünen ve \textcolor{vurguyesil}{yeşil} oluşan bölüm bitleridir.

\begin{center}
\footnotesize
\resizebox{\linewidth}{!}{%
\begin{tabular}{c l c c c}
\toprule
\textbf{Adım} & \textbf{İşlem} & \textbf{\textcolor{vurgukirmizi}{Bölen}} & \textbf{\textcolor{vurgumavi}{Kalan[üst]}} & \textbf{Kalan[alt]/\textcolor{vurguyesil}{Bölüm}} \\
\midrule
Başlangıç & --- & \textcolor{vurgukirmizi}{0011} & \textcolor{gray}{0000} & \textcolor{vurguturuncu}{0111} \\
\midrule
1 & Sola kaydır & \textcolor{vurgukirmizi}{0011} & \textcolor{gray}{000}\textcolor{vurgumavi}{0} & \textcolor{vurguturuncu}{1110} \\
  & $\geq 0$: Çıkar $\textcolor{vurgumavi}{0000} - \textcolor{vurgukirmizi}{0011} = \textcolor{vurgumavi}{1101}$ ($<0$) & \textcolor{vurgukirmizi}{0011} & \textcolor{vurgumavi}{1101} & \textcolor{vurguturuncu}{111}\textcolor{vurguyesil}{0} \\
\midrule
2 & Sola kaydır & \textcolor{vurgukirmizi}{0011} & \textcolor{vurgumavi}{1011} & \textcolor{vurguturuncu}{11}\textcolor{vurguyesil}{00} \\
  & $< 0$: Topla $\textcolor{vurgumavi}{1011} + \textcolor{vurgukirmizi}{0011} = \textcolor{vurgumavi}{1110}$ ($<0$) & \textcolor{vurgukirmizi}{0011} & \textcolor{vurgumavi}{1110} & \textcolor{vurguturuncu}{1}\textcolor{vurguyesil}{000} \\
\midrule
3 & Sola kaydır & \textcolor{vurgukirmizi}{0011} & \textcolor{vurgumavi}{1101} & \textcolor{vurguyesil}{0000} \\
  & $< 0$: Topla $\textcolor{vurgumavi}{1101} + \textcolor{vurgukirmizi}{0011} = \textcolor{vurgumavi}{0000}$ ($\geq 0$) & \textcolor{vurgukirmizi}{0011} & \textcolor{vurgumavi}{0000} & \textcolor{vurguyesil}{0001} \\
\midrule
4 & Sola kaydır & \textcolor{vurgukirmizi}{0011} & \textcolor{vurgumavi}{0001} & \textcolor{vurguyesil}{0010} \\
  & $\geq 0$: Çıkar $\textcolor{vurgumavi}{0001} - \textcolor{vurgukirmizi}{0011} = \textcolor{vurgumavi}{1110}$ ($<0$) & \textcolor{vurgukirmizi}{0011} & \textcolor{vurgumavi}{1110} & \textcolor{vurguyesil}{0010} \\
\midrule
Son & $< 0$: Düzelt $\textcolor{vurgumavi}{1110} + \textcolor{vurgukirmizi}{0011} = \textcolor{vurgumavi}{0001}$ & \textcolor{vurgukirmizi}{0011} & \textcolor{vurgumavi}{0001} & \textcolor{vurguyesil}{0010} \\
\bottomrule
\end{tabular}}%
\end{center}

Bölüm $= \textcolor{vurguyesil}{(0010)_2 = 2}$, \quad Kalan $= \textcolor{vurgumavi}{(0001)_2 = 1}$. Doğrulama: $7 = \textcolor{vurguyesil}{2} \times \textcolor{vurgukirmizi}{3} + \textcolor{vurgumavi}{1}$.\;\checkmark

Son adımdaki düzeltmeye dikkat: geri yüklemesiz bölmede işlem sonunda kalan eksi çıkabilir; bu durumda bölen bir kez eklenerek kalan artı yapılır. Geri yüklemeli bölmede ise her adımda koşullu geri yükleme yapıldığından kalan hiçbir zaman eksi kalmaz.
\end{ornek}

\subsection{SRT Bölme}
\label{sec:srt-bolme}
\index{SRT bölme}\index{SRT division|see{SRT bölme}}

Yukarıdaki bölme algoritmalarının tümü her adımda yalnızca tek bir bölüm biti üretir. $n$ bitlik bir bölme işlemi için $n$ adım gerekir ve bu durum bölmeyi çarpmadan çok daha yavaş kılar.

\emph{SRT bölme}\index{Sweeney--Robertson--Tocher} (\emph{Sweeney--Robertson--Tocher}) algoritması bu sorunu çözmek için her adımda birden fazla bölüm biti üretir. Radix-4 SRT bölmede her adımda 2 bölüm biti hesaplanır ve adım sayısı $n/2$'ye düşer.

SRT algoritmasını anlamak için önce önceki algoritmalarla ortak noktalarına bakalım. Tüm bölme algoritmalarında her adımda bir \emph{kalan ara değeri} hesaplanır; bu ara değere \emph{kısmi kalan} (\emph{partial remainder}) denir. Önceki bölümlerde bu ara değere kısaca ``kalan'' demiştik; SRT bağlamında ``kısmi kalan'' terimi tercih edilir çünkü bu değer her adımda değişir ve son adımda gerçek kalana dönüşür. Tüm yöntemlerde temel yineleme aynıdır:
\[
P_{i+1} = r \cdot P_i - q_i \cdot d
\]
Burada $P_i$ kısmi kalan, $d$ bölen, $q_i$ o adımın bölüm rakamı ve $r$ tabandır (\emph{radix}). Geri yüklemeli ve geri yüklemesiz bölmede $r = 2$ ve $q_i \in \{0, 1\}$'dir; SRT bölmenin farkı \emph{fazlalıklı rakam kümesi} kullanmasıdır.

SRT algoritmasının temel farkları:
\begin{itemize}[nosep]
    \item \textbf{Fazlalıklı rakam kümesi:} Bölüm rakamları yalnızca $\{0, 1\}$ değil, daha geniş bir kümeden seçilir. Radix-2 SRT'de $q_i \in \{-1, 0, +1\}$, radix-4 SRT'de $q_i \in \{-2, -1, 0, +1, +2\}$. Bu fazlalıklı (\emph{redundant}) gösterim, bölüm rakamı tahmininde hata payına izin verir: seçim hatalı olsa bile bir sonraki adımda düzeltilebilir.
    \item \textbf{Arama tablosu:} Bölüm rakamı $q_i$, kısmi kalanın ve bölenin en anlamlı birkaç bitine bakılarak bir \emph{arama tablosu} (\emph{lookup table}) aracılığıyla belirlenir. Bu, koşullu çıkarma sonucunu beklemeye gerek kalmadan hızlı karar verilmesini sağlar.
    \item \textbf{Son dönüşüm:} Fazlalıklı rakamlarla ifade edilen bölüm, işlem sonunda standart ikili gösterime dönüştürülür.
\end{itemize}

Radix-4 SRT, çağdaş işlemcilerde tamsayı ve kayan nokta bölme birimlerinin standart yöntemidir. Daha yüksek tabanlı sürümler (radix-8, radix-16) daha büyük arama tablolarıyla adım sayısını daha da azaltabilir.

SRT bölme donanımının temel yapısı Şekil~\ref{fig:srt-donanim}'de gösterilmiştir.

\begin{figure}[htbp]
\centering
\begin{tikzpicture}[scale=0.82, transform shape,
    blok/.style={draw, rectangle, fill=blokmavi, thick, font=\footnotesize,
                 align=center, minimum height=1.2cm},
    tablo/.style={draw, rectangle, fill=bloksari, thick, font=\footnotesize,
                  align=center, minimum width=3.4cm, minimum height=1.2cm},
    okstil2/.style={-{Stealth[length=2.5mm]}, thick}]

    % K{\i}smi kalan yazma{\c{c}}{\i} (64 bit, birle{\c{s}}ik)
    \node[blok, minimum width=5.4cm] (kalan) at (0, 5)
          {Kalan Yazmac{\i} {\scriptsize(64 bit)}};
    % Kesikli {\c{c}}izgi (sol/sa{\u{g}} yar{\i})
    \draw[dashed, thick, gray] (kalan.north) -- (kalan.south);
    \node[font=\scriptsize, above=1pt of kalan.north west, anchor=south west] {[63:32]};
    \node[font=\scriptsize, above=1pt of kalan.north east, anchor=south east] {[31:0]};

    % B{\"o}len yazma{\c{c}}{\i} (32 bit)
    \node[blok, minimum width=3.4cm] (bolen) at (7, 5)
          {B{\"o}len Yazmac{\i} {\scriptsize(32 bit)}};

    % Arama tablosu
    \node[tablo] (arama) at (7, 2.8)
          {Arama Tablosu\\{\scriptsize(B{\"o}l{\"u}m rakam{\i} se{\c{c}}imi)}};

    % Toplama / {\c{C}}{\i}karma (32 bit, sol yar{\i}yla hizal{\i}, dar)
    \node[draw, trapezium, trapezium angle=70, fill=bloksari, thick,
          shape border rotate=180, font=\footnotesize, align=center,
          minimum width=1.2cm, minimum height=0.9cm] (toplama) at (-1.35, 2.8)
          {32 bit\\Toplama/{\c{C}}{\i}karma};

    % Sola Kayd{\i}r
    \node[blok, minimum width=2.6cm] (kaydir) at (-1.35, 0.8)
          {Sola Kayd{\i}r{\i}c{\i}};

    % B{\"o}l{\"u}m yazma{\c{c}}{\i} (32 bit, fazlal{\i}kl{\i})
    \node[blok, minimum width=3.4cm] (bolum) at (7, 0.8)
          {B{\"o}l{\"u}m Yazmac{\i} {\scriptsize(32 bit)}\\{\scriptsize(fazlal{\i}kl{\i} g{\"o}sterim)}};

    % ========== BA{\u{G}}LANTILAR ==========

    % Kalan [63:32] $\rightarrow$ Toplama/{\c{C}}{\i}karma
    \draw[okstil2] ([xshift=-1.35cm]kalan.south) -- (toplama.north)
          node[midway, left, font=\scriptsize] {[63:32]};

    % B{\"o}len $\rightarrow$ Arama tablosu
    \draw[okstil2] (bolen.south) -- (arama.north)
          node[midway, right, font=\scriptsize] {{\"u}st bitler};

    % Arama tablosu $\rightarrow$ B{\"o}l{\"u}m
    \draw[okstil2] (arama.south) -- (bolum.north)
          node[midway, right, font=\scriptsize] {$q_i$};

    % Arama tablosu $\rightarrow$ Toplama/{\c{C}}{\i}karma
    \draw[okstil2] (arama.west) -- (toplama.east)
          node[midway, above, font=\scriptsize] {$q_i \cdot d$};

    % Toplama $\rightarrow$ Kayd{\i}r
    \draw[okstil2] (toplama.south) -- (kaydir.north);

    % Geri besleme: kayd{\i}r $\rightarrow$ kalan [63:32]
    \draw[okstil2] (kaydir.west) -- ++(-1.2, 0) coordinate (fb-alt)
          -- ++(0, 4.2) coordinate (fb-ust) -- (kalan.west);
    \node[font=\scriptsize, above=2pt, rotate=90, anchor=south] at ($(fb-alt)!0.5!(fb-ust)$)
          {$2P_i - q_i \cdot d$};

    % Kalan {\"u}st bitleri $\rightarrow$ arama tablosu (kalan{\i}n sa{\u{g}} kenar{\i}ndan)
    \draw[okstil2] (kalan.east) -- ++(1.2, 0) |- ([yshift=2mm]arama.west)
          node[pos=0.0, above, font=\scriptsize] {{\"u}st bitler};

\end{tikzpicture}
\caption[SRT b{\"o}lme donan{\i}m{\i}]{SRT bölme donanımının blok şeması. Kalan yazmacı 64 bitlik birleşik yapıdadır; toplama/çıkarma yalnızca üst yarı [63:32] üzerinde 32 bit olarak yapılır.}
\label{fig:srt-donanim}
\end{figure}

Bu algoritmadaki arama tablosundaki küçük bir hata, bilgisayar tarihinin en ünlü donanım hatalarından birine yol açmıştır (bkz.\ Bölüm~\ref{sec:fdiv-hatasi}).

Radix-4 SRT'de her adımda 2 bölüm biti üretilir ve $q_i \in \{-2, -1, 0, +1, +2\}$ kümesinden seçilir. 32 bitlik bir bölme 16 adımda tamamlanır. Fazlalıklı küme sayesinde bölüm rakamı tahmininin kesin olması gerekmez; hatalı bir seçim bir sonraki adımda eksi bir rakamla düzeltilir. Bu esneklik, arama tablosunun kısmi kalanın yalnızca birkaç üst bitine bakarak hızlı karar vermesini mümkün kılar.

\begin{bilginot}[Bölmeyi Çarpmaya Dönüştürmek]
SRT bölme yinelemeli çıkarma üzerine kurulu olduğundan çarpmadan her zaman daha yavaştır. Bu nedenle bazı çağdaş tasarımlar bölmeyi çarpmaya dönüştüren yöntemler kullanır:

\textbf{Newton--Raphson yöntemi:}\index{Newton--Raphson bölme} $1/d$ değerini yinelemeli olarak yakınsar. $x_{i+1} = x_i \times (2 - d \times x_i)$ formülüyle her yinelemede duyarlık \emph{ikiye katlanır} (kuadratik yakınsama). Başlangıç tahmini bir arama tablosundan alınır ve 32 bit duyarlık için yalnızca 4--5 yineleme yeterlidir. Her yineleme iki çarpma ve bir çıkarma içerir.

\textbf{Goldschmidt yöntemi:}\index{Goldschmidt bölme} Benzer fikre dayanır ancak pay ve paydayı eş zamanlı olarak ölçekleyerek böleni 1'e yakınsatır. Goldschmidt'in üstünlüğü, her yinelemede çarpmaların \emph{bağımsız} olması ve koşut çarpıcılarla hızlandırılabilmesidir.

Bu yöntemler özellikle kayan nokta bölme birimlerinde ve GPU'larda tercih edilir.
\end{bilginot}

\subsection{Bölme Yöntemlerinin Karşılaştırması}

Geri yüklemeli bölmenin üç yolu, donanımı adım adım yalınlaştırarak aynı sonuca ulaşır: 1.\ yolda 64 bitlik AMB ve sağa kayan bölen kullanılırken, 2.\ yolda AMB 32 bite iner ve kalan sola kaydırılır; 3.\ yolda ise bölüm yazmacı ortadan kalkarak kalan yazmacının alt yarısına gömülür. Geri yüklemesiz bölme bu temelin üzerine koşullu geri yüklemeyi kaldırarak her adımı sabit süreli yapar. SRT bölme ise fazlalıklı rakam kümesi ve arama tablosuyla her adımda birden fazla bit üreterek adım sayısını yarıya indirir. Tüm bu yöntemlerin karşılaştırması Tablo~\ref{tab:bolme-yol-karsilastirma}'da özetlenmiştir.

\begin{table}[htbp]
\centering
\caption[Bölme yöntemlerinin karşılaştırması]{Bölme yöntemlerinin gecikme, donanım karmaşıklığı ve temel özellik açısından karşılaştırması ($n$ bit genişlik).}
\label{tab:bolme-yol-karsilastirma}
\footnotesize
\begin{tabularx}{\textwidth}{l c c l c X}
\toprule
\textbf{Yöntem} & \textbf{Adım} & \textbf{AMB} & \textbf{Yazmaç} & \textbf{Alan} & \textbf{Kayıp} \\
\midrule
Geri yükl. (1.\ yol) & $n{+}1$ & 64 bit & $2{\times}64 + 32$ bit & Büyük & 64 bit AMB yavaş \\
Geri yükl. (2.\ yol) & $n$ & 32 bit & $64 + 2{\times}32$ bit & Orta & --- \\
Geri yükl. (3.\ yol) & $n{+}1$ & 32 bit & $64 + 32$ bit & Küçük & --- \\
Geri yüklemesiz & $n{+}1$ & 32 bit & $64 + 32$ bit & Küçük & Son düzeltme adımı \\
SRT (radix-2) & $n$ & 32 bit & $64 + 32$ bit + tablo & Orta & Arama tablosu alanı \\
SRT (radix-4) & $n/2$ & 32 bit & $64 + 32$ bit + tablo & Orta & Büyük arama tablosu \\
\midrule
Newton--Raphson & 4--5 yineleme & 2 çarpıcı & Arama tablosu (ilk tahmin) & Büyük & Kuadratik yakınsama \\
Goldschmidt & 4--5 yineleme & 2 koşut çarpıcı & Arama tablosu (ilk tahmin) & Büyük & Koşut çarpma avantajı \\
\bottomrule
\end{tabularx}
\end{table}

Tüm yinelemeli yöntemlerde bölme çarpmadan yavaştır: $n$ bitlik bir çarpma 32 saat vuruşunda tamamlanırken, bölme en az $n$ vuruş gerektirir ve her vuruşta koşullu karar (çıkarma sonucunun işareti) boru hattını zorlaştırır. Bu nedenle derleyiciler mümkün olduğunca bölmeyi kaydırma veya çarpma ile değiştirmeye çalışır (örneğin $x / 4$ yerine $x \gg 2$).


\subsection{RISC-V Bölme Buyrukları}
\label{sec:riscv-bolme-buyruklar}

Çarpma bölümünde olduğu gibi, bölme de RISC-V'in temel buyruk kümesi RV32I'da yer almaz; M uzantısıyla eklenir. M uzantısı dört bölme buyruğu tanımlar:

\begin{table}[htbp]
\centering
\caption{RISC-V Bölme Buyrukları (M Uzantısı)}
\label{tab:bolme-buyruklar}
\small
\begin{tabularx}{\textwidth}{l l c X}
\toprule
\textbf{Buyruk} & \textbf{Söz Dizimi} & \textbf{Tür} & \textbf{Açıklama} \\
\midrule
\texttt{div} & \texttt{div rd, rs1, rs2} & İşaretli & Bölümü rd'ye yazar. Sıfıra doğru yuvarlanır. \\
\texttt{divu} & \texttt{divu rd, rs1, rs2} & İşaretsiz & İşaretsiz bölümü rd'ye yazar. \\
\texttt{rem} & \texttt{rem rd, rs1, rs2} & İşaretli & Kalanı rd'ye yazar. Kalanın işareti bölünenle aynıdır. \\
\texttt{remu} & \texttt{remu rd, rs1, rs2} & İşaretsiz & İşaretsiz kalanı rd'ye yazar. \\
\bottomrule
\end{tabularx}
\end{table}

Her dört buyruk da aynı donanım bölme birimini kullanır; aralarındaki tek fark, işlenenlerin ve sonucun işaretli mi yoksa işaretsiz mi yorumlanacağıdır. İşaretli bölmede (\texttt{div}) bölüm \emph{sıfıra doğru yuvarlanır}: örneğin $-7 / 2 = -3{,}5$ yuvarlanınca $-3$ olur, $-4$'e değil. Kalan ise bölünen $=$ bölüm $\times$ bölen $+$ kalan özdeşliğinden bulunur ve her zaman bölünenle aynı işareti taşır.

MIPS'ten farklı olarak RISC-V'te bölüm ve kalan ayrı buyrukla elde edilir. MIPS'te \texttt{div} buyruğu her ikisini birden özel yazmaçlara (\texttt{hi}/\texttt{lo}) yazardı; RISC-V'te ise programcı hem bölümü hem kalanı istiyorsa \texttt{div} ve \texttt{rem} buyrukları art arda çalıştırılmalıdır. Ancak çağdaş RISC-V işlemcileri bu örüntüyü tanır ve iki buyruğu tek bir bölme işlemiyle birleştirir (\emph{division fusion}).

Dikkat edilmesi gereken önemli bir nokta: \texttt{div}/\texttt{rem} ile \texttt{divu}/\texttt{remu} buyrukları aynı bit örüntüsünü tamamen farklı yorumlar. Yazmaçtaki bitler değişmez; değişen yalnızca donanımın bu bitleri okuma biçimidir. Bu nedenle programcının (veya derleyicinin) veri türüne uygun buyruğu seçmesi gerekir: C dilindeki \texttt{int} gibi işaretli değişkenler için \texttt{div}/\texttt{rem}, \texttt{unsigned int} gibi işaretsiz değişkenler için \texttt{divu}/\texttt{remu} kullanılmalıdır. Bu farkın ne denli büyük sonuçlar doğurabileceği aşağıdaki örnekte görülmektedir.

\begin{ornek}[5.14]
\textbf{Soru:} $-7 \div 2$ işlemini RISC-V'te hem işaretli hem işaretsiz bölme buyrukları ile gerçekleştirin. Sonuçlar neden farklıdır?

\textbf{Çözüm:} Önce yazmaçları hazırlayalım:

\begin{lstlisting}[style=riscv]
    li   t1, -7    # bölünen: -7 (ikiye tümleyen: 0xFFFFFFF9)
    li   t2, 2     # bölen: 2
\end{lstlisting}

\textbf{1) İşaretli bölme} (\texttt{div} ve \texttt{rem}): Bu buyruklar yazmaç içeriğini ikiye tümleyen işaretli sayı olarak yorumlar.

\begin{lstlisting}[style=riscv]
    div  t3, t1, t2   # t3 = bölüm: -7 / 2 = -3
    rem  t4, t1, t2   # t4 = kalan: (-7) mod 2 = -1
\end{lstlisting}

Burada \texttt{div} buyruğu bölümü \texttt{t3} yazmacına, \texttt{rem} buyruğu ise kalanı \texttt{t4} yazmacına yazar. Bölüm sıfıra doğru yuvarlanır: $-7 / 2 = -3{,}5$ yuvarlanınca $-3$ olur. Kalan ise bölüm $\times$ bölen $+$ kalan $=$ bölünen denkleminden bulunur: $(-3) \times 2 + (-1) = -7$.\;\checkmark\quad Dikkat: işaretli bölmede kalan her zaman bölünenle aynı işareti taşır; burada bölünen $-7$ olduğundan kalan da eksidir ($-1$).

\textbf{2) İşaretsiz bölme} (\texttt{divu} ve \texttt{remu}): Şimdi aynı yazmaçlarla işaretsiz buyrukları deneyelim:

\begin{lstlisting}[style=riscv]
    divu t5, t1, t2   # t5 = bölüm: 4.294.967.289 / 2 = 2.147.483.644
    remu t6, t1, t2   # t6 = kalan: 4.294.967.289 mod 2 = 1
\end{lstlisting}

Bu kez donanım \texttt{t1}'deki aynı bit örüntüsünü ($\texttt{0xFFFFFFF9}$) işaretsiz bir sayı olarak yorumlar. İkiye tümleyen gösterimde $-7$'yi temsil eden bu bitler, işaretsiz okunduğunda $2^{32} - 7 = 4\,294\,967\,289$ gibi çok büyük bir değere karşılık gelir. Dolayısıyla \texttt{divu} aslında $4\,294\,967\,289 \div 2$ hesaplar: bölüm $= 2\,147\,483\,644$ (\texttt{t5}'e yazılır), kalan $= 1$ (\texttt{t6}'ya yazılır). Doğrulama: $2\,147\,483\,644 \times 2 + 1 = 4\,294\,967\,289$.\;\checkmark

\textbf{Sonuçların karşılaştırması:}

\begin{center}
\small
\begin{tabular}{l c c l}
\toprule
\textbf{Buyruk} & \textbf{Bölüm} & \textbf{Kalan} & \textbf{Açıklama} \\
\midrule
\texttt{div}/\texttt{rem} & $-3$ & $-1$ & \texttt{t1}'deki bitleri $-7$ olarak yorumlar \\
\texttt{divu}/\texttt{remu} & $2\,147\,483\,644$ & $1$ & Aynı bitleri $4\,294\,967\,289$ olarak yorumlar \\
\bottomrule
\end{tabular}
\end{center}

Yazmaçtaki bitler iki durumda da aynıdır (\texttt{0xFFFFFFF9}); değişen yalnızca donanımın bu bitleri yorumlama biçimidir. Bu nedenle programcının veri türüne uygun buyruğu seçmesi gerekir: C dilindeki \texttt{int} gibi işaretli değişkenler için \texttt{div}/\texttt{rem}, \texttt{unsigned int} gibi işaretsiz değişkenler için \texttt{divu}/\texttt{remu} kullanılmalıdır.
\end{ornek}


\subsection{Sıfıra Bölme ve Taşma}
\label{sec:sifira-bolme}

Tam sayı aritmetiğinde sıfıra bölmenin anlamlı bir sonucu yoktur; bölüm veya kalan olarak yazılabilecek geçerli bir tam sayı üretilemez. (Kayan nokta aritmetiğinde ise durum farklıdır: Bölüm~\ref{sec:kayan-nokta}'da gördüğümüz gibi IEEE~754 standardı sıfırdan farklı bir sayının sıfıra bölünmesini $\pm\infty$ olarak tanımlar.) Peki tam sayı bölme donanımı bölen sıfır olduğunda ne yapar?

Farklı mimariler bu sorunu farklı biçimlerde ele alır. Bazı işlemciler (örneğin MIPS) sıfıra bölme durumunda bir \emph{kural dışı durum}\index{kural dışı durum}\index{exception|see{kural dışı durum}} (\emph{exception}) üreterek programın çalışmasını durdurur ve işletim sistemine denetimi devreder. RISC-V ise farklı bir yaklaşım benimser: sıfıra bölme durumunda \textbf{kural dışı durum üretmez}; bunun yerine sonuç yazmaçlarına belirli değerler yazar ve programın çalışmasına devam eder. Bu tasarım kararının arkasında RISC-V'in yalınlık ilkesi yatar; kural dışı durum mekanizması donanıma karmaşıklık katar ve bölme biriminin boru hattıyla uyumunu zorlaştırır.

RISC-V'te sıfıra bölme ve taşma durumlarında döndürülen değerler Tablo~\ref{tab:sifira-bolme}'de özetlenmiştir.

\begin{table}[htbp]
\centering
\caption[Sıfıra bölme ve taşma sonuçları]{RISC-V'te sıfıra bölme ve taşma durumlarında döndürülen değerler (32 bit).}
\label{tab:sifira-bolme}
\small
\begin{tabular}{l l c c}
\toprule
\textbf{Durum} & \textbf{Buyruk} & \textbf{Bölüm} & \textbf{Kalan} \\
\midrule
Sıfıra bölme (işaretli) & \texttt{div}/\texttt{rem} & $-1$ (\texttt{0xFFFFFFFF}) & bölünen \\
Sıfıra bölme (işaretsiz) & \texttt{divu}/\texttt{remu} & $2^{32}-1$ (\texttt{0xFFFFFFFF}) & bölünen \\
Taşma ($-2^{31} \div -1$) & \texttt{div}/\texttt{rem} & $-2^{31}$ & $0$ \\
\bottomrule
\end{tabular}
\end{table}

Sıfıra bölmede bölüm olarak tüm bitleri 1 olan değerin döndürülmesi, yazılımın bu durumu kolayca saptamasını sağlar: program bölme buyruğundan sonra sonucun \texttt{0xFFFFFFFF} olup olmadığını denetleyerek gerekli önlemi alabilir. Kalan olarak bölünenin kendisinin döndürülmesi ise bölünen $=$ bölüm $\times$ bölen $+$ kalan özdeşliğiyle tutarlıdır (çünkü bölüm $\times$ 0 $+$ bölünen $=$ bölünen).

İşaretli bölmedeki taşma durumu da dikkat gerektirir: $-2^{31} \div (-1) = +2^{31}$ olması gerekirdi, ancak $+2^{31}$ değeri 32 bitlik ikiye tümleyen ile gösterilemez (en büyük pozitif değer $2^{31}-1$'dir). RISC-V bu durumda bölümü bölünenin kendisi ($-2^{31}$) olarak döndürür; bu matematiksel olarak yanlış bir sonuçtur, ancak donanım karmaşıklığını artırmamak adına bilinçli bir tercih yapılmıştır. Programcının bu uç durumu yazılımda denetlemesi beklenir.


% =============================================
\section{Kayan Nokta İşlemleri}
\label{sec:kayan-nokta}

Şimdiye kadar yalnızca tam sayılarla çalıştık. Peki $3{,}14$, $0{,}001$ ya da $6{,}022 \times 10^{23}$ gibi ondalıklı ve çok büyük/küçük sayıları nasıl göstereceğiz? 32 bitle işaretsiz olarak en fazla $2^{32}-1$, işaretli olarak $-2^{31}$ ile $2^{31}-1$ arasındaki tam sayıları gösterebiliyoruz; ancak gerçek dünyanın sorunları tam sayılarla sınırlı değil. İşte bu ihtiyaç için \emph{kayan nokta}\index{kayan nokta}\index{floating point|see{kayan nokta}} gösterimi geliştirilmiştir. İkilik taban dönüşümleri ve sayı gösterimi hakkında ayrıntılı bilgi için Ek~\ref{ch:sayi-sistemleri}'da bakınız.

Burada önemli bir noktanın altını çizelim: bellekte ya da yazmaçta saklanan 32 bitlik bir değerin tam sayı mı yoksa kayan nokta mı olduğunu \textbf{donanım kendiliğinden bilmez}. Aynı bit örüntüsü, kullanılan buyruğa bağlı olarak tamamen farklı yorumlanır. Örneğin \texttt{add} buyruğu yazmaçtaki bitleri ikiye tümleyen tam sayı olarak ele alırken, \texttt{fadd.s} buyruğu aynı bitleri IEEE~754 kayan nokta sayısı olarak yorumlar. Bu ayrımı belirleyen programcı veya derleyicidir; donanım yalnızca kendisine verilen buyruğa uygun işlemi gerçekleştirir. Aynı ilkeyi daha önce işaretli ve işaretsiz tam sayılarda da görmüştük (Bölüm~\ref{sec:riscv-bolme-buyruklar}'deki \texttt{div}/\texttt{divu} örneğini hatırlayın); kayan nokta ile tam sayı arasındaki fark da bundan farklı değildir.

\subsection{Sabit Nokta ve Kayan Nokta}
\label{sec:sabit-kayan}

Kesirli sayıları ikilik tabanda göstermenin iki temel yolu vardır: \textbf{sabit nokta}\index{sabit nokta}\index{fixed point|see{sabit nokta}} ve \textbf{kayan nokta}\index{kayan nokta}.

\subsubsection{Sabit Nokta Gösterimi}

Sabit nokta gösteriminde virgülün yeri önceden sabittir: bit dizisinin belirli bir kısmı tam sayı bölümünü, geri kalanı kesir bölümünü temsil eder. Örneğin 16 bitlik bir sabit nokta biçiminde 8 bit tam kısma, 8 bit kesir kısmına ayrılabilir:

\begin{center}
\resizebox{0.85\linewidth}{!}{%
\begin{tikzpicture}[x=1cm, y=1cm]
  % Tam kısım
  \draw[fill=blokmavi, draw=black, thick]
      (0,0) rectangle (6.0,0.9);
  \node[font=\small\bfseries] at (3.0, 1.2)
      {Tam kısım (8 bit)};
  \node[font=\normalsize\ttfamily] at (3.0, 0.45)
      {TTTTTTTT};
  % Kesir kısmı
  \draw[fill=blokyesil, draw=black, thick]
      (6.0,0) rectangle (12.0,0.9);
  \node[font=\small\bfseries] at (9.0, 1.2)
      {Kesir kısmı (8 bit)};
  \node[font=\normalsize\ttfamily] at (9.0, 0.45)
      {KKKKKKKK};
  % Virgül
  \draw[very thick, red] (6.0, -0.15) -- (6.0, 1.05);
  \node[font=\footnotesize\bfseries, red] at (6.0, -0.35) {virgül};
\end{tikzpicture}}
\end{center}

Bu düzende gösterilebilecek en büyük sayı $255{,}996$ ($\approx 2^8 - 2^{-8}$), en küçük pozitif sayı ise $0{,}0039$ ($= 2^{-8}$) olur. Sabit nokta, donanım açısından çok yalındır; toplama ve çıkarma tam sayı devreleriyle aynıdır. Ancak ciddi bir sınırlaması vardır: \textbf{değer aralığı çok dardır}.

Bu sınırlamayı somut bir örnekle görelim. Bir mühendislik programında aşağıdaki üç değerle aynı anda çalışılması gerektiğini düşünelim:

\begin{center}
\small
\begin{tabularx}{\textwidth}{l r X l}
\toprule
\textbf{Büyüklük} & \textbf{Değer} & \textbf{8.8 sabit noktada} & \textbf{Sonuç} \\
\midrule
Köprü uzunluğu & $1\,250{,}75$ m & Tam kısım $1250 > 255$ & \textcolor{red}{Taşma} \\[3pt]
Vida çapı & $12{,}347$ mm & $12 + 0{,}34375$; $12{,}347 \approx 12{,}344$ & Duyarlık kaybı \\[3pt]
Mikrop boyutu & $0{,}00085$ mm & En küçük kesir $0{,}0039 > 0{,}00085$ & \textcolor{red}{Gösterilemez} \\
\bottomrule
\end{tabularx}
\end{center}

Sorun açıktır: virgülün yeri sabitlendikten sonra, tam kısmın kapasitesini aşan büyük sayılar \emph{taşar} (\emph{overflow}); kesir kısmının çözünürlüğünden küçük sayılar ise sıfıra yuvarlanarak kaybolur.

Üstelik burada gizli bir \emph{ödünleşme} de vardır. Köprünün uzunluğu ($1250{,}75$) tam kısma sığmadığı için gösterilemez; ancak bu sayının kesir kısmı yalnızca 2 ondalık basamak içerir ve 8 bitlik kesir alanına rahatlıkla sığardı. Sabit nokta biçiminde tam kısma ayrılan 8 bit bu sayı için yetersizken, kesir kısmına ayrılan 8 bit boşa harcanmaktadır. Tersine, mikrobun boyutu ($0{,}00085$) için tam kısım hiç gerekmezken kesir alanı yetersiz kalmaktadır. Virgülü sağa kaydırsak büyük sayıları gösterebiliriz ama küçük sayıları kaybederiz; sola kaydırsak tam tersi olur. Virgül nereye konursa konsun, 16 bitin tamamını verimli kullanan \emph{tek bir konum yoktur}; sorun virgülün sabit olmasının kendisidir.

Bit genişliğini artırmak (örneğin 16.16 biçimine geçmek) aralığı genişletir ancak sorunun özünü değiştirmez; çünkü virgül hâlâ sabittir ve yeterince büyük ya da yeterince küçük bir sayı yine gösterilemez.

\subsubsection{Kayan Nokta: Virgülü Serbest Bırakmak}

Kayan nokta gösterimi bu sorunu, virgülün yerini \emph{değişken} yaparak çözer. Sayı, bir \emph{anlamlı kısım} ve bir \emph{üs} olarak ikiye ayrılır; üs, virgülün nereye kayacağını belirler. Böylece aynı bit genişliğiyle hem çok büyük ($10^{38}$) hem de çok küçük ($10^{-38}$) sayılar ifade edilebilir, tabii ki sınırlı bir duyarlıkla.

Bu fikir, onluk tabandaki bilimsel gösterimin ikilik tabana uyarlanmasıdır. Nasıl çalıştığını anlamak için önce bilimsel gösterimi ve olağanlaştırmayı ele alalım.

\subsection{Bilimsel Gösterim ve Olağanlaştırma}

Bilimsel gösterim\index{bilimsel gösterim}\index{scientific notation|see{bilimsel gösterim}}, sayıda virgülün solunda tek basamak olmasıdır. Virgülden önceki basamak 0 değilse buna \emph{olağanlaştırılmış}\index{olağanlaştırma}\index{normalization|see{olağanlaştırma}} gösterim denir.

\begin{itemize}
\item $21{,}41239$ -- bilimsel gösterimde değildir.
\item $0{,}2141239 \times 10^2$ -- bilimsel gösterimdedir.
\item $2{,}141239 \times 10^1$ -- bilimsel gösterimdedir ve \emph{olağandır}.
\end{itemize}

Bilimsel gösterim 2'lik düzene uygulandığında:
\begin{align*}
(01{,}10)_2 &= 1{,}10 \times 2^0 \\
(10011010{,}110)_2 &= 1{,}0011010110 \times 2^7
\end{align*}

İkilik düzende olağanlaştırılmış gösterim $1{,}\ldots \times 2^e$ ile ifade edilir. Çünkü virgülün solundaki basamağın 0'dan farklı olması ve tek bir basamak olması zorunludur. 2'lik düzende 0'dan farklı tek basamak 1 olduğu için virgülden önceki basamağın 1 olması kaçınılmazdır.

\begin{terimnotu}
\textbf{olağanlaştırma} (\emph{normalization})\,: İngilizce \emph{normalize} fiilinden gelen bu terim, Türkçe bilişim yazınında ``normalleştirme'' ve ``normalizasyon'' biçimlerinde de geçer. Bu kitapta Türkçe kökenli \textbf{olağanlaştırma} tercih edilmiştir: \emph{olağan} sözcüğü ``normal'' anlamına gelir; olağanlaştırma da sayıyı \emph{olağan biçime} (virgülün solunda tek anlamlı basamak kalacak şekle) getirmek demektir. Aynı kökten türeyen \emph{olağanlaştırılmamış} (\emph{denormalized}) terimi ise IEEE~754'te sıfıra çok yakın özel değerleri tanımlar.

Not: Türkçe bilişim terminolojisinde terim birliği henüz tam sağlanamamıştır; ``normalleştirme'' veri tabanı ve veri bilimi alanlarında daha yaygınken, ``olağanlaştırma'' bilgisayar mimarisi ve sayısal çözümleme bağlamında tercih edilen biçimdir.
\end{terimnotu}

\subsection{IEEE 754 Kayan Nokta Standardı}
\label{sec:ieee754}

Kayan nokta sayılarının bellekte nasıl temsil edileceğini belirleyen evrensel standart \textbf{IEEE 754}\index{IEEE 754}'tür (\emph{Institute of Electrical and Electronics Engineers}, 1985). Bu standart öncesinde her işlemci üreticisi kendi kayan nokta biçimini kullanıyordu; bir makinede hesaplanan sonuç başka bir makinede farklı çıkabiliyordu. IEEE 754, işaret, üs ve kesir alanlarının bit genişliklerini, eklentili üs kodlamasını, özel değerleri (sıfır, sonsuz, NaN) ve yuvarlama kurallarını kesin biçimde tanımlayarak \emph{taşınabilir} ve \emph{öngörülebilir} kayan nokta aritmetiğini mümkün kıldı.

\subsubsection{Tek Duyarlı Gösterim (32 bit)}\index{kayan nokta!tek duyarlı}

Tek duyarlı (\emph{single precision}) biçimde 32 bit üç alana bölünür:

\begin{center}
\resizebox{\linewidth}{!}{%
\begin{tikzpicture}[x=1cm, y=1cm]
  % İşaret alanı
  \draw[fill=blokkirmizi, draw=black, thick]
      (0,0) rectangle (1.2,0.9);
  \node[font=\small\bfseries] at (0.6, 1.55) {İşaret};
  \node[font=\scriptsize] at (0.6, 1.2) {(1 bit)};
  \node[font=\scriptsize, anchor=north west] at (0.05, -0.15) {31};
  \node[font=\normalsize] at (0.6, 0.45) {$i$};
  % Üs alanı
  \draw[fill=blokmavi, draw=black, thick]
      (1.2,0) rectangle (5.0,0.9);
  \node[font=\small\bfseries] at (3.1, 1.55) {Üs (\emph{exponent})};
  \node[font=\scriptsize] at (3.1, 1.2) {(8 bit)};
  \node[font=\scriptsize, anchor=north west] at (1.25, -0.15) {30};
  \node[font=\scriptsize, anchor=north east] at (4.95, -0.15) {23};
  \node[font=\normalsize] at (3.1, 0.45) {$e_7 \;\ldots\; e_0$};
  % Kesir alanı
  \draw[fill=blokyesil, draw=black, thick]
      (5.0,0) rectangle (14.0,0.9);
  \node[font=\small\bfseries] at (9.5, 1.55) {Kesir (\emph{fraction})};
  \node[font=\scriptsize] at (9.5, 1.2) {(23 bit)};
  \node[font=\scriptsize, anchor=north west] at (5.05, -0.15) {22};
  \node[font=\scriptsize, anchor=north east] at (13.95, -0.15) {0};
  \node[font=\normalsize] at (9.5, 0.45) {$f_{22} \;\ldots\; f_0$};
\end{tikzpicture}}
\end{center}

\begin{itemize}
  \item \textbf{İşaret (bit 31):} 0 ise sayı artı, 1 ise eksidir.

  \item \textbf{Üs\index{üs}\index{exponent|see{üs}} (bit 30--23):} Gerçek üs, $e_{\text{gerçek}} = e - 127$ formülüyle hesaplanır. 8 bitlik alan 0--255 arası değer alır; ancak 0 ve 255 özel değerler için ayrıldığından kullanılabilir aralık 1--254'tür. Bu kodlamaya \textbf{eklentili gösterim}\index{eklentili gösterim}\index{biased notation|see{eklentili gösterim}}\index{excess-N notation|see{eklentili gösterim}} (\emph{biased} / \emph{excess-127}) denir: gerçek üsse sabit bir \emph{eklenti} (\emph{bias}) değeri (tek duyarlıda 127) eklenerek üs alanına yazılır. Negatif üslerin ayrı bir işaret bitine gerek kalmadan ifade edilmesini sağlar; ayrıca iki kayan nokta sayısının üs alanlarını işaretsiz tam sayı gibi karşılaştırarak büyüklük sıralamaya olanak tanır.

  \item \textbf{Kesir (bit 22--0):} Olağanlaştırılmış sayılarda virgülden önceki bit her zaman \texttt{1} olduğundan bu bit bellekte \emph{saklanmaz}; donanım tarafından otomatik olarak eklenir. Bu saklanmayan bite \textbf{gizli 1}\index{gizli 1}\index{hidden bit|see{gizli 1}} (\emph{hidden bit}) denir. Dolayısıyla 23 bit ile aslında 24 bitlik duyarlık elde edilir.
\end{itemize}

Burada iki terimi birbirinden ayırt etmek gerekir: bellekte saklanan 23 bitlik alana \textbf{kesir}\index{kesir}\index{fraction|see{kesir}} (\emph{fraction}) denir; bu, virgülden sonraki kısımdır. Gizli 1 bu alanın başına eklendiğinde elde edilen $1{,}\text{kesir}$ değerine ise \textbf{anlamlı kısım}\index{anlamlı kısım}\index{significand|see{anlamlı kısım}} (\emph{significand}) denir. Eski kaynaklarda anlamlı kısım için \emph{mantis} (\emph{mantissa}) terimini görebilirsiniz; ancak bu terim logaritma bağlamından geldiği için IEEE 754 standardı \emph{significand} terimini tercih eder. Özetle:
\[
\underbrace{1}_{\text{gizli 1}}\;.\;\underbrace{f_{22}\,f_{21}\,\ldots\,f_0}_{\text{kesir (23 bit, bellekte saklanır)}} \;=\; \underbrace{1{,}f_{22}\,f_{21}\,\ldots\,f_0}_{\text{anlamlı kısım (24 bit etkin duyarlık)}}
\]

\begin{bilginot}[Gizli 1 Her Yerde Var mı?]
Gizli 1, IEEE 754'ün akıllıca bir tasarım kararıdır; ancak tüm kayan nokta biçimleri bu yöntemi kullanmaz. IEEE 754 öncesinde ve sonrasında farklı yaklaşımlar benimsenmiştir:

\begin{center}
\small
\begin{tabularx}{\textwidth}{l c c X}
\toprule
\textbf{Biçim} & \textbf{Taban} & \textbf{Gizli 1?} & \textbf{Açıklama} \\
\midrule
IEEE 754 (FP32/FP64) & 2 & Evet & Olağanlaştırılmış sayılarda virgülden önceki 1 saklanmaz; 1 bit bedavaya gelir. \\[3pt]
Intel x87 genişletilmiş (80 bit) & 2 & Hayır & 64 bitlik anlamlı kısım tam olarak saklanır; baştaki 1 açıkça yazılır. Olağanlaştırılmamış sayıların ve ara hesap sonuçlarının özel işlem gerektirmemesi için bu tasarım tercih edilmiştir. \\[3pt]
IBM System/360 (HFP) & 16 & Hayır & Üs 16'nın kuvvetlerini gösterir; anlamlı kısım ikilik tabanda en az bir anlamlı onaltılık basamakla olağanlaştırılır. 16 tabanlı olağanlaştırmada baştaki bit her zaman 1 olmadığından gizli 1 kullanılamaz; bu yüzden 3 bite kadar duyarlık kaybı olabilir. \\[3pt]
DEC VAX (F/D/G/H) & 2 & Evet & IEEE 754'e benzer biçimde gizli 1 kullanır; ancak olağanlaştırılmamış sayıları ve NaN'ı desteklemez. \\
\bottomrule
\end{tabularx}
\end{center}

Intel x87'nin 80 bitlik genişletilmiş biçimi özellikle ilginçtir: bu biçim IEEE 754'ün bir parçası olmasına rağmen gizli 1 kullanmaz. Bunun nedeni, x87'nin kayan nokta hesaplamalarında \emph{ara sonuçları} tutmak için tasarlanmış olmasıdır; olağanlaştırılmamış ara değerlerde gizli 1in varsayılması hatalara yol açabilirdi. Açık bit sayesinde donanımın her durumda doğru davranması garanti edilir; bedeli ise anlamlı kısım için 1 bit fazla alan harcanmasıdır.

IBM System/360'ın onaltılık tabanlı kayan nokta biçimi ise tarihsel bir ders niteliğindedir: 16 tabanlı olağanlaştırma, ikilik tabana göre daha ``gevşek'' olduğundan anlamlı kısmın ilk 1--3 biti sıfır olabilir. Bu durum etkili duyarlığı düşürür ve sayısal kararlılık sorunlarına yol açar. IEEE 754'ün ikilik tabana ve gizli 1e geçişinin temel nedenlerinden biri de bu duyarlık kaybını ortadan kaldırmaktı.
\end{bilginot}

\begin{bilginot}[IEEE 754 Üsleri Neden Eklentili Saklar?]
Tam sayı aritmetiğinde ikiye tümleyen çok iyi çalışır; o zaman üs alanı için de aynı yöntemi kullanmak doğal görünür. Ancak kayan nokta sayılarında bir gereksinim daha vardır: \emph{büyüklük sıralaması}. Eğer üs alanı ikiye tümleyenle kodlansaydı, eksi bir üs (örneğin $-5$, ikiye tümleyende \texttt{11111011}) artı bir üsten (örneğin $+3$, ikiye tümleyende \texttt{00000011}) \emph{büyük} görünürdü; çünkü en anlamlı bit 1'dir. Bu durumda iki kayan nokta sayısını karşılaştırmak için üs alanının ikiye tümleyen mi işaretsiz mi olduğunu ayırt eden özel bir devre gerekirdi.

Eklentili gösterimde ise tüm üs değerleri 0'dan 255'e kadar artı sayılardır. En küçük gerçek üs ($-126$) en küçük kodlanmış değere ($1$), en büyük gerçek üs ($+127$) en büyük kodlanmış değere ($254$) karşılık gelir; sıralama doğal olarak korunur. Daha da önemlisi, işaret bitleri aynı olan iki kayan nokta sayısı, 32 bitlik gösterimleri \emph{tek bir işaretsiz tam sayı gibi} karşılaştırılarak sıralanabilir. Bu özellik donanımda karşılaştırma devresini büyük ölçüde yalınlaştırır.

\end{bilginot}

\begin{terimnotu}
\textbf{eklentili gösterim}\index{eklentili gösterim}\index{biased notation|see{eklentili gösterim}}\index{excess-N notation|see{eklentili gösterim}} (\emph{biased notation}, \emph{excess-N notation})\,: İngilizce \emph{biased} sözcüğü ``önyargılı'' anlamına gelir; ancak ``önyargılı üs'' sanki üste bir hata eklenmiş izlenimi verir. Türkçe kaynaklarda kullanılan ``sapmalı'' da istenmeyen bir kayma çağrışımı yapar. Oysa yapılan işlem yalındır: gerçek üsse sabit bir değer \emph{eklenir}. Bu nedenle bu kitapta \textbf{eklentili gösterim} terimi tercih edilmiştir.
\end{terimnotu}

Gerçek değer şu şekilde hesaplanır:
\[
  (-1)^{i} \times 1{,}f_{22}f_{21}\ldots f_0 \times 2^{e - 127}
\]

C dilinde \texttt{float} veri tipiyle karşılık gelir; yaklaşık 7 onluk basamak duyarlık ve $\approx 10^{\pm 38}$ aralık sunar. Örneğin aşağıdaki C kodunda tanımlanan her değişken bellekte 32 bitlik IEEE~754 tek duyarlı biçimde saklanır:

\begin{lstlisting}[language=C, xleftmargin=2em]
float sicaklik = 36.6;     /* 32 bit IEEE 754 */
float pi       = 3.14159;  /* 32 bit IEEE 754 */
float kutle    = 0.00085;  /* 32 bit IEEE 754 */
\end{lstlisting}

Derleyici bu tanımları gördüğünde her değeri IEEE~754 biçimine dönüştürür ve bellekte 32 bitlik bir alan ayırır. Programcı ondalıklı bir sayı yazmış olsa da bellekteki gösterim işaret, üs ve kesir bitlerinden oluşur; onluk tabandaki basamaklarla doğrudan bir ilişkisi yoktur.

\subsubsection{Çift Duyarlı Gösterim (64 bit)}\index{kayan nokta!çift duyarlı}

Çift duyarlı (\emph{double precision}) biçimde 64 bit yine üç alana bölünür:

\begin{center}
\resizebox{\linewidth}{!}{%
\begin{tikzpicture}[x=1cm, y=1cm]
  % İşaret alanı
  \draw[fill=blokkirmizi, draw=black, thick]
      (0,0) rectangle (0.7,0.9);
  \node[font=\small\bfseries] at (0.35, 1.55) {İşaret};
  \node[font=\scriptsize] at (0.35, 1.2) {(1 bit)};
  \node[font=\scriptsize, anchor=north west] at (0.05, -0.15) {63};
  \node[font=\normalsize] at (0.35, 0.45) {$i$};
  % Üs alanı
  \draw[fill=blokmavi, draw=black, thick]
      (0.7,0) rectangle (3.5,0.9);
  \node[font=\small\bfseries] at (2.1, 1.55) {Üs};
  \node[font=\scriptsize] at (2.1, 1.2) {(11 bit)};
  \node[font=\scriptsize, anchor=north west] at (0.75, -0.15) {62};
  \node[font=\scriptsize, anchor=north east] at (3.45, -0.15) {52};
  \node[font=\normalsize] at (2.1, 0.45) {$e_{10} \;\ldots\; e_0$};
  % Kesir alanı
  \draw[fill=blokyesil, draw=black, thick]
      (3.5,0) rectangle (14.0,0.9);
  \node[font=\small\bfseries] at (8.75, 1.55) {Kesir};
  \node[font=\scriptsize] at (8.75, 1.2) {(52 bit)};
  \node[font=\scriptsize, anchor=north west] at (3.55, -0.15) {51};
  \node[font=\scriptsize, anchor=north east] at (13.95, -0.15) {0};
  \node[font=\normalsize] at (8.75, 0.45) {$f_{51} \;\ldots\; f_0$};
\end{tikzpicture}}
\end{center}

Üs alanı 11 bit olup eklenti değeri \textbf{1023}'tür; gerçek üs $e - 1023$ ile bulunur. Kesir alanı 52 bit olup gizli 1 ile birlikte 53 bitlik duyarlık sağlar ($\approx 15{-}16$ onluk basamak, $\approx 10^{\pm 308}$ aralık). C dilinde \texttt{double} veri tipiyle ifade edilir:

\begin{lstlisting}[language=C, xleftmargin=2em]
double uzaklik  = 149597870.7;  /* 64 bit IEEE 754 */
double elektron = 9.10938e-31;  /* 64 bit IEEE 754 */
\end{lstlisting}

Daha yüksek duyarlık veya daha geniş değer aralığı gereken hesaplamalarda \texttt{float} yerine \texttt{double} tercih edilir; ancak bellek kullanımı iki katına çıkar.

\subsubsection{Genişletilmiş Duyarlı Gösterim: Intel 80 Bit Biçimi}\index{kayan nokta!genişletilmiş duyarlı}\index{x87}\index{long double|see{genişletilmiş duyarlı}}

IEEE 754 standardı tek ve çift duyarlının yanı sıra isteğe bağlı bir \emph{genişletilmiş duyarlı} (\emph{extended precision}) biçim kategorisi de tanımlar. Standart bu biçimin tam düzenini belirlemez; yalnızca asgari gereksinimleri koyar: en az 64 bit anlamlı kısım ve en az 15 bit üs. Intel'in x87 kayan nokta birimi bu gereksinimleri karşılayan 80 bitlik bir biçim tanımlamıştır:

\begin{center}
\resizebox{\linewidth}{!}{%
\begin{tikzpicture}[x=1cm, y=1cm]
  % İşaret alanı
  \draw[fill=blokkirmizi, draw=black, thick]
      (0,0) rectangle (0.7,0.9);
  \node[font=\small\bfseries] at (0.35, 1.55) {İşaret};
  \node[font=\scriptsize] at (0.35, 1.2) {(1 bit)};
  \node[font=\scriptsize, anchor=north west] at (0.05, -0.15) {79};
  \node[font=\normalsize] at (0.35, 0.45) {$i$};
  % Üs alanı
  \draw[fill=blokmavi, draw=black, thick]
      (0.7,0) rectangle (4.0,0.9);
  \node[font=\small\bfseries] at (2.35, 1.55) {Üs};
  \node[font=\scriptsize] at (2.35, 1.2) {(15 bit)};
  \node[font=\scriptsize, anchor=north west] at (0.75, -0.15) {78};
  \node[font=\scriptsize, anchor=north east] at (3.95, -0.15) {64};
  \node[font=\normalsize] at (2.35, 0.45) {$e_{14} \;\ldots\; e_0$};
  % Tam bit (açık 1)
  \draw[fill=bloksari, draw=black, thick]
      (4.0,0) rectangle (5.0,0.9);
  \node[font=\small\bfseries] at (4.5, 1.55) {T};
  \node[font=\scriptsize] at (4.5, 1.2) {(1 bit)};
  \node[font=\scriptsize, anchor=north] at (4.5, -0.15) {63};
  \node[font=\normalsize\bfseries] at (4.5, 0.45) {\texttt{1}};
  % Kesir alanı
  \draw[fill=blokyesil, draw=black, thick]
      (5.0,0) rectangle (14.0,0.9);
  \node[font=\small\bfseries] at (9.5, 1.55) {Kesir};
  \node[font=\scriptsize] at (9.5, 1.2) {(63 bit)};
  \node[font=\scriptsize, anchor=north west] at (5.05, -0.15) {62};
  \node[font=\scriptsize, anchor=north east] at (13.95, -0.15) {0};
  \node[font=\normalsize] at (9.5, 0.45) {$f_{62} \;\ldots\; f_0$};
\end{tikzpicture}}
\end{center}

Bu biçimin tek ve çift duyarlıdan en önemli farkı, çizimde sarıyla gösterilen \textbf{T} (\emph{integer bit}, tam kısım biti) alanıdır. Tek ve çift duyarlıda virgülden önceki \texttt{1} gizli 1tir; saklanmaz, donanım tarafından varsayılır. 80 bitlik genişletilmiş biçimde ise bu bit \textbf{açıkça saklanır}: olağanlaştırılmış sayılarda \texttt{1}, olağanlaştırılmamış sayılarda \texttt{0} değerini alır. Bu tasarım kararının nedeni, genişletilmiş biçimin çok adımlı hesaplamalarda \emph{ara sonuçları} tutmak için tasarlanmış olmasıdır; ara değerler olağanlaştırılmamış olabilir ve gizli 1in varsayılması hatalara yol açabilirdi. Açık bit sayesinde donanımın her durumda doğru davranması garanti edilir; bedeli yalnızca 1 bit fazla alan kullanılmasıdır.

Üs alanı 15 bit olup eklenti değeri \textbf{16\,383}'tür. Bu geniş üs alanı $\approx 10^{\pm 4932}$ aralık sağlar; 64 bitlik anlamlı kısım ise $\approx 18{-}19$ onluk basamak duyarlık verir.

\textbf{C dilinde karşılığı:} 80 bitlik genişletilmiş biçim, C dilinde \texttt{long double} veri tipiyle kullanılır; ancak bu davranış platforma bağlıdır:

\begin{lstlisting}[language=C, xleftmargin=2em]
long double pi = 3.14159265358979323846L;  /* x86'da 80 bit */
\end{lstlisting}

\begin{center}
\small
\begin{tabularx}{\textwidth}{l l X}
\toprule
\textbf{Platform} & \texttt{long double} & \textbf{Açıklama} \\
\midrule
x86 Linux (GCC) & 80 bit (x87) & Bellekte 12 veya 16 bayt hizalanır \\
x86 Windows (MSVC) & 64 bit (FP64) & \texttt{long double} = \texttt{double} ile aynı \\
ARM / Apple Silicon & 128 bit (FP128) & IEEE 754 dörtlü duyarlı (yazılım emülasyonu) \\
RISC-V & 128 bit (FP128) & Q uzantısı ile donanım desteği tanımlı (henüz yaygın değil) \\
\bottomrule
\end{tabularx}
\end{center}

Görüldüğü gibi \texttt{long double} yazıp ``yüksek duyarlık'' elde ettiğini düşünen bir programcı, kodunu farklı bir platformda derlediğinde beklenmedik sonuçlarla karşılaşabilir. Taşınabilir kod yazarken bu farkın bilincinde olmak gerekir.

\begin{bilginot}[RISC-V'te Dörtlü Duyarlık: Q Uzantısı]
RISC-V mimarisi 128 bitlik dörtlü duyarlı (\emph{quad precision}) kayan nokta desteğini isteğe bağlı \textbf{Q uzantısı}\index{Q uzantısı} ile tanımlar. Q uzantısı IEEE 754 dörtlü duyarlı biçimi kullanır: 1 işaret $+$ 15 üs $+$ 112 kesir $=$ 128 bit. Bu biçim $\approx 34$ onluk basamak duyarlık ve $\approx 10^{\pm 4932}$ aralık sağlar. Intel'in 80 bitlik genişletilmiş biçiminden farklı olarak gizli 1 kullanır ve tam IEEE 754 uyumludur.

Ancak Q uzantısı henüz yaygın donanım desteğine sahip değildir; günümüzde dörtlü duyarlık gerektiğinde çoğu RISC-V derleyicisi bunu yazılımla gerçekleştirir (\emph{soft float}). Buna karşın F (tek duyarlı) ve D (çift duyarlı) uzantıları neredeyse tüm uygulama sınıfı RISC-V işlemcilerinde donanımla desteklenir.
\end{bilginot}

IEEE 754 yalnızca tek, çift ve genişletilmiş duyarlı biçimleri tanımlamaz; günümüzde yapay zeka donanımları çok daha dar biçimler de kullanır. Bu güncel biçimler Bölüm~\ref{sec:yapay-zeka-kayan-nokta}'da ele alınmaktadır.

\subsubsection{Kayan Nokta Formülü}

Kayan nokta ile gösterilen bir sayı aşağıdaki gibi ifade edilir:
\begin{equation}
(-1)^i \times (1 + \text{kesir}) \times 2^{(e - \text{referans})}
\label{eq:kayan-nokta}
\end{equation}

Burada:
\begin{itemize}
\item $(-1)^i$: Sayının işareti ($i=0$ ise artı, $i=1$ ise eksi)
\item Kesir: Virgülden sonraki kısım (1 ile toplanarak gerçek anlamlı kısım\index{anlamlı kısım}\index{significand|see{anlamlı kısım}}\index{mantissa|see{anlamlı kısım}} bulunur)
\item $e$: Üst alanındaki değer
\item Referans: Tek duyarlıda 127, çift duyarlıda 1023
\end{itemize}

IEEE 754 standardında üs gösterimi için 2'ye tümleyen yerine eklentili (\emph{biased} / \emph{excess-N}) gösterim kullanılır. Tek duyarlıda eklenti değeri 127'dir. 8 bitlik üs alanı 0'dan 255'e kadar 256 farklı değer alabilir; ancak 0 ve 255 özel durumlar (sıfır, olağanlaştırılmamış sayılar, sonsuz, NaN) için ayrıldığından kullanılabilir üs değerleri 1'den 254'e kadardır. Gerçek üs $= e - 127$ olduğundan, gerçek üs aralığı $-126$'dan $+127$'ye kadardır.

\begin{ornek}[5.15]
\textbf{Soru:} 5,75 sayısını IEEE 754 standardında tek ve çift duyarlılıkta gösterin.

\textbf{Çözüm:} $5 = (101)_2$

Virgülden sonraki kısım:
\begin{align*}
0{,}75 \times 2 &= 1{,}5 \quad \rightarrow 1 \\
0{,}5 \times 2 &= 1{,}0 \quad \rightarrow 1
\end{align*}

$5{,}75 = (101{,}11)_2 = 1{,}0111 \times 2^2$

\textbf{Tek duyarlılık:} Üs $= 2 + 127 = 129 = (10000001)_2$

\begin{center}
\resizebox{\linewidth}{!}{%
\begin{tikzpicture}[x=1cm, y=1cm]
  \draw[fill=blokkirmizi, draw=black, thick] (0,0) rectangle (1.5,1.0);
  \node[font=\large\ttfamily\bfseries] at (0.75, 0.5) {0};
  \node[font=\scriptsize] at (0.75, -0.3) {İşaret};
  \draw[fill=blokmavi, draw=black, thick] (1.5,0) rectangle (5.5,1.0);
  \node[font=\large\ttfamily] at (3.5, 0.5) {10000001};
  \node[font=\scriptsize] at (3.5, -0.3) {Üs = 129};
  \draw[fill=blokyesil, draw=black, thick] (5.5,0) rectangle (16.0,1.0);
  \node[font=\large\ttfamily] at (10.75, 0.5) {01110000000000000000000};
  \node[font=\scriptsize] at (10.75, -0.3) {Kesir};
\end{tikzpicture}}
\end{center}

\textbf{Çift duyarlılık:} Üs $= 2 + 1023 = 1025 = (10000000001)_2$

\begin{center}
\resizebox{\linewidth}{!}{%
\begin{tikzpicture}[x=1cm, y=1cm]
  \draw[fill=blokkirmizi, draw=black, thick] (0,0) rectangle (0.8,1.0);
  \node[font=\large\ttfamily\bfseries] at (0.4, 0.5) {0};
  \node[font=\scriptsize] at (0.4, -0.3) {İşaret};
  \draw[fill=blokmavi, draw=black, thick] (0.8,0) rectangle (4.0,1.0);
  \node[font=\large\ttfamily] at (2.4, 0.5) {10000000001};
  \node[font=\scriptsize] at (2.4, -0.3) {Üs = 1025};
  \draw[fill=blokyesil, draw=black, thick] (4.0,0) rectangle (16.0,1.0);
  \node[font=\normalsize\ttfamily] at (10.0, 0.5) {0111000000000000000000000000000000000000000000000000};
  \node[font=\scriptsize] at (10.0, -0.3) {Kesir (52 bit)};
\end{tikzpicture}}
\end{center}
\end{ornek}

\begin{ornek}[5.16]
\textbf{Soru:} 2,1 sayısını IEEE 754 standardında tek duyarlılıkta gösterin.

\textbf{Çözüm:} $2 = (10)_2$

Virgülden sonraki kısım (devreden bir sayı dizisi):
\begin{align*}
0{,}1 \times 2 &= 0{,}2 \quad \rightarrow 0 \\
0{,}2 \times 2 &= 0{,}4 \quad \rightarrow 0 \\
0{,}4 \times 2 &= 0{,}8 \quad \rightarrow 0 \\
0{,}8 \times 2 &= 1{,}6 \quad \rightarrow 1 \\
0{,}6 \times 2 &= 1{,}2 \quad \rightarrow 1 \\
0{,}2 \times 2 &= 0{,}4 \quad \rightarrow 0 \quad \text{(tekrar başlar: } 0011\ldots\text{)}
\end{align*}

$2{,}1 \approx (10{,}00011001100110011001100\ldots)_2 = 1{,}000011001100110011001100 \times 2^1$

Üs: $1 + 127 = 128 = (10000000)_2$

\begin{center}
\resizebox{\linewidth}{!}{%
\begin{tikzpicture}[x=1cm, y=1cm]
  \draw[fill=blokkirmizi, draw=black, thick] (0,0) rectangle (1.5,1.0);
  \node[font=\large\ttfamily\bfseries] at (0.75, 0.5) {0};
  \node[font=\scriptsize] at (0.75, -0.3) {İşaret};
  \draw[fill=blokmavi, draw=black, thick] (1.5,0) rectangle (5.5,1.0);
  \node[font=\large\ttfamily] at (3.5, 0.5) {10000000};
  \node[font=\scriptsize] at (3.5, -0.3) {Üs = 128};
  \draw[fill=blokyesil, draw=black, thick] (5.5,0) rectangle (16.0,1.0);
  \node[font=\large\ttfamily] at (10.75, 0.5) {00001100110011001100110};
  \node[font=\scriptsize] at (10.75, -0.3) {Kesir};
\end{tikzpicture}}
\end{center}

Bu sayıda devreden ($1100$ tekrar eden) bir sayı dizisi bulunur; 2,1 sayısı ikilik tabanda \emph{tam olarak} gösterilemez.
\end{ornek}

\begin{ornek}[5.17]
\textbf{Soru:} Tek duyarlılıkta \texttt{11000001101110100000000000000000} şeklinde gösterilen sayı 10'luk tabanda kaçtır?

\textbf{Çözüm:} Sayı parçalara ayrılır:

\begin{center}
\resizebox{\linewidth}{!}{%
\begin{tikzpicture}[x=1cm, y=1cm]
  \draw[fill=blokkirmizi, draw=black, thick] (0,0) rectangle (1.5,1.0);
  \node[font=\large\ttfamily\bfseries] at (0.75, 0.5) {1};
  \node[font=\scriptsize] at (0.75, -0.3) {İşaret};
  \draw[fill=blokmavi, draw=black, thick] (1.5,0) rectangle (5.5,1.0);
  \node[font=\large\ttfamily] at (3.5, 0.5) {10000011};
  \node[font=\scriptsize] at (3.5, -0.3) {Üs = 131};
  \draw[fill=blokyesil, draw=black, thick] (5.5,0) rectangle (16.0,1.0);
  \node[font=\large\ttfamily] at (10.75, 0.5) {01110100000000000000000};
  \node[font=\scriptsize] at (10.75, -0.3) {Kesir};
\end{tikzpicture}}
\end{center}

\begin{itemize}
\item İşaret: 1 (eksi)
\item Üs: $(10000011)_2 = 131$, gerçek üs: $131 - 127 = 4$
\item Kesir: $0{,}011101$
\end{itemize}

$(-1)^1 \times (1 + 0{,}011101) \times 2^4 = -1{,}011101 \times 2^4 = -10111{,}01$

\begin{center}
$2^4 + 2^2 + 2^1 + 2^0 + 2^{-2} = 16 + 4 + 2 + 1 + 0{,}25 = 23{,}25$
\end{center}

Sonuç: $\mathbf{-23{,}25}$
\end{ornek}

\subsection{Kayan Nokta Özel Değerleri}\index{kayan nokta!özel değerler}

IEEE 754'te üs alanının tüm bitleri 0 ($00000000$) veya tüm bitleri 1 ($11111111$) olan durumlar olağan sayıları değil, \emph{özel değerleri} temsil eder. Bu özel değerler Tablo~\ref{tab:ozel-degerler-kodlama}'da özetlenmiştir.

\begin{table}[htbp]
\centering
\caption{IEEE 754 özel değerlerinin kodlanması (tek duyarlı)}
\label{tab:ozel-degerler-kodlama}
\small
\begin{tabularx}{\textwidth}{l c c X}
\toprule
\textbf{Değer} & \textbf{Üs} & \textbf{Kesir} & \textbf{Açıklama} \\
\midrule
$+0$ / $-0$ & $00000000$ & $= 0$ & İşaret bitine göre $+0$ veya $-0$. Karşılaştırmada eşit kabul edilir. \\[3pt]
Olağanlaştırılmamış\index{olağanlaştırılmamış sayı}\index{denormalized number|see{olağanlaştırılmamış sayı}} & $00000000$ & $\neq 0$ & Örtük ``1'' varsayılmaz; $(-1)^i \times 0{,}\text{Kesir} \times 2^{-126}$. Sıfıra çok yakın küçük sayıları gösterir. \\[3pt]
$+\infty$ / $-\infty$\index{sonsuz}\index{infinity|see{sonsuz}} & $11111111$ & $= 0$ & İşaret bitine göre artı veya eksi sonsuz. \\[3pt]
NaN\index{NaN} & $11111111$ & $\neq 0$ & ``Sayı değil'' (\emph{Not a Number}). Tanımsız işlem sonucu. \\
\bottomrule
\end{tabularx}
\end{table}

\subsubsection{Sıfır}

Kayan nokta formülünü hatırlayalım: $(-1)^i \times (1 + \text{kesir}) \times 2^{(e - 127)}$. Formüldeki $1 + \text{kesir}$ ifadesine dikkat edelim: kesir alanının tüm bitleri sıfır olsa bile anlamlı kısım $1{,}000\ldots0 = 1{,}0$ olur; sıfır değil. Üs alanını ne yaparsak yapalım, $1{,}0 \times 2^e$ asla sıfır olamaz; çünkü $2^e$ her zaman artı bir sayıdır. Kısacası, olağanlaştırılmış gösterimde sıfırı temsil etmenin \emph{hiçbir yolu yoktur}; gizli 1, sayının sıfır olmasını engeller.

Bu nedenle IEEE 754, sıfır için özel bir kodlama tanımlar: \textbf{üs alanı ve kesir alanının tamamı sıfırdır}. Donanım bu bit örüntüsünü gördüğünde olağan formülü uygulamaz; sonucu doğrudan sıfır olarak yorumlar.

İşaret bitine göre $+0$ ve $-0$ olmak üzere iki sıfır vardır. Matematiksel olarak $+0 = -0$'dır ve IEEE 754'e göre karşılaştırmada eşit kabul edilir. İki ayrı sıfırın bulunmasının nedeni, bazı sayısal çözümleme algoritmalarında sıfıra ``hangi yönden'' yaklaşıldığının önemli olmasıdır; örneğin $1 / (+0) = +\infty$ iken $1 / (-0) = -\infty$ olur.

\subsubsection{Olağanlaştırılmamış Sayılar}

Üs alanı 0 olduğunda ve kesir alanı 0 olmadığında sayı \emph{olağanlaştırılmamış} (\emph{denormalized} veya \emph{subnormal}) kabul edilir. Bu durumda gizli 1 varsayılmaz; sayı $(-1)^i \times 0{,}\text{Kesir} \times 2^{-126}$ biçiminde yorumlanır. Bu sayede sıfıra çok yakın değerler (duyarlık kaybına rağmen) gösterilebilir. Tek duyarlıda en küçük pozitif olağanlaştırılmamış sayı $2^{-149} \approx 1{,}4 \times 10^{-45}$'tir; olağanlaştırılmamış sayılar olmasaydı sıfır ile en küçük olağan sayı ($2^{-126} \approx 1{,}2 \times 10^{-38}$) arasında temsil edilemeyen büyük bir boşluk kalırdı.

\subsubsection{Sonsuz}

Üs alanının tüm bitleri 1 ve kesir alanı 0 olduğunda sayı sonsuz olarak yorumlanır. İşaret bitine göre $+\infty$ ve $-\infty$ vardır. Sonsuz, taşma durumlarında otomatik olarak üretilir; örneğin çok büyük iki sayının çarpımı gösterilebilir aralığı aştığında sonuç $\pm\infty$ olur. Ayrıca sıfırdan farklı bir sayının sıfıra bölünmesi de sonsuz üretir: $5{,}0 \div 0{,}0 = +\infty$.

Sonsuz değerlerle aritmetik işlemler tanımlıdır: $n \div \infty = 0$, $\infty + \infty = \infty$, $\infty \times \infty = \infty$ vb. Bu sayede programlar taşma sonrasında bile çalışmaya devam edebilir.

\subsubsection{NaN (Sayı Değil)}

$0 \div 0$ işleminin sonucu nedir? Sonsuz olamaz; çünkü sonsuz, sıfırdan \emph{farklı} bir sayının sıfıra bölünmesiyle tanımlanır. Sıfır da olamaz; çünkü $0 \times 0 = 0$ olsa da bölümün tek bir değeri yoktur; herhangi bir sayı çarpı sıfır sıfır eder. Sonuç matematiksel olarak \emph{belirsizdir}. IEEE 754 bu tür belirsiz sonuçları temsil etmek için özel bir değer tanımlar: \textbf{NaN}\index{NaN} (\emph{Not a Number}, sayı değil).

NaN, üs alanının tüm bitleri 1 ve kesir alanı 0'dan farklı olduğunda kodlanır. Matematiksel olarak tanımsız olan her işlem NaN üretir:

\begin{center}
\small
\begin{tabular}{l l}
\toprule
\textbf{İşlem} & \textbf{Neden tanımsız?} \\
\midrule
$0 \div 0$ & Bölüm belirsiz \\
$\infty - \infty$ & Fark belirsiz \\
$\infty \div \infty$ & Oran belirsiz \\
$\infty \times 0$ & Çarpım belirsiz \\
$\sqrt{-1}$ & Gerçek sayılarda tanımsız \\
\bottomrule
\end{tabular}
\end{center}

NaN'ın en dikkat çekici özelliği \emph{bulaşıcı} olmasıdır: NaN içeren herhangi bir aritmetik işlemin sonucu yine NaN olur. Örneğin $\text{NaN} + 5 = \text{NaN}$, $\text{NaN} \times 0 = \text{NaN}$. Bu bulaşıcılık, programda bir yerde oluşan tanımsız sonucun tüm hesaplamaya sessizce yayılmasına neden olabilir.

Bir diğer sürpriz: NaN kendisine eşit değildir. IEEE 754'e göre $\text{NaN} \neq \text{NaN}$ her zaman \emph{doğrudur} ve $\text{NaN} = \text{NaN}$ her zaman \emph{yanlıştır}. Bu kural birçok programcı için beklenmediktir; ancak NaN'ı saptamanın en kolay yolunu sağlar: bir değişken kendisine eşit değilse NaN'dır.

\begin{bilginot}[NaN ile Programlama Uygulaması]
NaN, program çökmesine veya kural dışı duruma yol açmaz; program çalışmaya devam eder. Bu durum onu tehlikeli kılar: hata sessizce oluşur ve sonuçlara yayılır. Aşağıdaki C ve Python örneklerinde NaN'ın nasıl ortaya çıktığı ve ekrana nasıl yazdırıldığı görülmektedir:

\begin{lstlisting}[language=C, xleftmargin=2em]
#include <math.h>
#include <stdio.h>
int main() {
    double x = 0.0 / 0.0;    // NaN üretilir
    printf("%f\n", x);        // Çıktı: NaN
    printf("%d\n", x == x);   // Çıktı: 0 (yanlış!)
    return 0;
}
\end{lstlisting}

\begin{lstlisting}[language=Python, xleftmargin=2em]
x = float('nan')
print(x)          # Çıktı: NaN
print(x == x)     # Çıktı: False
print(x + 5)      # Çıktı: NaN  (bulaşıcılık)
\end{lstlisting}

Görüldüğü gibi program hata vermez; \texttt{nan} değeri ekrana yazdırılır ve hesaplamalar devam eder. Bu nedenle sayısal hesaplamalar yapan programlarda NaN denetimi önemlidir; C'de \texttt{isnan()} işlevi, Python'da \texttt{math.isnan()} işlevi bu amaçla kullanılır.
\end{bilginot}

Bu özel değerlerle yapılan işlemlerin tam listesi Tablo~\ref{tab:ozel-degerler}'de özetlenmiştir.

\begin{table}[htbp]
\centering
\caption[Özel değerlerle yapılan işlemler]{Kayan nokta özel değerleriyle yapılan işlemlerin sonuçları.}
\label{tab:ozel-degerler}
\small
\begin{tabular}{l c l}
\toprule
\textbf{İşlem} & \textbf{Sonuç} & \textbf{Açıklama} \\
\midrule
$n \div \pm\infty$ & $0$ & Sonlu sayı sonsuza bölününce sıfıra yakınsar \\
$\pm\infty \times \pm\infty$ & $\pm\infty$ & İki sonsuzun çarpımı sonsuz \\
$n \div 0$ \; ($n \neq 0$) & $\pm\infty$ & Sıfıra bölme sonsuz üretir \\
$(+\infty) + (+\infty)$ & $+\infty$ & Aynı yönlü sonsuzlar toplanabilir \\
\midrule
$0 \div 0$ & NaN & Belirsiz \\
$(+\infty) - (+\infty)$ & NaN & Belirsiz \\
$\infty \div \infty$ & NaN & Belirsiz \\
$\infty \times 0$ & NaN & Belirsiz \\
\bottomrule
\end{tabular}
\end{table}

\subsection{Yuvarlama}

Kayan nokta aritmetiğinde her işlemin sonucu, saklama alanının izin verdiğinden daha fazla bit içerebilir. Örneğin iki 24 bitlik anlamlı kısmın çarpımı 48 bit üretir; ancak tek duyarlıda yalnızca 23 bit saklanabilir. Fazla bitlerin atılması gerekir; ama nasıl? İşte bu soruya yanıt veren mekanizma \emph{yuvarlama}\index{yuvarlama}\index{rounding|see{yuvarlama}}dır.

\subsubsection{IEEE 754 Yuvarlama Yöntemleri}

IEEE 754 dört yuvarlama yöntemi tanımlar:

\begin{enumerate}
\item \textbf{En yakın çifte yuvarlama\index{en yakın çifte yuvarlama}\index{round to nearest even|see{en yakın çifte yuvarlama}} (\emph{round to nearest, ties to even}):} \textbf{Varsayılan} yöntemdir. Sayı, en yakınındaki gösterilebilir değere yuvarlanır. Eğer sayı tam iki gösterilebilir değerin ortasındaysa, son biti \emph{çift} (0) olan değere yuvarlanır.

\item \textbf{Sıfıra yuvarlama (\emph{round toward zero}):} Sayının mutlak değeri küçültülerek sıfıra doğru yuvarlanır. Bu yönteme ``kırpma'' (\emph{truncation}) da denir; fazla bitler basitçe atılır.

\item \textbf{Yukarıya yuvarlama (\emph{round toward $+\infty$}):} Sayı $+\infty$ yönünde, her zaman büyük değere yuvarlanır.

\item \textbf{Aşağıya yuvarlama (\emph{round toward $-\infty$}):} Sayı $-\infty$ yönünde, her zaman küçük değere yuvarlanır.
\end{enumerate}

\subsubsection{Neden ``En Yakın Çifte'' Varsayılandır?}

Yuvarlama yönteminin seçimi, binlerce işlem art arda yapıldığında hataların birikmesini doğrudan etkiler. Bunu onluk tabanda basit bir örnekle görelim.

Diyelim ki onluk tabanda yalnızca 3 basamak saklayabiliyoruz ve aşağıdaki dört sayıyı toplamamız gerekiyor:

\begin{center}
$1{,}235 + 1{,}245 + 1{,}255 + 1{,}265$
\end{center}

Gerçek toplam $= 5{,}000$. Şimdi her sayıyı önce 3 basamağa yuvarlayıp sonra toplayalım:

\begin{center}
\small
\begin{tabular}{l c c c}
\toprule
\textbf{Yöntem} & \textbf{Yuvarlanan değerler} & \textbf{Toplam} & \textbf{Hata} \\
\midrule
Her zaman yukarı & $1{,}24 + 1{,}25 + 1{,}26 + 1{,}27$ & $5{,}02$ & $+0{,}02$ \\
Her zaman aşağı & $1{,}23 + 1{,}24 + 1{,}25 + 1{,}26$ & $4{,}98$ & $-0{,}02$ \\
En yakın çifte & $1{,}24 + 1{,}24 + 1{,}26 + 1{,}26$ & $5{,}00$ & $0$ \\
\bottomrule
\end{tabular}
\end{center}

``Her zaman yukarı'' yöntemi sistematik olarak yukarı yanlılık üretir; ``her zaman aşağı'' ise aşağı yanlılık. ``En yakın çifte'' yöntemi ise tam ortadaki değerleri bazen yukarı bazen aşağı yuvarlayarak \emph{yanlılığı sıfıra yakın tutar}. İşte IEEE 754'ün bu yöntemi varsayılan yapmasının nedeni budur: uzun hesaplamalarda hata birikimini en aza indirir.

\subsubsection{Donanımda Yuvarlama: Koruma, Yuvarlama ve Yapışkan Bitler}

Yuvarlama kararını verebilmek için donanımda hesaplama sonucunun saklanabilir bitlerinin ötesinde \textbf{üç ek bit} tutulur:

\begin{itemize}
\item \textbf{Koruma biti\index{koruma biti}\index{guard bit|see{koruma biti}} (\emph{guard}, G):} Saklanan anlamlı kısmın hemen sağındaki ilk ek bit.
\item \textbf{Yuvarlama biti\index{yuvarlama biti}\index{round bit|see{yuvarlama biti}} (\emph{round}, R):} Koruma bitinin sağındaki ikinci ek bit.
\item \textbf{Yapışkan bit\index{yapışkan bit}\index{sticky bit|see{yapışkan bit}} (\emph{sticky}, S):} Yuvarlama bitinin sağında kalan \emph{tüm bitlerin} VEYA'sı. Bu bitlerin herhangi biri 1 ise yapışkan bit 1 olur; kesin değerlerinin bilinmesine gerek yoktur.
\end{itemize}

Bu üç bitin ne anlama geldiğini somutlaştırmak için tek duyarlıklı bir çarpma düşünelim. Tek duyarlıda anlamlı kısım 23 bit saklar; dolayısıyla çarpma sonucunda 23 bitten sonra gelen ilk bit koruma (G), ikinci bit yuvarlama (R), geri kalanların VEYA'sı ise yapışkan bittir (S):

\[
\texttt{1.}\;\underbrace{b_1 \; b_2 \; \cdots \; b_{23}}_{\text{saklanan 23 bit}} \;\underbrace{b_{24}}_{\text{G}} \;\underbrace{b_{25}}_{\text{R}} \;\underbrace{b_{26} \, b_{27} \, \cdots}_{\text{S = hepsinin VEYA'sı}}
\]

Saklanacak son bit $L = b_{23}$'tür. Yuvarlama kararı, \texttt{G\,R\,S} kombinasyonuna bakarak verilir. Bu üç bitin görevi, atılacak kısmın son gösterilebilir değerin \emph{yarısından} küçük mü, büyük mü, yoksa tam yarısına eşit mi olduğunu belirlemektir:

\begin{itemize}
\item \textbf{G = 0} demek, atılacak kısım yarıdan küçüktür (R ve S ne olursa olsun). $\Rightarrow$ \textbf{Kırp} (aşağı yuvarla).
\item \textbf{G = 1, RS $\neq$ 00} demek, atılacak kısım yarıdan büyüktür. $\Rightarrow$ \textbf{Yukarı yuvarla} ($L$'ye 1 ekle).
\item \textbf{G = 1, R = 0, S = 0} demek, atılacak kısım \emph{tam yarıya} eşittir. $\Rightarrow$ \textbf{Berabere!} $L$ tek ise (1) yukarı yuvarla, $L$ çift ise (0) kırp. Bu ``en yakın \emph{çifte} yuvarlama'' kuralıdır ve istatistiksel yanlılığı ortadan kaldırır.
\end{itemize}

Bunu bir örnekle görelim. Tek duyarlıda bir çarpma sonucu, olağanlaştırmadan sonra şu bit dizisini üretsin (yalnızca son bitler gösterilmiştir):

\[
\texttt{1.}\;\ldots\;\underbrace{1}_{L=b_{23}} \;\underbrace{1}_{G} \;\underbrace{0}_{R} \;\underbrace{1}_{S}
\]

G~=~1 olduğundan atılacak kısım en az yarıdır. R ve S'ye bakıyoruz: RS~=~01, dolayısıyla $\neq$~00. O halde atılacak kısım yarıdan \emph{büyüktür} $\Rightarrow$ yukarı yuvarlama yapılır: $L$'ye 1 eklenir ve saklanan son bit $L = 1 + 1 = 0$ olur (elde bir sonraki bite yayılır).

Eğer bit dizisi \texttt{\ldots\,1\;1\;0\;0} (G~=~1, R~=~0, S~=~0) olsaydı tam yarı noktası söz konusu olurdu. Bu durumda $L = 1$ tek olduğundan yine yukarı yuvarlanır ve $L$ çift (0) yapılırdı. Ancak $L = 0$ olsaydı (çift), kırpma yapılır ve $L$ olduğu gibi kalırdı.

Yapışkan bitin neden gerekli olduğunu anlamak önemlidir: yuvarlama ve koruma bitleri yalnızca iki ek bit verir; ancak gerçek sonuç bunların da ötesinde bitlere sahip olabilir. Eğer bu uzak bitlerden herhangi biri 1 ise, atılacak kısım ``tam yarı'' olamaz; mutlaka yarıdan büyüktür. İşte yapışkan bit, tüm bu uzak bitlerin tek bir VEYA işlemiyle özetlenerek bu ayrımın yapılabilmesini sağlar. Donanımda bu, uzun bir VEYA kapısı zinciriyle gerçekleştirilir ve yalnızca bir bitlik yer kaplar.

\subsubsection{Yuvarlama Hatası ve Sonuçları}

Yuvarlama, her kayan nokta işlemine küçük bir hata ekler. Bu hatanın büyüklüğü, anlamlı kısmın kaç bit sakladığına bağlıdır: tek duyarlıda yuvarlama virgülden sonraki 23.\ bitte, çift duyarlıda ise 52.\ bitte gerçekleşir. Bunu somutlaştırmak için $1 \div 3$ işlemini düşünelim; sonuç $0{,}333\ldots$ olup sonsuz tekrarlıdır ve her duyarlıkta yuvarlanmak zorundadır:

\begin{center}
\small
\begin{tabular}{l l l}
\toprule
\textbf{Duyarlık} & \textbf{Saklanan değer} & \textbf{Hata} \\
\midrule
Tek (23 bit) & $0{,}333\,333\,34\ldots$ & $\sim 10^{-8}$ (virgülden sonra 7 doğru basamak) \\
Çift (52 bit) & $0{,}333\,333\,333\,333\,333\,3\ldots$ & $\sim 10^{-16}$ (virgülden sonra 15 doğru basamak) \\
\bottomrule
\end{tabular}
\end{center}

Görüldüğü gibi duyarlılığı artırmak, yuvarlama noktasını daha uzak bir bite taşır ve her işlemdeki hatayı doğrudan azaltır. Tek bir işlemde bu hata en kötü durumda son bitin yarısı ($\frac{1}{2}\,\text{ulp}$, \emph{unit in the last place}\index{ulp|see{son birim değeri}}) kadardır ve önemsiz görünebilir; ancak binlerce işlem art arda yapıldığında bu küçük hatalar birikebilir.

\begin{bilginot}[Kayan Nokta Duyarlığı ve Günlük Etkisi]
Kayan nokta aritmetiğinde her işlemde küçük yuvarlama hataları oluşabilir. Örneğin, çoğu programlama dilinde $0{,}1 + 0{,}2 \neq 0{,}3$ sonucu alınır; çünkü $0{,}1$ ve $0{,}2$ ikilik tabanda sonsuz tekrarlı kesir olduğundan tam olarak temsil edilemez:

\begin{lstlisting}[language=Python, xleftmargin=2em]
>>> 0.1 + 0.2
0.30000000000000004
>>> 0.1 + 0.2 == 0.3
False
\end{lstlisting}

Bu tür hatalar biriktikçe bilimsel benzetimlerde, finansal hesaplamalarda ve mühendislik uygulamalarında beklenmedik sonuçlara yol açabilir. Bu nedenle sayısal çözümleme alanında hata yayılımı (\emph{error propagation}) ve kararlılık (\emph{numerical stability}) kavramları büyük önem taşır. Finansal uygulamalarda kayan nokta yerine sabit noktalı veya ondalık (\emph{decimal}) aritmetik tercih edilmesinin nedeni de budur.

Özellikle tehlikeli bir durum \emph{yıkıcı yokoluş}\index{yıkıcı yokoluş}\index{catastrophic cancellation|see{yıkıcı yokoluş}} (\emph{catastrophic cancellation}) olarak bilinir: birbirine çok yakın iki kayan nokta sayısı çıkarıldığında, anlamlı kısımdaki üst bitler birbirini götürür ve geriye kalan bitler yuvarlama hatalarından ibarettir. Örneğin $a = 1{,}00000001 \times 2^{10}$ ile $b = 1{,}00000000 \times 2^{10}$ çıkarıldığında sonuç $0{,}00000001 \times 2^{10} = 1{,}0 \times 2^{2}$ olur; burada 23 bitlik anlamlı kısımdan yalnızca 1 bit gerçek bilgi taşır, geri kalanı güvenilmezdir.
\end{bilginot}

\subsection{Kayan Nokta İşlemlerinde Toplama, Çarpma ve Bölme}

Kayan nokta gösterimindeki sayılarla da tam sayılarla olduğu gibi dört işlem yapılabilir; ancak üs ve anlamlı kısmın ayrı ayrı ele alınması gerektiğinden bu işlemler tamsayı aritmetiğinden daha karmaşıktır.

\subsubsection{Toplama ve Çıkarma}

İki kayan nokta sayısını toplamak için anlamlı kısımlarını doğrudan toplayamayız; çünkü her sayının virgülü farklı bir konumdadır (üsleri farklıdır). Önce virgülleri aynı hizaya getirmemiz, ardından toplamayı yapıp sonucu yeniden olağan biçime getirmemiz gerekir. Bu süreç beş adımdan oluşur:

\begin{enumerate}
\item \textbf{Üsleri karşılaştır:} İki sayının üsleri arasındaki fark hesaplanır. Bu fark, virgülün ne kadar kaydırılması gerektiğini belirler.

\item \textbf{Anlamlı kısmı kaydır:} Küçük üslü sayının anlamlı kısmı, üs farkı kadar \emph{sağa} kaydırılır. Böylece her iki sayı da aynı üsse sahip olur ve virgüller hizalanır. Sağa kaydırma sırasında dışarı kayan bitler koruma, yuvarlama ve yapışkan bitlerinde tutulur.

\item \textbf{Anlamlı kısımları topla veya çıkar:} Üsleri eşitlenen sayıların anlamlı kısımları, işaret bitlerine göre toplanır ya da çıkarılır.

\item \textbf{Olağanlaştır:} Sonuç $1{,}xxx \times 2^e$ biçiminde olmayabilir. Örneğin toplam $0{,}001\ldots$ biçiminde çıkarsa sola kaydırılıp üs azaltılır; $1x{,}xxx$ biçiminde çıkarsa sağa kaydırılıp üs artırılır. Bu adımda taşma veya alttan taşma da denetlenir.

\item \textbf{Yuvarla:} Sonuç, hedef duyarlığa yuvarlanır. Yuvarlama sonucu olağan biçim bozulursa (örneğin yuvarlama bir elde üretirse) 4.~adıma geri dönülür.
\end{enumerate}

Bu işlemi gerçekleştiren donanımın blok şeması Şekil~\ref{fig:kayan-nokta-toplama-donanim}'de gösterilmiştir. Dikkat edilirse tamsayı toplamada tek bir toplayıcı yeterliyken, kayan nokta toplamada karşılaştırma, kaydırma, toplama, olağanlaştırma ve yuvarlama birimlerinin hepsine ihtiyaç vardır; bu nedenle kayan nokta toplayıcısı, tamsayı toplayıcısından çok daha fazla donanım alanı kaplar ve daha yavaş çalışır.

\begin{figure}[htbp]
\centering
\begin{tikzpicture}[scale=0.82, transform shape,
    blok/.style={draw, rectangle, fill=blokmavi, thick, font=\footnotesize,
                 align=center, minimum width=5.0cm, minimum height=0.9cm,
                 rounded corners=2pt},
    girdi/.style={draw, rectangle, fill=bloksari, thick, font=\scriptsize,
                  align=center, minimum width=2.8cm, minimum height=0.7cm,
                  rounded corners=2pt},
    sonuc/.style={draw, rectangle, fill=blokturuncu, thick, font=\footnotesize,
                  align=center, minimum width=5.0cm, minimum height=0.9cm,
                  rounded corners=2pt}
]

    % ---------- Sabit dikey aralık ----------
    \pgfmathsetmacro{\aralik}{1.8}

    % ---------- Girdiler ----------
    \node[girdi] (A) at (-3.0, 6.8*\aralik)
        {{\footnotesize İşlenen A}\\{\tiny İşaret $|$ Üs $|$ Kesir}};
    \node[girdi] (B) at ( 3.0, 6.8*\aralik)
        {{\footnotesize İşlenen B}\\{\tiny İşaret $|$ Üs $|$ Kesir}};

    % ---------- Blok 1 ----------
    \node[blok] (kars) at (0, 5*\aralik)
        {1. Üsleri karşılaştır\\(üs farkını hesapla)};

    % ---------- Blok 2 ----------
    \node[blok] (kaydir) at (0, 4*\aralik)
        {2. Küçük üslü sayının anlamlı kısmını\\üs farkı kadar sağa kaydır};

    % ---------- Blok 3 ----------
    \node[blok] (topla) at (0, 3*\aralik)
        {3. Anlamlı kısımları topla\,/\,çıkar\\(işarete göre seç)};

    % ---------- Blok 4 ----------
    \node[blok] (olag) at (0, 2*\aralik)
        {4. Olağanlaştır\\(sola\,/\,sağa kaydır, üsü ayarla)};

    % ---------- Blok 5 ----------
    \node[blok] (yuvarla) at (0, 1*\aralik)
        {5. Yuvarla};

    % ---------- Sonuç ----------
    \node[sonuc] (sonuc) at (0, 0)
        {Sonuç\\{\tiny İşaret $|$ Üs $|$ Kesir}};

    % ---------- Oklar: Girdiler → Blok 1 ----------
    \draw[okstil] (A.south) -- ++(0,-0.6) -| ([xshift=-9mm]kars.north);
    \draw[okstil] (B.south) -- ++(0,-0.6) -| ([xshift= 9mm]kars.north);

    % ---------- Blok zinciri ----------
    \draw[okstil] (kars.south)    -- (kaydir.north)
        node[midway, right=4pt, font=\scriptsize] {üs farkı};
    \draw[okstil] (kaydir.south)  -- (topla.north);
    \draw[okstil] (topla.south)   -- (olag.north);
    \draw[okstil] (olag.south)    -- (yuvarla.north);
    \draw[okstil] (yuvarla.south) -- (sonuc.north);

    % ---------- Taşma geri bildirimi ----------
    \draw[denetimstil] (olag.east) -- ++(1.2,0)
        node[right, font=\footnotesize, align=left] {Taşma varsa:\\üsü artır,\\sağa kaydır};

\end{tikzpicture}
\caption{Kayan nokta toplama/çıkarma donanımının blok şeması}
\label{fig:kayan-nokta-toplama-donanim}
\end{figure}

\FloatBarrier
\begin{ornek}[5.18]
IEEE~754 tek duyarlıklı gösterimde $0{,}5 + (-0{,}375)$ toplamını hesaplayalım.

\textbf{Çözüm 5.18:}

Önce her iki sayıyı ikilik kayan nokta biçiminde yazalım:
\begin{align*}
0{,}5  &= 1{,}0 \times 2^{-1}  &&\text{(üs} = -1\text{)} \\
-0{,}375 &= -1{,}1 \times 2^{-2} &&\text{(üs} = -2\text{)}
\end{align*}

İki sayının üsleri farklı olduğundan anlamlı kısımları doğrudan toplayamayız; önce virgülleri hizalamamız gerekir.

\textbf{Adım 1: Üsleri karşılaştır.} Üs farkı $= (-1) - (-2) = 1$. Küçük üslü sayı $-1{,}1 \times 2^{-2}$'dir.

\textbf{Adım 2: Anlamlı kısmı kaydır.} Küçük üslü sayının anlamlı kısmı, üs farkı kadar (1 bit) sağa kaydırılır. Böylece her iki sayı da aynı üsse ($-1$) sahip olur:
\[
-1{,}1 \times 2^{-2} \;\longrightarrow\; -0{,}11 \times 2^{-1}
\]

\textbf{Adım 3: Anlamlı kısımları topla.} Üsler artık eşit olduğundan anlamlı kısımları toplayabiliriz:
\[
1{,}00 + (-0{,}11) = 0{,}01 \quad (\text{üs} = -1)
\]

\textbf{Adım 4: Olağanlaştır.} Sonuç $0{,}01 \times 2^{-1}$ olağan biçimde değildir; baştaki bit 1 değil. Anlamlı kısmı 2 bit sola kaydırıp üsü 2 azaltarak olağanlaştırırız:
\[
0{,}01 \times 2^{-1} \;\longrightarrow\; 1{,}0 \times 2^{-3}
\]

\textbf{Adım 5: Yuvarla.} Bu örnekte fazla bit olmadığından yuvarlama gerekmez.

\textbf{Sonuç:} $1{,}0 \times 2^{-3} = 0{,}125_{10}$. Doğrulama: $0{,}5 + (-0{,}375) = 0{,}125$. \checkmark
\end{ornek}

\subsubsection{Çarpma}

Kayan nokta çarpması, toplamadan farklı olarak üslerin eşitlenmesini gerektirmez; bunun yerine üsler toplanır, anlamlı kısımlar çarpılır. Adımlar şu şekildedir:

\begin{enumerate}
\item \textbf{Özel durumları denetle:} İşlenenlerin sıfır, sonsuz veya NaN olup olmadığı denetlenir. Özel durum varsa sonuç doğrudan üretilir ve sonraki adımlar atlanır: $0 \times \infty = \text{NaN}$, $0 \times n = 0$, $\infty \times n = \pm\infty$ (sıfır hariç). Böylece gereksiz çarpma işlemi engellenir.

\item \textbf{İşareti belirle:} İki sayının işaret bitlerine Dışlayan VEYA uygulanır. İşaretler aynıysa sonuç artı, farklıysa eksidir.

\item \textbf{Üsleri topla:} Gerçek üsler toplanır. Eklentili gösterimde her iki üste de referans değeri (tek duyarlıda 127) eklenmiş olduğundan, eklentili üsler toplandığında referans iki kez sayılmış olur; bu nedenle referans bir kez çıkarılır:
\[
\text{Sonuç üssü} = (E_A + E_B) - 127
\]

\item \textbf{Anlamlı kısımları çarp:} İki sayının anlamlı kısımları ($1{,}f_A$ ve $1{,}f_B$) doğrudan çarpılır. $n$ bitlik iki anlamlı kısmın çarpımı $2n$ bit üretir; fazla bitler yuvarlama adımında kesilecektir.

\item \textbf{Olağanlaştır:} Çarpım $1x{,}xxx$ biçiminde çıkarsa (baştaki 1'in solunda fazladan bir bit varsa) anlamlı kısım sağa bir kaydırılıp üs 1 artırılır. Bu adımda taşma veya alttan taşma da denetlenir.

\item \textbf{Yuvarla:} Sonuç hedef duyarlığa yuvarlanır. Toplamada olduğu gibi, yuvarlama olağan biçimi bozarsa olağanlaştırma tekrarlanır.
\end{enumerate}

Kayan nokta çarpmasının akış şeması Şekil~\ref{fig:kayan-nokta-carpma-akis}'te gösterilmiştir.

\begin{figure}[htbp]
\centering
\begin{tikzpicture}[scale=0.72, transform shape,
    blok/.style={draw, rectangle, fill=blokmavi, thick, font=\footnotesize,
                 align=center, minimum width=3.2cm, minimum height=1cm},
    karar/.style={draw, diamond, fill=bloksari, thick, font=\footnotesize,
                  align=center, inner sep=2pt, aspect=2.2},
    okstil2/.style={-{Stealth[length=2.5mm]}, thick}]

    \node[blok] (isaret) at (0, 8) {\.{I}\c{s}aret bitlerini\\D{\i}\c{s}layan VEYA'la};
    \node[blok] (ust) at (0, 6) {\"{U}stleri topla,\\referans{\i} \c{c}{\i}kar};
    \node[blok] (kesir) at (0, 4) {Anlaml{\i} k{\i}s{\i}mlar{\i}\\\c{c}arp};
    \node[karar] (normal) at (0, 2) {Ola\u{g}an\\bi\c{c}imde mi?};
    \node[blok] (kaydir) at (4.5, 2) {Sa\u{g}a kayd{\i}r,\\\"{u}st{\"u} 1 art{\i}r};
    \node[blok] (yuvarla) at (0, 0) {Yuvarla};
    \node[karar] (tasma) at (0, -2) {Ta\c{s}ma?};
    \node[blok] (sonuc) at (0, -4) {Sonu\c{c}};
    \node[blok] (istisna) at (4.5, -2) {\.{I}stisna};

    \draw[okstil2] (isaret) -- (ust);
    \draw[okstil2] (ust) -- (kesir);
    \draw[okstil2] (kesir) -- (normal);
    \draw[okstil2] (normal) -- node[above, font=\scriptsize] {Hay{\i}r} (kaydir);
    \draw[okstil2] (normal) -- node[right, font=\scriptsize] {Evet} (yuvarla);
    \draw[okstil2] (kaydir) |- (yuvarla);
    \draw[okstil2] (yuvarla) -- (tasma);
    \draw[okstil2] (tasma) -- node[right, font=\scriptsize] {Hay{\i}r} (sonuc);
    \draw[okstil2] (tasma) -- node[above, font=\scriptsize] {Evet} (istisna);

\end{tikzpicture}
\caption{Kayan nokta \c{c}arpma i\c{s}leminin ak{\i}\c{s} \c{s}emas{\i}}
\label{fig:kayan-nokta-carpma-akis}
\end{figure}


\FloatBarrier
\begin{ornek}[5.19]
IEEE~754 tek duyarlıklı gösterimde $1{,}5 \times 2^{1}$ ile $1{,}25 \times 2^{1}$ sayılarını çarpalım.

\textbf{Çözüm 5.19:}

Önce her iki sayıyı ikilik kayan nokta biçiminde yazalım:
\begin{align*}
A &= 1{,}5 \times 2^{1} = 1{,}1_2 \times 2^{1} &&\text{(onluk değer: } 3\text{)} \\
B &= 1{,}25 \times 2^{1} = 1{,}01_2 \times 2^{1} &&\text{(onluk değer: } 2{,}5\text{)}
\end{align*}

Çarpmada toplamadan farklı olarak üsleri eşitlemeye gerek yoktur; üsler toplanır, anlamlı kısımlar ayrı ayrı çarpılır.

\textbf{Adım 1: Özel durumları denetle.} Her iki işlenen de olağan sayı; özel durum yok, devam.

\textbf{Adım 2: İşareti belirle.} Her iki sayının da işaret biti 0'dır (artı). Dışlayan VEYA: $0 \oplus 0 = 0$, dolayısıyla sonuç artıdır.

\textbf{Adım 3: Üsleri topla.} Gerçek üsleri doğrudan toplayabiliriz: $1 + 1 = 2$. Eklentili gösterimde ise her iki üste de referans (127) eklenmiş olduğundan, eklentili üsler toplandığında referans iki kez sayılır; bu nedenle bir kez çıkarılır:
\[
E_\text{sonuç} = (1 + 127) + (1 + 127) - 127 = 128 + 128 - 127 = 129 \quad (\text{gerçek üs} = 129 - 127 = 2)
\]

\textbf{Adım 4: Anlamlı kısımları çarp.} İkilik çarpmayı adım adım yapalım:
\[
\begin{array}{r}
  1{,}10 \\
\times\; 1{,}01 \\
\hline
  110 \quad\text{(}1{,}10 \times 1\text{)}\\
 0000 \quad\text{(}1{,}10 \times 0\text{, 1 bit sola)}\\
1100 \quad\text{(}1{,}10 \times 1\text{, 2 bit sola)}\\
\hline
1{,}1110
\end{array}
\]

\textbf{Adım 5: Olağanlaştır.} Çarpım $1{,}1110_2$ zaten $1{,}xxx$ biçimindedir; baştaki bit 1 olduğundan olağanlaştırma gerekmez. Eğer çarpım $1x{,}xxx$ biçiminde çıksaydı, sağa bir kaydırıp üsü 1 artıracaktık.

\textbf{Adım 6: Yuvarla.} Çarpım 4 bitlik anlamlı kısma tam olarak sığdığından yuvarlama gerekmez. Gerçek donanımda 23 bitin ötesindeki bitler koruma, yuvarlama ve yapışkan bitleri kullanılarak yuvarlanırdı.

\textbf{Sonuç:} $1{,}1110_2 \times 2^{2} = 111{,}10_2 = 7{,}5$. Doğrulama: $3 \times 2{,}5 = 7{,}5$. \checkmark
\end{ornek}

\subsubsection{Bölme}

Kayan nokta bölmesi çarpmaya benzer bir yapıda ilerler: üsler ayrı ayrı işlenir, anlamlı kısımlar bölünür. Ancak çarpmada üsler toplanırken bölmede \emph{çıkarılır}; anlamlı kısımlar ise çarpılmak yerine bölünür. Bölme adımları şu şekildedir:

\begin{enumerate}
\item \textbf{Özel durumları denetle:} İşlenenlerin sıfır, sonsuz veya NaN olup olmadığı denetlenir. Özel durum varsa sonuç doğrudan üretilir ve sonraki adımlar atlanır: $n \div 0 = \pm\infty$, $0 \div 0 = \text{NaN}$, $\infty \div \infty = \text{NaN}$. Bölen sıfırsa bölme işlemi hiç yapılmaz.

\item \textbf{İşareti belirle:} Çarpmadaki gibi, iki sayının işaret bitlerine Dışlayan VEYA uygulanır. İşaretler aynıysa sonuç artı, farklıysa eksidir.

\item \textbf{Üsleri çıkar:} Bölünenin üssünden bölenin üssü çıkarılır. Eklentili gösterimde referans değeri her iki üste de eklenmiş olduğundan, çıkarma sonucunda referans kaybolur ve geri eklenmesi gerekir:
\[
E_\text{sonuç} = (E_A - E_B) + 127
\]

\item \textbf{Anlamlı kısımları böl:} Bölünenin anlamlı kısmı ($1{,}f_A$) bölenin anlamlı kısmına ($1{,}f_B$) bölünür. Bu bölme, daha önce gördüğümüz tam sayı bölme algoritmalarından (geri yüklemeli, geri yüklemesiz veya SRT) biriyle gerçekleştirilir.

\item \textbf{Olağanlaştır:} Bölüm $1{,}xxx$ biçiminde olmayabilir; örneğin $1{,}0 \div 1{,}1 = 0{,}1\overline{10}$ gibi bir sonuç çıkarsa sola kaydırılıp üs 1 azaltılır.

\item \textbf{Yuvarla:} Sonuç hedef duyarlığa yuvarlanır. Yuvarlama olağan biçimi bozarsa olağanlaştırma tekrarlanır.
\end{enumerate}

\begin{ornek}[5.20]
IEEE~754 tek duyarlıklı gösterimde $-7{,}0 \div 2{,}0$ bölmesini hesaplayalım.

\textbf{Çözüm 5.20:}

Önce her iki sayıyı ikilik kayan nokta biçiminde yazalım:
\begin{align*}
A &= -7{,}0 = -1{,}11_2 \times 2^{2} &&\text{(üs} = 2\text{)} \\
B &= 2{,}0 = 1{,}0_2 \times 2^{1} &&\text{(üs} = 1\text{)}
\end{align*}

\textbf{Adım 1: Özel durumları denetle.} Her iki işlenen de olağan sayı; özel durum yok, devam.

\textbf{Adım 2: İşareti belirle.} $A$ eksi, $B$ artı. Dışlayan VEYA: $1 \oplus 0 = 1$, dolayısıyla sonuç eksidir.

\textbf{Adım 3: Üsleri çıkar.} Gerçek üsler: $2 - 1 = 1$. Eklentili gösterimde:
\[
E_\text{sonuç} = (2 + 127) - (1 + 127) + 127 = 129 - 128 + 127 = 128 \quad (\text{gerçek üs} = 128 - 127 = 1)
\]

\textbf{Adım 4: Anlamlı kısımları böl.}
\[
1{,}11_2 \div 1{,}0_2 = 1{,}11_2
\]
Bölen $1{,}0$ olduğundan bölüm doğrudan bölünenin anlamlı kısmıdır.

\textbf{Adım 5: Olağanlaştır.} Sonuç $1{,}11_2$ zaten $1{,}xxx$ biçiminde olduğundan olağanlaştırma gerekmez.

\textbf{Adım 6: Yuvarla.} Bu örnekte yuvarlama gerekmez.

\textbf{Sonuç:} $-1{,}11_2 \times 2^{1} = -11{,}1_2 = -3{,}5$. Doğrulama: $-7{,}0 \div 2{,}0 = -3{,}5$. \checkmark
\end{ornek}

% =============================================
\section{RISC-V Kayan Nokta Buyrukları}
\label{sec:riscv-kayan-nokta}

Kayan nokta gösterimini öğrendiğimize göre, şimdi RISC-V'in bu sayılarla nasıl işlem yaptığına bakalım. Kayan nokta desteğinin buyruk kümesine nasıl eklendiği, mimariden mimariye farklılık gösterir. x86 mimarisinde kayan nokta işlemleri başlangıçta ayrı bir yardımcı işlemci yongasıyla (Intel 8087) destekleniyordu; daha sonra SSE ve AVX uzantılarıyla ana işlemciye tümleştirildi. ARM mimarisinde ise kayan nokta başlangıçta isteğe bağlı bir VFP (\emph{Vector Floating-Point}) birimi olarak sunulmuş, ARMv8 ile birlikte zorunlu hale getirilmiştir; günümüzde NEON/SVE gibi SIMD uzantıları da kayan nokta işlemlerini destekler.

RISC-V bu yaklaşımlardan farklı olarak kayan nokta desteğini baştan modüler uzantılar biçiminde tasarlamıştır: \emph{F} uzantısı tek duyarlı, \emph{D} uzantısı çift duyarlı işlemleri ekler. Ayrı bir yardımcı işlemci yoktur; kayan nokta buyrukları doğrudan ana buyruk kümesinin uzantıları olarak tanımlanmıştır. İşlemci, tam sayı yazmaçlarının (\texttt{x0}--\texttt{x31}) yanı sıra \texttt{f0}--\texttt{f31} arasında 32 adet kayan nokta yazmacı içerir. F uzantısında bu yazmaçlar 32 bit, D uzantısında ise 64 bit genişliğindedir. Bu uzantıların sağladığı buyruklar Tablo~\ref{tab:kayan-nokta-buyruklar}'de verilmiştir.

\begin{table}[htbp]
\centering
\caption{RISC-V Kayan Nokta Buyrukları (F/D Uzantıları)}
\label{tab:kayan-nokta-buyruklar}
\small
\begin{tabularx}{\textwidth}{l l X}
\toprule
\textbf{Buyruk} & \textbf{Söz Dizimi} & \textbf{Açıklama} \\
\midrule
\multicolumn{3}{l}{\textit{Yükle/Sakla Buyrukları}} \\
\midrule
Tek Duyarlı Yükle & \texttt{flw fd, offset(rs1)} & Bellekten 4 bayt yükler (tek duyarlı) \\
Çift Duyarlı Yükle & \texttt{fld fd, offset(rs1)} & Bellekten 8 bayt yükler (çift duyarlı) \\
Tek Duyarlı Sakla & \texttt{fsw fs, offset(rs1)} & Yazmaçtan belleğe 4 bayt saklar (tek duyarlı) \\
Çift Duyarlı Sakla & \texttt{fsd fs, offset(rs1)} & Yazmaçtan belleğe 8 bayt saklar (çift duyarlı) \\
\midrule
\multicolumn{3}{l}{\textit{Hesaplama Buyrukları}} \\
\midrule
Topla & \texttt{fadd.s / fadd.d fd, fs1, fs2} & Kayan nokta toplama \\
Çıkar & \texttt{fsub.s / fsub.d fd, fs1, fs2} & Kayan nokta çıkarma \\
Çarp & \texttt{fmul.s / fmul.d fd, fs1, fs2} & Kayan nokta çarpma \\
Böl & \texttt{fdiv.s / fdiv.d fd, fs1, fs2} & Kayan nokta bölme \\
\midrule
\multicolumn{3}{l}{\textit{Karşılaştırma Buyrukları}} \\
\midrule
Eşitlik Karşılaştır & \texttt{feq.s rd, rs1, rs2} & Kayan nokta eşitlik karşılaştırması; eşitse rd=1, değilse rd=0 \\
Küçüktür Karşılaştır & \texttt{flt.s rd, rs1, rs2} & Kayan nokta küçüktür karşılaştırması; rs1 $<$ rs2 ise rd=1, değilse rd=0 \\
Küçük Eşittir Karşılaştır & \texttt{fle.s rd, rs1, rs2} & Kayan nokta küçük eşittir karşılaştırması; rs1 $\leq$ rs2 ise rd=1, değilse rd=0 \\
\midrule
\multicolumn{3}{l}{\textit{Dönüşüm ve Taşıma Buyrukları}} \\
\midrule
KN'den Tam Sayıya & \texttt{fcvt.w.s rd, rs1} & Kayan noktadan tam sayıya dönüşüm (yuvarlayarak) \\
Tam Sayıdan KN'ye & \texttt{fcvt.s.w rd, rs1} & Tam sayıdan kayan noktaya dönüşüm \\
FP Yazmacından Kopyala & \texttt{fmv.x.w rd, rs1} & Kayan nokta yazmacındaki bitleri tam sayı yazmacına kopyalar (değer dönüşümü yapmaz) \\
Tam Sayı Yazmacından Kopyala & \texttt{fmv.w.x rd, rs1} & Tam sayı yazmacındaki bitleri kayan nokta yazmacına kopyalar (değer dönüşümü yapmaz) \\
\bottomrule
\end{tabularx}
\end{table}


% =============================================
\FloatBarrier
\begin{bilginot}[Intel Pentium FDIV Hatası]
\label{sec:fdiv-hatasi}
Kayan nokta aritmetiği hassas bir konudur; küçük bir tasarım hatası bile büyük sonuçlara yol açabilir. Bilgisayar tarihinin en ünlü donanım hatalarından biri de tam olarak bu alandadır.

İlk dönem işlemcilerde kayan nokta birimi ayrı bir yonga olarak satılırdı. Intel, 486DX işlemcisiyle kayan nokta birimini ana işlemcinin içine yerleştirdi; ardından da Pentium\index{Intel Pentium} işlemcisini geliştirdi. İki işlemci de kayan nokta bölmesi yaparken SRT bölme algoritmasını (bkz.~Bölüm~\ref{sec:srt-bolme}) kullanarak her adımda 2 bölüm biti birden üretiyordu. Her adımda doğru bölüm rakamını seçebilmek için bir SRT arama tablosuna ihtiyaç vardı. Sorun, bu tablodaki 1066 hücreden beşinin yanlış doldurulmasıydı. Bu hatalı işlemlerden\index{FDIV hatası} en yaygını $4195835 / 3145727$ işleminin sonucunun $1{,}33382$ yerine $1{,}33374$ bulunmasıydı.

Bu durumu 1994 Kasım ayında ilk fark eden kişi Lynchburg Kolejinden Thomas Nicely\index{Nicely, Thomas} oldu. Intel'e haber vermesine rağmen yeterince ilgilenilmeyince haberi internette duyurdu. IBM, Pentium işlemcili bilgisayarlarının üretimini durdurdu. Intel yaklaşık 300 milyon dolar zarar gördü. Bölme buyruğu FDIV olduğu için bu hataya \emph{FDIV hatası}\index{FDIV hatası} denir.
\end{bilginot}

\begin{bilginot}[Patriot Füzesi Kazası: Yuvarlama Hatasının Bedeli]
\index{Patriot füzesi kazası}
FDIV hatası maddi zarara yol açtı; ancak kayan nokta hataları can kaybına da neden olabilir. 1991 Körfez Savaşı'nda ABD ordusunun Patriot hava savunma sistemi, Suudi Arabistan'ın Dhahran kentinde bir Irak Scud füzesini durduramadı; füze bir askeri barınağa isabet etti ve 28 Amerikan askeri hayatını kaybetti.

Sorunun kaynağı, sistemin dahili saatindeki kayan nokta yuvarlama hatasıydı. Patriot'un bilgisayarı zamanı 0{,}1 saniyelik artışlarla sayıyordu; ancak $0{,}1$ ikilik tabanda sonsuz tekrarlı bir kesirdir ($0{,}0\overline{0011}_2$) ve 24 bitlik sabit noktalı gösterimde her adımda yaklaşık $9{,}5 \times 10^{-8}$ saniyelik bir hata oluşuyordu. Sistem 100 saatten fazla kesintisiz çalıştığında bu küçük hatalar birikti ve toplam zamanlama kayması yaklaşık $0{,}34$ saniyeye ulaştı. Bu sürede bir Scud füzesi yaklaşık 600 metre yol alır; sistem, füzeyi aradığı yerde bulamadı ve önleme başarısız oldu.

Bu olay, kayan nokta yuvarlama hatalarının ``küçük'' ve ``önemsiz'' olmadığının trajik bir kanıtıdır. Biriken hatalar, yeterli süre geçtiğinde yaşamsal sonuçlara yol açabilir.
\end{bilginot}

% =============================================
\section{Yapay Zeka Çağında Kayan Nokta Biçimleri}
\label{sec:yapay-zeka-kayan-nokta}
\index{kayan nokta!yapay zeka biçimleri}

IEEE 754'ün tek duyarlı (FP32) biçimi on yıllarca kayan nokta hesaplamalarının standart birimi oldu. Ancak 2010'ların ortasından itibaren derin öğrenme modelleri büyüdükçe FP32'nin maliyeti sorun olmaya başladı: milyarlarca parametreli bir modelin her parametresi 32 bit kapladığında, hem bellek gereksinimi hem de veri aktarım bant genişliği devasa boyutlara ulaşır. Üstelik yapay zeka eğitiminde her parametrenin 7 onluk basamak duyarlıkla saklanması çoğu zaman gereksizdir; ağırlık güncellemeleri zaten gürültülüdür ve birkaç basamak duyarlık yeterlidir.

Bu gözlem, araştırmacıları daha dar kayan nokta biçimleri geliştirmeye yöneltti. Temel fikir yalındır: \emph{her parametrenin kapladığı bit sayısını yarıya indirirsen, aynı belleğe iki kat daha fazla parametre sığar ve aynı bant genişliğiyle iki kat daha fazla veri taşırsın.} Ayrıca dar çarpıcılar daha az alan kaplar ve daha az enerji tüketir; dolayısıyla aynı yonga üzerine çok daha fazla çarpma-toplama birimi yerleştirilebilir.

\subsection{FP16: IEEE 754 Yarım Duyarlı}
\index{FP16}\index{yarım duyarlı|see{FP16}}

IEEE 754-2008 ile standartlaştırılan yarım duyarlı (\emph{half precision}) biçim 16 bitten oluşur: 1 işaret $+$ 5 üs $+$ 10 kesir. Eklenti değeri 15'tir; gerçek üs aralığı $-14$ ile $+15$'tir.

FP16 başlangıçta grafik uygulamaları (doku sıkıştırma, renk hesaplama) için tasarlanmıştı. Yapay zeka alanında da ilk dar biçim olarak yaygınlaştı; NVIDIA'nın 2016'da çıkardığı Pascal mimarisi (P100 GPU) FP16 desteğini yapay zeka eğitimi için öne çıkaran ilk büyük ölçekli donanımdı.

Ancak FP16'nın ciddi bir sınırlaması vardır: 5 bitlik üs alanı yalnızca $-14$ ile $+15$ arasında gerçek üs değerlerine izin verir; gösterilebilen en büyük sonlu değer yaklaşık $65\,504$'tür. Derin öğrenme eğitiminde ağırlık güncellemeleri sırasında gradyanlar ve kayıp değerleri bu aralığı kolayca aşabilir; bu durumda taşma (\emph{overflow}) oluşur ve eğitim kararsızlaşır.

\subsection{BF16: Google Brain'den Doğan Biçim}
\index{BF16}\index{beyin kayan nokta|see{BF16}}\index{Google Brain}

2018 yılında Google Brain ekibi, kendi TPU (\emph{Tensor Processing Unit}) hızlandırıcıları için yeni bir 16 bitlik biçim tasarladı: \textbf{BF16} (\emph{Brain Floating Point 16}). Tasarımın çıkış noktası şu sorunun yanıtıydı: \emph{``FP32'nin bit bütçesini yarıya indirirken hangi alandan kısmalıyız: üsten mi, kesirden mi?''}

FP16, üs alanını 5 bite indirerek sayı aralığını daraltmıştı ve bu taşma sorunlarına yol açıyordu. Google mühendisleri tam tersi bir seçim yaptı: \textbf{üs alanını FP32 ile aynı} (8 bit) tutup \textbf{kesir alanını} 23 bitten 7 bite indirdiler (Şekil~\ref{fig:fp32-bf16-fp16}).

\begin{figure}[htbp]
\centering
\resizebox{\linewidth}{!}{%
\begin{tikzpicture}[x=1cm, y=1cm]
  % 32 bit oku (üstte)
  \draw[<->, thick] (0, 1.35) -- (16.0, 1.35);
  \node[font=\scriptsize] at (8.0, 1.6) {32 bit};

  % FP32
  \node[font=\small\bfseries, anchor=east] at (-0.3, 0.45) {FP32:};
  \draw[fill=blokkirmizi, draw=black, thick] (0,0) rectangle (0.5,0.9);
  \node[font=\tiny\bfseries] at (0.25, 0.45) {1};
  \draw[fill=blokmavi, draw=black, thick] (0.5,0) rectangle (4.5,0.9);
  \node[font=\scriptsize] at (2.5, 0.45) {8 bit üs};
  \draw[fill=blokyesil, draw=black, thick] (4.5,0) rectangle (16.0,0.9);
  \node[font=\scriptsize] at (10.25, 0.45) {23 bit kesir};

  % BF16
  \node[font=\small\bfseries, anchor=east] at (-0.3, -0.95) {BF16:};
  \draw[fill=blokkirmizi, draw=black, thick] (0,-1.4) rectangle (0.5,-0.5);
  \node[font=\tiny\bfseries] at (0.25, -0.95) {1};
  \draw[fill=blokmavi, draw=black, thick] (0.5,-1.4) rectangle (4.5,-0.5);
  \node[font=\scriptsize] at (2.5, -0.95) {8 bit üs};
  \draw[fill=blokyesil, draw=black, thick] (4.5,-1.4) rectangle (8.0,-0.5);
  \node[font=\scriptsize] at (6.25, -0.95) {7 bit kesir};

  % FP16
  \node[font=\small\bfseries, anchor=east] at (-0.3, -2.35) {FP16:};
  \draw[fill=blokkirmizi, draw=black, thick] (0,-2.8) rectangle (0.5,-1.9);
  \node[font=\tiny\bfseries] at (0.25, -2.35) {1};
  \draw[fill=blokmavi, draw=black, thick] (0.5,-2.8) rectangle (3.0,-1.9);
  \node[font=\scriptsize] at (1.75, -2.35) {5 bit üs};
  \draw[fill=blokyesil, draw=black, thick] (3.0,-2.8) rectangle (8.0,-1.9);
  \node[font=\scriptsize] at (5.5, -2.35) {10 bit kesir};

  % Oklar ve açıklamalar
  \draw[dashed, gray] (4.5, -0.1) -- (4.5, -0.4);
  \node[font=\scriptsize, gray, anchor=west] at (4.7, -0.25) {üs alanı aynı};

  % 16 bit oku (altta)
  \draw[<->, thick] (0, -3.15) -- (8.0, -3.15);
  \node[font=\scriptsize] at (4.0, -3.4) {16 bit};
\end{tikzpicture}}
\caption[FP32, BF16 ve FP16 bit düzenleri]{FP32, BF16 ve FP16 bit düzenlerinin karşılaştırması. BF16, FP32 ile aynı 8 bitlik üs alanını korurken kesir alanını 7 bite indirmiştir.}
\label{fig:fp32-bf16-fp16}
\end{figure}

Bu tasarım kararının sonuçları:
\begin{itemize}
  \item \textbf{Sayı aralığı FP32 ile aynı:} 8 bitlik üs alanı sayesinde $\approx 10^{\pm 38}$ aralık korunur. FP16'daki taşma sorunu ortadan kalkar.
  \item \textbf{Duyarlık düşer:} 7 bitlik kesir alanı yalnızca $\approx 2{-}3$ onluk basamak duyarlık sağlar (FP32'nin 7 basamağına karşı). Ancak yapay zeka ağırlıkları için bu çoğu zaman yeterlidir.
  \item \textbf{FP32'ye dönüşüm önemsiz:} BF16, FP32'nin üst 16 bitinin kesilmesiyle (\emph{truncation}) elde edilebilir; ek bir dönüşüm devresi gerekmez. Bu, karma duyarlık eğitiminde iki biçim arasında geçişi çok ucuz kılar.
\end{itemize}

BF16 ilk olarak Google'ın TPU v2 ve v3 hızlandırıcılarında kullanıldı. Kısa sürede NVIDIA (Ampere mimarisi, 2020), AMD, Intel ve ARM de BF16 desteği ekledi. Bugün yapay zeka eğitiminin varsayılan standart biçimi haline gelmiştir.

\subsection{TF32: NVIDIA'nın Uzlaşısı}
\index{TF32}

BF16 ve FP16 bellekte 16 bit kaplar ve FP32'ye göre iki kat hız kazandırır; ancak programcının veri tiplerini açıkça değiştirmesini gerektirir. NVIDIA, 2020'de Ampere mimarisini (A100 GPU) tanıtırken farklı bir yaklaşım izledi: mevcut FP32 kodunu hiç değiştirmeden hızlandırmak. Bunun için \textbf{TF32} (\emph{TensorFloat-32}) adında bir ara biçim tasarladı.

TF32, 1 işaret $+$ 8 üs $+$ 10 kesir $=$ 19 bitten oluşur. FP32'nin 8 bitlik üs alanını koruyarak aynı sayı aralığını ($\approx 10^{\pm 38}$) sunar; kesir alanını ise 23 bitten 10 bite indirerek çarpıcı devresini küçültür. 10 bitlik kesir, FP16 ile aynı duyarlığı ($\approx 3{-}4$ onluk basamak) sağlar ve yapay zeka eğitimi için yeterlidir.

Peki neden doğrudan FP32 kullanılmaz? Çarpma devresinin alanı, kesir genişliğinin \emph{karesiyle} orantılı büyür: $23 \times 23 = 529$ bitlik bir çarpıcı yerine $10 \times 10 = 100$ bitlik bir çarpıcı yaklaşık 5 kat daha küçüktür ve daha az enerji tüketir. Aynı yonga alanına çok daha fazla çarpma birimi yerleştirilebilir; A100'ün Tensör Çekirdeği birimleri bu sayede FP32'ye göre 8 kat daha yüksek işlem gücü sunar.

TF32'nin en önemli özelliği, bellekte ayrı bir veri tipi olarak saklanmamasıdır: veriler bellekte FP32 olarak durur, yalnızca Tensör Çekirdeği çarpma birimleri içinde sessizce TF32'ye kırpılarak çarpılır ve sonuç FP32 olarak döndürülür. Programcının kod değiştirmesine gerek yoktur; bu nedenle NVIDIA bunu ``FP32'nin hızlı modu'' olarak konumlandırmıştır.

\subsection{FP8: Çıkarım İçin En Dar Biçim}
\index{FP8}

Yapay zeka çıkarımında (\emph{inference}), yani eğitilmiş bir modelin yeni veriler üzerinde tahmin yapmasında, duyarlık gereksinimi eğitime göre daha da düşüktür. Bu gözlem, 8 bitlik kayan nokta biçimlerinin geliştirilmesine yol açtı.

FP8 standardı (2022'de NVIDIA, ARM ve Intel tarafından önerilen) iki alt biçim tanımlar:
\begin{itemize}
  \item \textbf{E4M3:} 1 işaret $+$ 4 üs $+$ 3 kesir. Daha yüksek duyarlık, daha dar aralık. Ağırlıklar ve etkinleştirmeler için.
  \item \textbf{E5M2:} 1 işaret $+$ 5 üs $+$ 2 kesir. Daha geniş aralık, daha düşük duyarlık. Gradyanlar için.
\end{itemize}

Bu iki alt biçimin bit düzenleri Şekil~\ref{fig:fp8-bit-duzeni}'nde gösterilmiştir. Her ikisi de 8 bit kullanır; aradaki fark üs ve kesir alanlarının paylaşımındadır.

\begin{figure}[htbp]
\centering
\resizebox{\linewidth}{!}{%
\begin{tikzpicture}[x=1cm, y=1cm]
  % 8 bit oku (üstte)
  \draw[<->, thick] (0, 1.35) -- (8.0, 1.35);
  \node[font=\scriptsize] at (4.0, 1.6) {8 bit};

  % E4M3
  \node[font=\small\bfseries, anchor=east] at (-0.3, 0.45) {E4M3:};
  \draw[fill=blokkirmizi, draw=black, thick] (0,0) rectangle (1.0,0.9);
  \node[font=\tiny\bfseries] at (0.5, 0.45) {1};
  \draw[fill=blokmavi, draw=black, thick] (1.0,0) rectangle (5.0,0.9);
  \node[font=\scriptsize] at (3.0, 0.45) {4 bit üs};
  \draw[fill=blokyesil, draw=black, thick] (5.0,0) rectangle (8.0,0.9);
  \node[font=\scriptsize] at (6.5, 0.45) {3 bit kesir};

  % Açıklama E4M3
  \node[font=\scriptsize, gray, anchor=west] at (8.3, 0.45) {duyarlık $\uparrow$\; aralık $\downarrow$};

  % E5M2
  \node[font=\small\bfseries, anchor=east] at (-0.3, -0.95) {E5M2:};
  \draw[fill=blokkirmizi, draw=black, thick] (0,-1.4) rectangle (1.0,-0.5);
  \node[font=\tiny\bfseries] at (0.5, -0.95) {1};
  \draw[fill=blokmavi, draw=black, thick] (1.0,-1.4) rectangle (6.0,-0.5);
  \node[font=\scriptsize] at (3.5, -0.95) {5 bit üs};
  \draw[fill=blokyesil, draw=black, thick] (6.0,-1.4) rectangle (8.0,-0.5);
  \node[font=\scriptsize] at (7.0, -0.95) {2 bit};

  % Açıklama E5M2
  \node[font=\scriptsize, gray, anchor=west] at (8.3, -0.95) {duyarlık $\downarrow$\; aralık $\uparrow$};

  % BF16 karşılaştırma (altta)
  \node[font=\small\bfseries, anchor=east] at (-0.3, -2.35) {BF16:};
  \draw[fill=blokkirmizi, draw=black, thick] (0,-2.8) rectangle (1.0,-1.9);
  \node[font=\tiny\bfseries] at (0.5, -2.35) {1};
  \draw[fill=blokmavi, draw=black, thick] (1.0,-2.8) rectangle (9.0,-1.9);
  \node[font=\scriptsize] at (5.0, -2.35) {8 bit üs};
  \draw[fill=blokyesil, draw=black, thick] (9.0,-2.8) rectangle (16.0,-1.9);
  \node[font=\scriptsize] at (12.5, -2.35) {7 bit kesir};

  % 16 bit oku
  \draw[<->, thick] (0, -3.15) -- (16.0, -3.15);
  \node[font=\scriptsize] at (8.0, -3.4) {16 bit};

\end{tikzpicture}}
\caption[FP8 bit düzenleri: E4M3 ve E5M2]{FP8'in iki alt biçiminin bit düzenleri, BF16 ile karşılaştırmalı. E4M3, kesir alanına daha fazla bit ayırarak duyarlığı artırır; E5M2 ise üs alanını genişleterek sayı aralığını büyütür.}
\label{fig:fp8-bit-duzeni}
\end{figure}

Bu dar biçimlerin neden önemli olduğunu somutlaştırmak için günümüz yapay zeka modellerinin boyutlarına bakalım. GPT-3 modeli 175 milyar, LLaMA~3 modeli 405 milyar, bazı büyük modeller ise 1 trilyonun üzerinde parametreye sahiptir. Her parametre bir kayan nokta sayısıdır ve bellekte saklanması gerekir. Şekil~\ref{fig:model-bellek}, farklı kayan nokta biçimlerinin bellek gereksinimini nasıl etkilediğini göstermektedir.

\begin{figure}[htbp]
\centering
\begin{tikzpicture}
  \begin{axis}[
    ybar,
    bar width=14pt,
    width=\linewidth,
    height=7cm,
    ylabel={Bellek gereksinimi (GB)},
    symbolic x coords={LLaMA 3.1 8B, LLaMA 3.1 70B, GPT-3 175B, LLaMA 3.1 405B},
    xtick=data,
    x tick label style={font=\footnotesize, align=center},
    y tick label style={font=\footnotesize},
    ylabel style={font=\small},
    ymin=0,
    ymax=1700,
    ytick={0, 200, 400, 600, 800, 1000, 1200, 1400, 1600},
    legend style={at={(0.02,0.98)}, anchor=north west, font=\footnotesize, draw=none, fill=white, fill opacity=0.8, text opacity=1},
    enlarge x limits=0.15,
    grid=major,
    grid style={gray!20},
    nodes near coords,
    nodes near coords style={font=\tiny, anchor=north, inner sep=1pt, yshift=-2pt},
  ]
  % FP32
  \addplot[fill=blokkirmizi!70, draw=blokkirmizi!90] coordinates {
    (LLaMA 3.1 8B, 32) (LLaMA 3.1 70B, 280) (GPT-3 175B, 700) (LLaMA 3.1 405B, 1620)
  };
  % BF16
  \addplot[fill=blokmavi!60, draw=blokmavi!80] coordinates {
    (LLaMA 3.1 8B, 16) (LLaMA 3.1 70B, 140) (GPT-3 175B, 350) (LLaMA 3.1 405B, 810)
  };
  % FP8
  \addplot[fill=blokyesil!60, draw=blokyesil!80] coordinates {
    (LLaMA 3.1 8B, 8) (LLaMA 3.1 70B, 70) (GPT-3 175B, 175) (LLaMA 3.1 405B, 405)
  };
  \legend{FP32 (32 bit), BF16 (16 bit), FP8 (8 bit)}
  \end{axis}
\end{tikzpicture}
\caption[Yapay zeka modellerinin kayan nokta biçimine göre bellek gereksinimleri]{Yapay zeka modellerinin farklı kayan nokta biçimlerindeki bellek gereksinimleri. Bit genişliğini yarıya indirmek, bellek gereksinimini de yarıya indirir.}
\label{fig:model-bellek}
\end{figure}

Örneğin 405 milyar parametreli LLaMA~3.1 modeli FP8 ile 405\,GB belleğe sığar; bu, 80\,GB bellekli beş GPU ile çalıştırılabilecek bir boyuttur. Aynı model FP32 ile 1\,620\,GB gerektirir ve 20'den fazla GPU'ya ihtiyaç duyulur. Bit genişliğini daraltmak yalnızca bellek tasarrufu sağlamaz; aynı zamanda bellekten işlemciye veri aktarımını da hızlandırır; çünkü aynı bant genişliğiyle iki veya dört kat daha fazla parametre taşınabilir.

NVIDIA H100 GPU, FP8 modunda FP16'ya göre iki kat daha yüksek işlem gücü sunar.

Peki bu kadar bit kısmak modelin doğruluğunu ne kadar etkiler? Daha önce gördüğümüz gibi bit sayısını azaltmak yuvarlama hatasını artırır; ancak yapay zeka modellerinde ağırlıklar zaten gürültülüdür ve küçük yuvarlama hataları modelin genel davranışını çoğu zaman değiştirmez. Araştırmalar şunu göstermiştir:

\begin{itemize}
\item \textbf{FP32 $\rightarrow$ BF16:} Model doğruluğunda kayıp neredeyse sıfırdır (\%0{,}1'in altında). Bunun nedenini anlamak için tipik bir yapay zeka modelinin ağırlık dağılımına bakmak yeterlidir. Şekil~\ref{fig:agirlik-dagilimi}'de görüldüğü gibi ağırlıkların büyük çoğunluğu sıfıra yakın dar bir aralıkta yoğunlaşır; aşırı büyük veya küçük değerler çok nadirdir. Bu dağılımda ağırlıkların \%95'inden fazlası yalnızca 2--3 onluk basamak duyarlıkla yeterince temsil edilebilir; BF16'nın sunduğu duyarlık tam da bu kadardır. Bu nedenle BF16, eğitimde FP32'nin yerini büyük ölçüde almıştır.

\begin{figure}[htbp]
\centering
\begin{tikzpicture}
  \begin{axis}[
    width=\linewidth,
    height=7cm,
    xlabel={Ağırlık değeri},
    ylabel={Bağıl sıklık},
    xlabel style={font=\small},
    ylabel style={font=\small},
    x tick label style={font=\footnotesize},
    y tick label style={font=\footnotesize},
    xmin=-0.4, xmax=0.4,
    ymin=0, ymax=42,
    xtick={-0.4, -0.3, -0.2, -0.1, 0, 0.1, 0.2, 0.3, 0.4},
    ytick={0, 10, 20, 30, 40},
    grid=major,
    grid style={gray!20},
    axis lines=left,
    legend style={at={(0.98,0.98)}, anchor=north east, font=\footnotesize, draw=none, fill=white, fill opacity=0.85, text opacity=1},
  ]
  % GPT-2 / Transformer (geniş kuyruklu, σ ≈ 0.08) — önce çiz (arkada kalır)
  \addplot[draw=blokkirmizi, very thick, smooth, domain=-0.4:0.4, samples=100]
    {14 * exp(-x^2 / (2*0.08^2))};
  \addlegendentry{GPT-2 Transformer ($\sigma \approx 0{,}08$)}

  % BERT (orta, σ ≈ 0.05)
  \addplot[draw=blokyesil!80!black, very thick, smooth, domain=-0.4:0.4, samples=100, densely dashed]
    {23 * exp(-x^2 / (2*0.05^2))};
  \addlegendentry{BERT ($\sigma \approx 0{,}05$)}

  % ResNet-50 (CNN, dar dağılım, σ ≈ 0.03) — en son çiz (önde görünsün)
  \addplot[draw=blokmavi!80!black, very thick, smooth, domain=-0.4:0.4, samples=100]
    {38 * exp(-x^2 / (2*0.03^2))};
  \addlegendentry{ResNet-50 ($\sigma \approx 0{,}03$)}

  % Açıklama
  \node[font=\scriptsize, align=center, text=gray] at (axis cs:0.28, 35) {Tüm modellerde\\ağırlıkların \%95'inden\\fazlası $\pm 2\sigma$ içinde};

  \end{axis}
\end{tikzpicture}
\caption[Farklı yapay zeka modellerinin ağırlık dağılımları]{Farklı yapay zeka modellerinin tipik ağırlık dağılımları. Tüm modellerde ağırlıklar sıfır etrafında çan eğrisi biçiminde yoğunlaşır; CNN modelleri (ResNet) daha dar, Transformer modelleri (GPT, BERT) daha geniş kuyruklu dağılım gösterir. Bu yoğunlaşma, düşük bitli kayan nokta biçimlerinin yeterli duyarlık sağlayabilmesinin temel nedenidir.}
\label{fig:agirlik-dagilimi}
\end{figure}
\item \textbf{BF16 $\rightarrow$ FP8:} Doğruluk kaybı genellikle \%0{,}5--1 aralığındadır. E4M3 biçiminde yalnızca 3 bitlik kesir alanı ($\approx 1{,}5$ onluk basamak) kaldığından, aşırı büyük veya küçük ağırlıklarda kırpma (\emph{clipping}) gerekebilir. Bunu telafi etmek için eğitim sırasında \emph{niceleme farkındalıklı eğitim} (\emph{quantization-aware training}) kullanılır: model, düşük duyarlıkta çalışacağını ``bilerek'' eğitilir ve ağırlıklarını bu sınırlara göre uyarlar.
\item \textbf{Çıkarımda daha aşırı niceleme:} Eğitilmiş bir modelin ağırlıkları, çıkarım sırasında 4 bitlik tam sayılara (INT4) kadar indirilebilir. Bu durumda doğruluk kaybı \%1--3'e çıkabilir; ancak bellek gereksinimi FP32'ye göre 8 kat azalır. Bu teknik, büyük dil modellerinin tek bir tüketici GPU'sunda çalıştırılabilmesini sağlamıştır.
\end{itemize}

\subsection{Karma Duyarlık Eğitimi}
\index{karma duyarlık}\index{mixed precision|see{karma duyarlık}}

Önceki alt bölümlerde gördüğümüz biçimlerin her birinin güçlü ve zayıf yanları vardır: FP32 yüksek duyarlık sunar ama yavaş ve bellek açısından pahalıdır; BF16 hızlı ve verimlidir ama duyarlığı düşüktür. Peki neden tek bir biçim seçmek yerine ikisini birden kullanmıyoruz?

Sorun şudur: eğitim sırasında ağırlık güncellemeleri (gradyanlar) genellikle çok küçük değerlerdir; örneğin $10^{-5}$ mertebesinde. Bir ağırlık $1{,}0$ iken $0{,}00001$ eklemek istediğimizde, BF16'nın yalnızca 2--3 onluk basamak duyarlığı bu küçük farkı temsil edemez ve güncelleme \emph{yuvarlanarak kaybolur}. Öte yandan tüm hesaplamayı FP32 ile yapmak, BF16'ya göre iki kat daha fazla bellek ve bant genişliği gerektirir; milyarlarca parametreli bir modelde bu fark devasa boyutlara ulaşır.

Çözüm, iki biçimin güçlü yanlarını birleştiren \emph{karma duyarlık} (\emph{mixed precision}) yöntemidir. Tipik bir karma duyarlık eğitim döngüsü şöyle çalışır:

\begin{enumerate}
  \item Ağırlıkların bir \textbf{ana kopyası FP32'de} tutulur (\emph{master weights}); küçük gradyan güncellemelerinin kaybolmaması için.
  \item \textbf{İleri ve geri yayılım BF16} (veya FP16) ile yapılır; hızlı ve bellek açısından verimlidir.
  \item Hesaplanan gradyanlar \textbf{FP32'ye dönüştürülüp} ana ağırlıklara uygulanır; yuvarlama hatası birikmesi önlenir.
\end{enumerate}

Bu döngü Şekil~\ref{fig:karma-duyarlik}'te gösterilmiştir.

\begin{figure}[htbp]
\centering
\begin{tikzpicture}[scale=0.85, transform shape,
    blok/.style={draw, rectangle, fill=blokmavi!15, thick, font=\footnotesize,
                 align=center, minimum width=4.2cm, minimum height=1cm,
                 rounded corners=2pt},
    fp32/.style={draw, rectangle, fill=blokturuncu!20, thick, font=\footnotesize,
                 align=center, minimum width=4.2cm, minimum height=1cm,
                 rounded corners=2pt},
    okstil/.style={-{Stealth[length=2.5mm]}, thick}
]
    % Bloklar
    \node[fp32] (master) at (0, 0)
        {Ana ağırlıklar\\{\scriptsize (FP32)}};
    \node[blok] (donustur) at (5.5, 0)
        {BF16'ya dönüştür};
    \node[blok] (ileri) at (11, 0)
        {İleri yayılım\\{\scriptsize (BF16)}};
    \node[blok] (geri) at (11, -2.5)
        {Geri yayılım\\{\scriptsize (gradyanlar, BF16)}};
    \node[fp32] (guncelle) at (0, -2.5)
        {Ağırlıkları güncelle\\{\scriptsize (FP32)}};

    % Oklar
    \draw[okstil] (master) -- (donustur);
    \draw[okstil] (donustur) -- (ileri);
    \draw[okstil] (ileri) -- (geri)
        node[midway, right=3pt, font=\scriptsize] {kayıp hesapla};
    \draw[okstil] (geri) -- (guncelle)
        node[midway, below=3pt, font=\scriptsize] {gradyanları FP32'ye dönüştür};
    \draw[okstil] (guncelle) -- (master)
        node[midway, left=3pt, font=\scriptsize] {güncelle};

    % Renk açıklamaları
    \node[font=\scriptsize, anchor=west] at (3.5, -4) {\colorbox{blokturuncu!20}{\strut FP32} yüksek duyarlık gereken adımlar};
    \node[font=\scriptsize, anchor=west] at (3.5, -4.5) {\colorbox{blokmavi!15}{\strut BF16} hız ve bellek gereken adımlar};

\end{tikzpicture}
\caption[Karma duyarlık eğitim döngüsü]{Karma duyarlık eğitim döngüsü. Ana ağırlıklar FP32'de tutularak küçük güncellemelerin kaybolması önlenir; hesaplama yoğun ileri ve geri yayılım ise BF16 ile hızlı yapılır.}
\label{fig:karma-duyarlik}
\end{figure}

Bu yaklaşım, FP32'nin duyarlığını korurken BF16'nın hız ve bellek avantajından yararlanır. Günümüzün neredeyse tüm büyük dil modelleri (GPT, Gemini, Claude, LLaMA vb.) karma duyarlık ile eğitilmektedir.

\subsection{Karşılaştırma}

Tüm bu biçimlerin bir arada karşılaştırması Tablo~\ref{tab:yapay-zeka-kn-bicimleri}'nde verilmiştir.

\begin{table}[htbp]
\centering
\caption[Kayan nokta biçimlerinin karşılaştırması]{Kayan nokta biçimlerinin karşılaştırması. Duyarlık sütunu yaklaşık onluk basamak sayısını gösterir.}
\label{tab:yapay-zeka-kn-bicimleri}
\small
\begin{tabularx}{\textwidth}{l c c c c c X}
\toprule
\textbf{Biçim} & \textbf{Bit} & \textbf{Üs} & \textbf{Kesir} & \textbf{Eklenti} & \textbf{Duyarlık} & \textbf{Kullanım} \\
\midrule
FP64 (double) & 64 & 11 & 52 & 1023 & $\sim$16 & Bilimsel hesaplama \\
FP32 (float)  & 32 & 8  & 23 & 127  & $\sim$7  & Genel amaçlı \\
TF32          & 19 & 8  & 10 & 127  & $\sim$3  & GPU Tensör Çekirdeği (iç biçim) \\
BF16          & 16 & 8  & 7  & 127  & $\sim$2  & YZ eğitimi \\
FP16          & 16 & 5  & 10 & 15   & $\sim$3  & Grafik, hafif YZ eğitimi \\
FP8 (E5M2)    & 8  & 5  & 2  & 15   & $\sim$1  & YZ gradyanları \\
FP8 (E4M3)    & 8  & 4  & 3  & 7    & $\sim$1  & YZ ağırlıkları, çıkarım \\
\bottomrule
\end{tabularx}
\end{table}

Görüldüğü gibi, yapay zeka donanımları duyarlıktan ödün vererek \emph{aynı yonga alanına daha fazla hesaplama birimi sığdırma} ve \emph{aynı bant genişliğiyle daha fazla veri taşıma} ilkesiyle çalışır. Bu bölümde öğrendiğimiz toplayıcı ve çarpıcı devreleri, bu biçimlerin her birinde aynı temel ilkelerle (yalnızca daha dar bit genişliğiyle) çalışır.


% =============================================
% Yanılgılar ve Tuzaklar
% =============================================
\section{Yanılgılar ve Tuzaklar}
\label{sec:aritmetik-yanilgilar}

Sayısal aritmetik, donanım düzeyinde çalıştığı için hatalar genellikle sessizce oluşur ve yazılım katmanında fark edilmesi güçtür.

\yanilgi{Kayan nokta aritmetiği gerçek aritmetik gibi çalışır}
Kayan nokta sayıları sınırlı duyarlığa sahiptir. Matematikteki toplama işleminin değişme ve birleşme özellikleri kayan noktada her zaman geçerli değildir: $(a + b) + c \neq a + (b + c)$ olabilir. Bu durum, büyük ölçekli bilimsel hesaplamalarda ve yapay zeka eğitiminde ciddi doğruluk sorunlarına yol açabilir.

\yanilgi{Taşma yalnızca büyük sayılarla ortaya çıkar}
İşaretli aritmetikte iki artı sayının toplamı eksi çıkabilir (örneğin 8 bitte $+120 + 10 = -126$). Bu durum yalnızca büyük sayılarda değil, veri tipi sınırlarına yaklaşıldığında her zaman oluşabilir. 2004'te Comair havayolu şirketinin tüm uçuşlarını iptal etmesi, bir 16 bitlik sayaç taşmasından kaynaklanmıştır.

\yanilgi{Çarpma ve bölme işlemleri aynı sürede tamamlanır}
Çarpma, çağdaş işlemcilerde genellikle 3--5 saat çevriminde tamamlanırken, bölme 20--40 çevrim sürebilir. Bu nedenle derleyiciler mümkün olduğunca bölmeyi kaydırma veya çarpma ile değiştirmeye çalışır. Örneğin $x / 4$ yerine $x \gg 2$ kullanmak çok daha hızlıdır.

\yanilgi{Daha dar kayan nokta biçimi her zaman daha hızlıdır}
BF16 ve FP8 gibi dar biçimler tek başına daha hızlı çarpma ve daha az bellek sunar; ancak duyarlık kaybı birikebilir. Yapay zeka eğitiminde yalnızca BF16 kullanmak, küçük gradyan güncellemelerinin yuvarlanarak kaybolmasına yol açar. Bu nedenle uygulamada tek bir dar biçim değil, karma duyarlık yöntemi kullanılır; hız gereken yerde dar biçim, duyarlık gereken yerde FP32.

\tuzak{İşaretsiz ve işaretli karşılaştırmayı karıştırmak}
C dilinde \texttt{-1 > 2u} ifadesi \emph{doğrudur} çünkü derleyici işaretsiz karşılaştırma yapar ve $-1$, ikiye tümleyende en büyük işaretsiz sayıya ($2^{32}-1$) dönüşür. Bu tür tür dönüşüm hataları güvenlik açıklarının yaygın kaynaklarından biridir.

\tuzak{IEEE 754 özel değerlerini göz ardı etmek}
Sıfıra bölme sonucu oluşan \texttt{Inf} (sonsuz) ve tanımsız işlemlerden kaynaklanan \texttt{NaN} (Not a Number) değerleri, denetlenmezse programda sessizce yayılarak tüm hesaplamaları geçersiz kılabilir. \texttt{NaN == NaN} ifadesinin \emph{yanlış} döndürmesi, yeni programcılar için sık karşılaşılan bir sürprizdir.

\tuzak{Kayan nokta sayılarını doğrudan eşitlik ile karşılaştırmak}
Yuvarlama hataları nedeniyle \texttt{if (x == 0.3)} gibi doğrudan eşitlik denetimleri neredeyse her zaman başarısız olur. Örneğin \texttt{0.1 + 0.2} ifadesinin sonucu \texttt{0.30000000000000004} olduğundan \texttt{0.3} ile eşit çıkmaz. Doğru yöntem, iki sayının farkının küçük bir eşik değerinden ($\varepsilon$) az olup olmadığını denetlemektir: $|a - b| < \varepsilon$.

\tuzak{Tam sayı ve kayan nokta sıfıra bölmeyi karıştırmak}
Kayan nokta aritmetiğinde sıfıra bölme sessizce $\pm\infty$ üretir ve program çalışmaya devam eder. Tam sayı aritmetiğinde ise sıfıra bölme bir kural dışı durum (\emph{exception}) tetikler; RISC-V'te sonuç olarak özel bir değer yazılır, x86'da ise program çökebilir. Bu iki davranışın karıştırılması, özellikle tam sayı ile kayan nokta arasında tür dönüşümü yapılan kodlarda tehlikeli hatalara yol açar.


% =============================================
\section{Gerçek Dünyada Aritmetik}
\label{sec:gercek-dunyada-aritmetik}

Bu bölümde öğrendiğimiz toplayıcı, çarpıcı ve kayan nokta devreleri bugün her işlemcinin ayrılmaz bir parçasıdır. Ancak bu her zaman böyle değildi. Kayan nokta biriminin ayrı bir yonga olarak satıldığı, hatta birçok bilgisayarda hiç bulunmadığı bir dönem yaşanmıştır.

\subsection{Yardımcı İşlemci Dönemi (1980'ler)}
Intel'in 8086 işlemcisi (1978) kayan nokta donanımına sahip değildi. Kayan nokta işlemleri gereken kullanıcılar, ayrı bir yonga olan \textbf{Intel 8087}\index{Intel 8087} yardımcı işlemcisini (\emph{coprocessor}) satın alıp ana karta takardı. 8087, 80 bitlik genişletilmiş duyarlıkta kayan nokta aritmetiği sunan ve kendi yazmaç yığıtına sahip bağımsız bir yongaydı. 8087 takılı değilse, aynı işlemler yazılımla (tamsayı buyrukları kullanılarak) taklit edilirdi; bu ise 10 ila 100 kat daha yavaştı.

Bu ayrık yapı 80286/80287 ve 80386/80387 çiftlerinde de sürdü. Bilgisayar reklamlarında ``matematik yardımcı işlemci yuvası mevcut'' ifadesi bir satış noktasıydı; çünkü 8087 tek başına birkaç yüz dolar tutuyordu ve birçok kullanıcı bu ek maliyetten kaçınıyordu. Sonuç olarak, yazılım geliştiricileri programlarının kayan nokta birimi olmadan da çalışacağını garanti etmek zorundaydı.

\subsection{Tümleşik Kayan Nokta Birimi (1989--)}
Intel, \textbf{80486DX}\index{Intel 80486DX} işlemcisiyle (1989) kayan nokta birimini ilk kez işlemci yongasının içine aldı. Bu geçiş, Moore yasasının\index{Moore yasası} sağladığı transistör bütçesi artışıyla mümkün oldu: 80386'daki 275\,000 transistörden 80486'daki 1,2 milyona atlama, kayan nokta birimini de içerecek yeterli alanı sağladı. O günden bu yana her genel amaçlı işlemci tümleşik kayan nokta birimleriyle donatılmaktadır.

\subsection{Günümüz: SIMD, Vektör ve Matris Birimleri}
Bugünün işlemcileri tek seferde yalnızca bir kayan nokta işlemi yapan birimlerle yetinmez. ARM mimarisindeki \textbf{NEON}\index{NEON (ARM)} ve \textbf{SVE/SVE2}\index{SVE (ARM)} uzantıları, tek bir buyrukla birden fazla kayan nokta işlemi gerçekleştiren SIMD (\emph{Single Instruction, Multiple Data}\index{SIMD}) birimleri sunar. Apple'ın M serisi\index{Apple M serisi} yongaları (M1--M5) bu ARM uzantılarını yüksek başarımlı çekirdeklerinde kullanır; ayrıca \textbf{AMX}\index{AMX (Apple)} (\emph{Apple Matrix coprocessor}) adlı özel bir matris çarpma birimi barındırır. AMX, bu bölümde öğrendiğimiz çarpma-toplama birimlerinden oluşan bir dizilimdir ve yapay öğrenme iş yüklerini hızlandırmak için tasarlanmıştır.

x86 dünyasında da benzer bir evrim yaşanmıştır: MMX (1997), SSE (1999), AVX (2011), AVX-512 (2017) ve en son AVX10/AMX (2023) ile aritmetik birimlerin genişliği ve sayısı sürekli artmıştır. NVIDIA ve Google gibi şirketler ise bu bölümde gördüğümüz Tensör Çekirdekleri\index{Tensör Çekirdeği} ve TPU\index{TPU} gibi özel amaçlı aritmetik hızlandırıcılarla bu evrimi daha da ileri taşımıştır.

\begin{terimnotu}
\textbf{yapay öğrenme} (\emph{machine learning})\index{yapay öğrenme}\index{machine learning|see{yapay öğrenme}}\,: Türkçeye sıklıkla ``makine öğrenmesi'' olarak aktarılır. Ancak bu çeviri dilbilgisel olarak hatalıdır: Türkçede ``makine öğrenmesi'' biçiminde bir isim tamlaması kurulmaz; tıpkı ``çocuk öğrenmesi'' demediğimiz gibi. Doğru tamlama ``makinenin öğrenmesi'' ya da ``makineyle öğrenme'' olurdu. TDK'nın önerdiği karşılık \emph{yapay öğrenme} olup ``yapay zekâ'' ile koşut bir yapıdadır.
\end{terimnotu}

Bu tarihsel yolculuk bir noktayı vurgular: bu bölümde öğrendiğimiz temel yapı taşları (toplayıcılar, çarpıcılar ve kayan nokta devreleri) on yıllar boyunca değişmeden kalmıştır. Değişen, bu yapı taşlarından kaç tanesinin tek bir yonga üzerine sığdırılabildiğidir.

% =============================================
\section{Özet}

Bu bölümde bilgisayarın ``hesap makinesi'' kısmını, yani aritmetik işlemleri nasıl gerçekleştirdiğini, adım adım gördük. Bölüme başlarken bu konunun günümüz yapay zeka donanımları için ne denli temel olduğunu vurgulamıştık; şimdi bunu somut sayılarla pekiştirelim.

Toplamadan başlayarak, basamaklar arasındaki elde geçişini dalga eldeli ve önceden elde üreten toplayıcılarla çözdük. Çıkarmanın ikiye tümleyen yardımıyla toplamaya dönüştürüldüğünü ve bu sayede aynı donanımın hem toplama hem de çıkarma için kullanılabildiğini öğrendik. Taşma gibi kural dışı durumları da ele aldık.

Mantık işlemlerini de ele aldıktan sonra tüm bu parçaları bir araya getirerek 1 bitlik AMB tasarladık. Bu 1 bitlik AMB'den 32 tanesini yan yana koyarak RISC-V'in temel buyruklarını yürütebilen 32 bitlik bir AMB elde ettik.

Çarpma işlemini tekrarlanan toplama ve kaydırmaya dayanan algoritmalarla gerçekleştirdik; işaretli sayılarla doğrudan çarpım yapabilen Booth algoritmasını inceledik. Bölme için ise geri yüklemeli, geri yüklemesiz ve SRT algoritmalarını gördük.

Tam sayıların ötesine geçerek kayan nokta gösterimini öğrendik. IEEE~754 standardı, eklentili üs gösterimi, özel değerler (sıfır, sonsuz, NaN) ve yuvarlama yöntemlerini gördük. Koruma, yuvarlama ve yapışkan bitlerinin yuvarlama kararını nasıl belirlediğini, kayan nokta toplama ve çarpmasının adım adım nasıl gerçekleştirildiğini inceledik. RISC-V'te kayan nokta işlemleri F ve D uzantılarıyla desteklenir.

Bölümün son kısmında yapay zeka çağının getirdiği yeni kayan nokta biçimlerini ele aldık: BF16, TF32, FP8 gibi dar biçimlerin neden geliştirildiğini, bu biçimlerin bellek ve hız avantajlarını, ağırlık dağılımının düşük bitli gösterimi neden mümkün kıldığını ve karma duyarlık eğitiminin bu biçimleri nasıl bir arada kullandığını gördük.

Bu bölümde tasarladığımız toplayıcı ve çarpıcı devreleri, günümüzün en güçlü yongalarının temel yapı taşlarıdır. Tablo~\ref{tab:yonga-aritmetik}'te bazı güncel yongalardaki aritmetik birim sayıları verilmiştir.

\begin{table}[htbp]
\centering
\caption{Güncel yongalardaki aritmetik birim sayıları}
\label{tab:yonga-aritmetik}
\begin{tabular}{l r r l}
\toprule
\textbf{Yonga} & \textbf{Çarpma-Toplama Birimi} & \textbf{Duyarlık} & \textbf{Kullanım} \\
\midrule
NVIDIA H100 GPU     & 16\,896 & FP32 & Yapay zeka eğitimi \\
NVIDIA B200 GPU     & 18\,432 & FP32 & Yapay zeka eğitimi \\
Google TPU v5e      & 16\,384 & BF16 & Yapay zeka eğitimi \\
Apple M4 GPU        & 1\,280  & FP32 & Taşınabilir cihazlar \\
Intel Gaudi 3       & 1\,536  & BF16 & Yapay zeka çıkarımı \\
\bottomrule
\end{tabular}
\end{table}

Görüldüğü gibi bir GPU veya yapay zeka hızlandırıcısı, bu bölümde öğrendiğimiz toplayıcı ve çarpıcı birimlerinden on binlercesini bir yonga üzerinde barındırır. Her bir çarpma-toplama birimi, çarpma ve toplama işlemini tek saat darbesinde gerçekleştiren birleşik bir devredir. Bu birimlerin hızı ve enerji verimliliği, doğrudan bu bölümde gördüğümüz tasarım kararlarına (önceden elde üreten toplayıcılar, Wallace ağacı çarpıcılar, kayan nokta yuvarlama devreleri) bağlıdır.

Artık elimizde çalışan bir AMB var. Bir sonraki bölümde bu AMB'yi kullanarak bir veri yolunun ve onu denetleyen birimin nasıl tasarlandığını göreceğiz. Bu işlemlerin RISC-V programlarındaki kullanımı için Bölüm~\ref{chap:programlama}'te bakınız.


% =============================================
\section*{Alıştırmalar}
\addcontentsline{toc}{section}{Alıştırmalar}

\begin{alistirma}[\orta]
Aşağıdaki işlemleri yapan 2 adet 4 bitlik veri girişi, 1 elde girişi ve 2 bitlik işlem girişi olan bir Aritmetik Mantık Birimi (AMB) tasarlayın: Toplama, Çıkarma, VE, VEYA. Tasarımınızı yaparken elinizde bir toplayıcı olduğunu varsayabilirsiniz.
\end{alistirma}

\begin{alistirma}[\zor]
8 bitlik, A ve B olarak iki girişi olan, 2'ye tümleyen yöntemiyle çıkarma (A$-$B), toplama, Dışlayan VEYA, VEYA, VE, A'nın değili, B'nin değili ve SIFIR (giriş ne olursa olsun sonuç sıfır), A, B, Bir (giriş ne olursa olsun sonuç 1), A'nın ikiye tümleyeni, B'nin ikiye tümleyeni işlemlerini yapan bir AMB tasarlayın, devrelerini çizin ve doğruluk tablolarını gösterin.
\end{alistirma}

\begin{alistirma}[\zor]
Birisi 4 bitlik diğeri 8 bitlik iki yazmaç ve 4 bitlik bir AMB kullanarak 4 bitlik işaretsiz iki sayıyı çarpan bir devrenin çizelgesini gösterin ve çarpma algoritmasını açıklayın. Algoritmanın çalıştığını $1001 \times 1100$ işleminin sonucunu adım adım yazarak gösterin.
\end{alistirma}

\begin{alistirma}[\orta]
Önceden elde üreten toplayıcının sıradan toplayıcılara göre ne farkı vardır? Çalışma mantığı nedir? Toplayıcının her bir bitindeki girişlere göre ilgili 1 bitlik toplayıcının yaptığı işlemi açıklayın.
\end{alistirma}

\begin{alistirma}[\orta]
Kayan nokta işlemlerinde toplama nasıl yapılır? $1{,}1001 \times 2^{-2}$ sayısı ile $1{,}1101 \times 2^{-4}$ sayısını toplayarak kayan nokta toplamasının aşamalarını açıklayın. Bu tür işlemlerde taşma olabilir mi? Olabilirse kaç tür taşma olabilir?
\end{alistirma}

\begin{alistirma}[\kolay]
Kayan nokta işlemleri yapılırken bir sayı sıfıra bölündüğünde ne sonucu çıkar? Sıfır sıfıra bölündüğünde (0/0) ise ne sonucu çıkar? Bu sonuçlar IEEE 754 biçiminde nasıl gösterilir?
\end{alistirma}

\begin{alistirma}[\orta]
Aşağıdaki onluk tabanda gösterilen sayıları 32 bitlik IEEE-754 biçiminde gösterin:
\begin{enumerate}
\item $135{,}8243$
\item $1{,}125$
\item $6{,}375$
\item $11{,}125$
\item $0{,}00025$
\item $17{,}0625$
\end{enumerate}
\end{alistirma}

\begin{alistirma}[\zor]
Eğer bir işlemci tasarımcısının elinde 16 bitlik bir alan varsa bu alanla gösterebileceği ondalıklı sayıların sınırı nedir? 16 bitlik alanda işaretli ondalıklı sayıların gösterilmesi için bu alanın nasıl bölünmesi gerekir? $135{,}8343$ sayısının, 16 bitlik alanda ikilik tabanda virgülden sonra en fazla kaç basamağı gösterilebilir?
\end{alistirma}

\begin{alistirma}[\zor]
Ondalıklı sayıların işlemci içinde saklanması için toplam 5 bitlik saklama birimi kullanılacaktır. Saklama biriminin en soldaki biti işaret için ayrılacak geri kalan 4 bit ise sayının gösterilmesi için kullanılacaktır.
\begin{enumerate}
\item Eğer 4 bitlik kısım yalnızca anlamlı kısmı göstermek için kullanılırsa gösterilebilecek sayılar hangi aralıkta olur?
\item Eğer 4 bitlik kısım yalnızca üst göstermek için kullanılırsa gösterilebilecek sayılar hangi aralıkta olur? Gösterilebilecek en küçük duyarlılık ne olur?
\item Eğer 2 bitlik kısım üst için 2 bitlik kısım da sayının anlamlı parçası için kullanılırsa gösterilebilecek sayılar hangi aralıkta olur?
\end{enumerate}
Aralıklarda hem gösterilebilecek en küçük ve en büyük sayıları hem de duyarlılık sınırlarını gösterin.
\end{alistirma}

\begin{alistirma}[\orta]
Kayan nokta işlemlerinde toplama nasıl yapılır? $1{,}0011 \times 2^{-3}$ sayısı ile $1{,}1101 \times 2^{2}$ sayısını toplayarak kayan nokta toplamasının aşamalarını açıklayın. Bu tür işlemlerde taşma olabilir mi? Olursa kaç tür taşma olabilir?
\end{alistirma}

\begin{alistirma}[\orta]
Booth algoritmasını kullanarak $(-13) \times (11)$ çarpma işlemini 8 bitlik ikiye tümleyen gösteriminde adım adım gerçekleştirin. Her adımda A yazmacının, Q yazmacının ve sayacın değerlerini gösteren bir tablo oluşturun.
\end{alistirma}

\begin{alistirma}[\orta]
Booth algoritmasının standart ve değiştirilmiş (radix-4) sürümlerini kullanarak $15 \times 7$ çarpma işlemini gerçekleştirin. Her iki yöntemde kaç toplama/çıkarma işlemi gerektiğini karşılaştırın.
\end{alistirma}

\begin{alistirma}[\zor]
4 bitlik iki işaretsiz sayının çarpılması sırasında oluşan kısmi çarpımları yazın. Bu kısmi çarpımları Wallace ağacı yöntemiyle indirgeyerek sonucu bulun. Wallace ağacının dizi çarpıcıya göre üstünlüğünü gecikme açısından açıklayın.
\end{alistirma}

\begin{alistirma}[\orta]
4$\times$4 bitlik bir dizi çarpıcının yapısını çizin. $A = 1011$ ve $B = 1101$ sayılarını bu çarpıcı üzerinde çarpın. Toplam kaç tam toplayıcı ve yarım toplayıcı gerektiğini hesaplayın.
\end{alistirma}

\begin{alistirma}[\orta]
Geri yüklemeli bölme algoritmasını kullanarak $(11010100)_2$ sayısını $(1011)_2$ sayısına bölün. 3. Yol donanımını kullanarak her adımda kalan yazmacının ve bölüm yazmacının değerlerini gösteren bir tablo oluşturun.
\end{alistirma}

\begin{alistirma}[\orta]
Geri yüklemesiz bölme algoritmasını kullanarak $(10110)_2$ sayısını $(101)_2$ sayısına bölün. Her adımda kalan yazmacının değerini ve yapılan işlemi (toplama/çıkarma) belirtin. Geri yüklemeli yönteme göre üstünlüğünü açıklayın.
\end{alistirma}

\begin{alistirma}[\kolay]
Aşağıdaki 32 bitlik IEEE 754 gösterimlerinin hangi özel değeri temsil ettiğini belirleyin:
\begin{enumerate}[(a)]
\item 0 11111111 00000000000000000000000
\item 1 11111111 00000000000000000000000
\item 0 11111111 10000000000000000000000
\item 0 00000000 00000000000000000000000
\item 1 00000000 00000000000000000000000
\item 0 00000000 00000000000000000000001
\end{enumerate}
\end{alistirma}

\begin{alistirma}[\kolay]
IEEE 754 tek duyarlı gösterimde (32 bit) gösterilebilen en küçük pozitif olağanlaştırılmamış sayının değerini hesaplayın. Bu değeri olağanlaştırılmış gösterimdeki en küçük pozitif sayıyla karşılaştırın.
\end{alistirma}

\begin{alistirma}[\orta]
IEEE 754 tek duyarlı gösterimde $1{,}101 \times 2^3$ ile $1{,}011 \times 2^2$ sayılarını çarpın. Üs toplama, anlamlı kısım çarpma, olağanlaştırma ve yuvarlama adımlarını ayrı ayrı gösterin.
\end{alistirma}

\begin{alistirma}[\orta]
IEEE 754 tek duyarlıklı gösterimde $1{,}110 \times 2^{5}$ sayısını $1{,}100 \times 2^{2}$ sayısına bölün. Üs çıkarma, anlamlı kısım bölme, olağanlaştırma ve yuvarlama adımlarını ayrı ayrı gösterin.
\end{alistirma}

\begin{alistirma}[\orta]
IEEE 754 tek duyarlı gösterimde $1{,}0001 \times 2^5$ ile $1{,}1100 \times 2^1$ sayılarını toplayın. Üs hizalama, toplama, olağanlaştırma ve yuvarlama adımlarını gösterin. Hizalama sırasında kaybedilen bitlerin sonucu nasıl etkilediğini tartışın.
\end{alistirma}

\begin{alistirma}[\orta]
4 bitlik iki sayı $A = 1011$ ve $B = 0110$ toplandığında, üretme ($G_i$) ve yayma ($P_i$) sinyallerini her bit için hesaplayın. Ardından elden önceden üretme formüllerini kullanarak $C_1$, $C_2$, $C_3$ ve $C_4$ elde sinyallerini doğrudan hesaplayın.
\end{alistirma}

\begin{alistirma}[\zor]
16 bitlik bir önceden elde üreten toplayıcı, dört adet 4 bitlik önceden elde üretici bloğu ve bir adet ikinci düzey önceden elde üretme birimi kullanılarak tasarlanır.
\begin{enumerate}[(a)]
\item Her 4 bitlik blok için grup üretme ($G^*$) ve grup yayma ($P^*$) sinyallerinin nasıl hesaplandığını açıklayın.
\item 16 bitlik toplayıcının toplam kapı gecikmesini hesaplayın (bir kapının gecikmesi $t_g$ olsun).
\item Aynı boyuttaki zincirleme elde toplayıcının gecikmesiyle karşılaştırın.
\end{enumerate}
\end{alistirma}

\begin{alistirma}[\orta]
Aşağıdaki 8 bitlik ikiye tümleyen sayı çiftleri için toplama ve çıkarma işlemlerini gerçekleştirin. Her işlemde taşma olup olmadığını belirleyin ve taşma tespit kuralını açıklayın.
\begin{enumerate}[(a)]
\item $A = 01100100$, $B = 00110010$ (toplama)
\item $A = 01110000$, $B = 01010000$ (toplama)
\item $A = 10010000$, $B = 10110000$ (toplama)
\item $A = 01100000$, $B = 10100000$ (çıkarma: $A - B$)
\end{enumerate}
\end{alistirma}

\begin{alistirma}[\zor]
Bir işlemci 16 bitlik IEEE 754 yarım duyarlı (half-precision) kayan nokta biçimi kullanır (1 işaret, 5 üs, 10 anlamlı kısım biti).
\begin{enumerate}[(a)]
\item Bu biçimde gösterilebilen en büyük sonlu artı sayıyı hesaplayın.
\item $(17{,}125)_{10}$ sayısını bu biçime çevirin.
\item Bu biçimde iki sayı toplanırken taşma (overflow) ve alttan taşma (underflow) durumlarını birer örnekle açıklayın.
\end{enumerate}
\end{alistirma}

\begin{alistirma}[\orta]
Dört bitlik $A = 1011$ ve $B = 0110$ işlenenleri için önceden elde üreten toplayıcı tasarımında kullanılan $P$ (yayılım) ve $G$ (üretim) değerlerini hesaplayın.
\begin{enumerate}[(a)]
\item Her bit konumu için $P_i = A_i \oplus B_i$ ve $G_i = A_i \cdot B_i$ değerlerini bulun.
\item $C_4$'ü doğrudan $P_i$, $G_i$ ve $C_0 = 0$ kullanarak (herhangi bir $C_i$ beklemeden) hesaplayın.
\item Bu sonucu dalgalanmalı taşımalı toplayıcıyla (RCA) karşılaştırın: kaç kapı gecikmesi farkı vardır?
\end{enumerate}
\end{alistirma}

\begin{alistirma}[\zor]
Booth kodlaması kullanarak $(-6) \times 5$ çarpma işlemini gerçekleştirin. İşlenenler 8 bitlik ikiye tümleyen biçimindedir ($(-6)_{10} = 11111010_2$, $5_{10} = 00000101_2$).
\begin{enumerate}[(a)]
\item Başlangıç değerlerini ayarlayın (çarpan, çarpılan, kısmi çarpım yazmacı).
\item Her adımda Booth çizelgesine göre toplama, çıkarma ya da işlem yapmama kararını verin ve adım adım tabloyu doldurun.
\item Sonucu 16 bitlik ikiye tümleyen biçiminde gösterin ve 10'lu tabanla doğrulayın.
\end{enumerate}
\end{alistirma}

\begin{alistirma}[\orta]
$(-13{,}625)_{10}$ sayısını IEEE 754 tek duyarlıklı (32 bit) kayan nokta biçimine dönüştürün.
\begin{enumerate}[(a)]
\item Sayının tam kısmını ve ondalık kısmını ayrı ayrı ikili tabana çevirin.
\item Olağanlaştırılmış biçimi ($1{,}\text{anlamlı\_kısım} \times 2^{\text{üs}}$) yazın.
\item İşaret, eklentili üs (bias = 127) ve anlamlı kısım alanlarını belirleyerek 32 bitlik ikili kalıbı oluşturun.
\end{enumerate}
\end{alistirma}

\begin{alistirma}[\orta]
Aşağıdaki 32 bitlik IEEE 754 tek duyarlıklı kayan nokta kalıbını 10'lu tabana çevirin:
\[
\texttt{1\;10000100\;10110000000000000000000}
\]
\begin{enumerate}[(a)]
\item İşaret, üs ve anlamlı kısım alanlarını belirleyin.
\item Gerçek üs değerini hesaplayın (önyargı = 127).
\item Sonucu ondalık sayı olarak yazın.
\end{enumerate}
\end{alistirma}

\begin{alistirma}[\zor]
IEEE 754 tek duyarlıklı biçimde $A = 1{,}5 \times 2^{3}$ ve $B = 1{,}25 \times 2^{1}$ sayılarını elle toplayın.
\begin{enumerate}[(a)]
\item Üsleri eşitlemek için küçük üslü sayının anlamlı kısmını kaydırın.
\item Anlamlı kısımları toplayın.
\item Gerekiyorsa sonucu olağanlaştırın ve en yakın çift (round-to-nearest-even) kuralıyla yuvarlayın.
\item Nihai sonucu 32 bitlik IEEE 754 kalıbı olarak yazın.
\end{enumerate}
\end{alistirma}

\begin{alistirma}[\kolay]
Bir RISC-V işlemcisinde 32 bitlik AMB (Aritmetik Mantık Birimi) aşağıdaki denetim sinyalleri için hangi işlemi gerçekleştirir? Her satır için işlem adını ve RISC-V buyruk karşılığını (mümkünse) yazın.

\begin{table}[htbp]
\centering
\small
\begin{tabular}{c c c c l l}
\toprule
\textbf{AMBİşlem} & \textbf{İnvert\,A} & \textbf{İnvert\,B} & \textbf{Taşıma\,Gir.} & \textbf{İşlem} & \textbf{Buyruk} \\
\midrule
AND  & 0 & 0 & 0 & ? & ? \\
OR   & 0 & 0 & 0 & ? & ? \\
TOPLA & 0 & 0 & 0 & ? & ? \\
TOPLA & 0 & 1 & 1 & ? & ? \\
\bottomrule
\end{tabular}
\end{table}
\end{alistirma}

\begin{alistirma}[\orta]
Aşağıdaki 4-bit işaretli (ikiye tümleyen) toplama ve çıkarma işlemlerinin her birinde taşma oluşup oluşmadığını belirleyin. Taşma tespiti için taşıma girişi ile çıkışı karşılaştırma kuralını ($C_{in} \neq C_{out}$) uygulayın.
\begin{enumerate}[(a)]
\item $0111 + 0001$ (artı + artı)
\item $1000 + 1111$ (eksi + eksi)
\item $0110 + 0101$ (artı + artı)
\item $1001 - 0111$ (çıkarma: ikiye tümleyeni ekle)
\end{enumerate}
\end{alistirma}

\begin{alistirma}[\zor]
Dört adet 4-bit sayıyı ($A$, $B$, $C$, $D$) koşut olarak çarpan bir Wallace ağacı çarpıcısının kısmi çarpım matrisini oluşturun. $A = 0101$, $B = 0011$, $C = 0110$, $D = 0001$ için:
\begin{enumerate}[(a)]
\item Her çift için ($A \times B$ ve $C \times D$) kısmi çarpım matrisini yazın.
\item Wallace ağacının ilk indirgeme katmanında hangi tam toplayıcılar (FA) ve yarım toplayıcılar (HA) kullanılır?
\item Toplam kaç FA ve HA gerekir? Bir dalgalanmalı çarpıcıyla karşılaştırın.
\end{enumerate}
\end{alistirma}

\begin{alistirma}[\orta]
$13 \div 3$ bölme işlemini 1.\ yol geri yüklemeli bölme algoritmasıyla adım adım gerçekleştirin. İşlenenler 4 bittir; bölen ve kalan yazmaçları 8 bittir.
\begin{enumerate}[(a)]
\item Başlangıç değerlerini tabloya yazın (bölen sola hizalanmış, kalan = bölünen, bölüm = 0).
\item Her adımda Kalan $-$ Bölen işlemini yapın; sonuç eksiyse geri yükleyin ve bölümü 0-bitleyin, artıysa kalan güncellenip bölümü 1-bitleyin.
\item Böleni sağa kaydırın; 4 adım sonra bölüm ve kalanı ondalıkla doğrulayın.
\end{enumerate}
\end{alistirma}

\begin{alistirma}[\orta]
8 bitlik elde ile seçmeli toplayıcıyı 12 bite genişletmek istiyorsunuz. Üç adet 4 bitlik blok kullanacaksınız: Blok~0 (bit [3:0]) normal, Blok~1 (bit [7:4]) ve Blok~2 (bit [11:8]) elde ile seçmeli yapıda.
\begin{enumerate}[(a)]
\item Blok diyagramını çizin. Kaç adet 4 bitlik toplayıcı ve çoklayıcı gerekir?
\item Dalga eldeli 12 bitlik toplayıcıya göre gecikme karşılaştırması yapın ($t_{mux} = t_g$ alın).
\item Hızlanma oranını ve alan artışını hesaplayın.
\end{enumerate}
\end{alistirma}

\begin{alistirma}[\kolay]
Yapay zeka uygulamalarında (derin öğrenme, sinir ağı eğitimi ve çıkarımı) IEEE 754 yarım duyarlı (FP16, 16 bit) kayan nokta biçimi neden tercih edilir?
\begin{enumerate}[(a)]
\item Tek duyarlı (FP32) ile yarım duyarlı (FP16) biçimlerini bellek kullanımı ve bant genişliği açısından karşılaştırın.
\item FP16'nın daha düşük duyarlığına karşın derin öğrenmede neden yeterli olduğunu açıklayın.
\item NVIDIA Tensör Çekirdeği gibi donanımların FP16 işlemlerini neden çok daha hızlı gerçekleştirebildiğini tartışın.
\end{enumerate}
\end{alistirma}
